\documentclass[11pt,a4paper,oneside,openany]{book}
\usepackage[utf8]{inputenc}
\usepackage[T1]{fontenc}
\usepackage[brazil]{babel}
\usepackage{amsmath,amsthm,mathtools}
\usepackage{newpxtext,newpxmath}
\usepackage[scaled=0.9]{helvet}  % sem serifa (a TeX Gyre Heros do newpxtext pode faltar)
\usepackage{microtype}
\usepackage[a4paper,top=2.6cm,bottom=2.6cm,inner=2.5cm,outer=2.5cm,headheight=14pt]{geometry}
\usepackage{enumitem}
\usepackage{booktabs,tabularx,multirow,array}
\usepackage[dvipsnames,table]{xcolor}
\usepackage[most]{tcolorbox}
\usepackage{tikz}
\usetikzlibrary{arrows.meta,positioning,shapes.geometric,shapes.multipart,calc,decorations.pathreplacing,
  matrix,fit,backgrounds,patterns,bending,quotes}
\usepackage{pgfplots}
\pgfplotsset{compat=1.18, /pgf/number format/.cd, use comma, 1000 sep={.}}
\usepackage{forest}
\usepackage{graphicx}
\usepackage{caption}
\captionsetup{font=small,labelfont={bf,color=MidnightBlue},labelsep=period,skip=4pt}
\usepackage[noend]{algpseudocode}
\usepackage{titlesec}
\usepackage{fancyhdr}
\usepackage[colorlinks=true,linkcolor=MidnightBlue,urlcolor=MidnightBlue]{hyperref}
\usepackage{bookmark}

% ------------------------------------------------------------------
% Aparência de livro
% ------------------------------------------------------------------
\definecolor{capcor}{named}{MidnightBlue}
\titleformat{\chapter}[display]
  {\normalfont\bfseries\color{capcor}}
  {\filleft\fontsize{48}{48}\selectfont\thechapter}
  {6pt}
  {\titlerule[1.2pt]\vspace{6pt}\filright\Huge}
  [{\vspace{2pt}\titlerule[0.4pt]}]
\titleformat{name=\chapter,numberless}[display]
  {\normalfont\bfseries\color{capcor}}{}{0pt}
  {\titlerule[1.2pt]\vspace{6pt}\filright\Huge}[{\vspace{2pt}\titlerule[0.4pt]}]
\titlespacing*{\chapter}{0pt}{-10pt}{24pt}
\titleformat{\section}{\normalfont\Large\bfseries\color{capcor}}{\thesection}{0.8em}{}
\titleformat{\subsection}{\normalfont\large\bfseries}{\thesubsection}{0.7em}{}
\titleformat{\part}[display]{\normalfont\bfseries\centering\color{capcor}}
  {\Large\partname~\thepart}{14pt}{\Huge}

\pagestyle{fancy}
\fancyhf{}
\makeatletter
\renewcommand{\chaptermark}[1]{\markboth{\if@mainmatter\thechapter.\ \fi #1}{}}
\makeatother
\renewcommand{\sectionmark}[1]{\markright{\thesection\ \ #1}}
\fancyhead[L]{\small\textsc{\nouppercase{\leftmark}}}
\fancyhead[R]{\small\thepage}
\renewcommand{\headrulewidth}{0.4pt}
\fancypagestyle{plain}{\fancyhf{}\fancyfoot[C]{\small\thepage}\renewcommand{\headrulewidth}{0pt}}

\setcounter{tocdepth}{1}

% ------------------------------------------------------------------
% Macros
% ------------------------------------------------------------------
\newcommand{\N}{\mathbb{N}}
\newcommand{\R}{\mathbb{R}}
\newcommand{\bal}{\operatorname{bal}}
\newcommand{\rank}{\operatorname{rank}}
\newcommand{\Find}{\textsc{Find}}
\newcommand{\Union}{\textsc{Union}}
\newcommand{\MakeSet}{\textsc{Make-Set}}
\newcommand{\lgs}{\log^{*}}
\DeclarePairedDelimiter{\floor}{\lfloor}{\rfloor}
\DeclarePairedDelimiter{\ceil}{\lceil}{\rceil}
\DeclarePairedDelimiter{\abs}{\lvert}{\rvert}
\newcommand{\am}{$^{*}$}   % marca de amortizado

% Pseudocódigo em português (algpseudocode)
\algrenewcommand\algorithmicif{\textbf{se}}
\algrenewcommand\algorithmicthen{\textbf{então}}
\algrenewcommand\algorithmicelse{\textbf{senão}}
\algrenewcommand\algorithmicwhile{\textbf{enquanto}}
\algrenewcommand\algorithmicdo{\textbf{faça}}
\algrenewcommand\algorithmicfor{\textbf{para}}
\algrenewcommand\algorithmicreturn{\textbf{retorna}}
\algrenewcommand\algorithmicprocedure{\textbf{Procedimento}}
\algrenewcommand\algorithmicfunction{\textbf{Função}}
\algrenewcommand\algorithmicrepeat{\textbf{repita}}
\algrenewcommand\algorithmicuntil{\textbf{até que}}
\algrenewcommand{\algorithmiccomment}[1]{\hfill{\color{gray}\small\% #1}}
\algnewcommand\algorithmicelsif{\textbf{senão se}}
\algdef{C}[IF]{IF}{ElsIf}[1]{\algorithmicelsif\ #1\ \algorithmicthen}
\newcommand{\ptr}{\mathbin{\uparrow}}
\newcommand{\nil}{\lambda}

% ------------------------------------------------------------------
% Teoremas e caixas
% ------------------------------------------------------------------
\theoremstyle{plain}
\newtheorem{teo}{Teorema}[chapter]
\newtheorem{prop}[teo]{Proposição}
\newtheorem{lema}[teo]{Lema}
\newtheorem{cor}[teo]{Corolário}
\theoremstyle{definition}
\newtheorem{defi}[teo]{Definição}
\newtheorem{exemplo}[teo]{Exemplo}
\theoremstyle{remark}
\newtheorem{obs}[teo]{Observação}

% Pseudocódigo
\newtcolorbox{pseudo}[1]{enhanced, breakable, colback=gray!3, colframe=gray!45, colbacktitle=gray!22, coltitle=black,
  fonttitle=\bfseries\small, title={#1}, left=3pt,right=3pt,top=2pt,bottom=2pt, fontupper=\small}
% Abertura de capítulo
\newtcolorbox{objetivos}{enhanced, breakable, colback=capcor!3, colframe=capcor!3, boxrule=0pt,
  left=8pt,right=6pt,top=5pt,bottom=5pt, borderline west={3pt}{0pt}{capcor!70},
  fontupper=\small, title={Neste capítulo}, coltitle=capcor!80!black, colbacktitle=capcor!3,
  fonttitle=\bfseries\small, attach title to upper={\par\smallskip}}
% O essencial de uma seção ou de um capítulo
\newtcolorbox{chave}[1][Pontos-chave]{enhanced, breakable, colback=MidnightBlue!4,
  colframe=MidnightBlue!70!black, fonttitle=\bfseries, title={#1}, left=4pt,right=4pt,top=3pt,bottom=3pt}
% Métodos e roteiros
\newtcolorbox{roteiro}[1][Roteiro]{enhanced, breakable, colback=ForestGreen!5,
  colframe=ForestGreen!70!black, fonttitle=\bfseries, title={#1}, left=4pt,right=4pt,top=3pt,bottom=3pt}
% Erros frequentes
\newtcolorbox{cuidado}[1][Armadilhas comuns]{enhanced, breakable, colback=BrickRed!4,
  colframe=BrickRed!75!black, fonttitle=\bfseries, title={#1}, left=4pt,right=4pt,top=3pt,bottom=3pt}
% Aprofundamento
\newtcolorbox{complemento}[1][Para saber mais]{enhanced, breakable, colback=orange!4, colframe=orange!70!black,
  fonttitle=\bfseries, title={#1}, left=4pt,right=4pt,top=3pt,bottom=3pt}
% A ideia por trás de um conceito, antes da formalização
\newtcolorbox{intuicao}[1][Intuição]{enhanced, breakable, colback=Goldenrod!7, colframe=Goldenrod!7, boxrule=0pt,
  left=8pt,right=6pt,top=4pt,bottom=4pt, borderline west={3pt}{0pt}{Goldenrod!85!black},
  title={#1}, coltitle=Goldenrod!40!black, colbacktitle=Goldenrod!7, fonttitle=\bfseries\small,
  attach title to upper={\par\smallskip}}
% Exemplos resolvidos ficam numa caixa leve, para se destacarem do texto corrido
\tcolorboxenvironment{exemplo}{enhanced, breakable, colback=ForestGreen!3, colframe=ForestGreen!3, boxrule=0pt,
  left=6pt,right=6pt,top=2pt,bottom=4pt, borderline west={2pt}{0pt}{ForestGreen!60!black},
  before skip=8pt, after skip=8pt}

% Figuras no lugar em que aparecem (não flutuam), com legenda numerada
% (o tipo é declarado no início, para que o link da legenda aponte para o topo da figura)
\NewDocumentEnvironment{figura}{o}
  {\par\addvspace{6pt}\noindent\begin{minipage}{\linewidth}\centering\IfValueT{#1}{\captionsetup{type=figure}}}
  {\IfValueT{#1}{\par\smallskip\caption{#1}}\end{minipage}\par\addvspace{6pt}}

% ------------------------------------------------------------------
% Estilos TikZ usados nas figuras
% ------------------------------------------------------------------
\colorlet{destaque}{Goldenrod!45}      % célula/vértice em foco
\colorlet{feito}{ForestGreen!25}       % já processado / definitivo
\colorlet{ruim}{BrickRed!25}           % violação / conflito
\colorlet{azul}{MidnightBlue!18}
\tikzset{
  cel/.style={draw, minimum width=7mm, minimum height=6mm, inner sep=0pt, font=\small},
  idx/.style={font=\scriptsize, text=gray},
  vtx/.style={circle, draw, thick, minimum size=6.5mm, inner sep=0pt, font=\small},
  vtxs/.style={circle, draw, thick, minimum size=5mm, inner sep=0pt, font=\scriptsize},
  aresta/.style={thick},
  arco/.style={thick, -{Stealth[length=2.2mm]}},
  peso/.style={font=\scriptsize, fill=white, inner sep=1pt},
  marca/.style={line width=2.2pt, draw=MidnightBlue!70, opacity=0.65},
  seta/.style={-{Stealth[length=2mm]}, thick},
  rot/.style={font=\scriptsize, text=gray!60!black},
  passo/.style={font=\small\bfseries, text=MidnightBlue},
}
% \vetor[posições em destaque]{a,b,c,...}: desenha um vetor indexado a partir de 1
\NewDocumentCommand{\vetor}{O{} m}{%
\begin{tikzpicture}[baseline=-0.6ex]
  \foreach \h in {#1} {\fill[destaque] ($(\h*7mm,0)+(-3.5mm,-3mm)$) rectangle ++(7mm,6mm);}
  \foreach \x [count=\i] in {#2} {\node[cel] at (\i*7mm,0) {$\x$}; \node[idx] at (\i*7mm,4.8mm) {\i};}
\end{tikzpicture}}
% Dentro de um tikzpicture: \pic at (x,y) {arr={k}{a,b,c}} desenha k células centradas em (x,y)
\tikzset{pics/arr/.style n args={2}{code={\foreach \arrv [count=\arri] in {#2}
  {\node[cel, minimum width=5mm, minimum height=4.8mm, font=\scriptsize] at ({(\arri-(#1+1)/2)*5mm},0) {$\arrv$};}}}}
% \vetorz: o mesmo, indexado a partir de 0
\NewDocumentCommand{\vetorz}{O{} m}{%
\begin{tikzpicture}[baseline=-0.6ex]
  \foreach \h in {#1} {\fill[destaque] ($(\h*7mm,0)+(-3.5mm,-3mm)$) rectangle ++(7mm,6mm);}
  \foreach \x [count=\i from 0] in {#2} {\node[cel] at (\i*7mm,0) {$\x$}; \node[idx] at (\i*7mm,4.8mm) {\i};}
\end{tikzpicture}}

% ------------------------------------------------------------------
% Exercícios com link para a resposta (Apêndice B)
% ------------------------------------------------------------------
% \begin{exercicio}{rótulo}[origem]<ref1,ref2>: os rótulos entre < > apontam a seção,
% o exemplo ou a figura que o exercício pratica e aparecem no cabeçalho como "ver §X".
\newcounter{exerc}[chapter]
\renewcommand{\theexerc}{\thechapter.\arabic{exerc}}
\makeatletter
\def\ap@sec{sec}\def\ap@ex{ex}\def\ap@fig{fig}\def\ap@teo{teo}\def\ap@lema{lema}
\def\ap@umaref#1:#2\@nil{\def\ap@tipo{#1}%
  \ifx\ap@tipo\ap@sec \S\ref{#1:#2}%
  \else\ifx\ap@tipo\ap@ex Exemplo~\ref{#1:#2}%
  \else\ifx\ap@tipo\ap@fig Figura~\ref{#1:#2}%
  \else\ifx\ap@tipo\ap@teo Teorema~\ref{#1:#2}%
  \else\ifx\ap@tipo\ap@lema Lema~\ref{#1:#2}%
  \else\ref{#1:#2}\fi\fi\fi\fi\fi}
\newcommand{\verref}[1]{ver~\def\ap@sep{}%
  \@for\ap@r:=#1\do{\ap@sep\def\ap@sep{ e~}\expandafter\ap@umaref\ap@r\@nil}}
\makeatother
\NewDocumentEnvironment{exercicio}{m o !d<>}{%
  \par\addvspace{7pt}\refstepcounter{exerc}\label{ex:#1}%
  \noindent{\color{capcor}\bfseries\theexerc}%
  \IfValueT{#2}{\enspace{\small\color{gray!75!black}\textsf{[#2]}}}%
  \IfValueT{#3}{\enspace{\small\color{gray!75!black}(\verref{#3})}}%
  \hfill{\scriptsize\hyperref[resp:#1]{resposta: p.~\pageref*{resp:#1}}}%
  \par\nopagebreak\noindent\ignorespaces}{\par}
\newcommand{\resp}[1]{\par\addvspace{6pt}\noindent\phantomsection\label{resp:#1}%
  \hyperref[ex:#1]{\textcolor{capcor}{\textbf{\ref*{ex:#1}}}}\enspace\ignorespaces}
\newcommand{\respcap}[1]{\subsection*{Capítulo~\ref*{#1}: \nameref*{#1}}}
\newcommand{\estrela}{$\bigstar$}
\newcommand{\excurso}{\subsection*{Exercícios do curso}}
\newcommand{\excompl}{\subsection*{Exercícios complementares}}

% Árvores (forest)
\forestset{
  bst/.style={for tree={circle, draw, minimum size=1.5em, inner sep=1pt, font=\small,
      s sep=6mm, l sep=4mm, edge={thick}}},
  bstp/.style={for tree={circle, draw, minimum size=1.4em, inner sep=0.5pt, font=\scriptsize,
      s sep=2.2mm, l sep=2.2mm, edge={thick}}},
  bb/.style={label={[font=\scriptsize,text=BrickRed,inner sep=1pt]right:{$#1$}}},
  feito/.style={fill=ForestGreen!25},
  aviol/.style={edge={very thick, BrickRed}},
  acam/.style={edge={very thick, MidnightBlue!80}},
  vazio/.style={phantom},
  tri/.style={draw, isosceles triangle, shape border rotate=90, isosceles triangle apex angle=60,
      inner sep=0.5pt, minimum height=7mm, font=\scriptsize, anchor=apex, child anchor=apex},
  novo/.style={fill=yellow!40},
  desreg/.style={fill=BrickRed!25},
}

\setlist{itemsep=1.5pt,topsep=3pt}
\renewcommand{\arraystretch}{1.15}
\numberwithin{equation}{chapter}

\begin{document}

% ------------------------------------------------------------------
% Folha de rosto
% ------------------------------------------------------------------
\frontmatter
\hypersetup{pageanchor=false}  % a folha de rosto não tem número; evita duas páginas "i"
\begin{titlepage}
\centering
\vspace*{2.5cm}
{\color{capcor}\rule{\linewidth}{1.5pt}\par\vspace{16pt}
{\fontsize{40}{46}\selectfont\bfseries Algoritmos\par}
\vspace{12pt}
{\Large Estruturas de dados, paradigmas de projeto\\[2pt] e algoritmos em grafos\par}
\vspace{12pt}\rule{\linewidth}{0.5pt}\par}
\vspace{2cm}
{\LARGE Anizio Silva Correia Júnior\par}
\vspace{14pt}
{\large com orientação e material didático do\\[2pt] Prof.\ Uéverton Souza\par}
\vspace{10pt}
{\normalsize IMPA Tech\par}
\vfill
\begin{tcolorbox}[enhanced, width=0.82\linewidth, colback=capcor!4, colframe=capcor!60, boxrule=0.6pt,
  halign=center, left=8pt, right=8pt, top=6pt, bottom=6pt]
\small\textbf{Produção auxiliada por inteligência artificial.}\\[2pt]
Este texto foi elaborado por Anizio Silva Correia Júnior com o auxílio de ferramentas de inteligência artificial,
a partir da orientação e do material didático proporcionados pelo Prof.\ Uéverton Souza.
\end{tcolorbox}
\vspace{1cm}
\end{titlepage}
\hypersetup{pageanchor=true}

\chapter{Prefácio}
Este livro apresenta os fundamentos de projeto e análise de algoritmos na sequência em que eles são estudados na disciplina de Algoritmos do IMPA Tech: começa pela linguagem da análise de complexidade, passa pelas estruturas de dados clássicas (filas de prioridades, conjuntos disjuntos, árvores de busca), pela ordenação, pela seleção e pela técnica de divisão e conquista, e chega à programação dinâmica, aos algoritmos gulosos e aos algoritmos em grafos.

Cada capítulo trata de um tópico. Ele começa com um quadro que diz o que será estudado, desenvolve o assunto com definições, algoritmos em pseudocódigo, demonstrações e exemplos resolvidos passo a passo, e termina com um resumo e uma lista de exercícios. Os exemplos numéricos foram conferidos por simulação.

\section*{Como o livro está organizado}
A Parte~I trata de análise de algoritmos e de estruturas de dados: complexidade (Capítulo~\ref{cap:complex}), heaps binários, binomiais e de Fibonacci com análise amortizada (Capítulos~\ref{cap:heap} e~\ref{cap:heapav}), conjuntos disjuntos (Capítulo~\ref{cap:uf}), árvores AVL (Capítulo~\ref{cap:avl}), ordenação (Capítulos~\ref{cap:sortcomp} e~\ref{cap:sortlin}), seleção (Capítulo~\ref{cap:selecao}) e divisão e conquista (Capítulo~\ref{cap:dc}). O Capítulo~\ref{cap:revisao} reúne os critérios para escolher entre essas ferramentas.

A Parte~II trata de paradigmas de projeto e de grafos: programação dinâmica (Capítulo~\ref{cap:pd}), o algoritmo de Floyd-Warshall (Capítulo~\ref{cap:fw}), algoritmos gulosos e matroides (Capítulo~\ref{cap:gulosos}), árvores geradoras mínimas (Capítulo~\ref{cap:agm}), busca em largura (Capítulo~\ref{cap:bfs}), o algoritmo de Dijkstra (Capítulo~\ref{cap:dijkstra}) e a busca em profundidade em grafos e digrafos (Capítulos~\ref{cap:dfs} e~\ref{cap:dfsdig}). O Apêndice~\ref{cap:hash} trata de tabelas hash e o Apêndice~\ref{ap:respostas} traz as respostas dos exercícios.

\section*{As caixas}
Ao longo do texto, caixas coloridas ajudam a situar cada trecho:
\begin{objetivos}
O que o capítulo vai apresentar.
\end{objetivos}
\begin{chave}[Pontos-chave]
O que precisa ficar claro ao fim de uma seção ou de um capítulo.
\end{chave}
\begin{roteiro}[Roteiro]
Um método passo a passo para resolver um tipo de problema ou escolher entre alternativas.
\end{roteiro}
\begin{cuidado}[Armadilhas comuns]
Erros frequentes e confusões que costumam aparecer.
\end{cuidado}
\begin{complemento}[Para saber mais]
Aprofundamentos que vão além do essencial.
\end{complemento}

\section*{Os exercícios}
Cada capítulo termina com pelo menos vinte exercícios, divididos em dois grupos. Os \textbf{exercícios do curso} vêm do material da disciplina; uma etiqueta indica a origem de cada um:
\begin{center}\small
\begin{tabular}{@{}ll@{}}
\textsf{[Slides]} & exercícios propostos nos slides da disciplina;\\
\textsf{[Lista]} & listas de exercícios da disciplina;\\
\textsf{[Monitoria]} & exercícios das monitorias;\\
\textsf{[Avaliação]} & questões das avaliações da disciplina;\\
\textsf{[Pesquisa]} & temas dos trabalhos de pesquisa;\\
\textsf{[Material de apoio]} & listas de outras disciplinas indicadas como material complementar.\\
\end{tabular}
\end{center}
Os \textbf{exercícios complementares}, etiquetados \textsf{[Complementar]}, foram escolhidos para fixar e aprofundar o conteúdo; os mais difíceis estão marcados com \estrela. Todos têm resposta, sugestão ou esboço de solução no Apêndice~\ref{ap:respostas}. O link ``resposta: p.~X'' ao lado de cada enunciado leva até ela, e o número da resposta leva de volta ao enunciado.

\section*{Convenções}
\begin{itemize}
  \item Heaps são \textbf{de máximo} e guardados em vetor indexado a partir de 1: o pai de $i$ é $\floor{i/2}$ e os filhos são $2i$ e $2i+1$.
  \item Em árvores AVL, o balanço é $\bal(v)=h_D(v)-h_E(v)$ (altura da direita menos a da esquerda).
  \item Em grafos, $n=|V|$ e $m=|E|$; $N(v)$ é a vizinhança de $v$ e, num digrafo, $N_{out}(v)$ é a vizinhança de saída.
  \item $\log$ sem base indica $\log_2$; o símbolo \am\ indica custo amortizado.
\end{itemize}

\section*{Nota sobre a produção}
Este livro foi produzido por \textbf{Anizio Silva Correia Júnior}, com o \textbf{auxílio de ferramentas de inteligência artificial}, a partir da orientação e do material didático proporcionados pelo \textbf{Prof.\ Uéverton Souza} na disciplina de Algoritmos do IMPA Tech. As fontes utilizadas estão listadas nas Referências bibliográficas, ao final do volume.

\tableofcontents

\mainmatter
\part{Análise de algoritmos e estruturas de dados}
% ================================================================
\chapter{Complexidade de algoritmos}\label{cap:complex}\label{sec:complex}
% ================================================================
\begin{objetivos}
\begin{itemize}
  \item Medir o custo de um algoritmo em função do tamanho da entrada: pior caso, caso médio e melhor caso.
  \item As notações assintóticas $O$, $\Omega$, $\Theta$, $o$ e $\omega$ e suas propriedades.
  \item Analisar laços na prática: como contar as iterações de um trecho de código.
  \item A diferença entre a complexidade de um \emph{algoritmo} e a de um \emph{problema}: limites superiores, inferiores e algoritmos ótimos.
  \item Por que a fronteira entre algoritmos polinomiais e exponenciais importa, e o que é um algoritmo pseudopolinomial.
\end{itemize}
\end{objetivos}

Dois programas que resolvem o mesmo problema podem levar tempos muito diferentes. Considere a tarefa de calcular a soma $1+2+\cdots+N$. Um programa faz um laço que soma os números um a um; outro usa a fórmula de Gauss, $N(N+1)/2$.
\begin{figura}
\begin{minipage}[t]{0.42\linewidth}
\begin{pseudo}{Soma com laço}
\begin{algorithmic}
\State $s:=0$
\For{$i:=1$ \textbf{até} $N$} \State $s:=s+i$ \EndFor
\State \Return $s$
\end{algorithmic}
\end{pseudo}
\end{minipage}\hspace{0.06\linewidth}%
\begin{minipage}[t]{0.42\linewidth}
\begin{pseudo}{Soma com a fórmula}
\begin{algorithmic}
\State \Return $N(N+1)/2$
\end{algorithmic}
\end{pseudo}
\end{minipage}
\end{figura}
Para $N=10$, ambos são instantâneos. Para $N=10^{12}$, o primeiro faz um trilhão de somas e o segundo continua fazendo três operações. A diferença não está na máquina nem na linguagem: está no \emph{número de passos} que cada algoritmo executa e em \emph{como esse número cresce} quando a entrada aumenta.

Este capítulo fixa a linguagem que será usada no resto do livro para fazer esse tipo de comparação. As técnicas para resolver recorrências aparecem no Capítulo~\ref{cap:sortcomp}, junto com o primeiro algoritmo recursivo estudado, e a análise amortizada aparece no Capítulo~\ref{cap:heapav}.

\section{Medidas de complexidade}\label{sec:medidas}
\begin{intuicao}
Medir um algoritmo pelo tempo de relógio é pouco confiável: o mesmo programa roda em tempos diferentes em computadores diferentes. Em vez disso, contamos \textbf{passos elementares} (uma comparação, uma soma, uma atribuição, um acesso a uma posição de vetor), cada um com custo constante, e expressamos o total como uma \emph{função do tamanho da entrada}.
\end{intuicao}

O problema é que entradas do mesmo tamanho podem dar trabalhos diferentes. Fixe um algoritmo $A$ e um tamanho de entrada $n$. Sejam $E_1,\dots,E_m$ todas as entradas possíveis de tamanho $n$ e $T_i$ o número de passos que $A$ executa sobre $E_i$. Há três maneiras naturais de resumir os números $T_1,\dots,T_m$ num só valor:
\begin{itemize}
  \item a complexidade de \textbf{pior caso}, $\max_iT_i$;
  \item a complexidade de \textbf{caso médio}, $\sum_ip_iT_i$, onde $p_i$ é a probabilidade de a entrada ser $E_i$;
  \item a complexidade de \textbf{melhor caso}, $\min_iT_i$.
\end{itemize}

\begin{exemplo}[Os três casos da ordenação por inserção]\label{ex:inscasos}
A ordenação por inserção percorre o vetor da esquerda para a direita e insere cada elemento na posição correta entre os anteriores, que já estão ordenados. Para isso, compara o elemento com os anteriores, do mais próximo para o mais distante, até achar um menor ou chegar ao início.
\begin{figura}[Ordenação por inserção sobre $[3,1,4,2]$. Em cada linha, o elemento destacado é inserido entre os anteriores (já ordenados); à direita, as comparações feitas.]
\begin{tikzpicture}[x=7mm,y=8mm]
\foreach \lin/\vals/\h/\cmt in {0/{3,1,4,2}/2/{1 comparação},1/{1,3,4,2}/3/{1 comparação},2/{1,3,4,2}/4/{3 comparações},3/{1,2,3,4}/0/{}}{
  \foreach \x [count=\i] in \vals {
    \ifnum\i=\h \fill[destaque] (\i-0.5,-\lin-0.35) rectangle (\i+0.5,-\lin+0.35); \fi
    \node[cel] at (\i,-\lin) {\x};}
  \node[rot,anchor=west] at (5,-\lin) {\cmt};}
\node[rot,anchor=east] at (0.3,0) {$j=2$};
\node[rot,anchor=east] at (0.3,-1) {$j=3$};
\node[rot,anchor=east] at (0.3,-2) {$j=4$};
\node[rot,anchor=east] at (0.3,-3) {fim};
\end{tikzpicture}
\end{figura}
Para $n=4$, as 24 ordens possíveis de quatro números distintos dão:
\begin{itemize}
  \item \textbf{melhor caso}: o vetor já ordenado, $[1,2,3,4]$; cada elemento é comparado só com o vizinho e não se move: $3$ comparações, e em geral $n-1$;
  \item \textbf{pior caso}: o vetor decrescente, $[4,3,2,1]$; cada elemento é comparado com todos os anteriores: $1+2+3=6$ comparações, e em geral $n(n-1)/2$;
  \item \textbf{caso médio}, supondo as 24 ordens igualmente prováveis: a média é $59/12\approx4{,}92$ comparações (a entrada $[3,1,4,2]$ da figura faz 5).
\end{itemize}
O mesmo algoritmo é linear em entradas favoráveis e quadrático nas desfavoráveis.
\end{exemplo}

O pior caso é a medida padrão: ele dá uma \emph{garantia} que vale para toda entrada, e não depende de hipóteses sobre a distribuição das entradas, que nem sempre são realistas. O caso médio é útil quando essa distribuição é conhecida (o quicksort, por exemplo, é rápido em média e lento no pior caso). O melhor caso raramente diz algo útil, mas aparece para mostrar que um algoritmo pode ser rápido em entradas favoráveis, como a ordenação por inserção num vetor já ordenado.

Uma quarta medida, a complexidade \textbf{amortizada}, divide o custo de uma \emph{sequência} de operações sobre uma estrutura de dados pelo número de operações, no pior caso da sequência. Ela não envolve probabilidade e será estudada no Capítulo~\ref{cap:heapav}.

\subsection*{Tamanho da entrada}
O tamanho de uma entrada é o comprimento da sua codificação, em bits ou em símbolos. Para uma lista de $n$ números de tamanho limitado, o tamanho é proporcional a $n$. Mas quando a entrada é um número $N$, ela ocupa apenas cerca de $\log_2 N$ bits (ou $\log_{10}N$ dígitos), e um algoritmo cujo custo cresce com o \emph{valor} de $N$ pode ser exponencial no tamanho da entrada.

\begin{exemplo}[A soma com laço é exponencial no número de dígitos]\label{ex:somalaco}
Volte à soma $1+2+\cdots+N$. O laço faz $N$ iterações, o que parece ``linear''. Mas a entrada é o próprio número $N$, escrito com $d\approx\log_{10}N$ dígitos:
\begin{center}\small
\begin{tabular}{@{}rrr@{}}
\toprule
dígitos de $N$ & iterações do laço & tempo a $10^9$ iterações/s\\
\midrule
$3$ & $10^3$ & $1$ microssegundo\\
$9$ & $10^9$ & $1$ segundo\\
$15$ & $10^{15}$ & $\approx 11{,}6$ dias\\
$30$ & $10^{30}$ & $\approx3\cdot10^{13}$ anos\\
\bottomrule
\end{tabular}
\end{center}
Cada dígito a mais multiplica o tempo por 10: o custo é $\Theta(10^d)$, exponencial no tamanho $d$ da entrada. A fórmula $N(N+1)/2$, por outro lado, faz uma multiplicação de números com $d$ dígitos, o que custa $O(d^2)$ pelo método escolar: polinomial.
\end{exemplo}

Um algoritmo é \textbf{pseudopolinomial} quando é polinomial no \emph{valor} dos números da entrada (isto é, quando eles são escritos em unário), mas não necessariamente no seu tamanho em bits. A soma com laço é um exemplo artificial; exemplos importantes são o \emph{counting sort} do Capítulo~\ref{cap:sortlin}, de custo $\Theta(n+k)$ para inteiros em $[0,k]$, e a programação dinâmica para a mochila do Capítulo~\ref{cap:pd}, de custo $\Theta(nW)$.

\section{Notações assintóticas}\label{sec:notacoes}
\begin{intuicao}
Quando $n$ é grande, só o termo que cresce mais rápido importa. Em $n^2+10n+20$, para $n=1000$ o termo $n^2$ vale $10^6$, enquanto $10n+20$ vale $10\,020$, cerca de $1\%$ do total. E constantes multiplicativas dependem da máquina: um computador duas vezes mais rápido transforma $2n^2$ em $n^2$. Por isso dizemos simplesmente que esse custo ``é da ordem de $n^2$''. As notações assintóticas tornam essa frase precisa.
\end{intuicao}

As notações abaixo valem para funções $f,g$ não negativas a partir de algum ponto. Leia $f=O(g)$ como ``$f$ cresce no máximo como $g$''.
\begin{center}
\begin{tabular}{@{}llll@{}}
\toprule
Notação & Definição & Via limite (se existir) & Leitura\\
\midrule
$f=O(g)$ & $\exists c>0,n_0:\ f(n)\le c\,g(n)\ \forall n\ge n_0$ & $\lim f/g<\infty$ & ``$\le$''\\
$f=\Omega(g)$ & $\exists c>0,n_0:\ f(n)\ge c\,g(n)\ \forall n\ge n_0$ & $\lim f/g>0$ & ``$\ge$''\\
$f=\Theta(g)$ & $f=O(g)$ e $f=\Omega(g)$ & $\lim f/g=c\in(0,\infty)$ & ``$=$''\\
$f=o(g)$ & $\forall c>0\ \exists n_0:\ f(n)<c\,g(n)\ \forall n\ge n_0$ & $\lim f/g=0$ & ``$<$''\\
$f=\omega(g)$ & $\forall c>0\ \exists n_0:\ f(n)>c\,g(n)\ \forall n\ge n_0$ & $\lim f/g=\infty$ & ``$>$''\\
\bottomrule
\end{tabular}
\end{center}
Formalmente, $O(g)$ é um \emph{conjunto} de funções, e $f=O(g)$ significa $f\in O(g)$; o sinal de igualdade é um abuso de notação consagrado, que deve ser lido sempre da esquerda para a direita ($n=O(n^2)$ está certo; ``$O(n^2)=n$'' não faz sentido).

\subsection*{As notações $O$, $\Omega$ e $\Theta$ numa figura}
A definição de $O$ tem duas constantes, e cada uma tem um papel: $c$ absorve as constantes multiplicativas, e $n_0$ permite ignorar os valores pequenos de $n$, onde termos de ordem menor ainda pesam. A Figura~\ref{fig:otheta} mostra as duas em ação.

\begin{figura}[A função $f(n)=n^2+10n+20$ fica ``ensanduichada'' entre $1\cdot n^2$ e $2\cdot n^2$ a partir de $n_0=12$. A curva de baixo prova $f=\Omega(n^2)$, a de cima prova $f=O(n^2)$ e as duas juntas provam $f=\Theta(n^2)$.\label{fig:otheta}]
\begin{tikzpicture}
\begin{axis}[width=0.78\linewidth, height=5.6cm, domain=0:25, samples=60, xmin=0, xmax=25, ymin=0, ymax=1000,
  axis lines=left, xlabel={$n$}, ylabel={passos}, ylabel style={font=\small}, xlabel style={font=\small},
  tick label style={font=\scriptsize}, legend style={font=\scriptsize, draw=none, at={(0.02,0.98)}, anchor=north west},
  legend cell align=left]
\addplot[thick, MidnightBlue] {x^2+10*x+20}; \addlegendentry{$f(n)=n^2+10n+20$}
\addplot[thick, BrickRed, dashed] {2*x^2}; \addlegendentry{$c_2\,g(n)=2n^2$}
\addplot[thick, ForestGreen, dashed] {x^2}; \addlegendentry{$c_1\,g(n)=n^2$}
\draw[gray, densely dotted] (axis cs:11.66,0) -- (axis cs:11.66,1000);
\node[font=\scriptsize, anchor=south west, fill=white, inner sep=1pt] at (axis cs:11.8,30) {$n_0\approx 12$};
\end{axis}
\end{tikzpicture}
\end{figura}

\begin{exemplo}[Provando $O$, $\Omega$ e $\Theta$ pela definição]\label{ex:otheta}
Seja $f(n)=n^2+10n+20$.
\begin{itemize}
  \item $f=O(n^2)$: queremos $c$ e $n_0$ com $n^2+10n+20\le c\,n^2$. Com $c=2$, a desigualdade vira $10n+20\le n^2$, que vale para $n\ge12$ (para $n=12$: $140\le144$). Logo $c=2$, $n_0=12$ servem. Também serviria $c=31$, $n_0=1$, pois $10n\le10n^2$ e $20\le20n^2$ para $n\ge1$: as constantes não são únicas, basta exibir \emph{uma} escolha.
  \item $f=\Omega(n^2)$: com $c=1$ e $n_0=1$, $f(n)\ge n^2$ é óbvio.
  \item $f=\Theta(n^2)$: segue dos dois itens. Note que $f$ também é $O(n^3)$ (um limite superior folgado), mas \emph{não} é $O(n)$: se fosse, existiria $c$ com $n^2\le f(n)\le c\,n$ para $n$ grande, ou seja, $n\le c$, o que é falso para $n>c$.
\end{itemize}
\end{exemplo}

\begin{exemplo}
Pela mesma técnica, a função $4n^3+5n$ é $O(n^4)$ e $O(n^3)$, mas não é $O(n^2)$; é $\Omega(n^2)$ e $\Omega(n^3)$, mas não é $\Omega(n^4)$; portanto é $\Theta(n^3)$. Para provar $4n^3+5n=O(n^3)$ pela definição, basta tomar $c=9$ e $n_0=1$: $4n^3+5n\le4n^3+5n^3=9n^3$ para todo $n\ge1$.
\end{exemplo}

\subsection*{As versões estritas $o$ e $\omega$}
Em $f=O(g)$, a função $f$ pode crescer \emph{tanto quanto} $g$; em $f=o(g)$, ela cresce \emph{estritamente menos}: para \emph{qualquer} constante $c$, por menor que seja, $f$ acaba ficando abaixo de $c\,g$. A diferença está no quantificador: em $O$ basta \emph{existir} um $c$, em $o$ a desigualdade vale \emph{para todo} $c$.

\begin{exemplo}[Usando limites]\label{ex:limites}
Quando o limite de $f/g$ existe, ele resolve tudo de uma vez:
\begin{itemize}
  \item $\displaystyle\lim_{n\to\infty}\frac{3n^2+5n}{n^2}=3$, um número em $(0,\infty)$: logo $3n^2+5n=\Theta(n^2)$.
  \item $\displaystyle\lim\frac{n}{n^2}=\lim\frac1n=0$: logo $n=o(n^2)$. Já $n^2\ne o(n^2)$, pois $n^2/n^2=1\not\to0$.
  \item $\displaystyle\lim\frac{\ln n}{\sqrt n}=\lim\frac{1/n}{1/(2\sqrt n)}=\lim\frac{2}{\sqrt n}=0$ (regra de L'Hôpital): logo $\log n=o(\sqrt n)$. Logaritmos crescem mais devagar que qualquer raiz.
  \item $\displaystyle\lim\frac{2^n}{n^3}=\infty$ (aplique L'Hôpital três vezes): logo $2^n=\omega(n^3)$. Exponenciais crescem mais rápido que qualquer polinômio.
\end{itemize}
\end{exemplo}

\subsection*{Propriedades}
As notações se comportam como as desigualdades que representam:
\begin{itemize}
  \item \textbf{transitividade}: $f=O(g)$ e $g=O(h)$ implicam $f=O(h)$ (e o mesmo para as outras notações);
  \item \textbf{simetria transposta}: $f=O(g)$ se e somente se $g=\Omega(f)$;
  \item \textbf{a soma é dominada pela maior parcela}: $f+g=\Theta(\max\{f,g\})$;
  \item \textbf{a base do logaritmo não importa}: $\log_an=\log_bn/\log_ba$ difere de $\log_bn$ pelo fator constante $1/\log_ba$.
\end{itemize}

\begin{exemplo}[Aplicando as propriedades]\label{ex:propriedades}\leavevmode
\begin{itemize}
  \item $5n^3+100\,n^2\log n+7=\Theta(n^3)$: a soma é dominada pela maior parcela, e $n^2\log n=o(n^3)$ porque $\log n=o(n)$.
  \item $\log_{10}n=\log_2n/\log_2 10\approx0{,}301\log_2n$; logo $\log_{10}n=\Theta(\log_2n)$, e podemos escrever simplesmente $\Theta(\log n)$.
  \item $\log(n^3)=3\log n=\Theta(\log n)$, mas $\log(3^n)=n\log3=\Theta(n)$: o expoente \emph{dentro} do logaritmo vira constante; o expoente que depende de $n$, não.
\end{itemize}
\end{exemplo}

Também valem $\log(n!)=\Theta(n\log n)$, $(\log n)^c=o(n^\varepsilon)$ para quaisquer $c>0$ e $\varepsilon>0$, e $n^c=o(a^n)$ para todo $a>1$. A escala abaixo resume as ordens de crescimento mais frequentes:
\[
1\ <\ \log\log n\ <\ \log n\ <\ \sqrt n\ <\ n\ <\ n\log n\ <\ n^2\ <\ n^3\ <\ n^{\log n}\ <\ 2^n\ <\ n!\ <\ n^n .
\]

\begin{figura}[As funções mais comuns para $n$ entre 1 e 20. Para $n$ pequeno, todas parecem próximas; logo a exponencial $2^n$ sai do gráfico, e $n^2$ se afasta de $n\log n$.\label{fig:crescimento}]
\begin{tikzpicture}
\begin{axis}[width=0.8\linewidth, height=6cm, domain=1:20, samples=80, xmin=1, xmax=20, ymin=0, ymax=150,
  axis lines=left, xlabel={$n$}, ylabel={$f(n)$}, tick label style={font=\scriptsize},
  xlabel style={font=\small}, ylabel style={font=\small}, clip=true,
  every axis plot/.append style={thick}]
\addplot[gray] {log2(x)} node[pos=0.97, below, font=\scriptsize] {$\log n$};
\addplot[ForestGreen] {x} node[pos=0.97, above=1pt, font=\scriptsize] {$n$};
\addplot[MidnightBlue] {x*log2(x)} node[pos=0.93, above left, font=\scriptsize] {$n\log n$};
\addplot[Plum, domain=1:12.2] {x^2} node[pos=1, left, font=\scriptsize] {$n^2$};
\addplot[BrickRed, domain=1:7.2] {2^x} node[pos=1, right, font=\scriptsize] {$2^n$};
\end{axis}
\end{tikzpicture}
\end{figura}

A Tabela~\ref{tab:tempos} traduz essas ordens em tempo real, supondo um computador que executa um bilhão de passos por segundo. Ela mostra por que a ordem de crescimento importa mais que qualquer constante.
\begin{figura}
\captionof{table}{Tempo para executar $f(n)$ passos a $10^9$ passos por segundo.}\label{tab:tempos}
\small
\begin{tabular}{@{}rcccccc@{}}
\toprule
$n$ & $\log n$ & $n$ & $n\log n$ & $n^2$ & $n^3$ & $2^n$\\
\midrule
$10$ & 3\,ns & 10\,ns & 33\,ns & 100\,ns & 1\,$\mu$s & 1\,$\mu$s\\
$100$ & 7\,ns & 100\,ns & 664\,ns & 10\,$\mu$s & 1\,ms & $4\cdot10^{13}$ anos\\
$1000$ & 10\,ns & 1\,$\mu$s & 10\,$\mu$s & 1\,ms & 1\,s & ---\\
$10^6$ & 20\,ns & 1\,ms & 20\,ms & 17\,min & 32 anos & ---\\
\bottomrule
\end{tabular}
\end{figura}

\section{Analisando laços na prática}\label{sec:lacos}
Na maior parte das vezes, analisar um algoritmo iterativo é contar quantas vezes o corpo dos seus laços é executado.

\begin{roteiro}[Como analisar um trecho com laços]
\begin{enumerate}
  \item Identifique o \textbf{corpo mais interno} e verifique que ele custa $O(1)$.
  \item Para cada laço, escreva quantas vezes ele itera \emph{em função das variáveis dos laços externos}.
  \item Some: laços aninhados viram somatórios aninhados.
  \item Resolva o somatório e simplifique com as notações assintóticas.
\end{enumerate}
\end{roteiro}

\begin{exemplo}[Três padrões de laço]\label{ex:lacos}
\textbf{(a) Laços aninhados independentes.} Em ``\textbf{para} $i:=1$ \textbf{até} $n$, \textbf{para} $j:=1$ \textbf{até} $n$: corpo'', o corpo executa $n\cdot n=n^2$ vezes: $\Theta(n^2)$.

\smallskip
\textbf{(b) Laço interno que depende do externo.} Em ``\textbf{para} $i:=1$ \textbf{até} $n$, \textbf{para} $j:=1$ \textbf{até} $i$: corpo'', o laço interno faz $i$ iterações, e o total é
\[
\sum_{i=1}^n i=\frac{n(n+1)}2=\Theta(n^2).
\]
A Figura~\ref{fig:escada} mostra a contagem como uma escada: as execuções ocupam cerca de metade do quadrado $n\times n$, e metade de $n^2$ ainda é $\Theta(n^2)$.

\smallskip
\textbf{(c) Laço que divide pela metade.} Em ``$i:=n$; \textbf{enquanto} $i>1$ \textbf{faça} $i:=\floor{i/2}$'', o valor de $i$ passa por $n,\ n/2,\ n/4,\dots$ até chegar a $1$. Depois de $k$ iterações, $i\approx n/2^k$, que chega a $1$ quando $k\approx\log_2n$: $\Theta(\log n)$. Para $n=16$, $i$ vale $16,8,4,2,1$: quatro iterações.
\end{exemplo}

\begin{figura}[O laço do item (b) com $n=6$: na linha $i$, o corpo executa $i$ vezes (células preenchidas). O total, $1+2+\cdots+6=21$, é pouco mais que a metade das 36 células.\label{fig:escada}]
\begin{tikzpicture}[x=6mm,y=6mm]
\foreach \i in {1,...,6}{
  \foreach \j in {1,...,6}{
    \ifnum\j>\i \draw[gray!40] (\j-1,-\i) rectangle (\j,-\i+1);
    \else \draw[fill=azul] (\j-1,-\i) rectangle (\j,-\i+1); \fi}
  \node[rot,anchor=east] at (-0.2,-\i+0.5) {$i=\i$};
  \node[rot,anchor=west] at (6.2,-\i+0.5) {\i\ execuções};}
\node[rot] at (3,0.5) {$j=1,\dots,6$};
\end{tikzpicture}
\end{figura}

Um laço cuja variável é multiplicada ou dividida por uma constante a cada passo dá $\Theta(\log n)$ iterações. É isso que torna a \textbf{busca binária} tão rápida: a cada comparação, o trecho do vetor ordenado onde o elemento procurado pode estar cai pela metade.

\begin{figura}[Busca binária de $x=27$ num vetor ordenado de 16 posições. O trecho candidato (em destaque) cai pela metade a cada comparação com o elemento do meio ($\blacktriangledown$).\label{fig:buscabin}]
\begin{tikzpicture}[x=6.4mm,y=7.5mm]
\foreach \lin/\a/\b/\m in {0/1/16/8,1/9/16/12,2/9/11/10,3/11/11/11}{
  \fill[destaque] (\a-0.5,-\lin-0.32) rectangle (\b+0.5,-\lin+0.32);
  \foreach \x [count=\i] in {2,5,8,12,15,16,19,21,23,25,27,30,34,40,41,47}{
    \node[cel, minimum width=6.4mm, font=\scriptsize] at (\i,-\lin) {\x};}
  \node[font=\tiny] at (\m,-\lin+0.5) {$\blacktriangledown$};}
\node[rot,anchor=west] at (16.8,0) {$27>21$: direita};
\node[rot,anchor=west] at (16.8,-1) {$27<30$: esquerda};
\node[rot,anchor=west] at (16.8,-2) {$27>25$: direita};
\node[rot,anchor=west] at (16.8,-3) {achou!};
\end{tikzpicture}
\end{figura}

\begin{cuidado}[Erros comuns na análise de laços]
\begin{itemize}
  \item Multiplicar os limites dos laços quando o interno depende do externo: isso sempre dá um limite \emph{superior}, mas ele pode ser folgado. No item (b) a ordem saiu certa; no Exercício~\ref{ex:cx-lacos}(a), em que $i$ dobra a cada volta, faça a soma com cuidado antes de concluir.
  \item Esquecer que $\Theta$ descarta constantes: $n(n+1)/2$ e $n^2$ são da mesma ordem.
  \item Confundir o número de iterações com o \emph{valor} das variáveis: em (c), $i$ começa em $n$, mas o laço só dá $\log_2 n$ voltas.
\end{itemize}
\end{cuidado}

\section{Complexidade de problemas e de algoritmos}\label{sec:problemas}
A complexidade de um \emph{algoritmo} é uma função que descreve o seu custo. A complexidade de um \emph{problema} é uma pergunta mais difícil: qual o menor custo possível entre \emph{todos} os algoritmos que o resolvem, inclusive os que ainda não foram inventados?

\begin{defi}\leavevmode
\begin{itemize}
  \item Uma função $f(n)$ é um \textbf{limite superior} para o problema $\Pi$ se \emph{existe} algum algoritmo que resolve $\Pi$ com no máximo $f(n)$ passos no pior caso. A complexidade de pior caso de qualquer algoritmo para $\Pi$ é um limite superior para $\Pi$.
  \item Uma função $f(n)$ é um \textbf{limite inferior} para $\Pi$ se \emph{todo} algoritmo que resolve $\Pi$ executa pelo menos $f(n)$ passos no pior caso.
  \item Um algoritmo é \textbf{ótimo} se sua complexidade de pior caso coincide com um limite inferior do problema.
\end{itemize}
\end{defi}

\begin{exemplo}[Encontrar o máximo é um problema de custo exatamente $n-1$]\label{ex:maximo}
Considere o problema de achar o maior de $n$ números distintos usando comparações.
\begin{itemize}
  \item \textbf{Limite superior}: o algoritmo que guarda o maior visto até agora e o compara com cada novo elemento faz $n-1$ comparações. Logo $n-1$ é um limite superior para o problema.
  \item \textbf{Limite inferior}: numa comparação ``$a<b$?'', diremos que o menor dos dois \emph{perdeu}. Um algoritmo só pode ter certeza de que $x$ é o máximo se cada um dos outros $n-1$ elementos já perdeu alguma comparação: se $y\ne x$ nunca perdeu, $y$ poderia ser o máximo sem contradizer nenhuma resposta obtida. Como cada comparação produz exatamente um perdedor, são necessárias pelo menos $n-1$ comparações, \emph{qualquer que seja o algoritmo}.
\end{itemize}
Os dois limites coincidem, e o algoritmo simples é ótimo.
\begin{figura}[O máximo de $[5,9,2,7]$. Cada comparação (seta) elimina um candidato, o perdedor (em vermelho). Com $n=4$ elementos, três perdedores exigem três comparações.]
\begin{tikzpicture}[x=1.6cm]
\node[vtx] (a) at (0,0) {5}; \node[vtx] (b) at (1,0) {9}; \node[vtx] (c) at (2,0) {2}; \node[vtx] (d) at (3,0) {7};
\node[vtx, fill=ruim] (a1) at (0.5,-1.1) {5}; \node[vtx] (b1) at (1.5,-1.1) {9};
\node[vtx, fill=ruim] (c2) at (2,-2.2) {2}; \node[vtx] (b2) at (1,-2.2) {9};
\node[vtx, fill=ruim] (d3) at (3,-3.3) {7}; \node[vtx, fill=feito] (b3) at (1.5,-3.3) {9};
\node[rot, anchor=west] at (3.3,-1.1) {$5<9$: 5 perde};
\node[rot, anchor=west] at (3.3,-2.2) {$2<9$: 2 perde};
\node[rot, anchor=west] at (3.3,-3.3) {$7<9$: 7 perde};
\draw[seta,gray] (a) -- (a1); \draw[seta,gray] (b) -- (b1);
\draw[seta,gray] (b1) -- (b2); \draw[seta,gray] (c) -- (c2);
\draw[seta,gray] (b2) -- (b3); \draw[seta,gray] (d) -- (d3);
\end{tikzpicture}
\end{figura}
\end{exemplo}

Limites superiores se provam exibindo um algoritmo; limites inferiores exigem um argumento sobre todos os algoritmos possíveis, como o argumento do exemplo acima ou o da árvore de decisão, que mostra que ordenar por comparações exige $\Omega(n\log n)$ comparações (Capítulo~\ref{cap:sortcomp}). O mergesort e o heapsort, que fazem $O(n\log n)$ comparações, são portanto ótimos.

As notações assintóticas são afirmações sobre funções, e não sobre ``limites'' no sentido acima: nada impede dizer que o pior caso de um algoritmo é $\Omega(n^3)$, ou que conhecemos para um problema um limite inferior que é $O(\log n)$.

\section{Algoritmos polinomiais e exponenciais}\label{sec:polinomiais}
Um algoritmo é \textbf{polinomial} se seu custo de pior caso é $O(n^c)$ para alguma constante $c$; a classe dos problemas que admitem algoritmos polinomiais é chamada de $P$. Custos como $2^n$, $k^n$, $n!$, $2^{n^2}$ e $n^{\log n}$ são \textbf{exponenciais} (ou superpolinomiais).

O consenso é que algoritmos polinomiais são eficientes e exponenciais não, por um motivo que vai além do tempo absoluto da Tabela~\ref{tab:tempos}: algoritmos polinomiais aproveitam o avanço tecnológico. Com uma máquina mais rápida, o maior tamanho de entrada resolvido num tempo fixo é \emph{multiplicado} por uma constante quando o custo é polinomial, mas apenas \emph{acrescido} de uma constante quando é exponencial.

\begin{exemplo}[O que uma máquina 100 vezes mais rápida compra]\label{ex:maquina}
Suponha que hoje, em uma hora, resolvemos entradas de tamanho até $N$. Com a máquina nova, no mesmo tempo, fazemos $100$ vezes mais passos:
\begin{center}\small
\begin{tabular}{@{}lll@{}}
\toprule
custo & novo tamanho máximo & por quê\\
\midrule
$n$ & $100N$ & $n'=100N$ dá $100N$ passos\\
$n^2$ & $\sqrt{100}\,N=10\,N$ & $(10N)^2=100N^2$\\
$n^3$ & $\sqrt[3]{100}\,N\approx4{,}64\,N$ & $(4{,}64N)^3\approx100N^3$\\
$2^n$ & $N+6{,}6$ & $2^{N+6{,}6}\approx100\cdot2^N$\\
\bottomrule
\end{tabular}
\end{center}
Se hoje resolvemos $N=1000$ com o algoritmo quadrático, passaremos a resolver $10\,000$. Se resolvemos $N=40$ com o exponencial, passaremos a resolver só $46$: cem vezes mais velocidade compram seis posições a mais.
\end{exemplo}

\begin{chave}
Pior caso é a medida padrão. $O$ é ``$\le$'', $\Omega$ é ``$\ge$'', $\Theta$ é ``$=$'' a menos de constantes, para $n$ grande; a base do logaritmo não importa. Laços aninhados viram somatórios; laços que dividem a variável por uma constante dão $\Theta(\log n)$ voltas. Limite superior de um problema: \emph{existe} um algoritmo; limite inferior: \emph{todo} algoritmo; ótimo: pior caso igual ao limite inferior. O tamanho da entrada se mede em bits, e algoritmos cujo custo cresce com o \emph{valor} de um número são pseudopolinomiais.
\end{chave}

% ----------------------------------------------------------------
\section{Exercícios}
\excurso
\begin{exercicio}{cx-slide}[Slides]<ex:otheta>
Diga, justificando pela definição, se $2n^2+1$ é: (a) $O(n)$; (b) $O(n^2)$; (c) $O(n^3)$; (d) $\Omega(n)$; (e) $\Omega(n^3)$; (f) $\Theta(n^2)$.
\end{exercicio}
\begin{exercicio}{cx-tec}[Slides]<ex:maquina>
Uma máquina nova é 10 vezes mais rápida que a antiga. Para algoritmos de custo $n$, $n^2$ e $2^n$, em quanto aumenta o maior tamanho de entrada que se resolve no mesmo tempo? O que isso diz sobre algoritmos polinomiais e exponenciais?
\end{exercicio}
\begin{exercicio}{cx-k}[Slides]<sec:medidas>
O \emph{counting sort} ordena $n$ inteiros em $[0,k]$ em tempo $\Theta(n+k)$. Explique por que ele \textbf{não} é polinomial no tamanho da entrada quando $k=2^n$, e por que é linear quando $k=O(n)$.
\end{exercicio}

\excompl
\begin{exercicio}{cx-ordem}[Complementar]<sec:notacoes>
Ordene as funções pela taxa de crescimento: $\sqrt n$,\ $\log^2 n$,\ $2^{\log_2 n}$,\ $n\log n$,\ $n^{\log_2 3}$,\ $1{,}01^n$,\ $n^{100}$,\ $n!$,\ $2^{2^n}$.
\end{exercicio}
\begin{exercicio}{cx-vf}[Complementar]<ex:propriedades>
Verdadeiro ou falso? Justifique.
(a) $2^{n+1}=O(2^n)$;\quad (b) $2^{2n}=O(2^n)$;\quad (c) $\log(n^2)=\Theta(\log n)$;\quad (d) $(\log n)^{100}=O(n^{0{,}01})$;\quad (e) $n^{\log n}=O(2^n)$;\quad (f) se $f=O(g)$ então $2^{f}=O(2^{g})$.
\end{exercicio}
\begin{exercicio}{cx-fat}[Complementar]<sec:notacoes>
Mostre que $\log(n!)=\Theta(n\log n)$ sem usar a fórmula de Stirling.
\end{exercicio}
\begin{exercicio}{cx-incomp}[Complementar]<sec:notacoes>
Dê funções positivas $f$ e $g$ tais que nem $f=O(g)$ nem $g=O(f)$.
\end{exercicio}
\begin{exercicio}{cx-lacos}[Complementar]<ex:lacos>
Dê a complexidade, em $\Theta$, de cada trecho (o corpo mais interno custa $O(1)$):
\begin{center}\small
\begin{minipage}{0.45\linewidth}
(a)\ $i:=1$\\
\textbf{enquanto} $i\le n$ \textbf{faça}\\
\hspace*{1em}\textbf{para} $j:=1$ \textbf{até} $i$ \textbf{faça} (corpo)\\
\hspace*{1em}$i:=2i$
\end{minipage}
\begin{minipage}{0.45\linewidth}
(b)\ \textbf{para} $i:=1$ \textbf{até} $n$ \textbf{faça}\\
\hspace*{1em}$j:=1$\\
\hspace*{1em}\textbf{enquanto} $j\le n$ \textbf{faça}\\
\hspace*{2em}(corpo);\ $j:=2j$
\end{minipage}
\end{center}
\end{exercicio}
\begin{exercicio}{cx-primo}[Complementar]<ex:somalaco>
O algoritmo que testa se $N$ é primo tentando dividir $N$ por $2,3,\dots,\floor{\sqrt N}$ faz $O(\sqrt N)$ divisões. Ele é polinomial? (Pense no tamanho da entrada.)
\end{exercicio}
\begin{exercicio}{cx-def3}[Complementar]<ex:otheta>
Prove pela definição, exibindo as constantes, que $3n^3-100n^2+7=\Theta(n^3)$.
\end{exercicio}
\begin{exercicio}{cx-nlogn}[Complementar]<ex:otheta>
Prove que $n^2\ne O(n\log n)$.
\end{exercicio}
\begin{exercicio}{cx-base}[Complementar]<ex:propriedades>
Prove que, para quaisquer constantes $a,b>1$, $\log_an=\Theta(\log_bn)$. Por que, por outro lado, $a^n$ e $b^n$ não são da mesma ordem quando $a\ne b$?
\end{exercicio}
\begin{exercicio}{cx-prop}[Complementar]<sec:notacoes>
Prove: (a) se $f=O(g)$ e $g=O(h)$, então $f=O(h)$; (b) $f+g=\Theta(\max\{f,g\})$; (c) $f=\Theta(g)$ se e somente se $g=\Theta(f)$.
\end{exercicio}
\begin{exercicio}{cx-somak}[Complementar]<ex:lacos>
Mostre que $\sum_{i=1}^ni^k=\Theta(n^{k+1})$ para toda constante inteira $k\ge1$.
\end{exercicio}
\begin{exercicio}{cx-harm}[Complementar]<sec:notacoes>
Mostre que $H_n=\sum_{i=1}^n\frac1i=\Theta(\log n)$, comparando a soma com a integral de $1/x$.
\end{exercicio}
\begin{exercicio}{cx-busca}[Complementar]<sec:medidas>
Na busca sequencial de $x$ num vetor de $n$ elementos distintos, quantas comparações são feitas no pior caso, no melhor caso e no caso médio, supondo que $x$ está no vetor e é igualmente provável em cada posição?
\end{exercicio}
\begin{exercicio}{cx-triplo}[Complementar]<ex:lacos>
Quantas vezes o corpo é executado no trecho ``\textbf{para} $i:=1$ \textbf{até} $n$, \textbf{para} $j:=i$ \textbf{até} $n$, \textbf{para} $k:=i$ \textbf{até} $j$: corpo''? Dê a fórmula exata e a ordem em $\Theta$.
\end{exercicio}
\begin{exercicio}{cx-loglog}[Complementar]<ex:lacos>
Qual a complexidade do trecho ``$j:=2$; \textbf{enquanto} $j\le n$ \textbf{faça} $j:=j^2$''?
\end{exercicio}
\begin{exercicio}{cx-fatexp}[Complementar]<ex:limites>
Mostre que $2^n=o(n!)$ e que $n!=o(n^n)$.
\end{exercicio}
\begin{exercicio}{cx-binaria}[Complementar]<fig:buscabin>
Quantas comparações a busca binária faz no pior caso num vetor ordenado com 1000 elementos? E com $10^6$? Dê a fórmula geral.
\end{exercicio}
\begin{exercicio}{cx-inf}[Complementar]<ex:maximo>
Mostre que todo algoritmo que decide se $x$ está num vetor \emph{não ordenado} de $n$ elementos, usando apenas comparações de $x$ com elementos do vetor, faz pelo menos $n$ comparações no pior caso. Conclua que a busca sequencial é ótima nesse modelo.
\end{exercicio}
\begin{exercicio}{cx-max}[Complementar]<ex:propriedades>
Verdadeiro ou falso? (a) $\max\{f,g\}=\Theta(f+g)$; (b) $f=O(g)$ implica $\log f=O(\log g)$, para $f,g\ge2$; (c) $f=O(g)$ implica $f^2=O(g^2)$; (d) $f(n)=O(f(n/2))$ para toda $f$ crescente.
\end{exercicio}
\begin{exercicio}{cx-insercao4}[Complementar]<ex:inscasos>
Conte as comparações que a ordenação por inserção do Exemplo~\ref{ex:inscasos} faz nas entradas $[2,4,1,3]$ e $[4,3,1,2]$, inserção por inserção. Por que a segunda faz tantas comparações quanto o vetor decrescente $[4,3,2,1]$, embora não seja decrescente?
\end{exercicio}
\begin{exercicio}{cx-tercos}[Complementar]<ex:lacos>
Quantas vezes o laço ``$i:=n$; \textbf{enquanto} $i>1$ \textbf{faça} $i:=\floor{i/3}$'' é executado para $n=100$ e para $n=1000$? Dê uma fórmula para $n\ge2$ e a ordem em $\Theta$.
\end{exercicio}

% ================================================================
\chapter{Listas de prioridades e heaps binários}\label{cap:heap}\label{sec:heap}
% ================================================================
\begin{objetivos}
\begin{itemize}
  \item O tipo abstrato lista de prioridades e o custo de implementá-lo com listas.
  \item O heap binário: definição, representação em vetor e as operações \emph{subir} e \emph{descer}.
  \item Inserção, remoção e alteração de prioridade em $O(\log n)$, e construção em $O(n)$.
  \item O heapsort e as técnicas para demonstrar propriedades de heaps.
\end{itemize}
\end{objetivos}

Muitas aplicações precisam, repetidamente, do elemento \emph{mais importante} de uma coleção que muda o tempo todo: o próximo processo a ser executado, o paciente mais grave de uma emergência, o próximo evento de uma simulação, o vértice mais próximo no algoritmo de Dijkstra. O tipo abstrato que atende a essas aplicações é a \textbf{lista de prioridades}, e sua implementação clássica é o \textbf{heap binário}: um vetor com uma ordem parcial que permite todas as operações em tempo logarítmico e a construção em tempo linear.

\section{Listas de prioridades}\label{sec:listaprior}
Uma lista de prioridades é uma tabela em que cada elemento tem uma \textbf{prioridade}, e que oferece as operações de
\begin{itemize}
  \item \emph{seleção} do elemento de maior prioridade;
  \item \emph{inserção} de um novo elemento;
  \item \emph{remoção} do elemento de maior prioridade;
  \item \emph{alteração} da prioridade de um elemento;
  \item \emph{construção} da tabela a partir de um conjunto inicial.
\end{itemize}
A operação de \emph{fusão} de duas listas será tratada no Capítulo~\ref{cap:heapav}.

\begin{exemplo}[A fila de um pronto-socorro]\label{ex:prontosocorro}
Num pronto-socorro, cada paciente recebe uma nota de gravidade de 0 a 100 na triagem. Uma sequência típica de eventos é:
\begin{center}\small
\begin{tabular}{@{}ll@{}}
\toprule
evento & operação\\
\midrule
chegam Ana (40), Bruno (75) e Carla (20) & construção com $\{40,75,20\}$\\
o médico fica livre & seleção e remoção: atende Bruno (75)\\
chega Davi (90) & inserção de 90\\
o estado de Carla piora para 95 & alteração: $20\to95$\\
o médico fica livre & seleção e remoção: atende Carla (95)\\
\bottomrule
\end{tabular}
\end{center}
Note que nunca precisamos da ordem completa dos pacientes, apenas de quem é o \emph{primeiro}. É essa observação que o heap explora.
\end{exemplo}

As implementações mais simples são listas. Numa lista não ordenada, inserir é imediato, mas achar o máximo exige percorrer tudo; numa lista ordenada, o máximo está no início, mas cada inserção precisa encontrar a posição certa e deslocar elementos. Por exemplo, inserir 50 na lista ordenada $[90,70,40,20]$ exige deslocar o 40 e o 20 uma posição para a direita, e no pior caso (inserir um novo máximo) todos os elementos se deslocam. A tabela compara essas opções com o heap.
\begin{center}
\small
\begin{tabular}{@{}lccc@{}}
\toprule
Operação & Lista não ordenada & Lista ordenada & Heap\\
\midrule
seleção & $\Theta(n)$ & $\Theta(1)$ & $O(1)$\\
inserção & $\Theta(1)$ & $\Theta(n)$ & $O(\log n)$\\
remoção & $\Theta(n)$ & $\Theta(1)$ & $O(\log n)$\\
alteração & $\Theta(n)$ & $\Theta(n)$ & $O(\log n)$\\
construção & $\Theta(n)$ & $\Theta(n\log n)$ & $O(n)$\\
\bottomrule
\end{tabular}
\end{center}
Cada lista tem alguma operação de atualização que custa $\Omega(n)$. O heap abre mão de manter a ordem total e mantém só o necessário para que o máximo esteja sempre num lugar conhecido; com isso, todas as operações ficam logarítmicas.

\section{Definição e representação}\label{sec:heapdef}
\begin{intuicao}
Pense num organograma em que cada chefe é mais importante que seus subordinados diretos. Quem está no topo é o mais importante de todos, mas nada se exige entre colegas do mesmo nível, nem entre pessoas de departamentos diferentes. Um heap é exatamente isso: uma hierarquia em que \emph{cada pai domina os seus filhos}, e nada mais.
\end{intuicao}

\begin{defi}[Heap]
Um \textbf{heap} (de máximo) é uma sequência $s_1,\dots,s_n$ tal que
\[
s_i\le s_{\floor{i/2}}\qquad\text{para todo }1<i\le n .
\]
Num heap \emph{de mínimo}, a desigualdade é invertida: $s_i\ge s_{\floor{i/2}}$.
\end{defi}
A definição fica clara quando se desenha a sequência como uma \textbf{árvore binária completa}, numerada por níveis e, em cada nível, da esquerda para a direita: o nó $i$ tem pai $\floor{i/2}$, filho esquerdo $2i$ e filho direito $2i+1$ (quando esses índices são $\le n$). A condição diz que \emph{cada nó tem prioridade maior ou igual à de seus filhos}. Não há ponteiros: a árvore é só uma maneira de olhar para o vetor (Figura~\ref{fig:heaprep}).

\begin{figura}[O mesmo heap visto como árvore e como vetor. Os números pequenos são os índices; os filhos do nó $i$ estão nas posições $2i$ e $2i+1$, como indicam as setas desenhadas junto ao vetor.\label{fig:heaprep}]
\begin{forest} bst
[90,label={[font=\tiny]above:1}, s sep=16mm
  [55,label={[font=\tiny]left:2} [41,label={[font=\tiny]left:4} [37,label={[font=\tiny]left:8}] [,vazio]] [30,label={[font=\tiny]right:5}]]
  [82,label={[font=\tiny]right:3} [64,label={[font=\tiny]left:6}] [75,label={[font=\tiny]right:7}]]]
\end{forest}
\qquad
\begin{tikzpicture}[baseline=-1.2cm,x=7mm]
\foreach \x [count=\i] in {90,55,82,41,30,64,75,37} {\node[cel] (c\i) at (\i,0) {\x}; \node[idx] at (\i,0.5) {\i};}
\draw[seta, MidnightBlue!70, bend right=50] (c1.south) to (c2.south);
\draw[seta, MidnightBlue!70, bend right=50] (c1.south) to (c3.south);
\draw[seta, BrickRed!70, bend right=45] (c2.south) to (c4.south);
\draw[seta, BrickRed!70, bend right=45] (c2.south) to (c5.south);
\draw[seta, ForestGreen!70, bend left=40] ($(c3.north)+(0,3.2mm)$) to ($(c6.north)+(0,3.2mm)$);
\draw[seta, ForestGreen!70, bend left=40] ($(c3.north)+(0,3.2mm)$) to ($(c7.north)+(0,3.2mm)$);
\end{tikzpicture}
\end{figura}

\begin{exemplo}[Reconhecendo um heap]\label{ex:ehheap}
Para verificar se uma sequência é um heap, basta comparar cada posição $i\ge2$ com o seu pai $\floor{i/2}$.
\begin{itemize}
  \item $[9,7,8,3,5,6]$: os pares (filho, pai) são $(7,9)$, $(8,9)$, $(3,7)$, $(5,7)$ e $(6,8)$; em todos o filho é menor. É um heap.
  \item $[9,5,8,6,4,7]$: a posição 4 guarda $6$ e o seu pai, a posição 2, guarda $5$. Como $6>5$, \textbf{não} é um heap (Figura~\ref{fig:naoheap}). Todos os outros pares estão corretos: basta uma violação.
\end{itemize}
Observe que no heap $[9,7,8,3,5,6]$ os irmãos 7 e 8 aparecem em ordem crescente, e o 6, à direita, é maior que o 3 e o 5, à esquerda: nada disso importa. Só se comparam pais e filhos.
\begin{figura}[À esquerda, um heap. À direita, uma sequência que não é heap: a aresta em vermelho liga o pai 5 a um filho maior, 6.\label{fig:naoheap}]
\begin{forest} bstp
[9 [7 [3] [5]] [8 [6] [,vazio]]]
\end{forest}
\hspace{2cm}
\begin{forest} bstp
[9 [5,desreg [6,aviol,desreg] [4]] [8 [7] [,vazio]]]
\end{forest}
\end{figura}
\end{exemplo}

Alguns fatos sobre os índices aparecem em todas as demonstrações deste capítulo:
\begin{itemize}
  \item o nó $i$ está no nível $\floor{\log_2i}+1$ (a raiz está no nível 1), e a árvore tem $\floor{\log_2n}+1$ níveis;
  \item o nó $i$ tem filho se e somente se $2i\le n$; logo as \textbf{folhas} são exatamente as posições $\floor{n/2}+1,\dots,n$, e há $\ceil{n/2}$ delas;
  \item neste capítulo, a altura de um nó é medida em níveis: \textbf{folhas têm altura 1}.
\end{itemize}

\begin{exemplo}[Fazendo contas com índices]\label{ex:indices}
Num heap com $n=10$ elementos: o pai do nó $7$ é $\floor{7/2}=3$; os filhos do nó $4$ são $8$ e $9$; o nó $5$ tem só o filho esquerdo $10$, pois $11>n$; as folhas são as posições $6,\dots,10$ (são $\ceil{10/2}=5$ folhas); o nó $9$ está no nível $\floor{\log_29}+1=4$, e a árvore tem $\floor{\log_210}+1=4$ níveis.
\end{exemplo}

\section{As operações subir e descer}\label{sec:subirdescer}
Toda operação sobre um heap altera uma posição e depois restaura a propriedade de heap, que pode ter sido violada em apenas um ponto:
\begin{itemize}
  \item se a prioridade de um nó \emph{aumentou}, ele pode ter ficado maior que o pai, e precisa \textbf{subir};
  \item se \emph{diminuiu}, pode ter ficado menor que um filho, e precisa \textbf{descer}.
\end{itemize}

\begin{pseudo}{subir$(i)$ -- após \emph{aumentar} a prioridade de $i$}
\begin{algorithmic}
\Procedure{subir}{$i$}
  \State $j := \floor{i/2}$
  \If{$j\ge1$}
    \If{$T[i].chave > T[j].chave$}
      \State $T[i]\Leftrightarrow T[j]$
      \State \Call{subir}{$j$}
    \EndIf
  \EndIf
\EndProcedure
\end{algorithmic}
\end{pseudo}
\begin{pseudo}{descer$(i,n)$ -- após \emph{diminuir} a prioridade de $i$}
\begin{algorithmic}
\Procedure{descer}{$i,n$}
  \State $j := 2i$
  \If{$j\le n$}
    \If{$j<n$ \textbf{e} $T[j+1].chave > T[j].chave$}
      \State $j := j+1$ \Comment{maior filho}
    \EndIf
    \If{$T[i].chave < T[j].chave$}
      \State $T[i]\Leftrightarrow T[j]$
      \State \Call{descer}{$j,n$}
    \EndIf
  \EndIf
\EndProcedure
\end{algorithmic}
\end{pseudo}
Ambos percorrem um único caminho da árvore, logo custam $O(\log n)$. Um detalhe essencial do \emph{descer} é que ele troca o nó com o \textbf{maior} dos filhos: se trocasse com o menor, o maior filho ficaria abaixo de um valor menor que ele, e a propriedade continuaria violada.

\begin{exemplo}[Por que trocar com o maior filho]
Suponha que a raiz de $[4,9,7]$ precise descer. Trocando com o maior filho, 9, obtemos $[9,4,7]$, um heap. Trocando com o menor, 7, obteríamos $[7,9,4]$, e agora o filho 9 é maior que o pai 7: a violação não foi corrigida, só mudou de lugar.
\end{exemplo}

\section{Inserção, remoção e alteração de prioridade}\label{sec:heapops}
Com \emph{subir} e \emph{descer}, as demais operações são imediatas:
\begin{itemize}
  \item para \textbf{inserir}, o novo elemento ocupa a primeira posição livre do vetor, a $n+1$, e sobe;
  \item para \textbf{remover o máximo}, o último elemento é copiado para a raiz e desce;
  \item para \textbf{alterar} uma prioridade, o elemento sobe ou desce conforme a prioridade tenha aumentado ou diminuído.
\end{itemize}
Em todos os casos o \emph{formato} da árvore continua o de uma árvore completa (só a última posição é criada ou apagada), e a ordem é consertada ao longo de um único caminho.
\begin{pseudo}{Inserção, remoção do máximo e alteração de prioridade}
\begin{algorithmic}
\Procedure{inserir}{$novo$} \Comment{$M$ é a capacidade do vetor}
  \If{$n<M$} \State $T[n+1]:=novo$;\ $n:=n+1$;\ \Call{subir}{$n$}
  \Else\ \textbf{overflow}
  \EndIf
\EndProcedure
\Procedure{remover}{}
  \If{$n\ne0$} \State processar$(T[1])$;\ $T[1]:=T[n]$;\ $n:=n-1$;\ \Call{descer}{$1,n$}
  \Else\ \textbf{underflow}
  \EndIf
\EndProcedure
\Procedure{alterar}{$i, p$}
  \State $antiga := T[i].chave$;\ $T[i].chave := p$
  \If{$p > antiga$} \State \Call{subir}{$i$} \Comment{aumentou: sobe}
  \Else \State \Call{descer}{$i,n$} \Comment{diminuiu: desce}
  \EndIf
\EndProcedure
\end{algorithmic}
\end{pseudo}

\begin{exemplo}[Inserção passo a passo]\label{ex:heapins}
Inserimos 68 no heap $[90,55,82,41,30,64,75,37]$ da Figura~\ref{fig:heaprep}. Ele entra na posição 9, como filho de 41. Como $68>41$, troca com o pai; na posição 4, $68>55$, e troca de novo; na posição 2, $68<90$, e para. O resultado é $[90,68,82,55,30,64,75,37,41]$. Foram duas trocas, ao longo do caminho de uma folha até a raiz.
\begin{figura}[Inserção de 68: o novo elemento (em destaque) entra na primeira posição livre e sobe enquanto for maior que o pai. A aresta azul é a comparação seguinte.\label{fig:heapins}]
\begin{tabular}{@{}c@{\hspace{6mm}}c@{\hspace{6mm}}c@{}}
\begin{forest} bstp
[90 [55 [41 [37] [68,novo,acam]] [30]] [82 [64] [75]]]
\end{forest} &
\begin{forest} bstp
[90 [55 [68,novo,acam [37] [41]] [30]] [82 [64] [75]]]
\end{forest} &
\begin{forest} bstp
[90 [68,novo,acam [55 [37] [41]] [30]] [82 [64] [75]]]
\end{forest}\\
\small (a) entra na posição 9 & \small (b) troca com 41 & \small (c) troca com 55; $68<90$: fim
\end{tabular}
\end{figura}
\end{exemplo}

\begin{exemplo}[Remoção passo a passo]\label{ex:heaprem}
Removemos o máximo de $H=[90,55,82,\allowbreak41,30,\allowbreak64,75,37]$. O último elemento, 37, vai para a raiz, formando $[37,55,82,41,30,64,75]$. Agora 37 desce: os filhos são 55 e 82, e ele troca com o maior, 82; na posição 3, os filhos são 64 e 75, e ele troca com 75; na posição 7, é folha, e para. Resultado: $[82,55,75,41,30,64,37]$.
\begin{figura}[Remoção do máximo: o último elemento ocupa a raiz e desce, sempre trocando com o maior filho.\label{fig:heaprem}]
\begin{tabular}{@{}c@{\hspace{8mm}}c@{\hspace{8mm}}c@{}}
\begin{forest} bstp
[37,novo [55 [41] [30]] [82,acam [64] [75]]]
\end{forest} &
\begin{forest} bstp
[82 [55 [41] [30]] [37,novo [64] [75,acam]]]
\end{forest} &
\begin{forest} bstp
[82 [55 [41] [30]] [75 [64] [37,novo]]]
\end{forest}\\
\small (a) 37 vai para a raiz & \small (b) troca com 82 & \small (c) troca com 75: fim
\end{tabular}
\end{figura}
\end{exemplo}

\begin{exemplo}[Alteração de prioridade]\label{ex:heapalt}
No mesmo heap $H=[90,55,82,41,30,64,75,37]$:
\begin{itemize}
  \item \textbf{aumentar} a prioridade do nó 6 de 64 para 95: ele sobe trocando com 82 (posição 3) e com 90 (posição 1), e o heap fica $[95,55,90,41,$\allowbreak$30,82,75,37]$;
  \item \textbf{diminuir} a raiz de 90 para 33: ela desce trocando com 82 e com 75: $[82,55,75,41,$\allowbreak$30,64,33,37]$.
\end{itemize}
\end{exemplo}

\begin{cuidado}[Subir ou descer?]
É comum inverter as duas operações ao alterar uma prioridade, escrevendo algo como ``\textbf{se} $H[i]>H[\floor{i/2}]$ \textbf{então} descer$(i)$ \textbf{senão} subir$(i)$''. O correto é o contrário: se o novo valor ficou maior que o pai, a violação é \emph{com o pai}, e o elemento precisa \textbf{subir}; se não ficou, a única violação possível é com os filhos, e ele precisa \textbf{descer}. Em pseudocódigo: \textbf{se} $i>1$ \textbf{e} $H[i]>H[\floor{i/2}]$ \textbf{então} subir$(i)$ \textbf{senão} descer$(i,n)$.
\end{cuidado}

\section{Construção em tempo linear}\label{sec:arranjar}
Para transformar um vetor arbitrário num heap, a primeira ideia é inseri-lo elemento a elemento: \textbf{para} $i:=2,\dots,n$ \textbf{faça} subir$(i)$. Isso custa $O(n\log n)$, e de fato $\Theta(n\log n)$ no pior caso, por exemplo quando a entrada é crescente e cada elemento sobe até a raiz.

Uma ideia melhor é trabalhar de baixo para cima. As folhas, isto é, as posições $i>\floor{n/2}$, já são heaps de um único elemento. Percorrendo os demais nós do último para o primeiro, cada um encontra as duas subárvores dos filhos já organizadas como heaps, e basta fazê-lo descer.
\begin{pseudo}{arranjar$(n)$: construção do heap}
\begin{algorithmic}
\Procedure{arranjar}{$n$}
  \For{$i := \floor{n/2}$ \textbf{decrescendo até} $1$} \State \Call{descer}{$i,n$} \EndFor
\EndProcedure
\end{algorithmic}
\end{pseudo}
A correção segue do invariante ``antes de processar $i$, as subárvores com raízes $i+1,\dots,n$ são heaps''.

\begin{exemplo}[O \emph{arranjar} passo a passo]\label{ex:arranjar}
Construímos um heap a partir de $[5,13,8,2,11,9,14]$, com $n=7$. As folhas são as posições $4,\dots,7$; processamos $i=3,2,1$ (Figura~\ref{fig:arranjar}).
\begin{itemize}
  \item $i=3$ (valor 8): os filhos são 9 e 14; troca com 14. Vetor: $[5,13,14,2,11,9,8]$.
  \item $i=2$ (valor 13): os filhos são 2 e 11, ambos menores; nada muda.
  \item $i=1$ (valor 5): os filhos são 13 e 14; troca com 14, e na posição 3 os filhos são 9 e 8; troca com 9. Vetor final: $[14,13,9,2,11,5,8]$.
\end{itemize}
\begin{figura}[O \emph{arranjar} sobre $[5,13,8,2,11,9,14]$. Em cada passo, o nó processado (em destaque) desce enquanto for menor que o maior filho; as subárvores abaixo dele já são heaps.\label{fig:arranjar}]
\begin{tabular}{@{}c@{\hspace{7mm}}c@{\hspace{7mm}}c@{\hspace{7mm}}c@{}}
\begin{forest} bstp
[5 [13 [2] [11]] [8,novo [9] [14]]]
\end{forest} &
\begin{forest} bstp
[5 [13,novo [2] [11]] [14 [9] [8]]]
\end{forest} &
\begin{forest} bstp
[5,novo [13 [2] [11]] [14 [9] [8]]]
\end{forest} &
\begin{forest} bstp
[14,feito [13,feito [2,feito] [11,feito]] [9,feito [5,feito] [8,feito]]]
\end{forest}\\
\small (a) $i=3$ & \small (b) $i=2$: nada muda & \small (c) $i=1$ & \small (d) heap pronto
\end{tabular}
\end{figura}
\end{exemplo}

O ponto surpreendente é o custo: a maioria dos nós está perto das folhas e desce pouco. A Figura~\ref{fig:alturas} mostra a conta numa árvore cheia com 15 nós.
\begin{figura}[Numa árvore completa, metade dos nós são folhas e não descem nada; um quarto desce no máximo um nível; e assim por diante. Com 15 nós, o \emph{arranjar} faz no máximo $4\cdot1+2\cdot2+1\cdot3=11$ trocas, menos que $n$.\label{fig:alturas}]
\begin{tikzpicture}[x=0.9cm,y=0.85cm]
\foreach \d/\cor in {0/BrickRed!35,1/Goldenrod!50,2/MidnightBlue!25,3/ForestGreen!20} {
  \pgfmathtruncatemacro{\m}{2^\d}
  \foreach \k in {1,...,\m} {\node[circle,draw,fill=\cor,minimum size=3.5mm,inner sep=0pt] (n\d-\k) at ({(\k-0.5)*8/\m},-\d) {};}}
\foreach \d in {1,2,3} {
  \pgfmathtruncatemacro{\p}{\d-1}\pgfmathtruncatemacro{\m}{2^\d}
  \foreach \k in {1,...,\m} {\pgfmathtruncatemacro{\q}{(\k+1)/2}\draw (n\p-\q) -- (n\d-\k);}}
\foreach \d/\txt in {0/{altura 4:\ \ 1 nó $\times$ até 3 trocas},1/{altura 3:\ \ 2 nós $\times$ até 2 trocas},2/{altura 2:\ \ 4 nós $\times$ até 1 troca},3/{altura 1:\ \ 8 folhas $\times$ 0 trocas}}
  {\node[anchor=west,font=\small] at (8.6,-\d) {\txt};}
\end{tikzpicture}
\end{figura}

\begin{lema}\label{lema:alturas}
Numa árvore binária completa com $n$ nós (folhas com altura 1), no máximo $\ceil{n/2^j}$ nós têm altura $j$.
\end{lema}
\begin{proof}
Numa árvore completa, o caminho mais longo que desce de $i$ é o da esquerda: $i\to2i\to4i\to\cdots$. Logo o nó $i$ tem altura $\ge j$ se e somente se o seu descendente $j-1$ níveis abaixo existe, isto é, $i\,2^{j-1}\le n$, ou $i\le\floor{n/2^{j-1}}$. Assim, com $x=n/2^{j-1}$ e $m=\floor{x}$,
\[
\#\{\text{altura}=j\}=\floor{x}-\floor{x/2}=m-\floor{m/2}=\ceil{m/2}\le\ceil{x/2}=\ceil{n/2^j}.\qedhere
\]
\end{proof}

\begin{teo}
O procedimento arranjar$(n)$ roda em tempo $O(n)$.
\end{teo}
\begin{proof}
Num nó de altura $j$, descer faz no máximo $j-1$ trocas. Pelo Lema~\ref{lema:alturas}, o total de trocas é
\[
\sum_{j=2}^{h}(j-1)\ceil*{\frac{n}{2^j}}\ \le\ \sum_{j\ge2}(j-1)\Bigl(\frac{n}{2^j}+1\Bigr)
\ =\ n\sum_{j\ge2}\frac{j-1}{2^j}+O(\log^2 n)\ =\ n+O(\log^2n)=O(n),
\]
pois $\sum_{j\ge1}\frac{j-1}{2^j}=\sum_{j\ge1}\frac{j}{2^j}-\sum_{j\ge1}\frac1{2^j}=2-1=1$.
\end{proof}
A intuição é a da Figura~\ref{fig:alturas}: metade dos nós (as folhas) não custa nada, um quarto custa no máximo uma troca, um oitavo no máximo duas, e assim por diante; só a raiz pode descer $\log n$ níveis. Na construção por inserções sucessivas acontece o contrário: os muitos nós de baixo são justamente os que podem subir mais.

\begin{exemplo}[Mais duas construções]\label{ex:construcoes}\leavevmode
\begin{itemize}
  \item $[12,30,21,8,$\allowbreak$27,40,16,35,19]$, com $n=9$, processa $i=4,3,2,1$. Em $i=4$, troca $8\leftrightarrow35$; em $i=3$, $21\leftrightarrow40$; em $i=2$, $30\leftrightarrow35$, e o 30 para na posição 4, pois os seus novos filhos, 8 e 19, são menores; em $i=1$, $12\leftrightarrow40$ e depois $12\leftrightarrow21$. Resultado: $[40,35,21,30,$\allowbreak$27,12,16,8,19]$.
  \item $[6,9,14,3,12,10,2]$ processa $i=3,2,1$: nada muda em $i=3$, pois 14 já é maior que 10 e 2; em $i=2$, $9\leftrightarrow12$; em $i=1$, $6\leftrightarrow14$ e depois $6\leftrightarrow10$. Resultado: $[14,12,10,3,9,6,2]$.
\end{itemize}
\end{exemplo}

\section{Heapsort}\label{sec:heapsort}
O heap leva a um algoritmo de ordenação simples: construa o heap e, repetidamente, troque a raiz (o maior elemento restante) com a última posição da parte ainda não ordenada e faça a nova raiz descer. O vetor fica dividido em duas partes: à esquerda, um heap que encolhe; à direita, a parte já ordenada, que cresce.
\begin{pseudo}{Heapsort}
\begin{algorithmic}
\Procedure{heapsort}{$T,n$}
  \State \Call{arranjar}{$n$} \Comment{$O(n)$}
  \For{$m := n$ \textbf{decrescendo até} $2$}
    \State $T[1]\Leftrightarrow T[m]$ \Comment{o máximo vai para o fim}
    \State \Call{descer}{$1, m-1$} \Comment{$O(\log m)$}
  \EndFor
\EndProcedure
\end{algorithmic}
\end{pseudo}

\begin{exemplo}[{Heapsort de $[5,2,9,4,7]$}]\label{ex:heapsort}
O \emph{arranjar} processa $i=2$ (troca $2\leftrightarrow7$) e $i=1$ (troca $5\leftrightarrow9$), e o heap é $[9,7,5,4,2]$. Em seguida, a cada passo o máximo vai para o fim e a nova raiz desce (Figura~\ref{fig:heapsort}).
\begin{figura}[Heapsort de $[5,2,9,4,7]$. A parte verde já está na posição final; a parte branca é um heap. Em cada linha, a raiz trocou de lugar com a última posição do heap e depois desceu.\label{fig:heapsort}]
\begin{tikzpicture}[x=7mm,y=7.5mm]
\foreach \lin/\vals/\k/\txt in {0/{9,7,5,4,2}/5/{heap inicial},1/{7,4,5,2,9}/4/{$9\leftrightarrow2$; o 2 desce: troca com 7 e com 4},2/{5,4,2,7,9}/3/{$7\leftrightarrow2$; o 2 desce: troca com 5},3/{4,2,5,7,9}/2/{$5\leftrightarrow2$; o 2 desce: troca com 4},4/{2,4,5,7,9}/0/{$4\leftrightarrow2$: ordenado}}{
  \foreach \x [count=\i] in \vals {
    \ifnum\i>\k \fill[feito] (\i-0.5,-\lin-0.3) rectangle (\i+0.5,-\lin+0.3); \fi
    \node[cel] at (\i,-\lin) {\x};}
  \node[rot,anchor=west] at (5.8,-\lin) {\txt};}
\end{tikzpicture}
\end{figura}
\end{exemplo}

O custo é $O(n)+\sum_mO(\log m)=\Theta(n\log n)$ em \emph{todos} os casos, mesmo com a entrada já ordenada. O heapsort é \textbf{in-place} (só faz trocas dentro do vetor) e \textbf{não é estável}, pois a troca da raiz com o último elemento pode inverter a ordem relativa de elementos iguais. Como faz $O(n\log n)$ comparações, é ótimo entre os algoritmos baseados em comparações (Capítulo~\ref{cap:sortcomp}). Por exemplo, ele transforma $[12,30,21,8,27,40,16,35,19]$ em $[8,12,16,19,21,27,30,35,40]$.

% ------------------------------------------------------------------
\section{Propriedades de heaps e suas demonstrações}\label{sec:heapprops}
\begin{roteiro}[Como demonstrar propriedades de heaps]
Quase toda demonstração sobre heaps usa só três ferramentas: (i) a definição $s_i\le s_{\floor{i/2}}$ combinada com a \textbf{transitividade ao longo de um caminho}; (ii) os \textbf{fatos sobre índices} (pai, filhos, folhas, níveis); (iii) \textbf{indução} na altura, no nível ou no número de nós. Para ``provar ou dar contraexemplo'', teste primeiro exemplos pequenos, com $n\le8$.
\end{roteiro}

\begin{prop}[O ancestral domina]\label{prop:ancestral}
Se $j$ é ancestral de $i$ num heap, então $s_i\le s_j$. Em particular, todo caminho da raiz até uma folha é não crescente, a \textbf{raiz é o máximo} e \textbf{toda subárvore é um heap}.
\end{prop}
\begin{proof}
O caminho de $i$ até $j$ é $i,\floor{i/2},\floor{i/4},\dots,j$. Aplicando a definição em cada aresta e a transitividade de $\le$, obtemos $s_i\le s_{\floor{i/2}}\le\cdots\le s_j$. A raiz é ancestral de todos os nós. A condição de heap numa subárvore é a mesma condição, restrita a ela.
\end{proof}
Por exemplo, no heap da Figura~\ref{fig:heaprep}, o caminho da raiz até a posição 8 é $90\ge55\ge41\ge37$.

\begin{prop}[O mínimo está numa folha]
Num heap de máximo com chaves distintas, o menor elemento está numa folha, isto é, numa posição maior que $\floor{n/2}$.
\end{prop}
\begin{proof}
Se $i$ não é folha, tem o filho $2i$ e $s_{2i}<s_i$, pois as chaves são distintas; logo $s_i$ não é o mínimo. Com repetições, \emph{algum} mínimo está numa folha.
\end{proof}
No heap da Figura~\ref{fig:heaprep}, as folhas são as posições $5,\dots,8$, com valores $30,64,75,37$, e o mínimo, $30$, está entre elas. Note que o mínimo pode estar em \emph{qualquer} folha, não necessariamente na última posição.

\begin{prop}[Onde está o $k$-ésimo maior]
Com chaves distintas, o $k$-ésimo maior elemento está num nível $\le k$, isto é, numa posição $\le 2^k-1$. Em particular, o segundo maior é filho da raiz (posição 2 ou 3), e o terceiro maior está entre as posições 2 e 7.
\end{prop}
\begin{proof}
Pela Proposição~\ref{prop:ancestral}, todos os ancestrais de $x$ são maiores que $x$. Se $x$ é o $k$-ésimo maior, existem só $k-1$ elementos maiores que ele, e portanto $x$ tem no máximo $k-1$ ancestrais: está no nível $\le k$. As posições dos níveis $1,\dots,k$ são $1,\dots,2^k-1$.
\end{proof}
No heap da Figura~\ref{fig:heaprep}, o segundo maior, 82, está na posição 3, e o terceiro maior, 75, na posição 7, no limite permitido.

\begin{prop}[Altura]
Um heap com $n$ elementos tem altura $\floor{\log_2 n}+1$ (em níveis), que é $\Theta(\log n)$. Por isso subir, descer, inserir e remover custam $O(\log n)$.
\end{prop}
\begin{proof}
O último nó, $n$, está no nível $\floor{\log_2n}+1$, e numa árvore completa todos os níveis anteriores estão cheios: $2^{h-1}\le n\le 2^h-1$, logo $h=\floor{\log_2n}+1$.
\end{proof}
Um heap com um milhão de elementos tem apenas $\floor{\log_2 10^6}+1=20$ níveis: nenhuma operação faz mais que 19 trocas.

\begin{prop}[Heaps e ordenação]
Toda sequência em ordem não crescente é um heap; a recíproca é falsa, como mostra $[3,1,2]$. Construir um heap é, portanto, mais fácil que ordenar: $O(n)$ contra $\Omega(n\log n)$. E, como o heapsort ordena com uma construção e $n$ remoções, nenhuma implementação de heap baseada em comparações pode ter construção \emph{e} remoção com custo amortizado $o(\log n)$; caso contrário, ordenaríamos em $o(n\log n)$.
\end{prop}

\begin{prop}[Troca de dois elementos]
Sejam $S$ um heap e $i<j$. Trocar $s_i$ com $s_j$ não produz necessariamente um heap, nem quando $s_i<s_j$ nem quando $s_i>s_j$.
\end{prop}
\begin{proof}[Contraexemplos]
Para $s_i<s_j$: $S=[10,9,1,8,7,0,0,6]$ é um heap; trocando $s_3=1$ com $s_4=8$, obtemos $[10,9,8,1,$\allowbreak$7,0,0,6]$, em que a posição 8 ($6$) tem pai na posição 4 ($1$). O valor pequeno foi para uma posição cujos filhos podem ser maiores que ele. Para $s_i>s_j$: em $S=[10,9]$, trocar as duas posições dá $[9,10]$, com o filho maior que o pai.
\end{proof}

\begin{prop}[Número de heaps]
Seja $H(n)$ o número de permutações de $1,\dots,n$ que são heaps (o mesmo vale para heaps de mínimo). Então $H(0)=H(1)=1$ e
\[
H(n)=\binom{n-1}{L}\,H(L)\,H(n-1-L),
\]
onde $L$ é o número de nós da subárvore esquerda da raiz: com $h=\floor{\log_2 n}$ e $u=n-(2^h-1)$ nós no último nível, $L=(2^{h-1}-1)+\min\{u,2^{h-1}\}$. Os valores para $n=1,\dots,10$ são $1,1,2,3,8,20,80,210,896,3360$.
\end{prop}
\begin{proof}
A raiz é necessariamente $n$. O formato da árvore é fixo, então $L$ é determinado por $n$. Escolhem-se quais $L$ dos $n-1$ valores restantes vão para a esquerda, de $\binom{n-1}{L}$ modos; cada lado precisa ser um heap sobre os seus valores, e o número de heaps depende só da quantidade de valores.
\end{proof}

\begin{exemplo}[Os três heaps de $\{1,2,3,4\}$]\label{ex:tresheaps}
Com $n=4$, a raiz é 4, a subárvore esquerda tem as posições 2 e 4 ($L=2$) e a direita, só a posição 3. Escolhemos quais 2 dos valores $\{1,2,3\}$ vão para a esquerda, de $\binom32=3$ modos, e cada escolha dá um único heap: $\{1,2\}$ dá $[4,2,3,1]$; $\{1,3\}$ dá $[4,3,2,1]$; $\{2,3\}$ dá $[4,3,1,2]$.
\end{exemplo}

\begin{prop}[Correção de subir e descer]
(a) Se $T$ era um heap e a chave de $i$ aumentou, após subir$(i)$ $T$ volta a ser um heap. (b) Se as subárvores dos filhos de $i$ são heaps, após descer$(i,n)$ a subárvore de $i$ é um heap.
\end{prop}
\begin{proof}[Esboço]
(a) Invariante: a única violação possível é entre o nó corrente e o seu pai. Ao trocar com o pai $j$, o antigo valor de $j$ desce para a posição corrente e é maior ou igual aos filhos dela, que já eram seus descendentes; o novo valor em $j$ é maior que o antigo, e não viola nada com os filhos de $j$. Repete-se para cima. (b) Invariante: a única violação é entre o nó corrente e os seus filhos. Trocando com o \emph{maior} filho, este passa a dominar o outro filho e o valor que desceu. A violação desce um nível e termina numa folha ou quando deixa de existir.
\end{proof}

\begin{chave}
Heap $=$ vetor com ``pai $\ge$ filho''. Sobe-se após um aumento e desce-se após uma diminuição, sempre trocando com o maior filho. Inserir $=$ colocar no fim e subir; remover $=$ levar o último à raiz e descer. A construção é linear porque há no máximo $\ceil{n/2^j}$ nós de altura $j$. O heapsort é $\Theta(n\log n)$, in-place e não estável.
\end{chave}

% ----------------------------------------------------------------
\section{Exercícios}
\excurso
\begin{exercicio}{hp-l1}[Lista]<ex:ehheap>
Verifique se as sequências correspondem a um heap binário (de máximo):\\
(a)~$33\ 32\ 28\ 31\ 26\ 29\ 25\ 30\ 27$;\qquad (b)~$33\ 32\ 28\ 31\ 29\ 26\ 25\ 30\ 27$.
\end{exercicio}
\begin{exercicio}{hp-l2}[Lista]<sec:heapprops>
Prove ou dê contraexemplo: se $S$ é um heap e $s_i,s_j$ são chaves com $i<j$ e $s_i<s_j$, então trocar $s_i$ com $s_j$ produz um heap.
\end{exercicio}
\begin{exercicio}{hp-l3}[Lista]<sec:heapprops>
Repita o exercício anterior com $i<j$ e $s_i>s_j$.
\end{exercicio}
\begin{exercicio}{hp-l4}[Lista]<sec:heapprops>
Determine o número de permutações de $1,\dots,n$ que são heaps de mínimo.
\end{exercicio}
\begin{exercicio}{hp-l5}[Lista]<ex:arranjar>
Aplique o algoritmo linear de construção à lista $18\ 25\ 41\ 34\ 14\ 10\ 52\ 50\ 48$.
\end{exercicio}
\begin{exercicio}{hp-l6}[Lista]<lema:alturas>
Prove que no máximo $\ceil{n/2^j}$ nós têm altura $j$ numa árvore binária completa com $n$ nós.
\end{exercicio}
\begin{exercicio}{hp-l16}[Lista]<sec:subirdescer>
Escreva os algoritmos de inserção, remoção, aumento e diminuição de prioridade e construção de heaps binários \textbf{sem recursão}.
\end{exercicio}
\begin{exercicio}{hp-r1}[Monitoria]<sec:heapdef>
(a) Defina heap e a propriedade que $s_1,\dots,s_n$ satisfazem. (b) Por que o elemento de maior prioridade é selecionado em $O(1)$? (c) Dê a complexidade de pior caso de seleção, inserção, remoção do máximo, alteração de prioridade e construção.
\end{exercicio}
\begin{exercicio}{hp-r2}[Monitoria]<sec:heapops,sec:arranjar>
No heap $H=[95,60,78,39,28,66,70,33]$: (a) a prioridade 66 passa a 98; (b) remove-se a raiz. Construa, pelo procedimento baseado em \emph{descer}, heaps a partir de (c) $[20,15,30,40,10,25,5]$ e (d) $[7,2,15,9,1,20,13]$. Mostre as trocas em cada caso.
\end{exercicio}
\begin{exercicio}{hp-r5}[Monitoria]<sec:heapops>
Um aluno escreveu: ``$H[i]:=nova$; \textbf{se} $H[i]>H[\floor{i/2}]$ \textbf{então} Descer$(i)$ \textbf{senão} Subir$(i)$''. Qual é o erro? Corrija e explique.
\end{exercicio}
\begin{exercicio}{hp-tele}[Slides]<sec:listaprior>
Problema \emph{Telemarketing} (Beecrowd 2330): $n$ vendedores, numerados de 1 a $n$, atendem $L$ ligações em ordem; cada ligação vai para o vendedor livre de menor número, que fica ocupado pela duração da ligação. Quantas ligações cada vendedor atende? Projete a solução com heap e dê a complexidade.
\end{exercicio}

\excompl
\begin{exercicio}{hp-dario}[Complementar]<sec:heapdef>
Num heap $d$-ário (cada nó tem até $d$ filhos), em vetor indexado a partir de 1: (a) quais são os índices do pai e dos filhos do nó $i$? (b) Qual a altura? (c) Quanto custam \emph{subir} e \emph{descer}? (d) Por que um $d$ maior pode ser bom quando há muito mais aumentos de prioridade do que remoções?
\end{exercicio}
\begin{exercicio}{hp-kmerge}[Complementar]<sec:heapops>
Dadas $k$ listas ordenadas com $n$ elementos no total, intercale-as numa única lista ordenada em $O(n\log k)$.
\end{exercicio}
\begin{exercicio}{hp-kmaiores}[Complementar]<sec:heapops>
Um fluxo de $n$ números chega um a um, e você tem memória $O(k)$. Como obter os $k$ maiores ao final? Qual o tempo?
\end{exercicio}
\begin{exercicio}{hp-verif}[Complementar]<ex:ehheap>
Escreva um algoritmo $O(n)$ que decide se um vetor é um heap.
\end{exercicio}
\begin{exercicio}{hp-estavel}[Complementar]<ex:heapsort>
Mostre, com um exemplo de 3 elementos, que o heapsort não é estável.
\end{exercicio}
\begin{exercicio}{hp-crescente}[Complementar]<sec:arranjar>
Mostre que construir um heap de máximo inserindo $1,2,\dots,n$, nessa ordem, por inserções sucessivas custa $\Theta(n\log n)$, enquanto o \emph{arranjar} sobre o mesmo vetor custa $O(n)$.
\end{exercicio}
\begin{exercicio}{hp-ins}[Complementar]<ex:heapins>
Insira $15,40,25,50,10,45$, nessa ordem, num heap de máximo inicialmente vazio. Mostre o vetor após cada inserção.
\end{exercicio}
\begin{exercicio}{hp-rem2}[Complementar]<ex:heaprem>
Remova o máximo duas vezes do heap $[90,70,80,30,60,50,40,10,20]$, mostrando o vetor após cada remoção.
\end{exercicio}
\begin{exercicio}{hp-sort5}[Complementar]<ex:heapsort>
Execute o heapsort sobre $[4,10,3,5,1]$, mostrando o vetor após o \emph{arranjar} e após cada passo do laço principal.
\end{exercicio}
\begin{exercicio}{hp-remi}[Complementar]<sec:heapops>
Escreva uma operação que remove o elemento da posição $i$ (não necessariamente a raiz) de um heap em $O(\log n)$. Por que, depois de pôr o último elemento na posição $i$, às vezes é preciso subir e às vezes descer?
\end{exercicio}
\begin{exercicio}{hp-mediana}[Complementar]<sec:listaprior>
Mostre como manter a mediana de um conjunto que recebe inserções, com inserção em $O(\log n)$ e consulta à mediana em $O(1)$, usando um heap de máximo e um heap de mínimo.
\end{exercicio}
\begin{exercicio}{hp-comp}[Complementar \estrela]<sec:arranjar>
Mostre que o \emph{arranjar} faz no máximo $2n$ comparações entre chaves (para $n$ grande), contando as comparações feitas em cada passo do \emph{descer}.
\end{exercicio}

% ================================================================
\chapter{Heaps binomiais, heaps de Fibonacci e análise amortizada}\label{cap:heapav}\label{sec:heapav}
% ================================================================
\begin{objetivos}
\begin{itemize}
  \item Árvores binomiais e o heap binomial, cuja fusão imita a soma de números binários.
  \item A análise amortizada pelos métodos do banqueiro e do potencial.
  \item O heap de Fibonacci: a ideia de adiar trabalho, as marcas e os cortes em cascata, e por que o grau dos nós é logarítmico.
  \item Como escolher entre os heaps binário, binomial e de Fibonacci.
\end{itemize}
\end{objetivos}

O heap binário do capítulo anterior tem um ponto fraco: \textbf{unir} duas listas de prioridades exige concatenar os vetores e reconstruir o heap, o que custa $\Theta(n)$. Imagine um sistema operacional com uma fila de processos por processador: quando um processador é desligado, a sua fila precisa ser fundida com a de outro. Em aplicações com muitas fusões, pagar $\Theta(n)$ a cada uma é inviável.

Este capítulo apresenta duas estruturas em que a fusão é barata. O \textbf{heap binomial} representa a lista de prioridades como uma coleção de árvores cujas ordens correspondem aos bits de $n$, e funde duas listas como quem soma dois números binários, em $O(\log n)$ no pior caso. O \textbf{heap de Fibonacci} vai além: adia toda a arrumação para o momento da remoção do máximo, e só é eficiente quando se considera o custo de uma \emph{sequência} de operações. Por isso a análise amortizada, que formaliza esse tipo de garantia, aparece aqui.

% ------------------------------------------------------------------
\section{Árvores binomiais}\label{sec:binomial}
\begin{defi}
A árvore binomial $B_0$ é um único nó. Para $k\ge1$, $B_k$ é obtida \textbf{subordinando} (operação \emph{link}) uma $B_{k-1}$ a outra $B_{k-1}$: a raiz de uma passa a ser filha da raiz da outra. Num heap, a raiz de \emph{maior} prioridade fica por cima.
\end{defi}

\begin{figura}[As árvores binomiais $B_0$ a $B_4$. Em cada $B_k$, a subárvore em destaque é a cópia de $B_{k-1}$ que foi subordinada à outra.\label{fig:binomiais}]
\forestset{bn/.style={for tree={circle, draw, minimum size=2.8mm, inner sep=0pt, s sep=1.6mm, l sep=1.6mm, edge={thick}}}}
\begin{tabular}{@{}c@{\hspace{7mm}}c@{\hspace{7mm}}c@{\hspace{7mm}}c@{\hspace{7mm}}c@{}}
\begin{forest} bn, baseline [] \end{forest} &
\begin{forest} bn, baseline [ [,fill=destaque] ] \end{forest} &
\begin{forest} bn, baseline [ [,fill=destaque [,fill=destaque] ] [] ] \end{forest} &
\begin{forest} bn, baseline [ [,fill=destaque [,fill=destaque [,fill=destaque]] [,fill=destaque]] [ [] ] [] ] \end{forest} &
\begin{forest} bn, baseline [ [,fill=destaque [,fill=destaque [,fill=destaque [,fill=destaque]] [,fill=destaque]] [,fill=destaque [,fill=destaque]] [,fill=destaque]] [ [ [] ] [] ] [ [] ] [] ] \end{forest}\\[2pt]
$B_0$ & $B_1$ & $B_2$ & $B_3$ & $B_4$
\end{tabular}
\end{figura}

\begin{prop}\label{prop:binom}
Para todo $k\ge0$:
\begin{enumerate}
  \item $B_k$ tem exatamente $2^k$ nós;
  \item o grau da raiz, chamado \textbf{ordem} de $B_k$, é $k$, e é o maior grau da árvore;
  \item a altura de $B_k$, medida em arestas, é $k$;
  \item há exatamente $\binom kj$ nós na profundidade $j$;
  \item os filhos da raiz de $B_k$ são raízes de árvores $B_{k-1},B_{k-2},\dots,B_1,B_0$.
\end{enumerate}
\end{prop}
\begin{proof}
Todas as propriedades seguem por indução em $k$, olhando para o \emph{link} que forma $B_k$ a partir de duas cópias de $B_{k-1}$: o número de nós dobra ($2\cdot2^{k-1}$); a raiz ganha um filho; a árvore subordinada fica um nível abaixo, e a altura passa a $\max\{k-1,\,1+(k-1)\}=k$; os nós na profundidade $j$ vêm da cópia de cima (profundidade $j$) e da de baixo (profundidade $j-1$), e $\binom{k-1}{j}+\binom{k-1}{j-1}=\binom kj$ --- daí o nome ``binomial''; e o novo filho da raiz é uma $B_{k-1}$, que se junta aos filhos $B_{k-2},\dots,B_0$ que a raiz já tinha.
\end{proof}

\begin{exemplo}[Conferindo as propriedades em $B_3$]
Olhe para $B_3$ na Figura~\ref{fig:binomiais}. Ele tem $8=2^3$ nós; a raiz tem 3 filhos e nenhum outro nó tem grau maior; o caminho mais longo da raiz até uma folha tem 3 arestas. Contando os nós por profundidade, obtemos $1,3,3,1$, que são os coeficientes $\binom30,\binom31,\binom32,\binom33$ --- a linha 3 do triângulo de Pascal. Por fim, os filhos da raiz, da esquerda para a direita, são raízes de uma $B_2$ (4 nós), uma $B_1$ (2 nós) e uma $B_0$ (1 nó).
\end{exemplo}

\begin{cuidado}[Duas convenções de altura]
A altura de $B_k$ é $k$ quando medida em arestas, e $k+1$ quando medida em níveis (com folhas de altura 1, como nos heaps binários do Capítulo~\ref{cap:heap}). É a mesma árvore. Em exercícios sobre alturas, declare a convenção usada. Neste livro, a ordem de $B_k$ é o grau da raiz, $k$, e a altura de $B_k$ é $k$ arestas.
\end{cuidado}

% ------------------------------------------------------------------
\section{Heap binomial}\label{sec:heapbinomial}
\begin{defi}
Um \textbf{heap binomial} é uma floresta $H$ de árvores binomiais em que (1) cada árvore satisfaz a condição de heap (filho $\le$ pai) e (2) as árvores têm \textbf{ordens distintas}, isto é, há no máximo uma árvore de cada ordem. As raízes ficam numa lista, em ordem crescente de ordem.
\end{defi}

\begin{prop}
(a) Para cada $n$ existe um único formato de heap binomial com $n$ nós: $B_k$ pertence a $H$ se e somente se o bit $k$ da representação binária de $n$ é 1. (b) Portanto $H$ tem no máximo $\floor{\log_2n}+1$ árvores, cada uma de altura $\le\floor{\log_2n}$.
\end{prop}
\begin{proof}
A árvore $B_k$ contribui com $2^k$ nós e as ordens são distintas, logo $n=\sum_{B_k\in H}2^k$: é a representação binária de $n$, que é única. Um número $n$ tem $\floor{\log_2n}+1$ bits.
\end{proof}

\begin{exemplo}[Um heap binomial com 13 elementos]
Como $13=1101_2=8+4+1$, um heap binomial com 13 elementos tem exatamente as árvores $B_0$, $B_2$ e $B_3$ (Figura~\ref{fig:binheap13}). Do mesmo modo, $26=11010_2$ dá $H=\{B_1,B_3,B_4\}$. O máximo do heap é a maior das raízes, aqui o 50.
\begin{figura}[Um heap binomial com $13=1101_2$ elementos. Cada árvore é um heap (pai $\ge$ filhos), e há uma árvore para cada bit 1 de $n$. A lista de raízes está em ordem crescente de ordem.\label{fig:binheap13}]
\begin{tabular}{@{}c@{\hspace{10mm}}c@{\hspace{10mm}}c@{}}
\begin{forest} bstp, baseline [12] \end{forest} &
\begin{forest} bstp, baseline [41 [33 [17]] [28]] \end{forest} &
\begin{forest} bstp, baseline [50,feito [44 [31 [10]] [25]] [37 [29]] [6]] \end{forest}\\[2pt]
\small $B_0$ (bit 0) & \small $B_2$ (bit 2) & \small $B_3$ (bit 3)
\end{tabular}\\[4pt]
{\small\color{MidnightBlue} $13=\underbrace{1}_{B_3}\underbrace{1}_{B_2}\underbrace{0}_{\text{sem }B_1}\underbrace{1}_{B_0}{}_2$}
\end{figura}
\end{exemplo}

\subsection{Operações}
A operação central é a \textbf{fusão}, e as demais são escritas em termos dela.
\begin{pseudo}{Link e fusão de heaps binomiais}
\begin{algorithmic}
\Procedure{Link}{$y,z$} \Comment{$y$ e $z$ são raízes de árvores $B_{k-1}$; $z$ fica por cima}
  \State $pai[y]\gets z$;\quad $\textit{irmão}[y]\gets filho[z]$;\quad $filho[z]\gets y$;\quad $grau[z]\gets grau[z]+1$
\EndProcedure
\Procedure{Fusão}{$H',H''$}
  \State $H := H'\cup H''$ \Comment{intercala as listas de raízes por ordem crescente: $O(\log n)$}
  \While{$H$ contém duas árvores $B,B'$ de mesma ordem} \Comment{menor ordem primeiro}
    \If{$B.raiz.chave \ge B'.raiz.chave$} \State subordinar $B'$ a $B$
    \Else \State subordinar $B$ a $B'$
    \EndIf
  \EndWhile
\EndProcedure
\end{algorithmic}
\end{pseudo}

\begin{intuicao}[Fundir é somar em binário]
Na soma $011_2+011_2$, a coluna das unidades tem dois 1, que viram um 0 e um ``vai um'' para a coluna seguinte. Na fusão, duas árvores $B_0$ viram uma $B_1$, que ``vai'' para a posição das árvores de ordem 1. Ali podem se encontrar até três árvores de mesma ordem (duas originais e a que veio da posição anterior): uma fica, e as outras duas se juntam e vão para a posição seguinte, exatamente como $1+1+1=11_2$.
\end{intuicao}

Há $O(\log n)$ posições e, em cada uma, no máximo três árvores de mesma ordem ao mesmo tempo, de modo que a fusão custa $O(\log n)$ no pior caso. Numa implementação com ponteiros, a lista de raízes é percorrida mantendo o nó anterior, o atual e o próximo, e há quatro casos: ordens diferentes, três árvores seguidas de mesma ordem, e duas árvores de mesma ordem em que vence a atual ou a seguinte.

\begin{exemplo}[Fusão passo a passo]\label{ex:fusaobin}
Sejam $H_1=\{7,\ 9\to3\}$ e $H_2=\{5,\ 8\to6\}$, cada um com $3=11_2$ elementos: uma $B_0$ e uma $B_1$ (Figura~\ref{fig:fusaobin}).
\begin{enumerate}
  \item Ordem 0: há duas $B_0$, com raízes 7 e 5. O link põe a maior por cima e forma a $B_1$ $7\to5$ (o ``vai um'').
  \item Ordem 1: agora há três $B_1$: $7\to5$, $9\to3$ e $8\to6$. A primeira fica; as outras duas se juntam, com 9 por cima, numa $B_2$.
  \item O resultado é $\{B_1, B_2\}$, com $6=110_2$ elementos: de fato, $11_2+11_2=110_2$.
\end{enumerate}
\begin{figura}[A fusão de $H_1$ com $H_2$ (Exemplo~\ref{ex:fusaobin}) imita a soma $11_2+11_2=110_2$.\label{fig:fusaobin}]
\begin{tabular}{@{}c@{\hspace{9mm}}c@{\hspace{9mm}}c@{}}
\begin{forest} bstp, baseline [7] \end{forest}\ \begin{forest} bstp, baseline [9 [3]] \end{forest}
\ \ \raisebox{0.5ex}{\small$+$}\ \
\begin{forest} bstp, baseline [5] \end{forest}\ \begin{forest} bstp, baseline [8 [6]] \end{forest}
&
\begin{forest} bstp, baseline [7,novo [5]] \end{forest}\ \begin{forest} bstp, baseline [9 [3]] \end{forest}\ \begin{forest} bstp, baseline [8 [6]] \end{forest}
&
\begin{forest} bstp, baseline [7 [5]] \end{forest}\ \begin{forest} bstp, baseline [9,novo [8 [6]] [3]] \end{forest}\\[3pt]
\small (a) $H_1$ e $H_2$ & \small (b) $B_0+B_0\to B_1$ & \small (c) duas das três $B_1$ viram $B_2$
\end{tabular}
\end{figura}
\end{exemplo}

\begin{exemplo}[{Fusão de $H_1=\{B_1,B_2,B_3\}$ com $H_2=\{B_0,B_1,B_3\}$}]\label{ex:fusaoformas}
Olhando apenas para os formatos: os heaps têm $14=1110_2$ e $11=1011_2$ elementos. A $B_0$ de $H_2$ fica sozinha na ordem 0. Na ordem 1, as duas $B_1$ formam uma $B_2$ (o ``vai um''); na ordem 2, essa $B_2$ se junta à $B_2$ de $H_1$ e forma uma $B_3$; na ordem 3 há, então, três árvores --- as $B_3$ de $H_1$ e de $H_2$ e a que acabou de se formar ---, e, como na soma $1+1+1=11_2$, uma delas fica e as outras duas formam uma $B_4$. São três links, e o resultado é $\{B_0,B_3,B_4\}$, isto é, $25=11001_2=1110_2+1011_2$.
\end{exemplo}

As demais operações são as seguintes.
\begin{itemize}
  \item \textbf{Inserção}: fusão com um heap $\{B_0\}$ que contém só o novo elemento. No pior caso custa $O(\log n)$, quando $n=2^k-1$, todos os bits são 1 e há $k$ links. Mas inserir equivale a somar 1 a um contador binário, e cada link corresponde a um ``vai um''. Se os bits de $n$ são equiprováveis, a inserção faz em média um link, pois a posição $j$ do primeiro 0 de $n$ (lido da direita para a esquerda) tem esperança $\sum_{j\ge1}j/2^j=2$. E, ao longo de uma sequência de inserções, o custo \textbf{amortizado} é $O(1)$, como mostra o exemplo do contador binário da §\ref{sec:amort}.
  \item \textbf{Seleção do máximo}: o máximo é uma das no máximo $\floor{\log n}+1$ raízes, e é encontrado em $O(\log n)$, ou em $O(1)$ se um ponteiro para a maior raiz for mantido.
  \item \textbf{Remoção do máximo}: encontre a raiz máxima $r$, numa árvore $B_k$, e retire $B_k$ de $H$, formando $H'=H-B_k$. Removendo $r$, seus filhos $B_{k-1},\dots,B_0$ formam um heap binomial $H''$. Faça a fusão $H'\cup H''$. Total: $O(\log n)$.
  \item \textbf{Aumento de prioridade}: o nó sobe dentro da própria árvore trocando de lugar com o pai, como no heap binário, em $O(\log n)$, pois a altura é $\le\log n$.
  \item \textbf{Diminuição de prioridade}: remova o nó (veja abaixo) e reinsira-o com a nova prioridade, em $O(\log n)$. Fazê-lo descer trocando com o maior filho também funciona, mas cada nível tem até $O(\log n)$ filhos a comparar, o que dá $O(\log^2n)$.
  \item \textbf{Remoção de um nó qualquer}: aumente sua prioridade para $+\infty$, para que ele suba até a raiz, e remova o máximo: $O(\log n)$.
  \item \textbf{Construção}: $n$ inserções sucessivas custam $O(n)$ no total, pelo mesmo argumento do contador binário.
\end{itemize}

\begin{exemplo}[Remoção do máximo]\label{ex:rembin}
No heap resultante do Exemplo~\ref{ex:fusaobin}, $\{7\to5,\ B_2 \text{ de raiz } 9\}$, o máximo é a raiz 9. Retirando a sua árvore, sobra $H'=\{7\to5\}$. Os filhos de 9 formam outro heap binomial, $H''=\{3,\ 8\to6\}$ (uma $B_0$ e uma $B_1$). A fusão $H'\cup H''$ encontra duas $B_1$, $7\to5$ e $8\to6$, e as junta com 8 por cima. O resultado tem $5=101_2$ elementos (Figura~\ref{fig:rembin}).
\begin{figura}[Remoção do máximo: os filhos da raiz removida formam um heap binomial, que é fundido com o resto.\label{fig:rembin}]
\begin{tabular}{@{}c@{\hspace{10mm}}c@{\hspace{10mm}}c@{}}
\begin{forest} bstp, baseline [7 [5]] \end{forest}\ \begin{forest} bstp, baseline [9,desreg [8 [6]] [3]] \end{forest}
&
\begin{forest} bstp, baseline [7 [5]] \end{forest}\ \ \raisebox{0.5ex}{\small$\cup$}\ \ \begin{forest} bstp, baseline [3] \end{forest}\ \begin{forest} bstp, baseline [8 [6]] \end{forest}
&
\begin{forest} bstp, baseline [3] \end{forest}\ \begin{forest} bstp, baseline [8,novo [7 [5]] [6]] \end{forest}\\[3pt]
\small (a) o máximo é 9 & \small (b) $H'$ e os filhos de 9 & \small (c) após a fusão
\end{tabular}
\end{figura}
\end{exemplo}

% ------------------------------------------------------------------
\section{Análise amortizada}\label{sec:amort}
\begin{intuicao}
Pense num cofrinho. Toda vez que você faz algo barato, paga um pouco a mais e guarda o troco no cofrinho. Quando chega uma despesa grande, ela é paga com o dinheiro guardado. Se o cofrinho nunca fica negativo, o total gasto nunca passou do que você ``cobrou'' de cada operação. A análise amortizada é a contabilidade desse cofrinho.
\end{intuicao}

Em várias estruturas, uma operação individual pode ser cara, mas operações caras são raras, porque cada uma delas desfaz um trabalho acumulado por muitas operações baratas. A \textbf{análise amortizada} mede o custo total de uma \emph{sequência} de operações, no pior caso, e o divide pelo número de operações. É uma garantia de pior caso para a sequência e \textbf{não envolve probabilidade}.

\paragraph{Método do banqueiro (créditos).} Cada operação recebe um número fixo de créditos, que é o seu custo amortizado. Operações baratas deixam créditos guardados na estrutura, e esses créditos pagam as operações caras mais tarde. Se o saldo de créditos nunca fica negativo, o custo real total é no máximo a soma dos custos amortizados.

\paragraph{Método do potencial (físico).} Escolhe-se uma função $\Phi\ge0$ sobre os estados da estrutura, que mede o trabalho ``armazenado'' (o dinheiro no cofrinho). O custo amortizado da $i$-ésima operação, de custo real $t_i$, é $a_i=t_i+\Phi(E_i)-\Phi(E_{i-1})$. Somando, os termos de potencial se cancelam:
\[
\sum_{i=1}^m t_i=\sum_{i=1}^m a_i-\Phi_m+\Phi_0\ \le\ \sum_{i=1}^m a_i\qquad\text{sempre que }\Phi_m\ge\Phi_0 .
\]

\begin{exemplo}[Contador binário]\label{ex:contador}
Incrementar um contador de $k$ bits pode trocar $k$ bits (de $0111$ para $1000$, por exemplo), mas isso é raro. A tabela mostra os oito primeiros incrementos a partir de zero:
\begin{center}\small
\begin{tabular}{@{}rcccccc@{}}
\toprule
incremento & contador & bits trocados ($t_i$) & $\Phi$ (bits 1) & $\Delta\Phi$ & $a_i=t_i+\Delta\Phi$ & total real\\
\midrule
--- & \texttt{0000} & --- & 0 & --- & --- & 0\\
1 & \texttt{000\textbf{1}} & 1 & 1 & $+1$ & 2 & 1\\
2 & \texttt{00\textbf{10}} & 2 & 1 & $0$ & 2 & 3\\
3 & \texttt{001\textbf{1}} & 1 & 2 & $+1$ & 2 & 4\\
4 & \texttt{0\textbf{100}} & 3 & 1 & $-1$ & 2 & 7\\
5 & \texttt{010\textbf{1}} & 1 & 2 & $+1$ & 2 & 8\\
6 & \texttt{01\textbf{10}} & 2 & 2 & $0$ & 2 & 10\\
7 & \texttt{011\textbf{1}} & 1 & 3 & $+1$ & 2 & 11\\
8 & \texttt{\textbf{1000}} & 4 & 1 & $-2$ & 2 & 15\\
\bottomrule
\end{tabular}
\end{center}
Os custos reais oscilam (1, 2, 1, 3, 1, 2, 1, 4), mas o custo amortizado é sempre 2. Em geral, com $\Phi=$ número de bits iguais a 1, um incremento que zera $t$ bits e liga um tem custo amortizado $a=(t+1)+(1-t)=2$. Logo $m$ incrementos trocam no máximo $2m$ bits: na tabela, $15\le16$. É o mesmo argumento que mostra que as inserções num heap binomial custam $O(1)$ amortizado: inserir é somar 1.
\end{exemplo}

\begin{exemplo}[Pilha com desempilhamento múltiplo]\label{ex:pilha}
Cada operação desempilha $k\ge0$ elementos e depois empilha um; uma operação isolada pode custar $\Theta(n)$. Com $\Phi=$ número de elementos na pilha, uma operação que encontra $p$ elementos tem custo amortizado $a=(1+k)+(p-k+1)-p=2$. Logo $m$ operações custam no máximo $2m=O(m)$.

Pelo método do banqueiro, a mesma conta fica assim: cada operação paga 2 créditos; um paga o empilhamento e o outro fica ``grudado'' no elemento empilhado, para pagar o seu futuro desempilhamento. Na sequência com $k=0,0,0,3,0$, os custos reais são $1,1,1,4,1$, o total é $8$, e os créditos pagos somam $2\cdot5=10$. Na quarta operação, os três créditos guardados nos elementos da pilha pagam os três desempilhamentos.
\end{exemplo}

\begin{cuidado}[Amortizado não é caso médio]
Uma garantia amortizada vale para \emph{toda} sequência de operações, sem hipóteses probabilísticas, mas não limita o custo de uma operação \emph{individual}, que ainda pode ser alto (o oitavo incremento da tabela custou 4). Isso é decisivo quando um sistema exige que cada operação termine num tempo previsível.
\end{cuidado}

% ------------------------------------------------------------------
\section{Heap de Fibonacci}\label{sec:fib}
O heap binomial é ``rigoroso'': cada inserção ou fusão resolve imediatamente as colisões de ordem. O heap de Fibonacci é \textbf{preguiçoso}: adia toda a arrumação até a remoção do máximo, quando ela é feita de uma só vez.

\begin{intuicao}
É como deixar a louça acumular na pia e lavar tudo de uma vez. Cada refeição (inserção) é rápida; a lavagem (consolidação) é demorada, mas só acontece de vez em quando e é paga pelas refeições que sujaram a louça. A remoção do máximo é o momento da lavagem: como é preciso achar o novo máximo de qualquer jeito, aproveita-se para arrumar a estrutura.
\end{intuicao}

\paragraph{Estrutura.} É uma floresta de árvores com a condição de heap. As raízes formam uma \textbf{lista circular duplamente encadeada}, e os filhos de cada nó também. Um ponteiro indica a raiz máxima. Cada nó guarda a chave, um ponteiro para o pai, um para um filho, dois para os irmãos esquerdo e direito, o grau e um bit de \textbf{marca}. As listas circulares duplas permitem inserir ou retirar um nó e concatenar duas listas em $O(1)$.

\begin{figura}[Um heap de Fibonacci. As raízes formam uma lista circular dupla (setas duplas; a última volta para a primeira) e o ponteiro \textsf{max} aponta a maior raiz. Os nós escuros estão marcados (já perderam um filho). Os filhos de cada nó também formam uma lista circular dupla, omitida no desenho.\label{fig:fibestrutura}]
\begin{tikzpicture}[x=1.25cm,y=0.95cm, every node/.append style={font=\scriptsize}]
\node[vtxs] (r1) at (0,0) {23};
\node[vtxs] (r2) at (1.3,0) {7};
\node[vtxs, fill=feito] (r3) at (2.6,0) {48};
\node[vtxs] (r4) at (4.4,0) {17};
\node[vtxs] (r5) at (5.7,0) {24};
\foreach \a/\b in {r1/r2,r2/r3,r3/r4,r4/r5} {\draw[{Stealth[length=1.5mm]}-{Stealth[length=1.5mm]}, gray] (\a) -- (\b);}
\draw[-{Stealth[length=1.5mm]}, gray, rounded corners] (r5.north) -- ++(0,0.45) -| (r1.north);
\node[font=\scriptsize\sffamily, MidnightBlue] (mx) at (2.6,1.25) {max};
\draw[seta, MidnightBlue] (mx) -- (r3);
\node[vtxs, fill=black!70, text=white] (a1) at (2.1,-1) {39};
\node[vtxs] (a2) at (3.1,-1) {41};
\node[vtxs, fill=black!70, text=white] (a3) at (2.1,-2) {18};
\node[vtxs] (a4) at (2.1,-3) {3};
\node[vtxs] (a5) at (3.1,-2) {38};
\draw[aresta] (r3)--(a1) (r3)--(a2) (a1)--(a3) (a3)--(a4) (a2)--(a5);
\node[vtxs] (b1) at (4.0,-1) {16};
\node[vtxs, fill=black!70, text=white] (b2) at (4.8,-1) {12};
\node[vtxs] (b3) at (4.8,-2) {5};
\draw[aresta] (r4)--(b1) (r4)--(b2) (b2)--(b3);
\node[vtxs] (c1) at (5.7,-1) {20};
\draw[aresta] (r5)--(c1);
\end{tikzpicture}
\end{figura}

\paragraph{Operações.} Descrevemos o heap de máximo; ``aumentar a prioridade'' corresponde à operação \emph{decrease-key} de um heap de mínimo.
\begin{itemize}
  \item \textbf{Inserção}: cria uma árvore de um nó, coloca-a na lista de raízes e atualiza o ponteiro do máximo. Custo real $O(1)$.
  \item \textbf{Fusão}: concatena as duas listas de raízes; o máximo é o maior dos dois. Custo real $O(1)$.
  \item \textbf{Seleção do máximo}: leitura do ponteiro, $O(1)$.
  \item \textbf{Remoção do máximo}: (1) retira a raiz máxima e promove todos os seus filhos a raízes, desmarcados; (2) \textbf{consolida}: percorre as raízes com um vetor $A[0..D]$ indexado pelo grau, onde $D=O(\log n)$ é o grau máximo possível, e, sempre que duas raízes têm o mesmo grau, subordina a de menor prioridade à de maior (o grau aumenta de 1) e repete; (3) reconstrói a lista de raízes a partir de $A$ e encontra o novo máximo. O custo real é $O(D+t)$, onde $t$ é o número de raízes; o amortizado é $O(\log n)$.
  \item \textbf{Aumento de prioridade de $x$}: aumenta a chave; se ela não ficou maior que a do pai, termina. Senão, \textbf{corta} $x$, com sua subárvore, e o coloca na lista de raízes, desmarcado. Em seguida, olhando o pai $y$: se $y$ está \emph{desmarcado}, ele é marcado (a menos que seja raiz) e o processo termina; se $y$ já estava \emph{marcado}, isto é, já tinha perdido um filho, $y$ também é cortado, desmarcado e levado à lista de raízes, e o processo continua com o pai de $y$. São os \textbf{cortes em cascata}. Custo amortizado $O(1)$.
  \item \textbf{Remoção de um nó qualquer}: aumenta a prioridade para $+\infty$ e remove o máximo: $O(\log n)$ amortizado.
  \item \textbf{Diminuição de prioridade}: não é a operação barata desta estrutura; faz-se removendo e reinserindo o nó, em $O(\log n)$ amortizado.
\end{itemize}

\begin{exemplo}[Inserções preguiçosas e uma consolidação]\label{ex:consolida}
Num heap de Fibonacci vazio, inserimos $3,8,5,12,6,9,2$. Cada inserção apenas acrescenta um nó à lista de raízes: são 7 árvores de um nó, e o ponteiro do máximo aponta para 12. Agora removemos o máximo. O 12 sai (não tem filhos) e a consolidação percorre as raízes $3,8,5,6,9,2$, usando o vetor $A$ indexado pelo grau:
\begin{itemize}
  \item $3$: $A[0]=3$. \quad $8$: colide com $A[0]=3$; o link põe 8 por cima, e a árvore $8\to3$ tem grau 1: $A[0]$ fica vazio e $A[1]=8\to3$.
  \item $5$: $A[0]=5$. \quad $6$: colide com 5, formando $6\to5$, de grau 1, que colide com $A[1]=8\to3$; o link põe 8 por cima, formando uma árvore de grau 2 em $A[2]$.
  \item $9$: $A[0]=9$. \quad $2$: colide com 9, formando $9\to2$ em $A[1]$.
\end{itemize}
Restam duas raízes, 8 (grau 2) e 9 (grau 1), e o novo máximo é 9 (Figura~\ref{fig:consolida}). A remoção custou bem mais que uma inserção, mas pagou de uma vez a ``preguiça'' das sete inserções.
\begin{figura}[A consolidação do Exemplo~\ref{ex:consolida}: antes, sete raízes; depois de remover o 12 e consolidar, duas árvores de graus distintos.\label{fig:consolida}]
\begin{tabular}{@{}c@{\hspace{14mm}}c@{}}
\begin{tikzpicture}[x=1.25cm]
\foreach \v [count=\i] in {3,8,5,12,6,9,2} {\node[vtx] (q\i) at (\i,0) {\v};}
\node[vtx, fill=ruim] at (q4) {12};
\foreach \i [evaluate=\i as \j using int(\i+1)] in {1,...,6} {\draw[{Stealth[length=1.5mm]}-{Stealth[length=1.5mm]}, gray] (q\i) -- (q\j);}
\node[font=\scriptsize\sffamily, MidnightBlue] (mx) at (4,0.9) {max};
\draw[seta, MidnightBlue] (mx) -- (q4);
\end{tikzpicture}
&
\begin{forest} bst, baseline [8 [6 [5]] [3]] \end{forest}\quad
\begin{forest} bst, baseline [9,feito [2]] \end{forest}\\[3pt]
\small (a) após 7 inserções & \small (b) após remover o máximo
\end{tabular}
\end{figura}
\end{exemplo}

\paragraph{Regras das marcas.} Todos os nós começam desmarcados. Um nó é marcado quando \emph{perde o primeiro filho} por um corte, e, ao perder o segundo, ele próprio é cortado. \textbf{Raízes nunca estão marcadas} e podem perder quantos filhos for preciso; quando um nó é subordinado a outro, fica desmarcado.

\begin{exemplo}[Cortes em cascata]\label{ex:cascata}
O nó interno $X$ tem pai $Y$, marcado, e avô $Z$, desmarcado, que não é raiz. A prioridade de $X$ aumenta e passa a superar a de $Y$.
\begin{enumerate}
  \item $X$ é cortado e vai para a lista de raízes, desmarcado.
  \item $Y$ acaba de perder o segundo filho: é cortado, desmarcado e levado à lista de raízes.
  \item $Z$, que não é raiz e estava desmarcado, perdeu o seu primeiro filho: passa a ser marcado, e a cascata termina.
\end{enumerate}
Ao final, $X$ e $Y$ estão na lista de raízes, desmarcados, e $Z$ continua no seu lugar, agora marcado: se ele perder mais um filho no futuro, será cortado também. A Figura~\ref{fig:cascata} mostra o exemplo com chaves concretas: a de $X$ sobe de 10 para 60.
\begin{figura}[Cortes em cascata. (a) $Y=30$ está marcado (escuro) e $Z=40$, não; a chave de $X$ passa de 10 a 60. (b) $X$ e $Y$ são cortados, um após o outro, e vão para a lista de raízes desmarcados; $Z$ perde um filho e passa a ser marcado. O novo máximo é 60.\label{fig:cascata}]
\forestset{marc/.style={fill=black!70, text=white}}
\begin{tabular}{@{}c@{\hspace{14mm}}c@{}}
\begin{forest} bst, baseline
[50 [45] [40,label={[font=\scriptsize]right:$Z$} [38] [30,marc,label={[font=\scriptsize]right:$Y$} [25] [60,novo,label={[font=\scriptsize]right:$X$}]]]]
\end{forest}
&
\begin{forest} bst, baseline [50 [45] [40,marc,label={[font=\scriptsize]right:$Z$} [38]]] \end{forest}\quad
\begin{forest} bst, baseline [60,novo] \end{forest}\quad
\begin{forest} bst, baseline [30 [25]] \end{forest}\\[3pt]
\small (a) antes dos cortes & \small (b) depois: três raízes, e $Z$ marcado
\end{tabular}
\end{figura}
\end{exemplo}

\subsection{\texorpdfstring{Por que funciona: o grau máximo é $O(\log n)$}{Por que funciona: o grau máximo é O(log n)}}
\begin{lema}
Seja $x$ um nó e sejam $y_1,\dots,y_k$ os seus filhos, na ordem em que foram ligados a $x$. Então $\operatorname{grau}(y_i)\ge i-2$.
\end{lema}
\begin{proof}
Quando $y_i$ foi ligado a $x$, este já tinha os filhos $y_1,\dots,y_{i-1}$, e portanto grau $\ge i-1$. Só se ligam nós de mesmo grau, logo $y_i$ tinha grau $\ge i-1$. Desde então, $y_i$ perdeu no máximo um filho, pois ao perder o segundo seria cortado. Logo $\operatorname{grau}(y_i)\ge i-2$.
\end{proof}
\begin{teo}
Um nó de grau $k$ tem uma subárvore com pelo menos $F_{k+2}\ge\varphi^k$ nós, onde $F$ é a sequência de Fibonacci e $\varphi=\frac{1+\sqrt5}2$. Logo o grau máximo é $D(n)\le\log_\varphi n=O(\log n)$, e depois de uma consolidação o heap tem $O(\log n)$ árvores.
\end{teo}
\begin{proof}[Esboço]
Seja $s_k$ o menor tamanho possível de uma subárvore cuja raiz tem grau $k$. Pelo lema, $s_k\ge 2+\sum_{i=2}^{k}s_{i-2}$, e por indução $s_k\ge F_{k+2}$, pois $F_{k+2}=1+\sum_{i=0}^{k}F_i$. É daí que vem o nome da estrutura: as menores árvores de posto $0,1,2,3,4,5$ têm $1,2,3,5,8,13$ nós.
\end{proof}

\begin{figura}[As menores árvores possíveis de grau $0$ a $4$ num heap de Fibonacci. Cada uma é uma raiz cujos filhos são as menores árvores de graus $0,0,1,2,\dots,k-2$ (o lema). Os tamanhos $1,2,3,5,8$ são números de Fibonacci.\label{fig:fibmin}]
\forestset{bn/.style={for tree={circle, draw, minimum size=2.8mm, inner sep=0pt, s sep=1.6mm, l sep=1.6mm, edge={thick}}}}
\begin{tabular}{@{}c@{\hspace{9mm}}c@{\hspace{9mm}}c@{\hspace{9mm}}c@{\hspace{9mm}}c@{}}
\begin{forest} bn, baseline [] \end{forest} &
\begin{forest} bn, baseline [ [] ] \end{forest} &
\begin{forest} bn, baseline [ [] [] ] \end{forest} &
\begin{forest} bn, baseline [ [] [] [ [] ] ] \end{forest} &
\begin{forest} bn, baseline [ [] [] [ [] ] [ [] [] ] ] \end{forest}\\[2pt]
\small grau 0: 1 nó & \small grau 1: 2 & \small grau 2: 3 & \small grau 3: 5 & \small grau 4: 8
\end{tabular}
\end{figura}

Sem as marcas e os cortes em cascata, seria possível cortar todos os netos de uma raiz e obter uma árvore de posto $k$ com apenas $k+1$ nós. O posto deixaria de ser $O(\log n)$ e a consolidação ficaria cara. Compare: uma $B_4$ do heap binomial tem 16 nós; a menor árvore de grau 4 de um heap de Fibonacci tem 8; sem marcas, teria só 5 (uma raiz e quatro folhas).

\subsection{Análise amortizada}
Use o potencial $\Phi(H)=t(H)+2\,m(H)$, onde $t$ é o número de raízes e $m$ o número de nós marcados.
\begin{itemize}
  \item \textbf{Inserção}: custo real $O(1)$ e $\Delta\Phi=+1$, logo $O(1)$ amortizado. \textbf{Fusão}: $\Delta\Phi=0$, logo $O(1)$.
  \item \textbf{Remoção do máximo}: custo real $O(D+t)$; depois dela há no máximo $D+1$ raízes, então $\Delta\Phi\le(D+1)-t$, e o custo amortizado é $O(D)=O(\log n)$. O trabalho de consolidar as $t$ raízes acumuladas foi pago pelas inserções e fusões preguiçosas que as criaram.
  \item \textbf{Aumento de prioridade} com $c$ cortes em cascata: custo real $O(c)$; $t$ aumenta de $c$ e $m$ diminui de pelo menos $c-2$, então $\Delta\Phi\le c+2(2-c)=4-c$, e o custo amortizado é $O(1)$. Na linguagem do banqueiro, o corte de um nó marcado é pago pelos créditos deixados pela operação que o marcou.
\end{itemize}

\begin{exemplo}[O potencial nos exemplos anteriores]\label{ex:potencial}
No Exemplo~\ref{ex:consolida}, depois das 7 inserções há $t=7$ raízes e nenhuma marca: $\Phi=7$. A remoção do máximo examinou 6 raízes e fez 4 links, mas terminou com $t=2$, e $\Phi$ caiu de 7 para 2: a queda de 5 unidades ``pagou'' quase todo o trabalho. No Exemplo~\ref{ex:cascata}, houve $c=2$ cortes; $t$ aumentou de 2, e o número de marcas não mudou (a de $Y$ sumiu e a de $Z$ apareceu), logo $\Delta\Phi=2+2\cdot0=2$, e o custo amortizado é $2+2=4$, constante. Cada corte a mais na cascata desmarcaria mais um nó, e o potencial cairia 2 para pagá-lo.
\end{exemplo}

\begin{teo}
A partir de um heap de Fibonacci vazio, qualquer sequência de $a_1$ inserções, $a_2$ remoções do máximo, $a_3$ aumentos de prioridade e $a_4$ fusões custa $O(a_1+a_3+a_4+a_2\log n)$.
\end{teo}

\begin{cuidado}
\begin{itemize}
  \item Os custos $O(1)$ do heap de Fibonacci são \textbf{amortizados}: \emph{uma} remoção do máximo depois de muitas inserções e fusões pode custar $\Theta(n)$, ao consolidar tudo de uma vez. Se cada operação precisa ter tempo previsível, essa estrutura é inadequada.
  \item As constantes são grandes e cada nó ocupa bastante memória (quatro ponteiros, grau e marca). Na prática, a estrutura só compensa com um número muito grande de aumentos de prioridade, como nos algoritmos de Dijkstra e Prim em grafos densos.
  \item Uma árvore de um heap de Fibonacci pode ter altura grande, até $\Theta(n)$ em sequências adversárias. O que é logarítmico é o \textbf{grau}, não a altura.
\end{itemize}
\end{cuidado}

% ------------------------------------------------------------------
\section{Comparação e escolha da estrutura}\label{sec:escolhaheap}
\begin{center}
\small
\begin{tabular}{@{}l>{\centering\arraybackslash}p{3.3cm}>{\centering\arraybackslash}p{4.2cm}>{\centering\arraybackslash}p{3.6cm}@{}}
\toprule
& Heap binário & Heap binomial & Heap de Fibonacci\\
\midrule
Seleção & $O(1)$ & $O(\log n)$ ou $O(1)$ com ponteiro & $O(1)$\\
Inserção & $O(\log n)$ & $O(\log n)$ pior; $O(1)$ amortizado & $O(1)$\am\\
Remoção do máximo & $O(\log n)$ & $O(\log n)$ & $O(\log n)$\am\\
Aumentar prioridade & $O(\log n)$ & $O(\log n)$ & $O(1)$\am\\
Fusão & $O(n)$ & $O(\log n)$ \textbf{pior caso} & $O(1)$\am\\
Construção & $O(n)$ & $O(n)$ & $O(n)$\\
Memória & vetor, sem ponteiros & 3 ponteiros $+$ grau por nó & 4 ponteiros $+$ grau $+$ marca por nó\\
Garantia & pior caso & pior caso & amortizada\\
Simplicidade & muito simples & média & complexa\\
\bottomrule
\end{tabular}
\end{center}

\begin{roteiro}[Como escolher a lista de prioridades]
\begin{enumerate}
  \item Não há fusões? Use o \textbf{heap binário}: é o mais simples, usa menos memória e tem boa localidade de acesso.
  \item Há fusões frequentes e \textbf{cada operação precisa ter tempo previsível} (limite de tempo por operação, sistemas de tempo real)? Use o \textbf{heap binomial}, com fusão $O(\log n)$ no pior caso.
  \item Há uma quantidade muito grande de \textbf{aumentos de prioridade} (ou de fusões) e o que importa é o \textbf{custo total} da execução? Use o \textbf{heap de Fibonacci}.
  \item Há uma restrição forte de \textbf{memória}? Use o heap binário; as outras estruturas pagam vários ponteiros por elemento.
\end{enumerate}
\end{roteiro}

\begin{exemplo}[Aplicando o roteiro]\leavevmode
\begin{itemize}
  \item Uma fila de impressão que só recebe trabalhos e despacha o mais urgente: não há fusões, e o heap binário basta.
  \item Um escalonador de tempo real em que filas de processadores são fundidas quando um processador falha, e cada operação tem prazo: heap binomial.
  \item O algoritmo de Dijkstra num grafo denso, com $n$ remoções e até $m\approx n^2$ aumentos de prioridade: o heap de Fibonacci dá $O(m+n\log n)$, contra $O(m\log n)$ do binário (Capítulo~\ref{cap:dijkstra}).
\end{itemize}
\end{exemplo}

\begin{chave}
Heap binomial $=$ árvores $B_k$ de ordens distintas $=$ bits de $n$; fusão $=$ soma binária, $O(\log n)$ no pior caso; inserção $O(1)$ amortizado. Heap de Fibonacci $=$ preguiça $+$ marcas; tudo $O(1)$ amortizado, exceto remoções, $O(\log n)$ amortizado; grau $\le\log_\varphi n$. Amortizado é garantia para a sequência, não para cada operação, e não é caso médio.
\end{chave}

% ----------------------------------------------------------------
\section{Exercícios}
\excurso
\begin{exercicio}{hb-l7}[Lista]<sec:heapbinomial>
Desenvolva um algoritmo $O(n)$ para construir um heap binomial com $n$ elementos.
\end{exercicio}
\begin{exercicio}{hb-l8}[Lista]<sec:heapbinomial>
Um heap binomial tem $k$ árvores. Que alturas elas devem ter para que a inserção seja o pior caso, isto é, faça o número máximo de passos?
\end{exercicio}
\begin{exercicio}{hb-l10}[Lista]<ex:rembin>
$H$ é um heap binomial com $k$ árvores, e a árvore cuja raiz tem a maior prioridade tem altura $h$. Determine as alturas das árvores de $H$ e o valor de $h$ para que a remoção do máximo faça o número máximo de passos.
\end{exercicio}
\begin{exercicio}{hb-l11}[Lista]<ex:fusaobin>
Prove ou dê contraexemplo: se $H_1$ e $H_2$ têm $k_1$ e $k_2$ árvores e sua fusão $H$ tem $k$ árvores, então $k\ge\min\{k_1,k_2\}$.
\end{exercicio}
\begin{exercicio}{hb-l12}[Lista]<ex:fusaobin>
Prove ou dê contraexemplo: nas mesmas condições, $k\le k_1+k_2$.
\end{exercicio}
\begin{exercicio}{hb-l13}[Lista]<sec:fib>
Determine a altura mínima e a máxima de uma árvore de um heap de Fibonacci em função do posto $k$.
\end{exercicio}
\begin{exercicio}{hb-l14}[Lista]<sec:heapbinomial,sec:fib>
Escreva os algoritmos de seleção, inserção, remoção, aumento e diminuição de prioridade, construção e fusão para heaps binomiais e para heaps de Fibonacci.
\end{exercicio}
\begin{exercicio}{hb-r6}[Monitoria]<sec:binomial>
(a) O que caracteriza uma árvore binomial $B_k$? Quantos nós tem? Qual a relação entre o grau da raiz e $k$? (b) Como um heap binomial representa uma coleção de árvores binomiais? (c) Duas árvores $B_2$ de um heap binomial de mínimo têm raízes de chaves 10 e 15: qual é o resultado do \emph{link}, sua ordem, seu número de nós e o grau da raiz? (d) Como achar o elemento de maior prioridade e qual o custo? (e) Qual o custo da união e por que ela não percorre os $n$ elementos?
\end{exercicio}
\begin{exercicio}{hb-r9}[Monitoria]<sec:fib>
(a) Qual a principal diferença estrutural entre heap de Fibonacci e heap binomial? (b) Quais operações são $O(1)$ amortizado? Quanto custa remover o mínimo (\emph{Extract-Min})? (c) Na \emph{Decrease-Key}$(x,k)$ de um heap de mínimo, por que alterar a chave pode violar a propriedade de heap e o que se faz? (d) Para que servem os cortes em cascata, e por que a operação é $O(1)$ amortizado mesmo fazendo vários cortes?
\end{exercicio}
\begin{exercicio}{hb-m1}[Monitoria]<ex:fusaoformas>
Faça a fusão de $H_1=\{B_0,B_2,B_4\}$ com $H_2=\{B_0,B_1\}$ como uma soma binária: quais subordinações acontecem, em que ordem, e qual o resultado em binário?
\end{exercicio}
\begin{exercicio}{hb-m2}[Monitoria]<ex:cascata>
Num heap de Fibonacci de máximo, o nó interno $X$ tem pai $Y$ marcado e avô $Z$ marcado. A prioridade de $X$ aumenta e viola a ordem. Descreva os cortes em cascata, os nós movidos para a lista de raízes e as mudanças nas marcas.
\end{exercicio}
\begin{exercicio}{hb-m3}[Monitoria]<sec:fib>
Por que não podemos simplesmente cortar um nó do pai sem marcar nem cortar em cascata? O que aconteceria com a relação entre a ordem de uma árvore e o seu número de nós?
\end{exercicio}
\begin{exercicio}{hb-t1}[Avaliação]<sec:escolhaheap>
Um aeroporto mantém uma fila de prioridade por setor de controle de tráfego. Fusões de filas (transferências entre setores, unificação de setores) são muito frequentes, com filas de centenas de milhares de aeronaves; alterações de prioridade são raras; e o tempo de \emph{cada} fusão precisa ser previsível, pois as janelas de processamento são rígidas. Qual heap você escolheria? Justifique e explique por que os outros dois são menos adequados.
\end{exercicio}
\begin{exercicio}{hb-t2}[Avaliação]<sec:escolhaheap>
Uma central de atendimento médico mantém cerca de $2\cdot10^6$ chamadas numa fila de prioridade, com $2\cdot10^6$ inserções, $2\cdot10^6$ remoções da chamada mais urgente e poucas alterações de prioridade. Cada chamada ocupa 16 bytes, endereços ocupam 8 bytes e a estrutura pode usar no máximo $40$ MB $=41\,943\,040$ bytes. Qual estrutura você usaria? Considere a complexidade das operações, a necessidade de fusão, a memória por elemento e a simplicidade.
\end{exercicio}
\begin{exercicio}{hb-t3}[Avaliação]<sec:escolhaheap>
Um supercomputador tem cerca de $5\cdot10^6$ tarefas numa fila de prioridade. Num período típico há $5\cdot10^6$ inserções, $5\cdot10^6$ remoções da tarefa mais prioritária, $8\cdot10^7$ aumentos de prioridade (feitos sem retirar e reinserir a tarefa) e poucas uniões de filas. O custo total da execução é o que importa. Qual heap é mais adequado? Por que os outros dois são menos vantajosos?
\end{exercicio}

\excompl
\begin{exercicio}{hb-vetor}[Complementar]<sec:amort>
Um vetor dinâmico dobra de capacidade quando fica cheio, copiando todos os elementos. Mostre, com o método do potencial e $\Phi=2\cdot(\text{número de elementos})-\text{capacidade}$, que cada inserção custa $O(1)$ amortizado.
\end{exercicio}
\begin{exercicio}{hb-contador}[Complementar]<ex:contador>
Um contador binário de $k$ bits passa a aceitar \emph{incremento} e \emph{decremento}. Mostre uma sequência de $m$ operações que custa $\Theta(mk)$. Por que a análise amortizada do contador só com incrementos não se aplica?
\end{exercicio}
\begin{exercicio}{hb-fila}[Complementar]<sec:amort>
Implemente uma fila com duas pilhas (uma de entrada e outra de saída) e mostre que cada operação custa $O(1)$ amortizado.
\end{exercicio}
\begin{exercicio}{hb-forma}[Complementar]<sec:heapbinomial>
Quais árvores binomiais compõem um heap binomial com 100 elementos? E depois de mais uma inserção? E depois de mais quatro inserções a partir de 100? Quantos \emph{links} cada uma dessas inserções faz?
\end{exercicio}
\begin{exercicio}{hb-fibbin}[Complementar]<sec:fib>
Mostre que, se um heap de Fibonacci sofre apenas inserções, fusões e remoções do máximo (sem alterações de prioridade), toda árvore dele é uma árvore binomial (não ordenada) logo após cada consolidação.
\end{exercicio}
\begin{exercicio}{hb-fibcaro}[Complementar]<ex:consolida>
Dê uma sequência de operações, a partir de um heap de Fibonacci vazio, em que \emph{uma única} remoção do máximo tem custo real $\Theta(n)$. Por que isso não contradiz o custo amortizado $O(\log n)$?
\end{exercicio}
\begin{exercicio}{hb-binsim}[Complementar]<ex:fusaobin>
Insira $4,9,2,7,5$, nessa ordem, num heap binomial de máximo inicialmente vazio. Mostre as árvores depois de cada inserção e conte os \emph{links}. Compare o número de links de cada inserção com o número de ``vai um'' ao somar 1 em binário.
\end{exercicio}
\begin{exercicio}{hb-fibsim}[Complementar]<ex:consolida>
Num heap de Fibonacci de máximo vazio, insira $6,2,9,4,8,1,7$ e depois remova o máximo, consolidando as raízes na ordem em que foram inseridas, como no Exemplo~\ref{ex:consolida}. Quais árvores restam? Quantos links são feitos? Quanto vale o potencial $\Phi=t+2m$ antes e depois da remoção?
\end{exercicio}

% ================================================================
\chapter{Conjuntos disjuntos: Union-Find}\label{cap:uf}\label{sec:uf}
% ================================================================
\begin{objetivos}
\begin{itemize}
  \item O problema de manter uma partição sob as operações \MakeSet, \Find{} e \Union.
  \item Implementações por vetor de representantes, por listas encadeadas e por florestas.
  \item As heurísticas de união por rank e de compressão de caminhos, e a análise $O(m\log^*n)$.
  \item Como escolher a implementação de acordo com as operações exigidas.
\end{itemize}
\end{objetivos}

Quando objetos vão sendo agrupados ao longo do tempo --- usuários que formam comunidades, computadores ligados por cabos, vértices ligados pelas arestas escolhidas pelo algoritmo de Kruskal --- e é preciso responder rapidamente se dois deles já estão no mesmo grupo, a estrutura adequada é a de \textbf{conjuntos disjuntos}, também chamada \emph{Union-Find}. Este capítulo mostra que, com duas heurísticas simples, cada operação passa a custar, na prática, tempo constante.

\begin{figura}[Uma rede de computadores que recebe cabos aos poucos. Depois dos cabos $1$--$2$, $3$--$4$, $2$--$4$ e $5$--$6$, há três grupos (componentes): $\{1,2,3,4\}$, $\{5,6\}$ e $\{7\}$. A pergunta ``1 e 3 se comunicam?'' tem resposta sim; ``1 e 6?'', não.\label{fig:ufrede}]
\begin{tikzpicture}[x=1.2cm,y=1cm]
\begin{scope}[on background layer]
\fill[azul, rounded corners=8pt] (-0.45,-1.45) rectangle (1.45,0.45);
\fill[feito, rounded corners=8pt] (2.55,-0.45) rectangle (4.45,0.45);
\fill[destaque, rounded corners=8pt] (5.55,-0.45) rectangle (6.45,0.45);
\end{scope}
\node[vtx] (1) at (0,0) {1}; \node[vtx] (2) at (1,0) {2};
\node[vtx] (3) at (0,-1) {3}; \node[vtx] (4) at (1,-1) {4};
\node[vtx] (5) at (3,0) {5}; \node[vtx] (6) at (4,0) {6};
\node[vtx] (7) at (6,0) {7};
\draw[aresta] (1)--(2) (3)--(4) (2)--(4) (5)--(6);
\end{tikzpicture}
\end{figura}

\section{O problema}
Queremos manter uma partição de um universo de $n$ elementos em conjuntos disjuntos, cada um identificado por um elemento chamado \textbf{representante}, com três operações:
\begin{itemize}
  \item \MakeSet$(x)$ cria o conjunto unitário $\{x\}$;
  \item \Find$(x)$ devolve o representante do conjunto que contém $x$; a pergunta ``$x$ e $y$ estão juntos?'' equivale a testar \Find$(x)=$ \Find$(y)$;
  \item \Union$(x,y)$ une o conjunto de $x$ ao de $y$.
\end{itemize}
Conjuntos \textbf{nunca se separam}: não há uma operação de divisão. Numa sequência típica, há $n$ operações \MakeSet{} seguidas de $m$ operações \Find{} e \Union{} intercaladas; como cada união efetiva diminui o número de conjuntos em um, há no máximo $n-1$ delas.

\begin{exemplo}[A rede da Figura~\ref{fig:ufrede} em operações]\label{ex:ufseq}
Cada computador começa sozinho, e cada cabo é uma união:
\begin{center}\small
\begin{tabular}{@{}ll@{}}
\toprule
operação & partição depois dela\\
\midrule
\MakeSet$(1),\dots,$\MakeSet$(7)$ & $\{1\},\{2\},\{3\},\{4\},\{5\},\{6\},\{7\}$\\
\Union$(1,2)$ & $\{1,2\},\{3\},\{4\},\{5\},\{6\},\{7\}$\\
\Union$(3,4)$ & $\{1,2\},\{3,4\},\{5\},\{6\},\{7\}$\\
\Union$(2,4)$ & $\{1,2,3,4\},\{5\},\{6\},\{7\}$\\
\Union$(5,6)$ & $\{1,2,3,4\},\{5,6\},\{7\}$\\
\bottomrule
\end{tabular}
\end{center}
Agora \Find$(1)$ e \Find$(3)$ devolvem o mesmo representante (qual deles depende da implementação), e \Find$(1)\ne$ \Find$(6)$. Um novo cabo $1$--$3$ não muda nada: \Union$(1,3)$ encontra os dois já no mesmo conjunto.
\end{exemplo}

As aplicações incluem o cálculo de componentes conexas de um grafo que recebe arestas, o algoritmo de Kruskal para árvores geradoras mínimas (Capítulo~\ref{cap:agm}), a equivalência de nomes em compiladores, grupos em redes sociais e problemas de percolação.

\section{Implementações}
\subsection*{Vetor de representantes}
Um vetor $rep[x]$ guarda o representante de cada elemento. \Find{} custa $O(1)$, mas \Union{} precisa trocar o representante de todos os elementos de um dos conjuntos, o que custa $O(n)$ por união.

\begin{exemplo}[Vetor de representantes]
Depois das uniões do Exemplo~\ref{ex:ufseq}, uma possibilidade é $rep=[1,1,1,1,5,5,7]$: os elementos $1,\dots,4$ têm representante 1, os elementos 5 e 6 têm representante 5, e o 7 é o seu próprio representante. Executar agora \Union$(1,6)$ exige percorrer o vetor e trocar por 1 o representante de todos os elementos cujo representante é 5 (as posições 5 e 6, em destaque); em geral, até $n-1$ posições podem mudar numa única união.
\begin{figura}
\small $rep$:\ \vetor{1,1,1,1,5,5,7}\quad$\xrightarrow{\ \Union(1,6)\ }$\quad\vetor[5,6]{1,1,1,1,1,1,7}
\end{figura}
\end{exemplo}

\subsection*{Listas encadeadas com ponteiro para o representante}
Cada conjunto é uma lista encadeada, e um objeto que representa o conjunto guarda a cabeça, a cauda e o tamanho da lista; todo nó tem um ponteiro direto para esse objeto. \MakeSet{} e \Find{} custam $O(1)$, pois \Find{} segue um único ponteiro. \Union{} concatena as listas em $O(1)$, ligando a cauda de uma à cabeça da outra, mas precisa atualizar o ponteiro de todos os nós da lista anexada (Figura~\ref{fig:uflistas}).

\begin{figura}[Conjuntos como listas encadeadas. Cada nó aponta para o próximo (setas pretas) e para o objeto do seu conjunto (setas cinzas), que guarda cabeça, cauda e tamanho. Na união, a lista menor, $\{d,e\}$, é anexada ao fim da maior, e só os seus dois nós mudam de objeto.\label{fig:uflistas}]
\begin{tikzpicture}[x=1cm,y=1cm, lnode/.style={draw, thick, minimum width=6mm, minimum height=5.5mm, font=\small}, obj/.style={draw, thick, rounded corners=2pt, fill=azul, font=\scriptsize, align=center}]
\node[obj] (S1) at (1,1.3) {$S_1$: tam $=3$};
\node[lnode] (a) at (0,0) {$a$}; \node[lnode] (b) at (1,0) {$b$}; \node[lnode] (c) at (2,0) {$c$};
\draw[seta] (a)--(b); \draw[seta] (b)--(c);
\foreach \x in {a,b,c} \draw[-{Stealth[length=1.5mm]}, gray] (\x.north) -- (S1.south);
\node[obj] (S2) at (4.5,1.3) {$S_2$: tam $=2$};
\node[lnode] (d) at (4,0) {$d$}; \node[lnode] (e) at (5,0) {$e$};
\draw[seta] (d)--(e);
\foreach \x in {d,e} \draw[-{Stealth[length=1.5mm]}, gray] (\x.north) -- (S2.south);
\node[font=\small] at (6.3,0.6) {$\xrightarrow{\ \Union\ }$};
\begin{scope}[xshift=7.6cm]
\node[obj] (T) at (2,1.3) {$S_1$: tam $=5$};
\node[lnode] (a) at (0,0) {$a$}; \node[lnode] (b) at (1,0) {$b$}; \node[lnode] (c) at (2,0) {$c$};
\node[lnode, fill=destaque] (d) at (3,0) {$d$}; \node[lnode, fill=destaque] (e) at (4,0) {$e$};
\draw[seta] (a)--(b); \draw[seta] (b)--(c); \draw[seta] (c)--(d); \draw[seta] (d)--(e);
\foreach \x in {a,b,c} \draw[-{Stealth[length=1.5mm]}, gray] (\x.north) -- (T.south);
\foreach \x in {d,e} \draw[-{Stealth[length=1.5mm]}, BrickRed, thick] (\x.north) -- (T.south);
\end{scope}
\end{tikzpicture}
\end{figura}

Sem cuidado, isso custa $\Theta(n)$ por união e $\Theta(n^2)$ numa sequência ruim, como anexar sempre a lista grande a um elemento unitário.

\begin{exemplo}[A sequência ruim]
Com $n=5$ elementos $x_1,\dots,x_5$, faça \Union$(x_2,x_1)$, \Union$(x_3,x_1)$, \Union$(x_4,x_1)$, \Union$(x_5,x_1)$, anexando sempre a lista de $x_1$ (a grande) à do novo elemento. As uniões atualizam $1,2,3,4$ ponteiros: 10 no total, e, em geral, $1+2+\cdots+(n-1)=\Theta(n^2)$. Anexando a menor à maior, cada uma dessas uniões atualizaria um único ponteiro, o do elemento novo: 4 no total.
\end{exemplo}

A solução é a \textbf{união ponderada}: anexa-se sempre a lista \emph{menor} à maior. Cada vez que o ponteiro de um elemento $x$ é atualizado, o conjunto que contém $x$ pelo menos dobra de tamanho; como nenhum conjunto tem mais de $n$ elementos, o ponteiro de $x$ é atualizado no máximo $\log_2n$ vezes. Assim, $n$ operações \MakeSet{} e $m$ outras operações custam $O(m+n\log n)$: \Union{} custa $O(\log n)$ \textbf{amortizado}. Esta é a única implementação que permite \textbf{listar os membros} de um conjunto em tempo $O(k)$, proporcional ao tamanho $k$ do conjunto, simplesmente percorrendo a lista.

\begin{intuicao}[Por que ``menor na maior'' dá $\log n$]
Acompanhe um único elemento $x$. Toda vez que o seu ponteiro muda, ele estava na lista menor, e a lista resultante tem pelo menos o dobro do tamanho da que ele estava: $1\to\ge2\to\ge4\to\ge8\to\cdots$. Depois de $k$ mudanças, $x$ está num conjunto com pelo menos $2^k$ elementos, e como $2^k\le n$, são no máximo $\log_2n$ mudanças. Com $n=1000$, cada elemento muda de ponteiro no máximo 9 vezes.
\end{intuicao}

\subsection*{Florestas}
Cada conjunto é uma árvore em que cada nó aponta para o pai, $pai[x]$, e a raiz é o representante ($pai[r]=r$). \MakeSet{} faz $pai[x]:=x$; \Find{} sobe até a raiz, em tempo proporcional à altura; \Union{} encontra as duas raízes e faz uma apontar para a outra, em $O(1)$ depois de encontradas.

\begin{figura}[Uma floresta e o seu vetor $pai$. As setas apontam do filho para o pai; as raízes (1 e 5) apontam para si mesmas e são os representantes. \Find$(4)$ percorre $4\to3\to1$.\label{fig:ufflor}]
\forestset{uf/.style={for tree={circle, draw, thick, minimum size=1.5em, inner sep=0.5pt, font=\scriptsize, s sep=3mm, l sep=3mm, edge={thick, {Stealth[length=1.6mm]}-}}}}
\begin{forest} uf, baseline [1 [2] [3 [4]]] \end{forest}\quad
\begin{forest} uf, baseline [5 [6]] \end{forest}\quad
\begin{forest} uf, baseline [7] \end{forest}
\hspace{1.2cm}
\raisebox{-0.6cm}{\small $pai$:\ \vetor{1,1,1,3,5,5,7}}
\end{figura}

O problema é que, sem nenhum critério, uma sequência de uniões pode formar uma \textbf{cadeia de altura $n-1$}, e então cada \Find{} custa $\Theta(n)$. Por exemplo, \Union$(1,2)$, \Union$(1,3)$, \Union$(1,4)$, \dots, sempre pondo a raiz atual sob o novo elemento, produz a cadeia $1\to2\to3\to\cdots\to n$, e \Find$(1)$ percorre todos os nós.

\subsection*{União por tamanho ou por rank}
A primeira heurística liga a raiz da árvore \emph{menor} à raiz da maior, medindo o tamanho em número de nós ou por um valor chamado \textbf{rank}. As regras do rank são:
\begin{itemize}
  \item todo nó começa com rank 0;
  \item só o rank de raízes é atualizado: na união de duas raízes de mesmo rank $r$, a nova raiz passa a ter rank $r+1$; com ranks diferentes, a de menor rank vai para baixo e nada muda;
  \item o rank nunca diminui, e, quando um nó deixa de ser raiz, seu rank nunca mais muda;
  \item sem compressão de caminhos, o rank é igual à altura; com compressão, é apenas uma cota superior para ela.
\end{itemize}

\begin{figura}[União por rank. À esquerda, ranks diferentes: a árvore de rank 1 vai para baixo da de rank 2, e a altura não muda. À direita, ranks iguais: a nova raiz ganha um nível, e o seu rank passa de 1 a 2. Os números em vermelho são os ranks.\label{fig:ufrank}]
\forestset{uf/.style={for tree={circle, draw, thick, minimum size=1.5em, inner sep=0.5pt, font=\scriptsize, s sep=3mm, l sep=3mm, edge={thick, {Stealth[length=1.6mm]}-}}}}
\begin{tabular}{@{}c@{\hspace{12mm}}c@{}}
\begin{forest} uf, baseline [$a$,bb=2 [$b$] [$c$ [$d$]]] \end{forest}\ \raisebox{0.3ex}{$\cup$}\
\begin{forest} uf, baseline [$e$,bb=1 [$f$]] \end{forest}
\ $\Rightarrow$\
\begin{forest} uf, baseline [$a$,bb=2 [$b$] [$c$ [$d$]] [$e$,novo [$f$]]] \end{forest}
&
\begin{forest} uf, baseline [$g$,bb=1 [$h$]] \end{forest}\ \raisebox{0.3ex}{$\cup$}\
\begin{forest} uf, baseline [$i$,bb=1 [$j$]] \end{forest}
\ $\Rightarrow$\
\begin{forest} uf, baseline [$g$,bb=2 [$h$] [$i$,novo [$j$]]] \end{forest}
\end{tabular}
\end{figura}

\begin{teo}[Altura logarítmica]\label{teo:ufaltura}
Com união por rank (ou por tamanho), toda árvore de rank $h$ tem pelo menos $2^h$ nós. Logo $h\le\floor{\log_2 n}$, e \Find{} custa $O(\log n)$.
\end{teo}
\begin{proof}
Indução no número de operações. Uma árvore de um nó tem rank 0 e $2^0=1$ nó. O rank de uma árvore só aumenta quando se unem duas árvores de mesmo rank $h$; pela hipótese de indução, cada uma tem pelo menos $2^h$ nós, e a nova árvore tem pelo menos $2\cdot2^h=2^{h+1}$ nós e rank $h+1$. Quando a menor é ligada sob a maior, o rank da maior não muda. Como $2^h\le n$, temos $h\le\log_2n$.
\end{proof}
Por exemplo, com $n=1000$ elementos nenhuma árvore passa de altura $\floor{\log_21000}=9$, por pior que seja a sequência de uniões; sem a heurística, a altura poderia chegar a $999$.

\begin{cor}
(i) O número de nós de rank exatamente $k$ é no máximo $n/2^k$: cada um deles é, ou foi, raiz de uma subárvore com pelo menos $2^k$ nós, e essas subárvores são disjuntas. (ii) Todo rank é no máximo $\floor{\log_2n}$. (iii) $\rank(x)<\rank(pai[x])$ para todo $x$ que não é raiz.
\end{cor}

\subsection*{Compressão de caminhos}
A segunda heurística aproveita cada \Find: depois de encontrar a raiz, faz-se \emph{todo nó do caminho} apontar diretamente para ela. As buscas futuras por esses nós passam a custar $O(1)$. Há variantes que fazem o trabalho numa só passada: na \emph{separação de caminhos}, cada nó do caminho passa a apontar para o avô; na \emph{divisão de caminhos}, isso é feito em nós alternados. Todas têm a mesma complexidade assintótica.

\begin{figura}[\Find$(E)$ na cadeia $E\to D\to C\to B\to A$ com as três variantes. Todas encurtam o caminho; a compressão completa deixa todos os nós a um passo da raiz.\label{fig:ufcompr}]
\forestset{uf/.style={for tree={circle, draw, thick, minimum size=1.5em, inner sep=0.5pt, font=\scriptsize, s sep=3mm, l sep=3mm, edge={thick, {Stealth[length=1.6mm]}-}}}}
\begin{tabular}{@{}c@{\hspace{9mm}}c@{\hspace{9mm}}c@{\hspace{9mm}}c@{}}
\begin{forest} uf, baseline [$A$ [$B$ [$C$ [$D$ [$E$,novo]]]]] \end{forest} &
\begin{forest} uf, baseline [$A$ [$B$] [$C$] [$D$] [$E$,novo]] \end{forest} &
\begin{forest} uf, baseline [$A$ [$B$ [$D$]] [$C$ [$E$,novo]]] \end{forest} &
\begin{forest} uf, baseline [$A$ [$B$] [$C$ [$D$] [$E$,novo]]] \end{forest}\\[3pt]
\small (a) antes & \small (b) compressão & \small (c) separação & \small (d) divisão\\[-1pt]
& \scriptsize todos $\to$ raiz & \scriptsize cada nó $\to$ avô & \scriptsize alternados $\to$ avô
\end{tabular}
\end{figura}

\begin{pseudo}{Floresta com união por rank e compressão de caminhos}
\begin{algorithmic}
\Procedure{Make-Set}{$x$} \State $pai[x]:=x$;\ $rank[x]:=0$ \EndProcedure
\Function{Find}{$x$}
  \If{$pai[x]\ne x$} \State $pai[x]:=$ \Call{Find}{$pai[x]$} \Comment{compressão}
  \EndIf
  \State \Return $pai[x]$
\EndFunction
\Procedure{Union}{$x,y$}
  \State $r:=$ \Call{Find}{$x$};\ $s:=$ \Call{Find}{$y$}
  \If{$r=s$} \Return \EndIf
  \If{$rank[r]<rank[s]$} \State $pai[r]:=s$
  \ElsIf{$rank[r]>rank[s]$} \State $pai[s]:=r$
  \Else \State $pai[s]:=r$;\ $rank[r]:=rank[r]+1$
  \EndIf
\EndProcedure
\end{algorithmic}
\end{pseudo}
O \Find{} recursivo acima comprime o caminho na volta da recursão: a chamada sobre $E$ chama a de $D$, que chama a de $C$, e assim por diante até a raiz $A$; na volta, cada chamada faz o seu nó apontar para o valor devolvido, que é $A$. Uma versão iterativa faz duas passadas: a primeira sobe até encontrar a raiz $y$; a segunda percorre o mesmo caminho fazendo $pai[\cdot]:=y$.

\section{União por rank com compressão: $O(m\log^*n)$}
\begin{intuicao}
O logaritmo iterado $\lgs n$ responde à pergunta: ``quantas vezes preciso apertar a tecla $\log_2$ da calculadora, começando em $n$, até o visor mostrar um número $\le1$?''. Para $n=1000$: $1000\to9{,}97\to3{,}32\to1{,}73\to0{,}79$, quatro vezes. Essa função cresce tão devagar que, para efeitos práticos, é uma constante.
\end{intuicao}

Formalmente, seja $\lgs n$ o número de vezes que é preciso aplicar $\log_2$ a $n$ até chegar a um valor $\le1$. Definindo $F(0)=1$ e $F(k)=2^{F(k-1)}$, temos $F=1,2,4,16,65536,2^{65536},\dots$, e $\lgs n$ é o menor $k$ com $F(k)\ge n$:
\begin{center}\small
\begin{tabular}{@{}lccccccc@{}}
\toprule
$n$ & $1$ & $2$ & $3,4$ & $5,\dots,16$ & $17,\dots,65536$ & $65537,\dots,2^{65536}$\\
\midrule
$\lgs n$ & 0 & 1 & 2 & 3 & 4 & 5\\
\bottomrule
\end{tabular}
\end{center}
Como $2^{65536}$ é muito maior que o número de átomos do universo observável (cerca de $10^{80}\approx2^{266}$), para qualquer $n$ que apareça na prática, $\lgs n\le5$.

Dividimos os ranks em \textbf{grupos}: o grupo $B\ge1$ contém os ranks no intervalo $(F(B-1),F(B)]$, e o grupo 0 contém os ranks $\le F(0)=1$. Assim, o grupo 0 é $\{0,1\}$, o grupo 1 é $\{2\}$, o grupo 2 é $\{3,4\}$, o grupo 3 é $\{5,\dots,16\}$ e o grupo 4 é $\{17,\dots,65536\}$.
\begin{itemize}
  \item \textbf{Propriedade A}: como os ranks são no máximo $\log n$, há no máximo $\lgs n$ grupos.
  \item \textbf{Propriedade B}: para $B\ge1$, o número de nós no grupo $B$ é no máximo $\sum_{r=F(B-1)+1}^{F(B)}\frac{n}{2^r}\le\frac{n}{2^{F(B-1)}}=\frac{n}{F(B)}$; o grupo 0 tem trivialmente no máximo $n=n/F(0)$ nós.
\end{itemize}
Num \Find, distribua o custo de cada nó $x$ do caminho. Se $x$ é a raiz ou filho da raiz, ou se $pai[x]$ está num \emph{grupo diferente} do de $x$, o custo é cobrado do próprio \Find; como os ranks crescem ao longo do caminho, há $O(\lgs n)$ mudanças de grupo por \Find, e o total cobrado dessa forma é $O(m\lgs n)$. Caso contrário, quando $pai[x]$ está no mesmo grupo de $x$, o custo é cobrado do nó $x$. A compressão dá a $x$ um novo pai de rank \emph{estritamente maior}, e um nó do grupo $B$ só pode ser cobrado assim no máximo $F(B)$ vezes antes que seu pai passe a outro grupo, depois do que nunca mais é cobrado. O total cobrado dos nós de um grupo é, pela Propriedade B, no máximo $\frac{n}{F(B)}\cdot F(B)=n$, e, somando sobre os $\lgs n$ grupos, $O(n\lgs n)$. O custo total é $O(m\lgs n)$, para $m\ge n$.

\begin{exemplo}[Onde cada passo de um \Find{} é cobrado]
Suponha um caminho de \Find{} com ranks $0,1,3,4,7,12,20$ (da folha para a raiz; os ranks crescem ao longo do caminho). Os grupos são $0,0,2,2,3,3,4$. Os passos que saem de um nó para um pai de \emph{outro grupo} ($1\to3$, $4\to7$ e $12\to20$), mais os dois últimos nós, são cobrados do \Find: no máximo $\lgs n+2$ por consulta. Os demais ($0\to1$, $3\to4$, $7\to12$) são cobrados dos próprios nós; depois da compressão, todos passam a apontar para a raiz, de rank 20, e esses nós ficam mais perto de ``mudar de grupo''.
\end{exemplo}

\begin{complemento}
A análise acima não é a melhor possível. O limite justo, devido a Tarjan, é $O(m\,\alpha(m,n))$, onde $\alpha$ é a inversa da função de Ackermann, que cresce ainda mais devagar que $\lgs$ e não passa de 4 para qualquer entrada concebível. Os dois limites mostram que a estrutura é ``quase linear''.
\end{complemento}

% ------------------------------------------------------------------
\section{Qual implementação usar?}\label{sec:ufescolha}
\begin{table}[h]
\centering\small
\caption{Implementações de Union-Find e quando usar cada uma.}\label{tab:uf}
\begin{tabularx}{\linewidth}{@{}p{3.3cm}cccX@{}}
\toprule
Implementação & \MakeSet & \Find & \Union & Quando é a melhor escolha\\
\midrule
Vetor de representantes & $O(1)$ & $O(1)$ & $O(n)$ & Quase só consultas e pouquíssimas uniões, com $n$ pequeno.\\
Listas encadeadas com ponteiro ao representante e união ponderada & $O(1)$ & $O(1)$ \textbf{pior caso} & $O(\log n)$\am & As consultas dominam e cada uma precisa custar $O(1)$ \emph{garantido}; é preciso \textbf{listar os membros} de um grupo em $O(k)$; há memória para cabeça, cauda, tamanho e ponteiros.\\
Floresta sem critério & $O(1)$ & $O(n)$ & $O(1)$ $+$ \Find s & Nunca; serve apenas para motivar as heurísticas.\\
Floresta com união por rank ou tamanho, sem compressão & $O(1)$ & $O(\log n)$ & $O(\log n)$ & Quando é preciso \textbf{desfazer uniões}: sem compressão, cada união muda só um ponteiro (e talvez um rank), que pode ser empilhado e restaurado em $O(1)$.\\
Floresta com rank e compressão de caminhos & $O(1)$ & \multicolumn{2}{c}{$O(\lgs n)$\am\ (ou $\alpha$)} & Caso geral, com muitas operações misturadas e só \Find{} e \Union: praticamente $O(1)$ por operação e memória mínima (vetores $pai$ e $rank$).\\
\bottomrule
\end{tabularx}
\end{table}

\begin{roteiro}[Cenários típicos]
\begin{itemize}
  \item \textbf{Listar os membros do grupo de $u$.} A floresta só tem ponteiros do filho para o pai; listar exigiria testar \Find$(i)=$ \Find$(u)$ para todos os $n$ elementos, em $O(n)$. Com \textbf{listas encadeadas}, percorre-se a lista em $O(k)$; ``mesmo grupo?'' custa $O(1)$ e a união custa $O(\log n)$ amortizado, anexando a menor à maior.
  \item \textbf{Conectividade enquanto arestas são destruídas.} Union-Find não desfaz uniões, e a compressão destrói a forma das árvores. A solução é \emph{offline}: leia todas as destruições, construa o estado final (sem as arestas que serão destruídas) e processe as destruições \textbf{de trás para frente}, cada uma como uma \Union; guarde as respostas e apresente-as na ordem original.
  \item \textbf{Operações online com desfazer} (\emph{snapshots}). Use floresta \textbf{sem compressão}, só com união por rank ou tamanho (altura $O(\log n)$ e \Find{} em $O(\log n)$). Cada \Union{} empilha o ponteiro e o rank que alterou; desfazer é desempilhar e restaurar, em $O(1)$.
  \item \textbf{Algoritmo de Kruskal.} Para cada aresta $uv$, em ordem crescente de peso, testa-se \Find$(u)\ne$ \Find$(v)$ e, em caso afirmativo, faz-se \Union$(u,v)$. A floresta com rank e compressão é a escolha natural.
\end{itemize}
\end{roteiro}

\begin{exemplo}[A solução offline, de trás para frente]
Uma rede tem os cabos $1$--$2$, $2$--$3$ e $3$--$4$, que serão destruídos nessa ordem; queremos o número de componentes após cada destruição. Em vez de destruir, começamos do fim: sem nenhum cabo, há 4 componentes. Reinserimos os cabos na ordem inversa: \Union$(3,4)$ recupera o estado logo após a segunda destruição (3 componentes); \Union$(2,3)$, o estado após a primeira (2 componentes); \Union$(1,2)$, o estado inicial (1 componente). Assim, as respostas, na ordem original das destruições, são $2$, $3$ e $4$ componentes, e cada passo custou só uma \Union.
\end{exemplo}

\section{Exemplos resolvidos}\label{sec:ufexemplos}
\begin{exemplo}[Uniões por rank]\label{ex:ufrank}
Partindo de $\{A\},\dots,\{F\}$, executam-se \Union$(A,B)$, \Union$(C,D)$, \Union$(A,C)$ e \Union$(E,F)$, com união por rank e, em caso de empate, a raiz do primeiro argumento como nova raiz.
\begin{itemize}
  \item \Union$(A,B)$: ranks $0=0$, a raiz é $A$ e $\rank(A)=1$.
  \item \Union$(C,D)$: a raiz é $C$ e $\rank(C)=1$.
  \item \Union$(A,C)$: ranks $1=1$, a raiz é $A$ e $\rank(A)=2$; a árvore tem $A$ com os filhos $B$ e $C$, e $C$ com o filho $D$.
  \item \Union$(E,F)$: a raiz é $E$ e $\rank(E)=1$.
\end{itemize}
Ao final há \textbf{dois conjuntos}, com representantes $A$ (de $\{A,B,C,D\}$) e $E$ (de $\{E,F\}$).
\begin{figura}[As quatro uniões do Exemplo~\ref{ex:ufrank}, com os ranks das raízes em vermelho.]
\forestset{uf/.style={for tree={circle, draw, thick, minimum size=1.5em, inner sep=0.5pt, font=\scriptsize, s sep=3mm, l sep=3mm, edge={thick, {Stealth[length=1.6mm]}-}}}}
\begin{tabular}{@{}c@{\hspace{8mm}}c@{\hspace{8mm}}c@{}}
\begin{forest} uf, baseline [$A$,bb=1 [$B$]] \end{forest}\
\begin{forest} uf, baseline [$C$,bb=1 [$D$]] \end{forest}
&
\begin{forest} uf, baseline [$A$,bb=2 [$B$] [$C$,novo [$D$]]] \end{forest}
&
\begin{forest} uf, baseline [$A$,bb=2 [$B$] [$C$ [$D$]]] \end{forest}\
\begin{forest} uf, baseline [$E$,bb=1 [$F$]] \end{forest}\\[3pt]
\small após \Union$(A,B)$ e \Union$(C,D)$ & \small após \Union$(A,C)$ & \small após \Union$(E,F)$
\end{tabular}
\end{figura}
\end{exemplo}
\begin{exemplo}[Compressão numa cadeia]
Na cadeia $pai[B]=A$, $pai[C]=B$, $pai[D]=C$, $pai[E]=D$, o \Find$(E)$ com compressão devolve $A$ e deixa $B,C,D,E$ apontando diretamente para $A$ (Figura~\ref{fig:ufcompr}(b)): a árvore passa a ter altura 1, e as próximas consultas a esses nós custam $O(1)$. Se depois, com $\rank(A)=\rank(D)=1$, for executada \Union$(A,D)$, há empate: a nova raiz é $A$ e $\rank(A)=2$. O rank só aumenta no empate porque, ligando a árvore menor sob a maior, a altura da maior não muda; só com alturas iguais a nova árvore fica um nível mais alta.
\end{exemplo}
\begin{exemplo}[Por que usar as duas heurísticas juntas]
Só com união por rank, \Find{} custa $\Theta(\log n)$ no pior caso, pois há sequências que atingem essa altura. Só com compressão, uma árvore pode ficar alta (uma cadeia) antes de ser comprimida: o custo total de $m$ operações é $O(m\log n)$, e um \Find{} individual pode custar $\Theta(n)$. Juntas, as duas dão $O(m\lgs n)$ no total, ou $O(m\,\alpha)$, praticamente constante por operação.
\end{exemplo}

\begin{chave}
\textbf{Listas encadeadas} dão \Find{} $O(1)$ no pior caso, \Union{} $O(\log n)$ amortizado com ``menor na maior'' e listagem de membros em $O(k)$. \textbf{Florestas} com \textbf{rank e compressão} dão $O(\lgs n)$ (ou $\alpha$) amortizado por operação, com memória mínima; \textbf{sem compressão}, quando é preciso desfazer uniões. Sempre o menor sob o maior; sem isso, listas e árvores degeneram.
\end{chave}

% ----------------------------------------------------------------
\section{Exercícios}
\excurso
\begin{exercicio}{uf-r1}[Monitoria]
(a) Qual o objetivo da estrutura Union-Find? (b) Explique \MakeSet, \Find{} e \Union. (c) Que informação cada elemento armazena para representar os conjuntos? (d) Como a união por rank evita árvores muito altas?
\end{exercicio}
\begin{exercicio}{uf-r2}[Monitoria]
Refaça, sem consultar o texto, os exemplos da §\ref{sec:ufexemplos}: (a) as uniões \Union$(A,B)$, \Union$(C,D)$, \Union$(A,C)$, \Union$(E,F)$ com união por rank; (b) \Find$(E)$ com compressão na cadeia $A\leftarrow B\leftarrow C\leftarrow D\leftarrow E$; (c) \Union$(A,D)$ com $\rank(A)=\rank(D)=1$; (d) por que usar rank e compressão juntos.
\end{exercicio}
\begin{exercicio}{uf-l1}[Lista]
Considere união por rank sem compressão de caminhos, sempre ligando a árvore de menor rank à de maior. Prove que uma árvore com $n$ elementos tem altura no máximo $\floor{\log n}$ e mostre que a cota é justa, com uma sequência de operações que produz altura $\Theta(\log n)$.
\end{exercicio}
\begin{exercicio}{uf-l2}[Lista]
Considere apenas compressão de caminhos, sem união por rank ou tamanho. Construa uma sequência de operações que produza uma árvore alta antes da compressão e discuta o custo $O(m\log n)$ de $m$ operações.
\end{exercicio}
\begin{exercicio}{uf-t4}[Avaliação]
Uma plataforma tem $10^7$ usuários, cada um inicialmente num grupo só seu. Num dia há $10^5$ uniões de grupos e $10^8$ consultas ``$u$ e $v$ estão no mesmo grupo?''; grupos nunca se separam e o histórico das uniões não importa. (a) Que estrutura usar, e como representar cada grupo: como uma árvore em que cada elemento aponta para outro até o representante, ou como uma lista encadeada em que cada elemento aponta diretamente para o representante? (b) Ao unir dois grupos, que estratégia decide quem é anexado a quem, e por que ela importa?
\end{exercicio}
\begin{exercicio}{uf-doc}[Lista]
Explique, para cada problema, por que o Union-Find tradicional (floresta com compressão) falha e qual adaptação resolve: (a) estradas são destruídas uma a uma e, após cada destruição, pergunta-se se o país continua conexo; (b) uma rede recebe operações online de adicionar cabo, consultar conectividade, salvar um \emph{snapshot} e desfazer tudo até o último \emph{snapshot}; (c) uma rede social precisa de \emph{Juntar}, \emph{MesmoGrupo} e \emph{ListarMembros}$(u)$.
\end{exercicio}
\begin{exercicio}{uf-cf}[Lista]
Problema \emph{Asya and Kittens} (Codeforces 1131F): $n$ gatinhos ocupavam $n$ células enfileiradas, um por célula. Em $n-1$ dias, uma divisória entre duas células vizinhas foi removida por dia, juntando os gatinhos $x_i$ e $y_i$, que estavam em células vizinhas de blocos diferentes. Dados os pares, reconstrua uma ordem inicial possível dos gatinhos.
\end{exercicio}

\excompl
\begin{exercicio}{uf-sim}[Complementar]
Com floresta, união por rank (no empate, a raiz do primeiro argumento vence) e compressão de caminhos, execute em $\{1,\dots,8\}$: \Union$(1,2)$, \Union$(3,4)$, \Union$(5,6)$, \Union$(7,8)$, \Union$(1,3)$, \Union$(5,7)$, \Union$(1,5)$, \Find$(8)$. Dê o vetor $pai$ e os ranks antes e depois do \Find.
\end{exercicio}
\begin{exercicio}{uf-ciclo}[Complementar]
Use Union-Find para decidir se um grafo não direcionado tem ciclo. Qual a complexidade?
\end{exercicio}
\begin{exercicio}{uf-rankpai}[Complementar]
Prove que, com união por rank e compressão de caminhos, $\rank(x)<\rank(pai[x])$ para todo $x$ que não é raiz.
\end{exercicio}
\begin{exercicio}{uf-max}[Complementar]
Estenda a estrutura para responder, em tempo $O(\alpha(n))$ amortizado, ``qual o maior elemento do conjunto de $x$?'' e ``quantos elementos tem o conjunto de $x$?''.
\end{exercicio}
\begin{exercicio}{uf-tamanho}[Complementar]
Mostre que a \emph{união por tamanho}, que pendura a árvore com menos nós sob a outra, também garante altura no máximo $\floor{\log_2n}$.
\end{exercicio}
\begin{exercicio}{uf-ingenua}[Complementar]
Na floresta sem critério, em que \Union$(x,y)$ faz a raiz do conjunto de $x$ apontar para a raiz do conjunto de $y$, dê uma sequência de $n-1$ uniões que produz uma árvore de altura $n-1$, e uma sequência de $n$ operações \Find{} que custa $\Theta(n^2)$ no total.
\end{exercicio}
\begin{exercicio}{uf-listas}[Complementar]
Com listas encadeadas e união ponderada, una os elementos $1,\dots,8$ em rodadas: primeiro aos pares ($\{1,2\},\{3,4\},\dots$), depois os pares aos pares, e por fim os dois grupos de 4. Quantas atualizações de ponteiro ao representante são feitas no total? Generalize para $n=2^k$.
\end{exercicio}
\begin{exercicio}{uf-listasprova}[Complementar]
Prove que, com listas encadeadas e união ponderada, qualquer sequência de $n-1$ uniões faz no máximo $n\log_2n$ atualizações de ponteiro.
\end{exercicio}
\begin{exercicio}{uf-semponderada}[Complementar]
Mostre que, com listas encadeadas \emph{sem} união ponderada, existe uma sequência de $n-1$ uniões que faz $\Theta(n^2)$ atualizações de ponteiro.
\end{exercicio}
\begin{exercicio}{uf-logstar}[Complementar]
Calcule $\lgs n$ para $n=2$, $16$, $65536$, $2^{65536}$ e $10^{80}$ (uma estimativa do número de átomos do universo observável).
\end{exercicio}
\begin{exercicio}{uf-comp}[Complementar]
Um grafo tem os vértices $1,\dots,10$ e as arestas $12$, $34$, $23$, $56$, $78$, $89$, $65$, $97$, processadas nessa ordem por Union-Find. Quais arestas provocam uniões efetivas e quais fecham ciclos? Quantas componentes conexas há ao final?
\end{exercicio}
\begin{exercicio}{uf-percol}[Complementar]
Num tabuleiro $n\times n$, casas vão sendo abertas uma a uma. Queremos saber, após cada abertura, se existe um caminho de casas abertas (vizinhas na horizontal ou na vertical) ligando a primeira linha à última. Como resolver com Union-Find usando dois vértices virtuais? Qual o custo total?
\end{exercicio}
\begin{exercicio}{uf-equiv}[Complementar]
Um compilador recebe declarações de equivalência entre variáveis, como $a\equiv b$, $c\equiv d$, $b\equiv c$, e depois consultas ``$x$ e $y$ ocupam a mesma posição de memória?''. Mostre como resolver com Union-Find e diga quais variáveis ficam equivalentes no exemplo, se também houver $e\equiv f$.
\end{exercicio}
\begin{exercicio}{uf-kruskal}[Complementar]
Quantas operações \Find{} e quantas \Union{} o algoritmo de Kruskal faz num grafo conexo com $n$ vértices e $m$ arestas, no pior caso? Qual o custo dessas operações com rank e compressão, comparado ao da ordenação das arestas?
\end{exercicio}
\begin{exercicio}{uf-rankcomp}[Complementar]
Por que a compressão de caminhos não atualiza os ranks? O que o rank passa a representar quando há compressão, e por que ele ainda serve para decidir as uniões?
\end{exercicio}

% ================================================================
\chapter{Árvores binárias de busca e árvores AVL}\label{cap:avl}\label{sec:avl}
% ================================================================
\begin{objetivos}
\begin{itemize}
  \item Árvores binárias de busca: busca, inserção e remoção em tempo proporcional à altura.
  \item Árvores AVL: a condição de balanceamento e a prova de que a altura é logarítmica.
  \item As rotações simples e duplas, a inserção com no máximo uma rotação e a remoção.
  \item Uma visão geral das árvores rubro-negras e critérios para escolher entre heap, ABB e AVL.
\end{itemize}
\end{objetivos}

Heaps são ótimos para encontrar o máximo, mas não sabem procurar uma chave qualquer. Para um \textbf{dicionário ordenado} --- buscar, inserir e remover por chave, listar as chaves em ordem, encontrar sucessor e predecessor --- usamos árvores binárias de busca. Seu custo é proporcional à altura, que pode degenerar até $n$. As árvores AVL impedem a degeneração com uma condição local de balanceamento, restaurada após cada atualização por meio de rotações que custam $O(1)$.

\section{Árvores binárias de busca}
Uma lista sequencial ordenada permite busca binária em $\Theta(\log n)$, mas inserções e remoções custam $\Omega(n)$, pois exigem deslocar elementos; uma lista encadeada insere rapidamente, mas a busca custa $\Omega(n)$. A árvore binária de busca combina busca eficiente com uma estrutura dinâmica.
\begin{defi}
Uma \textbf{árvore binária de busca} (ABB) é uma árvore binária em que, para todo nó $p$, a chave de $p$ é maior ou igual às chaves da subárvore esquerda de $p$ e menor ou igual às da subárvore direita. Equivalentemente, o percurso \emph{em ordem} lista as chaves em ordem crescente.
\end{defi}

\begin{intuicao}
Uma ABB é uma busca binária ``congelada'' numa árvore. A raiz faz o papel do elemento do meio: tudo o que é menor está à esquerda, tudo o que é maior, à direita, e o mesmo vale recursivamente em cada subárvore. Procurar uma chave é responder a uma sequência de perguntas ``é menor ou maior?'', descendo um nível a cada resposta.
\end{intuicao}

\begin{itemize}
  \item \textbf{Busca}: desce a partir da raiz, indo para a esquerda ou para a direita conforme a comparação; custa $\Theta(h)$, onde $h$ é a altura.
  \item \textbf{Inserção}: busca a posição onde a chave deveria estar e a insere como uma nova \textbf{folha}; custa $\Theta(h)$.
  \item \textbf{Remoção} de um nó $p$: se $p$ é folha, basta retirá-lo; se tem um único filho, esse filho ocupa o lugar de $p$; se tem dois filhos, o \textbf{maior elemento da subárvore esquerda} (o predecessor de $p$, encontrado descendo uma vez à esquerda e depois sempre à direita) ocupa o lugar de $p$ e é removido da sua posição original, que tem no máximo um filho. Simetricamente, pode-se usar o menor da subárvore direita. Custa $\Theta(h)$.
\end{itemize}

\begin{exemplo}[Busca e inserção]\label{ex:abbbusca}
Na ABB da Figura~\ref{fig:abbbusca}(a), a busca por 25 compara com 50 (menor: esquerda), com 30 (menor: esquerda), com 20 (maior: direita) e encontra 25: quatro comparações. A busca por 65 compara com 50 (direita), com 70 (esquerda) e com 60 (direita), e encontra um ponteiro nulo: 65 não está na árvore. Inserir 65 é exatamente pendurar uma nova folha nesse ponteiro nulo, como filho direito de 60 (Figura~\ref{fig:abbbusca}(b)). Percorrendo a árvore em ordem (subárvore esquerda, raiz, subárvore direita), obtemos $10,20,25,30,40,50,60,65,70,80$, em ordem crescente.
\begin{figura}[(a) As buscas por 25 (caminho azul, sucesso) e por 65 (caminho vermelho, fracasso). (b) A inserção de 65 pendura uma folha onde a busca fracassou.\label{fig:abbbusca}]
\begin{tabular}{@{}c@{\hspace{14mm}}c@{}}
\begin{forest} bstp
[50,fill=azul [30,fill=azul,acam [20,fill=azul,acam [10] [25,acam,fill=destaque]] [40]] [70,aviol [60,aviol [,vazio] [?,draw=BrickRed,dashed,aviol,text=BrickRed]] [80]]]
\end{forest}
&
\begin{forest} bstp
[50 [30 [20 [10] [25]] [40]] [70 [60 [,vazio] [65,novo]] [80]]]
\end{forest}\\[3pt]
\small (a) buscas por 25 e por 65 & \small (b) depois de inserir 65
\end{tabular}
\end{figura}
\end{exemplo}

\begin{exemplo}[Os três casos da remoção]\label{ex:abbrem}
Partimos da árvore da Figura~\ref{fig:abbbusca}(b).
\begin{itemize}
  \item \textbf{Remover 10} (folha): basta retirá-la.
  \item \textbf{Remover 60} (um filho): o filho 65 sobe e ocupa o lugar de 60.
  \item \textbf{Remover 30} (dois filhos): o predecessor de 30 é o maior da subárvore esquerda, encontrado descendo uma vez à esquerda (20) e depois sempre à direita (25). A chave 25 sobe para o lugar de 30, e o nó onde ela estava, uma folha, é retirado. A ordem continua correta: 25 é maior que tudo o que ficou à esquerda e menor que tudo à direita.
\end{itemize}
\begin{figura}[Os três casos da remoção numa ABB, cada um aplicado à árvore da Figura~\ref{fig:abbbusca}(b).\label{fig:abbrem}]
\begin{tabular}{@{}c@{\hspace{8mm}}c@{\hspace{8mm}}c@{}}
\begin{forest} bstp
[50 [30 [20 [,vazio] [25]] [40]] [70 [60 [,vazio] [65]] [80]]]
\end{forest}
&
\begin{forest} bstp
[50 [30 [20 [10] [25]] [40]] [70 [65,novo] [80]]]
\end{forest}
&
\begin{forest} bstp
[50 [25,novo [20 [10] [,vazio]] [40]] [70 [60 [,vazio] [65]] [80]]]
\end{forest}\\[3pt]
\small (a) sem 10 (folha) & \small (b) sem 60 (um filho) & \small (c) sem 30 (dois filhos)
\end{tabular}
\end{figura}
\end{exemplo}

A altura de uma ABB com $n$ nós está entre $\floor{\log_2n}+1$ e $n$ (medida em níveis). Inserir as chaves em ordem crescente ou decrescente produz uma ``lista'' de altura $n$, e todas as operações passam a custar $\Theta(n)$ (Figura~\ref{fig:abbdeg}).

\begin{figura}[As mesmas cinco chaves inseridas em duas ordens. Em ordem crescente, a ABB degenera numa lista de altura 5; na ordem $30,20,40,10,50$, fica com altura 3.\label{fig:abbdeg}]
\begin{tabular}{@{}c@{\hspace{18mm}}c@{}}
\begin{forest} bstp
[10 [,vazio] [20 [,vazio] [30 [,vazio] [40 [,vazio] [50]]]]]
\end{forest}
&
\begin{forest} bstp
[30 [20 [10] [,vazio]] [40 [,vazio] [50]]]
\end{forest}\\[3pt]
\small inserção de $10,20,30,40,50$ & \small inserção de $30,20,40,10,50$
\end{tabular}
\end{figura}

Árvores completas minimizam a altura, mas mantê-las completas depois de cada inserção exige reorganizações globais de custo $\Omega(n)$. A saída é exigir algo mais fraco: que a árvore seja \textbf{balanceada}, isto é, que tenha altura $O(\log n)$ e que o mesmo valha para cada subárvore, e que o balanceamento possa ser restaurado em tempo $O(\log n)$.

\section{Árvores AVL}
\begin{defi}
Uma árvore binária é \textbf{AVL} se, em todo nó $v$, as alturas das subárvores esquerda e direita diferem de no máximo 1. O \textbf{balanço} de $v$ é
\[
\bal(v)=h_D(v)-h_E(v),
\]
a altura da subárvore direita menos a da esquerda; um balanço $-1$ indica que $v$ ``pesa'' à esquerda. O nó $v$ é \emph{regulado} se $\bal(v)\in\{-1,0,1\}$ e \emph{desregulado} caso contrário. Folhas têm balanço 0, e toda árvore completa é AVL (a recíproca é falsa).
\end{defi}

\begin{exemplo}[Calculando balanços]\label{ex:balancos}
Na árvore da esquerda da Figura~\ref{fig:avlsimnao}, o nó 10 não tem filho esquerdo (altura 0) e tem o filho direito 12 (altura 1): $\bal(10)=1-0=+1$. Na raiz, as duas subárvores têm altura 2, e $\bal(8)=0$. Todos os balanços estão em $\{-1,0,1\}$: é uma AVL, embora não seja completa. Na árvore da direita, a subárvore direita de 10 tem altura 2 (o caminho $12\to14$) e a esquerda é vazia: $\bal(10)=+2$, e a árvore \textbf{não} é AVL, ainda que todos os outros nós, inclusive a raiz, estejam regulados.
\begin{figura}[Balanços em vermelho. À esquerda, uma AVL; à direita, uma ABB que não é AVL por causa do nó 10.\label{fig:avlsimnao}]
\begin{tabular}{@{}c@{\hspace{18mm}}c@{}}
\begin{forest} bstp
[8,bb=0 [4,bb=0 [2,bb=0] [6,bb=0]] [10,bb=+1 [,vazio] [12,bb=0]]]
\end{forest}
&
\begin{forest} bstp
[8,bb=0 [4,bb=-1 [2,bb=-1 [1,bb=0] [,vazio]] [6,bb=0]] [10,desreg,bb=+2 [,vazio] [12,bb=+1 [,vazio] [14,bb=0]]]]
\end{forest}
\end{tabular}
\end{figura}
\end{exemplo}

\begin{teo}\label{teo:avlaltura}
Toda árvore AVL com $n$ nós tem altura $O(\log n)$; mais precisamente, $h<1{,}44\log_2(n+2)$.
\end{teo}
\begin{proof}
Seja $N(h)$ o número mínimo de nós de uma AVL de altura $h$, medida em níveis. Então $N(1)=1$, $N(2)=2$ e, para $h\ge3$, uma AVL mínima de altura $h$ tem a raiz, uma subárvore de altura $h-1$ e a outra com a menor altura permitida, $h-2$, ambas mínimas:
\[
N(h)=1+N(h-1)+N(h-2).
\]
Essas árvores mínimas são chamadas \textbf{árvores de Fibonacci}, e $N(h)=F_{h+2}-1$. Como $N$ é crescente, $N(h)>2N(h-2)$, de onde $N(h)\ge2^{\floor{h/2}}$, ou seja, $h\le2\log_2N(h)+1$. Como qualquer AVL de altura $h$ com $n$ nós tem $n\ge N(h)$, segue que $h=O(\log n)$; o cálculo exato com a fórmula de $F_k$ dá a constante $1{,}44$.
\end{proof}

\begin{figura}[As árvores de Fibonacci: as AVL com o menor número de nós para cada altura. Cada uma é uma raiz com as duas anteriores penduradas.\label{fig:avlfib}]
\forestset{bn/.style={for tree={circle, draw, minimum size=2.8mm, inner sep=0pt, s sep=1.8mm, l sep=1.8mm, edge={thick}}}}
\begin{tabular}{@{}c@{\hspace{9mm}}c@{\hspace{9mm}}c@{\hspace{9mm}}c@{\hspace{9mm}}c@{}}
\begin{forest} bn, baseline [] \end{forest} &
\begin{forest} bn, baseline [ [] [,vazio] ] \end{forest} &
\begin{forest} bn, baseline [ [ [] [,vazio] ] [] ] \end{forest} &
\begin{forest} bn, baseline [ [ [ [] [,vazio] ] [] ] [ [] [,vazio] ] ] \end{forest} &
\begin{forest} bn, baseline [ [ [ [ [] [,vazio] ] [] ] [ [] [,vazio] ] ] [ [ [] [,vazio] ] [] ] ] \end{forest}\\[2pt]
\small $h=1$: 1 nó & \small $h=2$: 2 & \small $h=3$: 4 & \small $h=4$: 7 & \small $h=5$: 12
\end{tabular}
\end{figura}

\begin{exemplo}[Quão alta pode ser uma AVL?]
A recorrência dá $N(1),\dots,N(5)=1,2,4,7,12$, os tamanhos das árvores da Figura~\ref{fig:avlfib}. Como uma AVL de altura 5 tem pelo menos 12 nós, uma AVL com 10 nós tem altura no máximo 4, e essa altura é possível, pois $N(4)=7\le10$; uma árvore completa com 10 nós também tem altura 4. A diferença entre a melhor e a pior AVL é pequena, e continua pequena para $n$ grande: com um milhão de nós, a altura de uma AVL é no máximo 28, contra 20 de uma árvore completa e até $10^6$ de uma ABB degenerada.
\end{exemplo}

Assim, a busca numa AVL custa $O(\log n)$. Resta mostrar que a inserção e a remoção, incluindo o rebalanceamento, também custam $O(\log n)$.

% ------------------------------------------------------------------
\section{Inserção e atualização dos balanços}
A inserção começa como numa ABB: a nova chave vira uma folha (se a chave já existe, nada é feito). Em seguida, percorre-se o caminho de volta, do pai do novo nó em direção à raiz, atualizando os balanços. Suponha que a inserção ocorreu na subárvore \textbf{esquerda} de um nó $p$; o caso direito é simétrico, com os sinais trocados.
\begin{center}
\small
\begin{tabular}{@{}cccl@{}}
\toprule
$\bal(p)$ antes & depois & altura de $T(p)$ & ação\\
\midrule
$+1$ & $0$ & não mudou & \textbf{para}: os ancestrais não mudam\\
$0$ & $-1$ & aumentou 1 & continua com o pai de $p$ (se $p$ é a raiz, para)\\
$-1$ & $-2$ & --- & $p$ está desregulado: \textbf{rotação}, e depois \textbf{para}\\
\bottomrule
\end{tabular}
\end{center}
Uma inserção aumenta alturas em no máximo 1, então o nó desregulado tem balanço $\pm2$. O nó $p$ a rotacionar é o \textbf{primeiro nó desregulado encontrado ao subir da folha para a raiz}, isto é, o mais profundo. Como a rotação devolve à subárvore a altura que ela tinha antes da inserção, nenhum ancestral é afetado: \textbf{uma inserção faz no máximo uma rotação}, simples ou dupla.

\begin{exemplo}[As três linhas da tabela]\label{ex:avlupd}
Na árvore $20^{+1}(10,\ 30^{+1}(\cdot,40))$ (Figura~\ref{fig:avlupd}(a)):
\begin{itemize}
  \item \textbf{Inserir 5}, à esquerda de 10. Em 10, o balanço passa de $0$ a $-1$: a subárvore de 10 ficou mais alta, e subimos. Em 20, a inserção veio da esquerda e o balanço passa de $+1$ a $0$: a altura de $T(20)$ não mudou, e paramos (Figura~\ref{fig:avlupd}(b)).
  \item \textbf{Inserir 50}, à direita de 40. Em 40, $0\to+1$, e subimos; em 30, $+1\to+2$: 30 é o primeiro desregulado. O filho do lado pesado, 40, tem balanço $+1$ (mesmo sinal), e uma rotação simples à esquerda em 30 produz $40(30,50)$, com a altura de antes. Paramos sem tocar em 20 (Figura~\ref{fig:avlupd}(c)).
\end{itemize}
\begin{figura}[Atualização de balanços na inserção (Exemplo~\ref{ex:avlupd}). O nó novo está em destaque.\label{fig:avlupd}]
\begin{tabular}{@{}c@{\hspace{12mm}}c@{\hspace{12mm}}c@{}}
\begin{forest} bstp
[20,bb=+1 [10,bb=0] [30,bb=+1 [,vazio] [40,bb=0]]]
\end{forest}
&
\begin{forest} bstp
[20,bb=0 [10,bb=-1 [5,novo,bb=0] [,vazio]] [30,bb=+1 [,vazio] [40,bb=0]]]
\end{forest}
&
\begin{forest} bstp
[20,bb=0 [10,bb=-1 [5,bb=0] [,vazio]] [40,bb=0 [30,bb=0] [50,novo,bb=0]]]
\end{forest}\\[3pt]
\small (a) inicial & \small (b) após inserir 5 & \small (c) após inserir 50
\end{tabular}
\end{figura}
\end{exemplo}

\section{As rotações}
Seja $p$ o nó desregulado mais profundo. Há quatro casos, determinados pelo sinal do balanço de $p$ e pelo do filho do lado mais pesado.
\begin{center}
\small
\begin{tabularx}{\linewidth}{@{}lcllX@{}}
\toprule
Caso & $\bal(p)$ & filho do lado pesado & sinais & rotação\\
\midrule
esquerda-esquerda & $-2$ & $u$ (esq.), $\bal(u)=-1$ & iguais & \textbf{simples à direita} em $p$\\
esquerda-direita & $-2$ & $u$ (esq.), $\bal(u)=+1$ & opostos & \textbf{dupla à direita}: esquerda em $u$, direita em $p$\\
direita-direita & $+2$ & $z$ (dir.), $\bal(z)=+1$ & iguais & \textbf{simples à esquerda} em $p$\\
direita-esquerda & $+2$ & $z$ (dir.), $\bal(z)=-1$ & opostos & \textbf{dupla à esquerda}: direita em $z$, esquerda em $p$\\
\bottomrule
\end{tabularx}
\end{center}
Uma regra fácil de lembrar: \textbf{balanço positivo pede rotações à esquerda e negativo, à direita; sinais iguais pedem rotação simples e sinais opostos, dupla}. Um filho com balanço 0 só ocorre na remoção, e então se usa a rotação simples.

\begin{exemplo}[Os quatro casos com três chaves]\label{ex:quatrocasos}
A menor situação de cada caso aparece ao inserir 10, 20 e 30 numa árvore vazia, em ordens diferentes (Figura~\ref{fig:quatrocasos}). Nos quatro casos o resultado é o mesmo, $20(10,30)$: a rotação sempre leva para a raiz a chave do meio.
\begin{figura}[Os quatro casos de rotação. Na linha de cima, a árvore logo após a terceira inserção, com o nó desregulado em vermelho; na de baixo, depois da rotação.\label{fig:quatrocasos}]
\begin{tabular}{@{}c@{\hspace{9mm}}c@{\hspace{9mm}}c@{\hspace{9mm}}c@{}}
\begin{forest} bstp [30,desreg,bb=-2 [20,bb=-1 [10,novo] [,vazio]] [,vazio]] \end{forest} &
\begin{forest} bstp [30,desreg,bb=-2 [10,bb=+1 [,vazio] [20,novo]] [,vazio]] \end{forest} &
\begin{forest} bstp [10,desreg,bb=+2 [,vazio] [20,bb=+1 [,vazio] [30,novo]]] \end{forest} &
\begin{forest} bstp [10,desreg,bb=+2 [,vazio] [30,bb=-1 [20,novo] [,vazio]]] \end{forest}\\
\small $30,20,10$: esq.-esq. & \small $30,10,20$: esq.-dir. & \small $10,20,30$: dir.-dir. & \small $10,30,20$: dir.-esq.\\
\small simples à direita & \small dupla à direita & \small simples à esquerda & \small dupla à esquerda\\[4pt]
\begin{forest} bstp [20,bb=0 [10,bb=0] [30,bb=0]] \end{forest} &
\begin{forest} bstp [20,bb=0 [10,bb=0] [30,bb=0]] \end{forest} &
\begin{forest} bstp [20,bb=0 [10,bb=0] [30,bb=0]] \end{forest} &
\begin{forest} bstp [20,bb=0 [10,bb=0] [30,bb=0]] \end{forest}
\end{tabular}
\end{figura}
\end{exemplo}

Em geral, as subárvores penduradas nos nós envolvidos precisam ser redistribuídas. As figuras abaixo mostram as rotações à direita com subárvores genéricas $T_1,\dots,T_4$.

\begin{figura}[Rotação simples à direita (caso esquerda-esquerda), com $h(T_1)=h+1$ e $h(T_2)=h(T_3)=h$. Depois da rotação, $\bal(u)=\bal(p)=0$ e a altura volta a $h+2$.\label{fig:rotsimples}]
\begin{forest} bst
[$p$ [$u$ [$T_1$,tri] [$T_2$,tri]] [$T_3$,tri]]
\end{forest}
\quad\raisebox{3em}{$\xrightarrow{\text{à direita em }p}$}\quad
\begin{forest} bst
[$u$ [$T_1$,tri] [$p$ [$T_2$,tri] [$T_3$,tri]]]
\end{forest}
\end{figura}
\begin{figura}[Rotação dupla à direita (caso esquerda-direita): $v$ sobe para a raiz e fica com $\bal(v)=0$. Se $\bal(v)$ era $-1$, ficam $\bal(u)=0$ e $\bal(p)=+1$; se era $+1$, $\bal(u)=-1$ e $\bal(p)=0$; se era $0$ ($v$ é o nó inserido), ambos ficam 0.\label{fig:rotdupla}]
\begin{forest} bst
[$p$ [$u$ [$T_1$,tri] [$v$ [$T_2$,tri] [$T_3$,tri]]] [$T_4$,tri]]
\end{forest}
\quad\raisebox{3em}{$\xrightarrow{\text{dupla à direita}}$}\quad
\begin{forest} bst
[$v$ [$u$ [$T_1$,tri] [$T_2$,tri]] [$p$ [$T_3$,tri] [$T_4$,tri]]]
\end{forest}
\end{figura}
Os casos direita-direita e direita-esquerda são os espelhos desses: na rotação simples à esquerda, o filho direito $z$ sobe, $p$ passa a ser seu filho esquerdo e o antigo filho esquerdo de $z$ passa a ser filho direito de $p$. \textbf{Rotações preservam a ordem da ABB}, pois o percurso em ordem $T_1,u,T_2,p,T_3$ é o mesmo antes e depois, e custam $O(1)$, porque só alteram alguns ponteiros.

\begin{intuicao}[Por que a rotação dupla?]
No caso esquerda-direita, uma rotação simples à direita em $p$ não resolve: a subárvore pesada, a de $v$, mudaria de lado junto com $u$ e continuaria pendurada no lugar errado, só que agora do lado direito. A primeira rotação (à esquerda em $u$) transforma o ``zigue-zague'' $p\to u\to v$ numa ``linha reta'' $p\to v\to u$, e a segunda (à direita em $p$) é o caso esquerda-esquerda comum. Teste com $30,10,20$: rotacionar só em 30 daria $10(\cdot,30(20,\cdot))$, ainda desregulado.
\end{intuicao}

\section{Pseudocódigo da inserção}
Cada nó tem os campos \texttt{esq}, \texttt{dir}, \texttt{chave} e \texttt{bal}; $\nil$ é o ponteiro nulo e $pt\ptr.campo$ acessa um campo do nó apontado por $pt$. A chamada externa é \textsc{Ins-AVL}$(x,ptraiz)$. A função devolve um valor lógico $h$ que indica se \textbf{a altura da subárvore aumentou}; só nesse caso o pai precisa atualizar o seu balanço.

\begin{pseudo}{Inserção em árvore AVL}
\begin{algorithmic}
\Procedure{Iniciar-nó}{$pt, x$}
  \State alocar$(pt)$;\ $pt\ptr.esq:=\nil$;\ $pt\ptr.dir:=\nil$;\ $pt\ptr.chave:=x$;\ $pt\ptr.bal:=0$
\EndProcedure
\Function{Ins-AVL}{$x$, \textbf{ref} $pt$} \Comment{verdadeiro se a altura de $T(pt)$ aumentou}
  \If{$pt=\nil$} \State \Call{Iniciar-nó}{$pt,x$};\ \Return verdadeiro \EndIf
  \If{$x=pt\ptr.chave$} \State \Return falso \Comment{chave repetida} \EndIf
  \If{$x<pt\ptr.chave$}
    \State $h:=$ \Call{Ins-AVL}{$x, pt\ptr.esq$}
    \If{$h$} \State $(pt,h):=$ \Call{AtualizaBalanço}{$pt$, ``esquerda''} \EndIf
  \Else
    \State $h:=$ \Call{Ins-AVL}{$x, pt\ptr.dir$}
    \If{$h$} \State $(pt,h):=$ \Call{AtualizaBalanço}{$pt$, ``direita''} \EndIf
  \EndIf
  \State \Return $h$
\EndFunction
\Function{AtualizaBalanço}{$pt, a$}
  \State $h:=$ falso
  \If{$a=$ ``esquerda''}
    \If{$pt\ptr.bal=1$} \State $pt\ptr.bal:=0$
    \ElsIf{$pt\ptr.bal=0$} \State $pt\ptr.bal:=-1$;\ $h:=$ verdadeiro \Comment{a altura aumentou}
    \Else \State $pt:=$ \Call{RotaçõesDireita}{$pt$} \Comment{era $-1$ e ficaria $-2$}
    \EndIf
  \Else
    \If{$pt\ptr.bal=-1$} \State $pt\ptr.bal:=0$
    \ElsIf{$pt\ptr.bal=0$} \State $pt\ptr.bal:=1$;\ $h:=$ verdadeiro
    \Else \State $pt:=$ \Call{RotaçõesEsquerda}{$pt$}
    \EndIf
  \EndIf
  \State \Return $(pt,h)$ \Comment{após uma rotação, $pt$ aponta para a nova raiz da subárvore}
\EndFunction
\end{algorithmic}
\end{pseudo}

\begin{pseudo}{Rotações}
\begin{algorithmic}
\Function{RotaçõesDireita}{$pt$}
  \State $ptu:=pt\ptr.esq$
  \If{$ptu\ptr.bal=-1$} \Comment{simples à direita}
    \State $pt\ptr.esq:=ptu\ptr.dir$;\ $ptu\ptr.dir:=pt$;\ $pt\ptr.bal:=0$;\ $ptu\ptr.bal:=0$;\ $pt:=ptu$
  \Else \Comment{dupla à direita}
    \State $ptv:=ptu\ptr.dir$
    \State $ptu\ptr.dir:=ptv\ptr.esq$;\ $ptv\ptr.esq:=ptu$;\ $pt\ptr.esq:=ptv\ptr.dir$;\ $ptv\ptr.dir:=pt$
    \If{$ptv\ptr.bal=-1$} $pt\ptr.bal:=1$ \Else\ $pt\ptr.bal:=0$ \EndIf
    \If{$ptv\ptr.bal=1$} $ptu\ptr.bal:=-1$ \Else\ $ptu\ptr.bal:=0$ \EndIf
    \State $ptv\ptr.bal:=0$;\ $pt:=ptv$
  \EndIf
  \State \Return $pt$
\EndFunction
\Function{RotaçõesEsquerda}{$pt$}
  \State $ptz:=pt\ptr.dir$
  \If{$ptz\ptr.bal=1$} \Comment{simples à esquerda}
    \State $pt\ptr.dir:=ptz\ptr.esq$;\ $ptz\ptr.esq:=pt$;\ $pt\ptr.bal:=0$;\ $ptz\ptr.bal:=0$;\ $pt:=ptz$
  \Else \Comment{dupla à esquerda}
    \State $pty:=ptz\ptr.esq$
    \State $ptz\ptr.esq:=pty\ptr.dir$;\ $pty\ptr.dir:=ptz$;\ $pt\ptr.dir:=pty\ptr.esq$;\ $pty\ptr.esq:=pt$
    \If{$pty\ptr.bal=1$} $pt\ptr.bal:=-1$ \Else\ $pt\ptr.bal:=0$ \EndIf
    \If{$pty\ptr.bal=-1$} $ptz\ptr.bal:=1$ \Else\ $ptz\ptr.bal:=0$ \EndIf
    \State $pty\ptr.bal:=0$;\ $pt:=pty$
  \EndIf
  \State \Return $pt$
\EndFunction
\end{algorithmic}
\end{pseudo}

\begin{cuidado}[Detalhes do pseudocódigo]
\begin{itemize}
  \item \textbf{Para que serve $h$?} Ele indica se a altura da subárvore cresceu. Se não cresceu, nenhum ancestral muda de balanço; por isso \textsc{AtualizaBalanço} só é chamada quando $h$ é verdadeiro.
  \item \textbf{Por que atualizar $pt$ depois da rotação?} A raiz da subárvore muda ($u$, $v$, $z$ ou $y$ sobe). Como $pt$ é passado por referência, a atribuição religa a nova raiz ao pai, ou atualiza a raiz da árvore.
  \item Os três testes de \textsc{AtualizaBalanço} precisam ser \textbf{mutuamente exclusivos} (``senão se''). Escritos como ``se'' consecutivos, um nó com balanço 1 viraria 0 e, em seguida, o teste ``balanço $=0$'' dispararia de novo, deixando o nó com $-1$ e $h$ verdadeiro.
  \item Depois de uma rotação na inserção, $h$ é falso: a altura voltou ao valor antigo.
\end{itemize}
\end{cuidado}

% ------------------------------------------------------------------
\section{Remoção}
A remoção começa como numa ABB (se o nó tem dois filhos, o maior elemento da subárvore esquerda toma o seu lugar) e sobe a partir do nó \emph{fisicamente removido}, atualizando os balanços. Agora a pergunta é se a altura \textbf{diminuiu}. Para uma remoção na subárvore \textbf{direita} de $p$ (o caso esquerdo é simétrico):
\begin{center}
\small
\begin{tabular}{@{}cccl@{}}
\toprule
$\bal(p)$ antes & depois & altura de $T(p)$ & ação\\
\midrule
$+1$ & $0$ & \textbf{diminuiu} & continua com o pai\\
$0$ & $-1$ & não mudou & \textbf{para}\\
$-1$ & $-2$ & --- & rotação à direita (ver abaixo)\\
\bottomrule
\end{tabular}
\end{center}
Na rotação, olha-se o filho esquerdo $u$ de $p$:
\begin{itemize}
  \item se $\bal(u)=-1$, faz-se a rotação simples à direita; ambos ficam com balanço 0 e a altura \textbf{diminuiu}, de modo que o processo continua subindo;
  \item se $\bal(u)=0$, caso que só ocorre na remoção, faz-se a rotação simples à direita; ficam $\bal(p)=-1$ e $\bal(u)=+1$, a altura \textbf{não muda} e o processo para;
  \item se $\bal(u)=+1$, faz-se a rotação dupla à direita; a altura \textbf{diminuiu} e o processo continua subindo.
\end{itemize}
Por isso uma remoção pode exigir uma rotação em \textbf{cada} ancestral: são $O(\log n)$ rotações de custo $O(1)$, e o total continua $O(\log n)$.

\begin{intuicao}[Por que a remoção pode propagar e a inserção não]
Na inserção, o nó desregulado ficou desregulado porque a sua subárvore \emph{cresceu}; a rotação a encolhe de volta à altura original, e os ancestrais nem percebem. Na remoção, a subárvore desregulada \emph{encolheu} de um lado; a rotação conserta o nó, mas em geral a subárvore resultante fica um nível mais baixa do que era antes da remoção, e esse encolhimento pode desregular o ancestral seguinte.
\end{intuicao}

\begin{cuidado}[Reaproveitar as rotações da inserção]
As funções \textsc{RotaçõesDireita} e \textsc{RotaçõesEsquerda} acima testam apenas ``$\bal(u)=-1$'' (ou ``$=1$'') para decidir entre a rotação simples e a dupla. Na remoção, o caso $\bal(u)=0$ cairia na rotação \emph{dupla}, o que está errado: com $\bal(u)=0$ deve-se fazer a rotação \textbf{simples} (o teste passa a ser ``$\le0$''), ajustar $\bal(p)=-1$ e $\bal(u)=+1$ e informar que a altura não diminuiu.
\end{cuidado}

\begin{exemplo}[Remoções pequenas]\label{ex:avlrem}\leavevmode
\begin{itemize}
  \item Em $20^{+1}(10,\ 40^{0}(30,50))$, remover $10$ deixa $\bal(20)=+2$, com o filho $40$ de balanço 0: rotação simples à esquerda em 20, resultando em $40^{-1}(20^{+1}(\cdot,30),\ 50)$, com a altura inalterada.
  \item Em $20^{-1}(10^{+1}(\cdot,15),\ 30)$, remover $30$ deixa $\bal(20)=-2$, com o filho $10$ de balanço $+1$: rotação dupla à direita, resultando em $15(10,20)$.
\end{itemize}
\begin{figura}[As duas remoções do Exemplo~\ref{ex:avlrem}. O nó removido está tracejado; o desregulado, em vermelho.\label{fig:avlrem}]
\forestset{rem/.style={draw=gray, dashed, text=gray}}
\begin{tabular}{@{}c@{\hspace{4mm}}c@{\hspace{4mm}}c@{\hspace{14mm}}c@{\hspace{4mm}}c@{\hspace{4mm}}c@{}}
\begin{forest} bstp [20,desreg,bb=+2 [10,rem] [40,bb=0 [30] [50]]] \end{forest} &
\raisebox{-4ex}{$\Rightarrow$} &
\begin{forest} bstp [40,bb=-1 [20,bb=+1 [,vazio] [30]] [50]] \end{forest} &
\begin{forest} bstp [20,desreg,bb=-2 [10,bb=+1 [,vazio] [15]] [30,rem]] \end{forest} &
\raisebox{-4ex}{$\Rightarrow$} &
\begin{forest} bstp [15,bb=0 [10] [20]] \end{forest}\\[3pt]
\multicolumn{3}{c}{\small filho com balanço 0: simples, altura mantida} & \multicolumn{3}{c}{\small sinais opostos: dupla, altura diminui}
\end{tabular}
\end{figura}
\end{exemplo}

% ------------------------------------------------------------------
\section{Exemplos resolvidos}\label{sec:avlexemplos}
\begin{exemplo}[Inserções com rotação simples e dupla]\label{ex:teste1q5}
Partimos da árvore abaixo (balanços em vermelho) e inserimos $5$, $15$, $35$ e $45$.
\begin{center}
\begin{forest} bst
[50,bb=-1 [30,bb=-1 [20,bb=0 [10,bb=0] [25,bb=0]] [40,bb=0]] [70,bb=0 [60,bb=0] [80,bb=0]]]
\end{forest}
\end{center}
\textbf{Inserir 5} (filho esquerdo de 10). Subindo: $\bal(10)=-1$, $\bal(20)=-1$, $\bal(30)=-2$. O primeiro nó desregulado é \textbf{30}, e o filho esquerdo $20$ tem balanço $-1$ (sinais iguais): caso esquerda-esquerda, \textbf{rotação simples à direita em 30}. A altura da subárvore volta a 3, e 50 continua com $-1$.
\begin{center}
\begin{forest} bst
[50,bb=-1 [20,bb=0 [10,bb=-1 [5,novo,bb=0] [,vazio]] [30,bb=0 [25,bb=0] [40,bb=0]]] [70,bb=0 [60,bb=0] [80,bb=0]]]
\end{forest}
\end{center}
\textbf{Inserir 15} (filho direito de 10): $\bal(10)$ passa de $-1$ a $0$ e a altura de $T(10)$ não muda; o processo para, sem rotação.
\begin{center}
\begin{forest} bst
[50,bb=-1 [20,bb=0 [10,bb=0 [5,bb=0] [15,novo,bb=0]] [30,bb=0 [25,bb=0] [40,bb=0]]] [70,bb=0 [60,bb=0] [80,bb=0]]]
\end{forest}
\end{center}
\textbf{Inserir 35} (filho esquerdo de 40). Subindo: $\bal(40)=-1$, $\bal(30)=+1$, $\bal(20)=+1$, $\bal(50)=-2$. O primeiro nó desregulado é \textbf{50}, e o filho esquerdo $20$ tem balanço $+1$ (sinais opostos): caso esquerda-direita, \textbf{rotação dupla à direita} --- à esquerda em 20 e depois à direita em 50. O nó $v=30$ sobe para a raiz; como $\bal(30)$ era $+1$, ficam $\bal(20)=-1$ e $\bal(50)=0$.
\begin{center}
\begin{forest} bst
[30,bb=0 [20,bb=-1 [10,bb=0 [5,bb=0] [15,bb=0]] [25,bb=0]] [50,bb=0 [40,bb=-1 [35,novo,bb=0] [,vazio]] [70,bb=0 [60,bb=0] [80,bb=0]]]]
\end{forest}
\end{center}
\textbf{Inserir 45} (filho direito de 40): $\bal(40)$ passa de $-1$ a $0$ e a altura de $T(40)$ não muda. A árvore final é:
\begin{center}
\begin{forest} bst
[30,bb=0 [20,bb=-1 [10,bb=0 [5,bb=0] [15,bb=0]] [25,bb=0]] [50,bb=0 [40,bb=0 [35,bb=0] [45,novo,bb=0]] [70,bb=0 [60,bb=0] [80,bb=0]]]]
\end{forest}
\end{center}
Numa ABB comum, inserções em ordem desfavorável (crescente, decrescente ou em zigue-zague) produziriam altura $\Theta(n)$ e operações de custo $\Theta(n)$. A AVL garante altura no máximo $1{,}44\log_2n$ (Teorema~\ref{teo:avlaltura}), qualquer que seja a ordem de chegada, pagando só $O(1)$ rotações por inserção.
\end{exemplo}

\begin{exemplo}[Inserindo $20,40,60,30,25$ numa árvore vazia]\label{ex:avl4020}
Depois de $20$ e $40$, temos $20^{+1}(\cdot,40)$. Inserir $60$ deixa $\bal(20)=+2$ com o filho $40$ de balanço $+1$: rotação simples à esquerda em 20, resultando em $40^{0}(20,60)$. Inserir $30$ dá $40^{-1}(20^{+1}(\cdot,30),\ 60)$. Inserir $25$ (à esquerda de 30) deixa $\bal(30)=-1$ e $\bal(20)=+2$: o primeiro desregulado é \textbf{20}, com filho direito $30$ de balanço $-1$, e a rotação é \textbf{dupla à esquerda} (à direita em 30, à esquerda em 20). A árvore final é $40^{-1}(25^{0}(20^{0},30^{0}),\ 60^{0})$.
\begin{figura}[As inserções do Exemplo~\ref{ex:avl4020}: uma rotação simples e uma dupla, ambas à esquerda.\label{fig:avl4020}]
\begin{tabular}{@{}c@{\hspace{9mm}}c@{\hspace{9mm}}c@{\hspace{9mm}}c@{}}
\begin{forest} bstp [20,desreg,bb=+2 [,vazio] [40,bb=+1 [,vazio] [60,novo]]] \end{forest} &
\begin{forest} bstp [40,bb=-1 [20,bb=+1 [,vazio] [30,novo]] [60]] \end{forest} &
\begin{forest} bstp [40 [20,desreg,bb=+2 [,vazio] [30,bb=-1 [25,novo] [,vazio]]] [60]] \end{forest} &
\begin{forest} bstp [40,bb=-1 [25,bb=0 [20] [30]] [60]] \end{forest}\\[3pt]
\small (a) 60 desregula 20 & \small (b) após a rotação e o 30 & \small (c) 25 desregula 20 & \small (d) após a dupla
\end{tabular}
\end{figura}
\end{exemplo}

\begin{exemplo}[Inserindo $50,40,30,20,10$ e depois $11$]
Inserir $30$ deixa $\bal(50)=-2$ com o filho $40$ de balanço $-1$: rotação simples à direita em 50, resultando em $40(30,50)$. Inserir $20$ dá $40^{-1}(30^{-1}(20,\cdot),\ 50)$. Inserir $10$: subindo, $\bal(20)=-1$ e $\bal(30)=-2$, de modo que o \textbf{primeiro desregulado é 30}, e não 40; faz-se a rotação simples à direita em 30, e a árvore fica $40^{-1}(20^{0}(10^{0},30^{0}),\ 50^{0})$.
\begin{center}
\begin{forest} bst
[40,bb=-1 [20,bb=0 [10,bb=0] [30,bb=0]] [50,bb=0]]
\end{forest}
\end{center}
Inserir $11$ (à direita de 10) deixa $\bal(10)=+1$, $\bal(20)=-1$ e $\bal(40)=-2$: o desregulado é 40, com filho esquerdo $20$ de balanço $-1$, e a rotação é \textbf{simples à direita em 40}, resultando em $20^{0}(10^{+1}(\cdot,11),$\allowbreak$\ 40^{0}(30,50))$.
\begin{center}
\begin{forest} bst
[20,bb=0 [10,bb=+1 [,vazio] [11,novo,bb=0]] [40,bb=0 [30,bb=0] [50,bb=0]]]
\end{forest}
\end{center}
Compare com a Figura~\ref{fig:abbdeg}: numa ABB comum, as chaves de 10 a 50 em ordem decrescente formariam uma lista, como em ordem crescente; aqui formam uma árvore de altura 3.
\end{exemplo}

\begin{cuidado}[Rotacionar no nó certo]
No exemplo anterior, ao inserir 10, é tentador rotacionar em 40, a raiz. Mas a rotação deve ser feita no desregulado \emph{mais profundo}, 30; depois dela, o nó 40 fica com balanço $-1$ e nem chega a se desregular. Rotacionar em 40 produziria a árvore $30(20(10,\cdot),40(\cdot,50))$, e a inserção seguinte de 11 levaria, erradamente, a uma rotação dupla.
\end{cuidado}

% ------------------------------------------------------------------
\section{Heap, ABB ou AVL?}
\begin{roteiro}[Como escolher]
\begin{itemize}
  \item \textbf{Fila de prioridade pura} (sempre atender o elemento de maior prioridade): \textbf{heap}. O máximo sai em $O(1)$, inserir e remover custam $O(\log n)$, e a estrutura é um vetor simples. A AVL também faria tudo em $O(\log n)$, mas não dá o máximo em $O(1)$ sem um ponteiro extra e é mais pesada.
  \item \textbf{Busca por chave arbitrária, remoção de um elemento específico, percurso em ordem, sucessor e predecessor, consultas por intervalo}: \textbf{AVL}. O heap não tem ordem entre irmãos, e buscar uma chave nele custa $O(n)$.
  \item \textbf{ABB simples ou AVL}: com inserções e remoções frequentes e ordem de chegada imprevisível, \textbf{AVL}, pois a ABB pode degenerar para altura $\Theta(n)$.
\end{itemize}
\end{roteiro}

\begin{exemplo}[Aplicando o roteiro]
Uma agenda de compromissos precisa listar os eventos entre duas datas e remover um evento cancelado qualquer: consultas por intervalo e remoções arbitrárias pedem uma \textbf{AVL} indexada pela data. Já o atendimento de chamados por urgência, em que só interessa o mais urgente, pede um \textbf{heap}. Um dicionário que é carregado uma vez a partir de um arquivo ordenado e depois só consultado degeneraria numa ABB simples (as chaves chegam em ordem crescente), mas pode ser um vetor ordenado com busca binária, ou uma AVL.
\end{exemplo}

% ------------------------------------------------------------------
\section{Árvores rubro-negras}\label{sec:rn}
As árvores rubro-negras são outra forma clássica de manter uma ABB balanceada, e são a escolha mais comum em bibliotecas de linguagens de programação.
\begin{defi}
Uma \textbf{árvore rubro-negra} é uma ABB em que cada nó é \emph{vermelho} ou \emph{preto} (um bit extra por nó) e: (1) a raiz é preta; (2) as folhas nulas ($\nil$) são pretas; (3) um nó vermelho tem os dois filhos pretos, isto é, não há dois vermelhos seguidos; (4) todo caminho de um nó até as folhas nulas abaixo dele tem o mesmo número de nós pretos, chamado \emph{altura negra} do nó, $bh$.
\end{defi}

\begin{exemplo}[Conferindo as propriedades]\label{ex:rn}
Na árvore da Figura~\ref{fig:rn}, a raiz 13 é preta; os vermelhos 8, 17, 6, 22 e 27 têm apenas filhos pretos (inclusive as folhas nulas, não desenhadas). Contando os pretos de 13 até as folhas nulas, sem contar o próprio 13 e contando a folha nula, todo caminho encontra 2: por exemplo, $13\to8\to1\to6\to\nil$ passa pelos pretos 1 e $\nil$, e $13\to17\to25\to22\to\nil$ passa por 25 e $\nil$. Logo $bh(13)=2$.
\begin{figura}[Uma árvore rubro-negra com altura negra 2. As folhas nulas, todas pretas, foram omitidas.\label{fig:rn}]
\forestset{rb/.style={fill=black!80, text=white}, rr/.style={fill=BrickRed!85, text=white, draw=BrickRed!60!black}}
\begin{forest} bst
[13,rb [8,rr [1,rb [,vazio] [6,rr]] [11,rb]] [17,rr [15,rb] [25,rb [22,rr] [27,rr]]]]
\end{forest}
\end{figura}
\end{exemplo}

\textbf{Por que é balanceada.} Pela propriedade (4), a subárvore de um nó $x$ tem pelo menos $2^{bh(x)}-1$ nós internos (por indução). Pela propriedade (3), pelo menos metade dos nós de qualquer caminho da raiz até uma folha é preta, então $bh(\text{raiz})\ge h/2$. Logo $n\ge2^{h/2}-1$, isto é, $h\le2\log_2(n+1)$. Dito de outro modo: como todos os caminhos têm o mesmo número de pretos e não há dois vermelhos seguidos, o caminho mais longo da raiz até uma folha nula tem no máximo cerca do dobro de nós do mais curto. Na Figura~\ref{fig:rn}, os caminhos vão de 3 nós ($13\to8\to11$) a 4 nós ($13\to17\to25\to22$).

\textbf{Inserção.} Insere-se como numa ABB e pinta-se o novo nó $z$ de \emph{vermelho}, o que não altera alturas negras; se o pai de $z$ também é vermelho, a propriedade (3) é violada. A correção olha o \emph{tio} $y$ de $z$: (caso 1) se o tio é vermelho, pintam-se o pai e o tio de preto e o avô de vermelho, e o processo continua a partir do avô; (caso 2) se o tio é preto e $z$ está em zigue-zague com o pai e o avô, uma rotação no pai o transforma no caso 3; (caso 3) se o tio é preto e $z$ está ``em linha'', faz-se uma rotação no avô e trocam-se as cores do pai e do avô. Ao final, a raiz é pintada de preto. São no máximo \textbf{duas rotações} por inserção, e $O(\log n)$ recolorações.

\begin{exemplo}[Inserindo $10,20,30$ numa rubro-negra vazia]
O 10 entra como raiz e é pintado de preto. O 20 entra vermelho, à direita de 10: o pai é preto, e nada é violado. O 30 entra vermelho, à direita de 20: pai e filho vermelhos violam a regra (3). O tio de 30 (filho esquerdo de 10) é uma folha nula, preta, e $10\to20\to30$ estão em linha: é o caso 3. A rotação à esquerda no avô 10 leva 20 para a raiz, e a troca de cores deixa 20 preto e 10 vermelho (Figura~\ref{fig:rnins}). É a mesma rotação que a AVL faria (Figura~\ref{fig:quatrocasos}), seguida de uma recoloração.
\begin{figura}[Inserção de 30 na rubro-negra $10(\cdot,20)$: o caso 3 resolve-se com uma rotação e uma troca de cores.\label{fig:rnins}]
\forestset{rb/.style={fill=black!80, text=white}, rr/.style={fill=BrickRed!85, text=white, draw=BrickRed!60!black}}
\begin{tabular}{@{}c@{\hspace{6mm}}c@{\hspace{6mm}}c@{}}
\begin{forest} bstp [10,rb [,vazio] [20,rr [,vazio] [30,rr]]] \end{forest} &
\raisebox{-5ex}{$\xrightarrow{\text{rotação à esquerda em 10, troca de cores}}$} &
\begin{forest} bstp [20,rb [10,rr] [30,rr]] \end{forest}
\end{tabular}
\end{figura}
\end{exemplo}

\textbf{Remoção.} Remove-se como numa ABB; se o nó retirado era preto, surge um ``preto duplo'', resolvido olhando o irmão, em quatro casos, com no máximo \textbf{três rotações}.

\begin{complemento}[AVL ou rubro-negra?]
A AVL é mais rigidamente balanceada ($h\le1{,}44\log n$, contra $2\log n$), o que torna as buscas um pouco mais rápidas. A rubro-negra faz menos rotações nas atualizações --- $O(1)$ rotações na remoção, contra $O(\log n)$ na AVL ---, o que a torna preferida em bibliotecas e sistemas com muitas atualizações.
\end{complemento}

\begin{chave}
AVL: $\bal=h_D-h_E\in\{-1,0,1\}$ e altura no máximo $1{,}44\log n$ (árvores de Fibonacci). Inserção: sobe atualizando os balanços; $\pm1\to0$ para, $0\to\pm1$ continua, $\pm1\to\pm2$ rotaciona no desregulado \textbf{mais profundo} e para (no máximo uma rotação). Sinais iguais: simples; opostos: dupla. Remoção: pode rotacionar em todos os ancestrais; filho com balanço 0 pede rotação simples, e a altura não muda.
\end{chave}

% ----------------------------------------------------------------
\section{Exercícios}
\excurso
\begin{exercicio}{avl-r1}[Monitoria]
(a) Defina árvore binária de busca balanceada. (b) Defina árvore AVL. (c) Qual a relação entre árvores AVL e árvores balanceadas? (d) Defina o fator de balanço de um nó.
\end{exercicio}
\begin{exercicio}{avl-r2}[Monitoria]
Um nó é inserido na subárvore esquerda de $p$. (a) Para $\bal(p)=1$, $0$ e $-1$ antes da inserção, diga o novo balanço, se a altura de $T(p)$ mudou e se é preciso continuar subindo. (b) Se $\bal(p)=-2$ e o filho esquerdo $u$ tem $\bal(u)=-1$, qual rotação deve ser feita? E se $\bal(u)=+1$? Qual a diferença entre os dois casos?
\end{exercicio}
\begin{exercicio}{avl-r4}[Monitoria]
No pseudocódigo da inserção: (a) qual a finalidade da variável lógica $h$? (b) Por que \textsc{AtualizaBalanço} só é chamada quando $h$ é verdadeiro? (c) Por que o ponteiro $pt$ precisa ser atualizado após uma rotação? (d) Um aluno escreveu os três testes de \textsc{AtualizaBalanço} como ``se'' consecutivos: o que dá errado?
\end{exercicio}
\begin{exercicio}{avl-r5}[Monitoria]
Insira (a) $40,20,10,30,35$ e (b) $10,20,30,40,50$ numa AVL vazia, indicando os balanços, as rotações (simples ou dupla, e a direção) e a árvore final. (c) Em seguida, insira 49 na árvore de (b).
\end{exercicio}
\begin{exercicio}{avl-t5}[Avaliação]
Na AVL $50(30(20(10,25),40),\,70(60,80))$, que armazena códigos de produtos, insira $5,15,35,45$, nessa ordem. (a) Após cada inserção, verifique os balanços e desenhe a árvore após cada rotação. (b) Quando a árvore deixar de ser AVL, identifique o primeiro nó desbalanceado encontrado ao subir da folha para a raiz, o caso e a rotação necessária. (c) Por que uma ABB comum poderia ter desempenho muito pior que a AVL se os códigos chegassem numa ordem desfavorável?
\end{exercicio}
\begin{exercicio}{avl-l3}[Lista]
Construa uma AVL em que remover um único nó exija rotações em \emph{todos} os ancestrais; o nó removido deve ter pelo menos três ancestrais.
\end{exercicio}
\begin{exercicio}{avl-l4}[Lista]
Descreva em alto nível uma estratégia para a fusão de duas árvores AVL.
\end{exercicio}
\begin{exercicio}{avl-l5}[Lista]
É possível usar árvores AVL como estrutura de conjuntos disjuntos? Qual seria a complexidade de \Find{} e \Union?
\end{exercicio}
\begin{exercicio}{avl-r6}[Monitoria]
(a) Uma base de registros tem inserções e remoções frequentes e consultas rápidas por chave: ABB simples ou AVL? (b) Uma fila de impressão sempre processa a tarefa de maior prioridade: heap ou AVL? Justifique.
\end{exercicio}
\begin{exercicio}{avl-pseudo}[Slides]
Escreva o pseudocódigo das rotações, da inserção e da remoção em árvores AVL.
\end{exercicio}
\begin{exercicio}{avl-juiz}[Slides]
Resolva os problemas 1195 (\emph{Árvore Binária de Busca}), 1200 e 1201 (\emph{Operações em ABP I e II}) e 1191 (\emph{Recuperação da Árvore}) do Beecrowd, e o problema SID do SPOJ. No 1191, dados os percursos em pré-ordem e em ordem de uma árvore binária, deve-se imprimir o percurso em pós-ordem.
\end{exercicio}
\begin{exercicio}{avl-rn}[Pesquisa]
Apresente a definição de árvore rubro-negra e suas propriedades; justifique por que elas garantem altura $O(\log n)$; explique os algoritmos de inserção e remoção, com rotações e recolorações.
\end{exercicio}

\excompl
\begin{exercicio}{avl-1a7}[Complementar]
Insira $1,2,3,4,5,6,7$ numa AVL vazia. Em quais inserções ocorrem rotações, e de que tipo? Qual a árvore final?
\end{exercicio}
\begin{exercicio}{avl-3rot}[Complementar]
Insira $10,20,30,25,28,27$ numa AVL vazia, indicando cada rotação e a árvore final com os balanços.
\end{exercicio}
\begin{exercicio}{avl-rem}[Complementar]
Na árvore final do Exercício~\ref{ex:avl-1a7}, remova $1$, depois $3$, depois $2$ (na remoção de nó com dois filhos, use o maior da subárvore esquerda). Indique balanços e rotações.
\end{exercicio}
\begin{exercicio}{avl-altmax}[Complementar]
Qual a altura máxima (em níveis) de uma AVL com 20 nós? E com 32 nós? Use $N(h)=1+N(h-1)+N(h-2)$.
\end{exercicio}
\begin{exercicio}{avl-inordem}[Complementar]
Prove que as rotações simples e duplas preservam o percurso em ordem e, portanto, a propriedade de ABB.
\end{exercicio}
\begin{exercicio}{avl-catalan}[Complementar]
Quantas ABBs distintas armazenam as chaves $\{1,2,3\}$? E $\{1,2,3,4\}$? Dê a recorrência.
\end{exercicio}
\begin{exercicio}{avl-rem30}[Complementar]
Na árvore final do Exemplo~\ref{ex:teste1q5}, remova a raiz 30 (substituindo-a pelo maior elemento da subárvore esquerda). Indique os balanços e a rotação necessária.
\end{exercicio}
\begin{exercicio}{avl-seq}[Complementar]
Insira $3,2,1,4,5,6,7,16,15,14$, nessa ordem, numa AVL vazia. Em quais inserções há rotação, de que tipo e em que nó? Qual a árvore final?
\end{exercicio}
\begin{exercicio}{avl-hmin}[Complementar]
Mostre que toda árvore binária com $n$ nós tem altura (em níveis) pelo menos $\ceil{\log_2(n+1)}$. Conclua que a altura de uma AVL é $\Theta(\log n)$.
\end{exercicio}
\begin{exercicio}{avl-verif}[Complementar]
Escreva um algoritmo $O(n)$ que decide se uma árvore binária é uma árvore AVL (isto é, uma ABB que satisfaz a condição de balanceamento).
\end{exercicio}
\begin{exercicio}{avl-suc}[Complementar]
Escreva um algoritmo que encontra o sucessor de um nó numa ABB em tempo $O(h)$, supondo que cada nó tem um ponteiro para o pai.
\end{exercicio}

% ================================================================
\chapter{Ordenação por comparação e limite inferior}\label{cap:sortcomp}\label{sec:ordenacao}
% ================================================================
\begin{objetivos}
\begin{itemize}
  \item Ferramentas para resolver recorrências: árvore de recursão, substituição e Teorema Mestre.
  \item Os algoritmos clássicos de ordenação por comparação: mergesort, quicksort e heapsort.
  \item Estabilidade e ordenação in-place.
  \item A prova de que toda ordenação por comparação faz $\Omega(n\log n)$ comparações.
\end{itemize}
\end{objetivos}

Ordenar é o problema algorítmico mais estudado: aparece sozinho e como passo de quase tudo, da busca binária à eliminação de repetidos e ao algoritmo de Kruskal. Este capítulo apresenta os algoritmos que ordenam \emph{comparando} pares de elementos e prova que nenhum algoritmo desse tipo pode fazer menos de $\Omega(n\log n)$ comparações no pior caso. Como o mergesort e o quicksort são recursivos, começamos pelas ferramentas para analisar algoritmos recursivos.

\section{O problema da ordenação}\label{sec:ordproblema}
Dada uma sequência $S$, queremos uma permutação $s_1,\dots,s_n$ de seus elementos com $s_i\le s_{i+1}$ para todo $i$. Duas propriedades de um algoritmo de ordenação importam além do tempo:
\begin{itemize}
  \item ele é \textbf{estável} se elementos de chaves iguais aparecem na saída na mesma ordem relativa da entrada. A estabilidade é essencial para ordenar por várias chaves (primeiro pela secundária, depois pela principal) e para o radix sort do Capítulo~\ref{cap:sortlin};
  \item ele é \textbf{in-place} se usa memória adicional que não depende de $n$, ou no máximo $O(\log n)$ para a pilha de recursão. O quicksort e o heapsort são in-place; o mergesort não é.
\end{itemize}

\begin{exemplo}[Estabilidade e ordenação por duas chaves]\label{ex:estavel}
Uma turma está em ordem alfabética: (Ana, 8), (Bia, 7), (Caio, 8), (Davi, 7). Ordenando por nota com um algoritmo \textbf{estável}, os empates mantêm a ordem de entrada, e o resultado é (Bia, 7), (Davi, 7), (Ana, 8), (Caio, 8): dentro de cada nota, os nomes continuam em ordem alfabética. Um algoritmo não estável poderia devolver (Davi, 7), (Bia, 7), (Caio, 8), (Ana, 8), que também está ordenado por nota, mas perdeu a ordem dos nomes.

É assim que se ordena por duas chaves: para obter a lista por nota e, nos empates, por nome, ordena-se \emph{primeiro pelo nome} (a chave secundária) e \emph{depois, de forma estável, pela nota} (a principal).
\end{exemplo}

% ------------------------------------------------------------------
\section{Recorrências}\label{sec:recorrencias}
O custo de um algoritmo recursivo é descrito por uma \textbf{recorrência}. Um algoritmo que divide uma entrada de tamanho $n$ em $a$ subproblemas de tamanho $n/b$ e gasta $f(n)$ para dividir e combinar tem custo
\[
T(n)=a\,T(n/b)+f(n).
\]
Por exemplo, o mergesort (§\ref{sec:sortcomp}) resolve $a=2$ subproblemas de tamanho $n/2$ e gasta $f(n)=\Theta(n)$ para intercalar: $T(n)=2T(n/2)+\Theta(n)$. A busca binária resolve $a=1$ subproblema de tamanho $n/2$ e gasta $O(1)$: $T(n)=T(n/2)+O(1)$.

\paragraph{Árvore de recursão.} Desenha-se a árvore das chamadas e soma-se o custo nível a nível. Para $T(n)=2T(n/2)+cn$, há $\log_2n$ níveis, e cada um custa $cn$ no total, pois os subproblemas de um nível somam $n$; logo $T(n)=\Theta(n\log n)$ (Figura~\ref{fig:arvrec}). É o custo do mergesort e de vários algoritmos de divisão e conquista do Capítulo~\ref{cap:dc}.

\begin{figura}[Árvore de recursão de $T(n)=2T(n/2)+cn$. Cada nó mostra o custo da chamada fora das chamadas recursivas; à direita, a soma de cada nível. Há $\log_2 n$ níveis de custo $cn$ e mais $n$ folhas de custo $c$.\label{fig:arvrec}]
\begin{tikzpicture}[x=1cm,y=1.05cm, every node/.style={font=\small}]
\node (r) at (0,0) {$cn$};
\node (a1) at (-2.4,-1) {$c\frac n2$}; \node (a2) at (2.4,-1) {$c\frac n2$};
\foreach \x [count=\i] in {-3.6,-1.2,1.2,3.6} {\node (b\i) at (\x,-2) {$c\frac n4$};}
\draw (r)--(a1) (r)--(a2) (a1)--(b1) (a1)--(b2) (a2)--(b3) (a2)--(b4);
\foreach \i in {1,...,4} {\draw[densely dotted] (b\i) -- ++(-0.35,-0.6); \draw[densely dotted] (b\i) -- ++(0.35,-0.6);}
\node at (0,-3) {$\vdots$};
\foreach \x in {-4.2,-3.6,-3,-2.4,-1.8,-1.2,-0.6,0,0.6,1.2,1.8,2.4,3,3.6,4.2} {\node[font=\scriptsize] at (\x,-3.8) {$c$};}
\foreach \y/\t in {0/{$cn$},-1/{$cn$},-2/{$cn$},-3.8/{$cn$}} {\node[anchor=west, MidnightBlue] at (5.2,\y) {\t};}
\draw[gray] (5.05,0.3) -- (5.05,-4.1);
\draw[decorate, decoration={brace, amplitude=5pt, mirror}, gray] (-4.8,0.2) -- (-4.8,-2.3) node[midway, left=6pt, font=\scriptsize, text=black] {$\log_2 n$ níveis};
\node[anchor=west, MidnightBlue, font=\small\bfseries] at (5.2,-4.6) {total $\Theta(n\log n)$};
\end{tikzpicture}
\end{figura}

\paragraph{Substituição.} Chuta-se a forma da solução, por exemplo $T(n)\le c\,g(n)$, e prova-se o chute por indução forte, ajustando a constante. É assim que se mostra, no Capítulo~\ref{cap:selecao}, que a seleção pela mediana das medianas é linear.

\begin{exemplo}[Substituição: $T(n)=T(n/2)+T(n/4)+n$ é linear]\label{ex:substituicao}
Chutamos $T(n)\le c\,n$ e supomos, por indução, que isso vale para todos os valores menores que $n$. Então
\[
T(n)=T(n/2)+T(n/4)+n\ \le\ c\,\frac n2+c\,\frac n4+n=\Bigl(\frac{3c}4+1\Bigr)n .
\]
Para fechar a indução, precisamos de $\frac{3c}4+1\le c$, isto é, $c\ge4$. Com $c=4$ (e escolhendo a base de modo que $T(n)\le4n$ para $n$ pequeno), temos $T(n)\le4n=O(n)$. A ideia por trás: os dois subproblemas somam $\frac34$ da entrada, menos que o todo, e os custos dos níveis da árvore de recursão formam uma progressão geométrica decrescente, $n+\frac34n+\frac9{16}n+\cdots\le4n$.
\end{exemplo}

\begin{teo}[Teorema Mestre]
Seja $T(n)=a\,T(n/b)+f(n)$, com $a\ge1$ e $b>1$, e seja $E=\log_b a$ o \emph{expoente crítico}.
\begin{enumerate}
  \item Se $f(n)=O(n^{E-\varepsilon})$ para algum $\varepsilon>0$, então $T(n)=\Theta(n^{E})$: as folhas dominam.
  \item Se $f(n)=\Theta(n^{E})$, então $T(n)=\Theta(n^{E}\log n)$: todos os níveis pesam o mesmo.
  \item Se $f(n)=\Omega(n^{E+\varepsilon})$ para algum $\varepsilon>0$ e $a f(n/b)\le c f(n)$ para algum $c<1$, então $T(n)=\Theta(f(n))$: a raiz domina.
\end{enumerate}
\end{teo}
A intuição vem da árvore de recursão: o nível $j$ tem $a^j$ subproblemas de tamanho $n/b^j$, e há $n^E$ folhas. O teorema compara o custo da raiz, $f(n)$, com o custo total das folhas, $n^E$ (Figura~\ref{fig:mestre}).

\begin{figura}[Os três casos do Teorema Mestre vistos pelo custo de cada nível da árvore de recursão (barras, de cima para baixo). No caso 1 os custos crescem e o último nível domina; no caso 2 são iguais e o total é o custo de um nível vezes o número de níveis; no caso 3 decrescem e a raiz domina.\label{fig:mestre}]
\begin{tikzpicture}[x=1cm,y=0.42cm]
\foreach \col/\tit/\lens/\labs in {
  0/{caso 1: $8T(\frac n2)+n^2$}/{0.35,0.7,1.4,2.8}/{$n^2$,$2n^2$,$4n^2$,$8n^2$},
  4.6/{caso 2: $2T(\frac n2)+n$}/{1.4,1.4,1.4,1.4}/{$n$,$n$,$n$,$n$},
  9.2/{caso 3: $2T(\frac n2)+n^2$}/{2.8,1.4,0.7,0.35}/{$n^2$,$\frac{n^2}2$,$\frac{n^2}4$,$\frac{n^2}8$}}{
  \node[font=\small, anchor=west] at (\col,1.2) {\tit};
  \foreach \l [count=\k] in \lens {\fill[azul, draw=MidnightBlue] (\col,-\k*1.25+1.25) rectangle ++(\l,-0.85);}
  \foreach \t [count=\k] in \labs {\node[font=\scriptsize, anchor=west] at (\col+3.0,-\k*1.25+0.85) {\t};}
}
\node[font=\scriptsize, MidnightBlue] at (1.6,-5.2) {$\Rightarrow\Theta(n^3)$};
\node[font=\scriptsize, MidnightBlue] at (6.2,-5.2) {$\Rightarrow\Theta(n\log n)$};
\node[font=\scriptsize, MidnightBlue] at (10.8,-5.2) {$\Rightarrow\Theta(n^2)$};
\end{tikzpicture}
\end{figura}

\begin{roteiro}[Como aplicar o Teorema Mestre]
\begin{enumerate}
  \item Identifique $a$ (número de subproblemas), $b$ (fator de redução do tamanho) e $f(n)$ (custo fora das chamadas).
  \item Calcule o expoente crítico $E=\log_ba$.
  \item Compare $f(n)$ com $n^E$: se $f$ é \emph{polinomialmente} menor, caso 1; se é da mesma ordem, caso 2; se é polinomialmente maior (e satisfaz a condição de regularidade), caso 3.
  \item Se $f$ e $n^E$ diferem só por um fator logarítmico, o teorema, nesta forma, não se aplica: use a árvore de recursão.
\end{enumerate}
\end{roteiro}

\begin{exemplo}[Aplicando o roteiro]\label{ex:mestre}\leavevmode
\begin{itemize}
  \item $T(n)=3T(n/2)+n$ (o algoritmo de Karatsuba para multiplicar inteiros): $a=3$, $b=2$, $E=\log_23\approx1{,}585$. Como $f(n)=n=O(n^{1{,}585-\varepsilon})$ com $\varepsilon=0{,}5$, é o caso 1: $T(n)=\Theta(n^{\log_23})$.
  \item $T(n)=T(2n/3)+1$: $a=1$, $b=3/2$, $E=\log_{3/2}1=0$, e $f(n)=1=\Theta(n^0)$: caso 2, $T(n)=\Theta(\log n)$.
  \item $T(n)=3T(n/4)+n\log n$: $E=\log_43\approx0{,}79$, e $f(n)=n\log n=\Omega(n^{0{,}79+\varepsilon})$ com $\varepsilon=0{,}2$. A regularidade vale: $3\cdot\frac n4\log\frac n4\le\frac34\,n\log n$, com $c=\frac34<1$. Caso 3: $T(n)=\Theta(n\log n)$.
\end{itemize}
\end{exemplo}

\begin{center}
\small
\begin{tabular}{@{}lll@{}}
\toprule
Recorrência & Exemplo & Solução\\
\midrule
$T(n)=T(n/2)+O(1)$ & busca binária & $\Theta(\log n)$\\
$T(n)=2T(n/2)+O(1)$ & máximo por divisão e conquista & $\Theta(n)$\\
$T(n)=2T(n/2)+\Theta(n)$ & mergesort, contagem de inversões & $\Theta(n\log n)$\\
$T(n)=8T(n/2)+\Theta(n^2)$ & produto de matrizes por blocos & $\Theta(n^3)$\\
$T(n)=7T(n/2)+\Theta(n^2)$ & algoritmo de Strassen & $\Theta(n^{\log_2 7})\approx\Theta(n^{2{,}807})$\\
$T(n)=T(n/5)+T(3n/4)+\Theta(n)$ & mediana das medianas & $\Theta(n)$, pois $\tfrac15+\tfrac34<1$\\
$T(n)=T(n-1)+\Theta(n)$ & quicksort no pior caso & $\Theta(n^2)$\\
$T(n)=T(9n/10)+T(n/10)+\Theta(n)$ & quicksort com divisão desbalanceada fixa & $\Theta(n\log n)$\\
\bottomrule
\end{tabular}
\end{center}

\begin{teo}[Divisão e conquista balanceada]
Se um algoritmo divide a instância em $k$ partes de tamanho no máximo $n/c$, com $c>1$ e tamanhos que diferem por uma constante, e gasta tempo linear para dividir e para combinar, então roda em $O(n\log n)$.
\end{teo}
\begin{proof}[Ideia]
Na árvore de recursão, os subproblemas de um mesmo nível são partes disjuntas da entrada, então cada nível custa $O(n)$; como o tamanho cai por um fator $c$ a cada nível, há $O(\log_c n)$ níveis.
\end{proof}

\paragraph{Somatórios úteis.} $\sum_{i=1}^n i=\frac{n(n+1)}2$; $\sum_{k\ge0}x^k=\frac1{1-x}$ para $|x|<1$; $\sum_{k\ge0}\frac{k}{2^k}=2$; $\sum_{k=0}^{\log n}\frac{n}{2^k}\le 2n$; $\sum_{i=1}^n\frac1i=\Theta(\log n)$; $\log(n!)=\sum_{i=1}^n\log i=\Theta(n\log n)$.

% ------------------------------------------------------------------
\section{Ordenação por comparação}\label{sec:sortcomp}
\subsection*{Mergesort}
O mergesort divide o vetor ao meio, ordena recursivamente cada metade e \emph{intercala} as duas metades ordenadas, percorrendo-as com dois índices e copiando sempre o menor dos elementos apontados.

\begin{intuicao}
Intercalar é como juntar dois montes de provas, cada um já em ordem alfabética, num monte só: olhe a prova do topo de cada monte, pegue a que vem primeiro e repita. Cada comparação põe uma prova no lugar definitivo, então juntar montes com $m$ provas no total custa no máximo $m-1$ comparações.
\end{intuicao}

\begin{pseudo}{MergeSort e Intercala}
\begin{algorithmic}
\Procedure{MergeSort}{$V, ini, fim$}
  \If{$ini<fim$}
    \State $meio:=\floor{(ini+fim)/2}$
    \State \Call{MergeSort}{$V,ini,meio$};\ \Call{MergeSort}{$V,meio+1,fim$}
    \State \Call{Intercala}{$V,ini,meio,fim$}
  \EndIf
\EndProcedure
\Procedure{Intercala}{$V,ini,meio,fim$}
  \State $i:=ini$;\ $j:=meio+1$;\ $k:=1$ \Comment{$W$ é um vetor auxiliar}
  \While{$i\le meio$ \textbf{e} $j\le fim$}
    \If{$V[i]\le V[j]$} \State $W[k]:=V[i]$;\ $i:=i+1$ \Comment{``$\le$'' mantém a estabilidade}
    \Else \State $W[k]:=V[j]$;\ $j:=j+1$
    \EndIf
    \State $k:=k+1$
  \EndWhile
  \State copie o que sobrou de $V[i..meio]$ ou de $V[j..fim]$ para o fim de $W$
  \State copie $W[1..fim-ini+1]$ de volta para $V[ini..fim]$
\EndProcedure
\end{algorithmic}
\end{pseudo}

\begin{exemplo}[Mergesort de $(5,2,4,7,1,3,2,6)$]\label{ex:mergesort}
A Figura~\ref{fig:mergesort} mostra a execução: a metade de cima divide o vetor até sobrarem elementos isolados (que já estão ordenados), e a de baixo intercala os pedaços, dois a dois. Na última intercalação, de $(2,4,5,7)$ com $(1,2,3,6)$, as comparações são: $2>1$ (sai 1); $2\le2$ (sai o 2 da esquerda, o que preserva a estabilidade); $4>2$ (sai 2); $4>3$ (sai 3); $4\le6$ (sai 4); $5\le6$ (sai 5); $7>6$ (sai 6); e a metade direita acabou, de modo que o 7 é copiado sem comparação. Foram 7 comparações para 8 elementos.
\begin{figura}[Mergesort de $(5,2,4,7,1,3,2,6)$. As setas cinzas dividem; as azuis intercalam.\label{fig:mergesort}]
\begin{tikzpicture}[x=1mm,y=1mm]
\pic at (0,0) {arr={8}{5,2,4,7,1,3,2,6}};
\pic at (-25,-10) {arr={4}{5,2,4,7}}; \pic at (25,-10) {arr={4}{1,3,2,6}};
\pic at (-37.5,-20) {arr={2}{5,2}}; \pic at (-12.5,-20) {arr={2}{4,7}}; \pic at (12.5,-20) {arr={2}{1,3}}; \pic at (37.5,-20) {arr={2}{2,6}};
\foreach \x/\v in {-43.75/5,-31.25/2,-18.75/4,-6.25/7,6.25/1,18.75/3,31.25/2,43.75/6} {\pic at (\x,-30) {arr={1}{\v}};}
\pic at (-37.5,-40) {arr={2}{2,5}}; \pic at (-12.5,-40) {arr={2}{4,7}}; \pic at (12.5,-40) {arr={2}{1,3}}; \pic at (37.5,-40) {arr={2}{2,6}};
\pic at (-25,-50) {arr={4}{2,4,5,7}}; \pic at (25,-50) {arr={4}{1,2,3,6}};
\pic[fill=feito] at (0,-60) {arr={8}{1,2,2,3,4,5,6,7}};
\begin{scope}[gray, -{Stealth[length=1.4mm]}]
\draw (-6,-3) -- (-20,-7.5); \draw (6,-3) -- (20,-7.5);
\foreach \p/\q in {-25/-37.5,-25/-12.5,25/12.5,25/37.5} {\draw (\p,-13) -- (\q,-17.5);}
\foreach \p/\q in {-37.5/-43.75,-37.5/-31.25,-12.5/-18.75,-12.5/-6.25,12.5/6.25,12.5/18.75,37.5/31.25,37.5/43.75} {\draw (\p,-23) -- (\q,-27.5);}
\end{scope}
\begin{scope}[MidnightBlue!80, -{Stealth[length=1.4mm]}]
\foreach \p/\q in {-43.75/-37.5,-31.25/-37.5,-18.75/-12.5,-6.25/-12.5,6.25/12.5,18.75/12.5,31.25/37.5,43.75/37.5} {\draw (\p,-33) -- (\q,-37.5);}
\foreach \p/\q in {-37.5/-25,-12.5/-25,12.5/25,37.5/25} {\draw (\p,-43) -- (\q,-47.5);}
\draw (-20,-53) -- (-6,-57.5); \draw (20,-53) -- (6,-57.5);
\end{scope}
\node[rot, anchor=west] at (52,-15) {dividir};
\node[rot, anchor=west] at (52,-45) {intercalar};
\end{tikzpicture}
\end{figura}
\end{exemplo}

A intercalação de duas metades que somam $m$ elementos custa $\Theta(m)$, e portanto $T(n)=2T(n/2)+\Theta(n)=\Theta(n\log n)$ em todos os casos. O mergesort é \textbf{estável}, desde que, no empate, o elemento da metade esquerda seja copiado primeiro, e \textbf{não é in-place}, pois o vetor auxiliar ocupa $\Theta(n)$. Existe também uma versão iterativa, de baixo para cima, que intercala blocos de tamanho $1,2,4,\dots$ ao longo do vetor.

\subsection*{Quicksort}
O quicksort faz o trabalho na divisão, e não na combinação: escolhe um \textbf{pivô} $x$, particiona os elementos em $S_1=\{w<x\}$ e $S_2=\{w\ge x\}$, ordena cada parte recursivamente e concatena. Depois da partição, o pivô já está na sua posição definitiva: todos à sua esquerda são menores e todos à sua direita, maiores ou iguais.
\begin{pseudo}{Quicksort com a partição de Lomuto}
\begin{algorithmic}
\Procedure{QuickSort}{$V,ini,fim$}
  \If{$ini<fim$}
    \State $q:=$ \Call{Particiona}{$V,ini,fim$}
    \State \Call{QuickSort}{$V,ini,q-1$};\ \Call{QuickSort}{$V,q+1,fim$}
  \EndIf
\EndProcedure
\Function{Particiona}{$V,ini,fim$}
  \State $x:=V[fim]$;\ $i:=ini-1$ \Comment{o pivô fica no fim; troque-o para lá se escolher outro}
  \For{$j:=ini$ \textbf{até} $fim-1$}
    \If{$V[j]<x$} \State $i:=i+1$;\ $V[i]\Leftrightarrow V[j]$ \EndIf
  \EndFor
  \State $V[i+1]\Leftrightarrow V[fim]$;\ \Return $i+1$
\EndFunction
\end{algorithmic}
\end{pseudo}

\begin{exemplo}[A partição de Lomuto passo a passo]\label{ex:lomuto}
Particionamos $(6,3,9,1,8,2,5)$ com pivô $x=5$, o último elemento. O índice $j$ percorre o vetor; as posições $ini..i$ guardam os elementos já vistos que são menores que o pivô (em verde), e as posições $i+1..j$, os já vistos maiores ou iguais (em amarelo). Quando $V[j]<x$, o índice $i$ avança e $V[j]$ é trocado para dentro da região verde (Figura~\ref{fig:lomuto}). No fim, o pivô é trocado com o primeiro elemento da região amarela, e fica na posição 4, que é devolvida.
\begin{figura}[A partição de Lomuto sobre $(6,3,9,1,8,2,5)$. Verde: menores que o pivô; amarelo: maiores ou iguais; azul: o pivô.\label{fig:lomuto}]
\begin{tikzpicture}[x=7mm,y=6.6mm]
\foreach \lin/\vals/\i/\j/\txt in {
  0/{6,3,9,1,8,2,5}/0/1/{$j=1$: $6\ge5$, nada muda},
  1/{3,6,9,1,8,2,5}/1/2/{$j=2$: $3<5$: $i=1$, troca $V[1]\leftrightarrow V[2]$},
  2/{3,6,9,1,8,2,5}/1/3/{$j=3$: $9\ge5$},
  3/{3,1,9,6,8,2,5}/2/4/{$j=4$: $1<5$: $i=2$, troca $V[2]\leftrightarrow V[4]$},
  4/{3,1,9,6,8,2,5}/2/5/{$j=5$: $8\ge5$},
  5/{3,1,2,6,8,9,5}/3/6/{$j=6$: $2<5$: $i=3$, troca $V[3]\leftrightarrow V[6]$},
  6/{3,1,2,5,8,9,6}/3/7/{fim: troca o pivô com $V[4]$; devolve $q=4$}}{
  \foreach \x [count=\k] in \vals {
    \ifnum\k>\j \else
      \ifnum\k>\i \fill[destaque] (\k-0.5,-\lin-0.3) rectangle (\k+0.5,-\lin+0.3);
      \else \fill[feito] (\k-0.5,-\lin-0.3) rectangle (\k+0.5,-\lin+0.3); \fi
    \fi
    \node[cel] at (\k,-\lin) {\x};}
  \node[rot,anchor=west] at (7.8,-\lin) {\txt};}
\foreach \lin in {0,...,5} {\fill[azul] (6.5,-\lin-0.3) rectangle (7.5,-\lin+0.3); \node[cel] at (7,-\lin) {5};}
\fill[azul] (3.5,-6.3) rectangle (4.5,-5.7); \node[cel] at (4,-6) {5};
\fill[destaque] (6.5,-6.3) rectangle (7.5,-5.7); \node[cel] at (7,-6) {6};
\end{tikzpicture}
\end{figura}
O resultado, $(3,1,2,\,5,\,8,9,6)$, tem o pivô no lugar certo; resta ordenar $(3,1,2)$ e $(8,9,6)$ recursivamente.
\end{exemplo}

\begin{itemize}
  \item \textbf{Pior caso}, $\Theta(n^2)$: quando as partições têm sempre tamanhos $0$ e $n-1$, como ao usar o primeiro elemento como pivô num vetor já ordenado ou invertido. A recorrência é $T(n)=T(n-1)+\Theta(n)$.
  \item \textbf{Melhor caso e caso médio}, $\Theta(n\log n)$. Com um pivô aleatório, em metade das vezes a parte maior tem tamanho no máximo $3n/4$. Se isso acontecesse sempre, no nível $t$ da recursão os subproblemas teriam tamanho no máximo $(3/4)^tn$, e haveria $O(\log n)$ níveis de custo $O(n)$. Como acontece com probabilidade $1/2$, em média o tamanho cai por $3/4$ a cada dois níveis, e o custo esperado continua $O(n\log n)$.
  \item \textbf{Escolha do pivô}: o primeiro elemento é ruim para entradas já ordenadas; boas opções são um elemento aleatório ou a \textbf{mediana de três} (primeiro, meio e último). Usar a \textbf{mediana exata}, calculada em tempo linear (Capítulo~\ref{cap:selecao}), garante $O(n\log n)$ no \emph{pior} caso, mas a constante é alta e na prática não compensa.
  \item O quicksort é \textbf{in-place}: a partição é feita por trocas, e a pilha de recursão tem $O(\log n)$ níveis se a recursão for feita primeiro na parte menor. \textbf{Não é estável}, porque as trocas de longa distância embaralham elementos iguais; pode ser tornado estável com uma partição estável, ao custo de $O(n)$ de memória extra.
\end{itemize}

\begin{figura}[O trabalho do quicksort, nível a nível. Cada barra é um subproblema; a largura é o seu tamanho, e a área total é proporcional ao custo. À esquerda, o pivô é sempre o maior: $n$ níveis e área $\approx n^2/2$. À direita, o pivô divide ao meio: $\log_2n$ níveis de largura $n$ e área $\approx n\log n$.\label{fig:quickcasos}]
\begin{tikzpicture}[x=4.5mm,y=3.6mm]
\foreach \k in {0,...,7} {\fill[ruim, draw=BrickRed!70] (0,-\k) rectangle (8-\k,-\k-0.75);}
\node[rot] at (4,1) {pivô sempre o maior};
\draw[decorate, decoration={brace, amplitude=4pt, mirror}, gray] (-0.4,0) -- (-0.4,-7.75) node[midway, left=5pt, font=\scriptsize, text=black] {$n$ níveis};
\begin{scope}[xshift=6.5cm]
\foreach \k/\m in {0/1,1/2,2/4,3/8} {
  \pgfmathsetmacro{\w}{8/\m}
  \foreach \t in {1,...,\m} {\fill[feito, draw=ForestGreen!70!black] ({(\t-1)*\w+0.08},-\k) rectangle ({\t*\w-0.08},-\k-0.75);}}
\node[rot] at (4,1) {pivô sempre a mediana};
\draw[decorate, decoration={brace, amplitude=4pt}, gray] (8.4,0) -- (8.4,-3.75) node[midway, right=5pt, font=\scriptsize, text=black] {$\log_2 n$ níveis};
\end{scope}
\end{tikzpicture}
\end{figura}

\begin{exemplo}[O quicksort não é estável]\label{ex:quickinst}
Três alunos com a mesma nota, Ana, Bruno e Carla, nessa ordem. Na primeira partição, o pivô é Carla (o último); ninguém é estritamente menor, e a troca final leva Carla para a primeira posição: (Carla, Bruno, Ana). A partição seguinte, sobre (Bruno, Ana), troca os dois. A saída é (Carla, Ana, Bruno), embora a entrada estivesse na ordem (Ana, Bruno, Carla): elementos iguais mudaram de ordem relativa.
\end{exemplo}

\subsection*{Heapsort}
O heapsort (Capítulo~\ref{cap:heap}) constrói um heap em $O(n)$ e remove o máximo $n$ vezes, em $\Theta(n\log n)$ em todos os casos. É in-place e não é estável.

\subsection*{Algoritmos quadráticos}
A \textbf{ordenação por inserção} insere $V[i]$ na parte já ordenada $V[1..i-1]$, deslocando os elementos maiores (veja o Exemplo~\ref{ex:inscasos}). Ela custa $\Theta(n^2)$ no pior caso e no caso médio, mas $\Theta(n)$ se o vetor já está ordenado; é estável e in-place, e é a melhor escolha para vetores pequenos ou quase ordenados. A \textbf{ordenação por bolha} também é $\Theta(n^2)$ e estável.

\begin{center}\small
\begin{tabular}{@{}lccccc@{}}
\toprule
algoritmo & melhor caso & caso médio & pior caso & estável? & in-place?\\
\midrule
inserção & $\Theta(n)$ & $\Theta(n^2)$ & $\Theta(n^2)$ & sim & sim\\
mergesort & $\Theta(n\log n)$ & $\Theta(n\log n)$ & $\Theta(n\log n)$ & sim & não\\
quicksort & $\Theta(n\log n)$ & $\Theta(n\log n)$ & $\Theta(n^2)$ & não & sim\\
heapsort & $\Theta(n\log n)$ & $\Theta(n\log n)$ & $\Theta(n\log n)$ & não & sim\\
\bottomrule
\end{tabular}
\end{center}

% ------------------------------------------------------------------
\section{Limite inferior para ordenação por comparação}\label{sec:limiteinf}
\begin{intuicao}
Ordenar por comparações é como jogar ``adivinhe a permutação'' fazendo só perguntas de sim ou não do tipo ``$a_i\le a_j$?''. Cada resposta, no máximo, divide ao meio o conjunto de permutações ainda possíveis. Começando com $n!$ possibilidades, são necessárias pelo menos $\log_2(n!)$ perguntas para chegar a uma só.
\end{intuicao}

\begin{figura}[Árvore de decisão da ordenação por inserção para $n=3$. Cada nó interno compara dois elementos (``$i{:}j$'' pergunta se $a_i\le a_j$); cada folha é a ordem encontrada. Há $3!=6$ folhas e altura 3: no pior caso, três comparações.\label{fig:arvdecisao}]
\forestset{dec/.style={for tree={draw, rounded corners=2pt, font=\scriptsize, inner sep=2pt, s sep=3mm, l sep=4mm, edge={thick}},
  where n children=0{fill=feito, rounded corners=0pt}{fill=azul}}}
\begin{forest} dec
[1{:}2
  [2{:}3, edge label={node[midway, left, font=\tiny]{$\le$}}
    [{$\langle1,2,3\rangle$}, edge label={node[midway, left, font=\tiny]{$\le$}}]
    [1{:}3, edge label={node[midway, right, font=\tiny]{$>$}}
      [{$\langle1,3,2\rangle$}, edge label={node[midway, left, font=\tiny]{$\le$}}]
      [{$\langle3,1,2\rangle$}, edge label={node[midway, right, font=\tiny]{$>$}}]]]
  [1{:}3, edge label={node[midway, right, font=\tiny]{$>$}}
    [{$\langle2,1,3\rangle$}, edge label={node[midway, left, font=\tiny]{$\le$}}]
    [2{:}3, edge label={node[midway, right, font=\tiny]{$>$}}
      [{$\langle2,3,1\rangle$}, edge label={node[midway, left, font=\tiny]{$\le$}}]
      [{$\langle3,2,1\rangle$}, edge label={node[midway, right, font=\tiny]{$>$}}]]]]
\end{forest}
\end{figura}

\begin{teo}
Todo algoritmo de ordenação baseado em comparações faz $\Omega(n\log n)$ comparações no pior caso.
\end{teo}
\begin{proof}
Para cada $n$, o algoritmo pode ser descrito por uma \textbf{árvore de decisão}: cada nó interno é uma comparação ``$a_i\le a_j$?'', com dois filhos, um para cada resposta; cada folha é um estado final, em que o algoritmo já sabe qual permutação ordena a entrada (Figura~\ref{fig:arvdecisao}). O número de comparações no pior caso é a altura $h$ da árvore. Como qualquer uma das $n!$ permutações pode ser a resposta, há pelo menos $n!$ folhas, e uma árvore binária de altura $h$ tem no máximo $2^h$ folhas. Logo
\[
2^h\ge n!\ge n(n-1)\cdots\bigl(\tfrac n2\bigr)\ge\bigl(\tfrac n2\bigr)^{n/2}
\ \Rightarrow\ h\ge\tfrac n2\log_2\tfrac n2=\tfrac n2\log n-\tfrac n2=\Omega(n\log n).\qedhere
\]
\end{proof}

\begin{exemplo}[O limite para $n=3$ e $n=4$]\label{ex:limite34}
Com $n=3$, há $3!=6$ ordens possíveis; como $2^2=4<6$, duas comparações não bastam, e todo algoritmo precisa de pelo menos $\ceil{\log_26}=3$ comparações no pior caso. A árvore da Figura~\ref{fig:arvdecisao} atinge esse limite. Com $n=4$, há $24$ ordens, e $2^4=16<24\le32=2^5$: são necessárias pelo menos 5 comparações. O mergesort sobre 4 elementos faz no máximo $1+1+3=5$: é ótimo para $n=4$.
\end{exemplo}

\begin{cor}
O heapsort e o mergesort são \textbf{ótimos} entre os algoritmos baseados em comparações, supondo que cada comparação custa $O(1)$. Para ordenar em $o(n\log n)$ é preciso usar a estrutura das chaves --- inteiros pequenos, dígitos, distribuição conhecida ---, e não apenas compará-las. É o assunto do próximo capítulo.
\end{cor}

\begin{chave}
Recorrências se resolvem por árvore de recursão, substituição ou Teorema Mestre ($E=\log_ba$). Qualquer ordenação por comparação faz $\Omega(n\log n)$ comparações (árvore de decisão, $2^h\ge n!$). Mergesort: $\Theta(n\log n)$ sempre, estável, $\Theta(n)$ de memória extra. Quicksort: in-place, $\Theta(n\log n)$ em média e $\Theta(n^2)$ no pior caso, não estável. Heapsort: in-place, $\Theta(n\log n)$ sempre, não estável. Mergesort e heapsort são ótimos.
\end{chave}

% ----------------------------------------------------------------
\section{Exercícios}
\excurso
\begin{exercicio}{sc-merge}[Slides]<sec:sortcomp>
Escreva o pseudocódigo do procedimento de intercalação (\emph{Merge}) do mergesort.
\end{exercicio}
\begin{exercicio}{sc-iter}[Slides]<sec:sortcomp>
Escreva uma versão iterativa, não recursiva, do mergesort.
\end{exercicio}
\begin{exercicio}{sc-semcomp}[Slides]<sec:limiteinf>
Desenvolva uma estratégia para ordenar os elementos de um vetor que não seja baseada em comparações. Quais hipóteses sobre os elementos ela exige?
\end{exercicio}
\begin{exercicio}{sc-estavel}[Slides]<ex:estavel>
(a) É possível implementar o mergesort de forma estável? (b) E o quicksort? Justifique.
\end{exercicio}

\excompl
\begin{exercicio}{sc-mestre}[Complementar]<ex:mestre>
Resolva pelo Teorema Mestre: (a) $T(n)=4T(n/2)+n$; (b) $T(n)=4T(n/2)+n^2$; (c) $T(n)=4T(n/2)+n^3$; (d) $T(n)=2T(n/4)+\sqrt n$; (e) $T(n)=9T(n/3)+n$; (f) $T(n)=T(n/2)+n$.
\end{exercicio}
\begin{exercicio}{sc-arvore}[Complementar]<sec:recorrencias>
Use a árvore de recursão para resolver $T(n)=T(n/3)+T(2n/3)+n$.
\end{exercicio}
\begin{exercicio}{sc-pior8}[Complementar]<ex:mergesort>
Quantas comparações o mergesort faz no pior caso para $n=8$? Dê a recorrência para potências de 2 e a sua solução exata.
\end{exercicio}
\begin{exercicio}{sc-max}[Complementar]<sec:limiteinf>
Mostre que todo algoritmo que encontra o máximo de $n$ elementos por comparações faz pelo menos $n-1$ comparações. Por que o argumento da árvore de decisão dá um limite pior aqui?
\end{exercicio}
\begin{exercicio}{sc-iguais}[Complementar]<ex:lomuto>
Qual o custo do quicksort com a partição de Lomuto (teste ``$V[j]<x$'') quando todos os elementos são iguais? Como corrigir?
\end{exercicio}
\begin{exercicio}{sc-insercao}[Complementar]<sec:sortcomp>
Mostre que o número de deslocamentos feitos pela ordenação por inserção é exatamente o número de inversões do vetor (pares $i<j$ com $V[i]>V[j]$). Conclua que ela é $\Theta(n^2)$ no pior caso e no caso médio.
\end{exercicio}
\begin{exercicio}{sc-intercalar}[Complementar \estrela]<ex:mergesort>
Mostre que intercalar duas listas ordenadas de $n$ elementos cada exige $2n-1$ comparações no pior caso.
\end{exercicio}
\begin{exercicio}{sc-mtrace}[Complementar]<ex:mergesort>
Execute o mergesort sobre $[38,27,43,3,9,82,10]$, com $meio=\floor{(ini+fim)/2}$, e mostre o resultado de cada intercalação, na ordem em que acontecem.
\end{exercicio}
\begin{exercicio}{sc-lomuto}[Complementar]<ex:lomuto>
Execute a partição de Lomuto sobre $[2,8,7,1,3,5,6,4]$ (pivô 4) e dê o vetor resultante e a posição devolvida.
\end{exercicio}
\begin{exercicio}{sc-quickpior}[Complementar]<sec:sortcomp>
Quantas comparações exatamente o quicksort com a partição de Lomuto faz para ordenar um vetor já ordenado de $n$ elementos distintos?
\end{exercicio}
\begin{exercicio}{sc-subst}[Complementar]<ex:substituicao>
Prove por indução que a solução de $T(1)=0$, $T(n)=2T(n/2)+n$ para $n$ potência de 2 é $T(n)=n\log_2n$.
\end{exercicio}
\begin{exercicio}{sc-selecao}[Complementar]<ex:estavel>
A ordenação por seleção acha o menor elemento de $V[i..n]$ e o troca com $V[i]$, para $i=1,\dots,n-1$. Quantas comparações e quantas trocas ela faz? Ela é estável? E a ordenação por inserção?
\end{exercicio}
\begin{exercicio}{sc-cinco}[Complementar]<ex:limite34>
Qual o número mínimo de comparações que um algoritmo baseado em comparações precisa fazer, no pior caso, para ordenar 5 elementos, segundo o argumento da árvore de decisão? E para 12 elementos?
\end{exercicio}
\begin{exercicio}{sc-kdist}[Complementar \estrela]<sec:limiteinf>
Um vetor de $n$ elementos tem apenas $k$ valores distintos. Mostre que $\Omega(n\log k)$ comparações são necessárias para ordená-lo e dê um algoritmo com esse custo.
\end{exercicio}
\begin{exercicio}{sc-nlog2}[Complementar]<sec:recorrencias>
Resolva $T(n)=2T(n/2)+n\log n$ pela árvore de recursão. Por que o Teorema Mestre, na forma apresentada, não se aplica?
\end{exercicio}
\begin{exercicio}{sc-lineares}[Complementar]<sec:recorrencias>
Resolva: (a) $T(1)=1$, $T(n)=T(n-1)+n$; (b) $T(1)=1$, $T(n)=2T(n-1)+1$ (a recorrência das Torres de Hanói); (c) $T(1)=1$, $T(n)=T(n-1)+\log n$.
\end{exercicio}
\begin{exercicio}{sc-quicktrace}[Complementar]<ex:lomuto>
Execute o quicksort com a partição de Lomuto sobre $[5,8,1,9,3,7,2]$. Para cada chamada de \textsc{Particiona}, dê o trecho particionado, o pivô, o vetor depois da partição e a posição devolvida. Quantas comparações com o pivô são feitas no total?
\end{exercicio}
\begin{exercicio}{sc-instrace}[Complementar]<sec:sortcomp>
Execute a ordenação por inserção sobre $[6,2,5,1,4]$, mostrando o vetor depois de cada inserção e o número de deslocamentos de cada uma. Confira que o total é o número de inversões do vetor.
\end{exercicio}

% ================================================================
\chapter{Ordenação em tempo linear}\label{cap:sortlin}\label{sec:sortlin}
% ================================================================
\begin{objetivos}
\begin{itemize}
  \item Como escapar do limite $\Omega(n\log n)$ usando a estrutura das chaves.
  \item Counting sort, radix sort, bucket sort e American Flag sort.
  \item O papel da estabilidade no radix sort.
  \item Um guia para escolher o algoritmo de ordenação adequado a cada situação.
\end{itemize}
\end{objetivos}

O limite $\Omega(n\log n)$ do capítulo anterior vale para algoritmos que só \emph{comparam} elementos. Qualquer algoritmo de ordenação em tempo linear precisa, portanto, fazer no máximo $O(n)$ comparações e tomar suas decisões de outra forma. Se soubermos algo sobre as chaves --- que são inteiros num intervalo pequeno, números com poucos dígitos ou reais uniformemente distribuídos ---, podemos usá-las diretamente como \emph{índices} e ordenar em tempo linear. O preço é memória extra e hipóteses sobre os dados.

\section{Counting sort}
O counting sort ordena inteiros no intervalo $[0,k]$. Ele conta quantas vezes cada valor ocorre, acumula as contagens para saber quantos elementos são menores ou iguais a cada valor e, com isso, coloca cada elemento diretamente na sua posição final.
\begin{pseudo}{Counting-Sort$(A,n,k)$}
\begin{algorithmic}
\For{$i:=0$ \textbf{até} $k$} \State $C[i]:=0$ \EndFor
\For{$j:=1$ \textbf{até} $n$} \State $C[A[j]]:=C[A[j]]+1$ \EndFor \Comment{$C[i]=\#\{j: A[j]=i\}$}
\For{$i:=1$ \textbf{até} $k$} \State $C[i]:=C[i]+C[i-1]$ \EndFor \Comment{$C[i]=\#\{j: A[j]\le i\}$}
\For{$j:=n$ \textbf{decrescendo até} $1$} \Comment{de trás para frente: estável}
  \State $B[C[A[j]]]:=A[j]$;\ $C[A[j]]:=C[A[j]]-1$
\EndFor
\State \Return $B$
\end{algorithmic}
\end{pseudo}
O tempo e a memória são $\Theta(n+k)$, o que é \textbf{linear quando $k=O(n)$}. Quando $k$ é muito maior que $n$ (por exemplo, $k=\Theta(2^n)$), o algoritmo fica exponencial no tamanho da entrada: é \textbf{pseudopolinomial}. Ele é \textbf{estável}, porque, percorrendo a entrada de trás para frente, a última ocorrência de cada valor vai para a última posição do seu bloco; não é in-place, pois usa os vetores $B$ e $C$; e não faz nenhuma comparação entre elementos.

\section{Radix sort}
O radix sort ordena inteiros com $d$ dígitos na base $b$ \textbf{dígito a dígito, do menos significativo para o mais significativo}, usando em cada passada um método de ordenação \textbf{estável}. A versão abaixo usa uma fila para cada valor de dígito.
\begin{pseudo}{Radix sort com filas}
\begin{algorithmic}
\For{$i:=1$ \textbf{até} $d$}
  \For{$j:=1$ \textbf{até} $n$}
    \State $k:=$ $i$-ésimo dígito menos significativo de $L[j]$ na base $b$;\ insira $L[j]$ no fim da fila $F[k]$
  \EndFor
  \State $j:=1$
  \For{$k:=0$ \textbf{até} $b-1$}
    \While{$F[k]$ não está vazia} \State retire o início de $F[k]$ e coloque-o em $L[j]$;\ $j:=j+1$ \EndWhile
  \EndFor
\EndFor
\end{algorithmic}
\end{pseudo}
\textbf{Por que a estabilidade é necessária?} Depois da passada pelo dígito $i$, elementos com o mesmo dígito $i$ devem manter entre si a ordem obtida pelos dígitos $1,\dots,i-1$. É isso que garante, por indução, que após a passada $i$ a lista está ordenada pelos $i$ últimos dígitos. As filas preservam a ordem de chegada naturalmente; outra opção é usar o counting sort em cada dígito.

O tempo é $\Theta(d(n+b))$: com $b$ e $d$ constantes, $\Theta(n)$. Uma escolha útil é usar a própria base $b=n$: inteiros menores que $n^c$ têm $c$ dígitos nessa base, e o radix sort os ordena em $\Theta(cn)$.

\begin{exemplo}[Radix sort na base 10]
Ordenemos $102,313,690,131,222,670,690,130,079,102$.\\
Pelas unidades: $690\ 670\ 690\ 130\ 131\ 102\ 222\ 102\ 313\ 079$;\\
pelas dezenas: $102\ 102\ 313\ 222\ 130\ 131\ 670\ 079\ 690\ 690$;\\
pelas centenas: $079\ 102\ 102\ 130\ 131\ 222\ 313\ 670\ 690\ 690$.
\end{exemplo}

\paragraph{Chaves de tamanhos variados.} Se os números têm quantidades diferentes de dígitos, $d_1,\dots,d_n$, completar todos com zeros até o maior tamanho desperdiça trabalho. É possível ordenar em tempo $O(\sum_id_i)$, linear no tamanho da entrada, agrupando os números pelo número de dígitos (um número com menos dígitos é menor) e fazendo cada número participar apenas das passadas correspondentes aos seus próprios dígitos. Para \textbf{strings}, ordena-se da direita para a esquerda completando as menores com um símbolo ``menor que todos'', ou usa-se a versão que começa pelo caractere mais significativo, descrita a seguir.

\section{Bucket sort}
O bucket sort supõe que as chaves são reais \textbf{uniformemente distribuídos} em $[0,1)$. Divide-se o intervalo em $r$ baldes de mesmo tamanho, coloca-se $A[i]$ no balde $\floor{r\,A[i]}$, ordena-se cada balde (por exemplo, por inserção) e concatenam-se os baldes. Com entrada uniforme, cada balde recebe em média $n/r$ elementos; com $r=\Theta(n)$, o tempo \textbf{esperado} é $\Theta(n)$, porque, sendo $n_i$ o número de elementos no balde $i$, vale $E[n_i^2]=2-1/n$, e $\sum_iE[O(n_i^2)]=O(n)$. O pior caso, quando todos os elementos caem no mesmo balde, é o do algoritmo usado dentro dos baldes.

\section{American Flag sort}
O American Flag sort é um radix sort que começa pelo dígito \textbf{mais} significativo e trabalha \textbf{in-place}. Ele conta quantos elementos há com cada valor do dígito (um vetor $C$ de tamanho $b$), calcula onde começa o bloco de cada valor e \emph{permuta no próprio vetor} cada elemento para o seu bloco, por meio de trocas; depois, chama a si mesmo recursivamente em cada bloco, com o dígito seguinte. A memória auxiliar se resume aos contadores: $O(b)$ por nível de recursão, $O(b\cdot d)$ no total, independente de $n$. As trocas mudam a ordem relativa de elementos iguais, de modo que o algoritmo \textbf{não é estável} --- e não precisa ser, porque cada bloco é ordenado pelos dígitos seguintes de forma independente.

% ------------------------------------------------------------------
\section{Qual algoritmo de ordenação usar?}\label{sec:escolhasort}
\begin{center}
\small
\begin{tabularx}{\linewidth}{@{}lX@{}}
\toprule
Situação & Escolha\\
\midrule
Inteiros num intervalo pequeno ($k=O(n)$) & \textbf{Counting sort}, $\Theta(n+k)$\\
Inteiros ou identificadores com poucos dígitos ($d$ fixo) e $k$ enorme & \textbf{Radix sort}, $\Theta(d(n+b))$\\
Reais uniformemente distribuídos num intervalo & \textbf{Bucket sort}, $\Theta(n)$ esperado\\
Chaves arbitrárias, com garantia de tempo e/ou estabilidade & \textbf{Mergesort}, $\Theta(n\log n)$, memória $\Theta(n)$\\
Chaves arbitrárias, pouca memória, com garantia de tempo & \textbf{Heapsort}, $\Theta(n\log n)$, in-place\\
Chaves arbitrárias, desempenho médio na prática & \textbf{Quicksort} com pivô aleatório ou mediana de três\\
Vetor quase ordenado, ou $n$ pequeno & \textbf{Inserção}\\
Dados em lista encadeada ou em arquivos externos & \textbf{Mergesort}\\
\bottomrule
\end{tabularx}
\end{center}

\begin{exemplo}[Escolhendo o algoritmo]
(A) $10^7$ notas inteiras entre 0 e 100: \textbf{counting sort}, com 101 contadores e tempo $\Theta(n)$; uma ordenação por comparação faria cerca de $n\log_2n\approx2{,}3\cdot10^8$ comparações. (B) $10^6$ números reais uniformes em $[0,1)$: \textbf{bucket sort}, $\Theta(n)$ esperado; counting e radix não se aplicam a reais. (C) $10^5$ registros com chaves arbitrárias: só resta comparar, e o \textbf{mergesort} dá garantia de $\Theta(n\log n)$ e estabilidade.
\end{exemplo}

\begin{chave}
Counting sort: $\Theta(n+k)$, estável (último laço de trás para frente), linear se $k=O(n)$. Radix sort (do menos para o mais significativo): $\Theta(d(n+b))$, exige passadas estáveis. Bucket sort: $\Theta(n)$ esperado para entrada uniforme. American Flag sort: radix pelo dígito mais significativo, in-place e não estável. Para chaves arbitrárias, o limite $\Omega(n\log n)$ continua valendo.
\end{chave}

% ----------------------------------------------------------------
\section{Exercícios}
\excurso
\begin{exercicio}{sl-t2}[Avaliação]
Uma empresa precisa ordenar grandes volumes de dados com características diferentes. Para cada cenário, escolha o algoritmo mais adequado entre counting sort, radix sort, bucket sort, mergesort e quicksort, justificando por tempo, espaço e características dos dados: (A) $10^7$ notas entre 0 e 100; (B) $10^6$ identificadores numéricos de 12 dígitos; (C) $10^5$ números reais uniformemente distribuídos entre 0 e 1; (D) $10^6$ inteiros com valores entre $-10^9$ e $10^9$; (E) $10^5$ registros com chaves arbitrárias e sem restrição de domínio.
\end{exercicio}
\begin{exercicio}{sl-var}[Slides]
Modifique o radix sort para ordenar inteiros com quantidades diferentes de dígitos em tempo $O\big(\sum_{i=1}^nd_i\big)$, onde $d_i$ é o número de dígitos do $i$-ésimo inteiro. Faça o mesmo para o American Flag sort.
\end{exercicio}
\begin{exercicio}{sl-str}[Slides]
Que modificações são necessárias no radix sort e no American Flag sort para ordenar \emph{strings} em vez de inteiros?
\end{exercicio}

\excompl
\begin{exercicio}{sl-count}[Complementar]
Execute o counting sort em $A=[2,5,3,0,2,3,0,3]$ com $k=5$: mostre o vetor $C$ após a contagem e após as somas acumuladas, e a saída $B$.
\end{exercicio}
\begin{exercicio}{sl-radix}[Complementar]
Execute o radix sort na base 10 sobre $329,457,657,839,436,720,355$, mostrando a lista após cada passada.
\end{exercicio}
\begin{exercicio}{sl-n3}[Complementar]
Mostre como ordenar $n$ inteiros do intervalo $[0,n^3-1]$ em tempo $O(n)$.
\end{exercicio}
\begin{exercicio}{sl-frente}[Complementar]
No último laço do counting sort, o que muda se $j$ for de $1$ até $n$ (em vez de $n$ até $1$)? O resultado continua ordenado? E o radix sort que usa esse counting sort em cada passada?
\end{exercicio}
\begin{exercicio}{sl-bucket}[Complementar]
Qual o pior caso do bucket sort com ordenação por inserção nos baldes? Como garantir $O(n\log n)$ no pior caso sem perder o tempo esperado $\Theta(n)$?
\end{exercicio}
\begin{exercicio}{sl-msd}[Complementar]
Por que o radix sort que começa pelo dígito menos significativo precisa de passadas estáveis, enquanto o American Flag sort, que começa pelo mais significativo, funciona mesmo não sendo estável?
\end{exercicio}
\begin{exercicio}{sl-neg}[Complementar]
Como usar o counting sort para ordenar inteiros no intervalo $[-m,m]$? Qual o custo?
\end{exercicio}
\begin{exercicio}{sl-n2}[Complementar]
Para ordenar $n$ inteiros no intervalo $[0,n^2-1]$, compare o counting sort direto com o radix sort na base $n$.
\end{exercicio}
\begin{exercicio}{sl-strtrace}[Complementar]
Ordene pelo radix sort as palavras COW, DOG, SEA, RUG, ROW, MOB, BOX, TAB, BAR, EAR, TAR, DIG, BIG, TEA, NOW, FOX, mostrando a lista após cada uma das três passadas.
\end{exercicio}
\begin{exercicio}{sl-buckettrace}[Complementar]
Execute o bucket sort com 10 baldes sobre $0{,}78;\ 0{,}17;\ 0{,}39;\ 0{,}26;\ 0{,}72;\ 0{,}94;\ 0{,}21;\ 0{,}12;\ 0{,}23;\ 0{,}68$. Mostre o conteúdo de cada balde não vazio.
\end{exercicio}
\begin{exercicio}{sl-base}[Complementar]
Para ordenar $10^6$ inteiros de 32 bits com o radix sort, compare as bases $2^4$, $2^8$ e $2^{16}$: quantas passadas cada uma faz e qual o custo aproximado de cada passada? Qual parece a melhor escolha?
\end{exercicio}
\begin{exercicio}{sl-prova}[Complementar]
Prove por indução no número de passadas que o radix sort (do dígito menos significativo para o mais significativo, com passadas estáveis) ordena corretamente.
\end{exercicio}
\begin{exercicio}{sl-registros}[Complementar]
Uma lista de alunos já está em ordem alfabética e deve ser ordenada pela nota (um inteiro de 0 a 10), mantendo em ordem alfabética os alunos com a mesma nota. Que algoritmo usar e por quê?
\end{exercicio}
\begin{exercicio}{sl-datas}[Complementar]
Datas no formato (dia, mês, ano) devem ser ordenadas cronologicamente com o radix sort. Em que ordem devem ser feitas as passadas e qual a base de cada uma?
\end{exercicio}
\begin{exercicio}{sl-dup}[Complementar]
Dados $n$ inteiros no intervalo $[0,n^2]$, decida em tempo $O(n)$ se há elementos repetidos.
\end{exercicio}
\begin{exercicio}{sl-intervalo}[Complementar]
Dados $n$ inteiros no intervalo $[0,k]$, faça um pré-processamento em tempo $O(n+k)$ que permita responder em $O(1)$ a perguntas do tipo ``quantos elementos estão no intervalo $[a,b]$?''.
\end{exercicio}
\begin{exercicio}{sl-bucketruim}[Complementar]
Dê uma entrada de $n$ números em $[0,1)$ para a qual o bucket sort com $n$ baldes e inserção nos baldes gasta $\Theta(n^2)$.
\end{exercicio}
\begin{exercicio}{sl-flag}[Complementar]
Quanta memória auxiliar o American Flag sort usa para ordenar $n$ inteiros de $d$ dígitos na base $b$? Por que se diz que ele é in-place, se usa vetores de contadores?
\end{exercicio}

% ================================================================
\chapter{Seleção em tempo linear: a mediana das medianas}\label{cap:selecao}\label{sec:selecao}
% ================================================================
\begin{objetivos}
\begin{itemize}
  \item O problema da seleção do $k$-ésimo menor elemento e as soluções ingênuas.
  \item O quickselect e a sua análise de tempo esperado linear.
  \item O algoritmo da mediana das medianas, que seleciona em tempo linear no pior caso.
  \item Por que grupos de 5 funcionam e grupos de 3 não.
\end{itemize}
\end{objetivos}

Para achar a mediana de um conjunto --- ou, em geral, o $k$-ésimo menor elemento --- não é preciso ordenar. O quickselect resolve o problema em tempo linear \emph{esperado}, recursando em apenas um dos lados da partição. O algoritmo da mediana das medianas, de Blum, Floyd, Pratt, Rivest e Tarjan, garante tempo linear \emph{no pior caso}, escolhendo como pivô a mediana das medianas de grupos de cinco elementos. Sua análise é um belo exemplo de recorrência com subproblemas de tamanhos diferentes.

\section{O problema e o quickselect}
No \textbf{problema da seleção}, recebemos um conjunto $A$ de $n$ números e um inteiro $k$, e queremos o $k$-ésimo menor elemento de $A$. Casos particulares são o mínimo ($k=1$), o máximo ($k=n$) e as medianas inferior, $k=\floor{(n+1)/2}$, e superior, $k=\ceil{(n+1)/2}$.

Há várias maneiras de resolver o problema: (1) ordenar e tomar a posição $k$, em $O(n\log n)$; (2) retirar o mínimo $k$ vezes, em $O(kn)$, o que dá $O(n^2)$ para a mediana; (3) o \textbf{quickselect}, que particiona como o quicksort mas recursa em \emph{um} lado só, em $O(n)$ esperado e $O(n^2)$ no pior caso; (4) o algoritmo da \textbf{mediana das medianas}, em $O(n)$ no \textbf{pior caso}.

\begin{pseudo}{Quickselect: tempo esperado $O(n)$}
\begin{algorithmic}
\Function{QuickSelect}{$S,k$}
  \State escolha um pivô $m$ \textbf{aleatório} em $S$
  \State $S_1:=\{s<m\}$;\ $S_2:=\{s=m\}$;\ $S_3:=\{s>m\}$
  \If{$|S_1|\ge k$} \State \Return \Call{QuickSelect}{$S_1,k$}
  \ElsIf{$|S_1|+|S_2|\ge k$} \State \Return $m$
  \Else \State \Return \Call{QuickSelect}{$S_3,\,k-|S_1|-|S_2|$}
  \EndIf
\EndFunction
\end{algorithmic}
\end{pseudo}
\textbf{Análise.} Com probabilidade pelo menos $1/2$, o pivô cai entre o primeiro e o terceiro quartis, e então o lado mantido tem no máximo $3n/4$ elementos. Agrupe as chamadas em \emph{fases}: uma fase termina quando o tamanho cai para no máximo $3n/4$. O número de chamadas numa fase é dominado por uma variável geométrica de parâmetro $1/2$, cuja esperança é 2, e cada chamada custa no máximo $cn$. Assim, $E[T(n)]\le E[T(3n/4)]+2cn$, e, somando a série geométrica, $E[T(n)]\le 2cn\sum_{t\ge0}(3/4)^t=8cn=O(n)$.

\section{O algoritmo da mediana das medianas}
A ideia é garantir que o pivô seja \emph{sempre} bom: ele é a mediana das medianas de grupos de 5 elementos.
\begin{pseudo}{Seleção$(S,k)$ pela mediana das medianas}
\begin{algorithmic}
\Function{Seleção}{$S,k$}
  \If{$|S|\le49$} \State ordene $S$ e \Return o $k$-ésimo menor \EndIf
  \State divida $S$ em $\ceil{|S|/5}$ grupos de até 5 elementos
  \State determine a mediana de cada grupo \Comment{ordenar 5 elementos custa $O(1)$}
  \State $M:=$ conjunto das medianas;\quad $m:=$ \Call{Seleção}{$M,\ceil{|M|/2}$} \Comment{recursão sobre $n/5$ elementos}
  \State $S_1:=\{s\in S: s<m\}$;\ $S_2:=\{s\in S: s=m\}$;\ $S_3:=\{s\in S: s>m\}$
  \If{$|S_1|\ge k$} \State \Return \Call{Seleção}{$S_1,k$} \Comment{recursão sobre no máximo $3n/4$}
  \ElsIf{$|S_1|+|S_2|\ge k$} \State \Return $m$
  \Else \State \Return \Call{Seleção}{$S_3,\ k-|S_1|-|S_2|$}
  \EndIf
\EndFunction
\end{algorithmic}
\end{pseudo}
\begin{cuidado}
O teste que decide ir para $S_1$ é $|S_1|\ge k$, e não $|S_1|=k$: o $k$-ésimo menor está em $S_1$ sempre que há pelo menos $k$ elementos menores que o pivô. E, ao ir para $S_3$, o índice muda para $k-|S_1|-|S_2|$, pois os elementos de $S_1$ e $S_2$ ficaram para trás.
\end{cuidado}

\section{Análise}
\begin{lema}
$|S_1|\le n-3\floor{n/10}$ e $|S_3|\le n-3\floor{n/10}$. Para $n\ge50$, ambos são no máximo $3n/4$.
\end{lema}
\begin{proof}
Das $\floor{n/5}$ medianas de grupos completos, pelo menos metade, isto é, $\floor{\floor{n/5}/2}=\floor{n/10}$, são maiores ou iguais a $m$. Cada uma delas vem de um grupo que tem mais dois elementos maiores ou iguais a ela, e portanto a $m$. Assim, pelo menos $3\floor{n/10}$ elementos são maiores ou iguais a $m$, ou seja, $|S_2|+|S_3|\ge3\floor{n/10}$ e $|S_1|\le n-3\floor{n/10}$. O argumento para $S_3$ é simétrico.
Para que $n-3\floor{n/10}\le3n/4$ basta $12\floor{n/10}\ge n$, isto é, $2\floor{n/10}\ge n-10\floor{n/10}$. Como o lado direito é no máximo 9, basta $\floor{n/10}\ge5$, ou seja, $n\ge50$. Por isso o caso base é $n\le49$.
\end{proof}

\begin{teo}
$T(n)\le 20cn$; portanto, a seleção pela mediana das medianas custa $O(n)$ no pior caso.
\end{teo}
\begin{proof}
$T(n)=O(1)$ para $n\le49$, e $T(n)\le T(n/5)+T(3n/4)+cn$ para $n>49$, pois o cálculo das medianas dos grupos e a partição são lineares. Por indução forte, escolhendo $c$ grande o suficiente para que a base valha:
\[
T(n)\le 20c\,\frac n5+20c\,\frac{3n}{4}+cn=4cn+15cn+cn=20cn.\qedhere
\]
\end{proof}

\paragraph{Por que grupos de 5?} Para que o tempo seja linear, a soma dos tamanhos dos dois subproblemas precisa ser uma fração menor que 1 de $n$: com grupos de 5, $\frac n5+\frac{3n}4=\frac{19n}{20}<n$. Com grupos de 3, o pivô só garante cerca de $n/3$ elementos de cada lado, e a recorrência $T(n)=T(n/3)+T(2n/3)+cn$, com $\frac13+\frac23=1$, tem solução $\Theta(n\log n)$. Grupos de 7, 9 ou mais também funcionam, com outras constantes.

\paragraph{Aplicação.} Usando como pivô do quicksort a mediana exata, calculada pela mediana das medianas, obtém-se $T(n)=2T(n/2)+O(n)=O(n\log n)$ no \textbf{pior} caso. Na prática, porém, esse quicksort é mais lento que o quicksort com pivô aleatório, por causa das constantes.

\begin{chave}
Seleção em $O(n)$ no pior caso: grupos de 5, mediana das medianas como pivô, pelo menos $3\floor{n/10}$ elementos de cada lado, $T(n)\le T(n/5)+T(3n/4)+cn\le20cn$. O quickselect com pivô aleatório é mais simples e linear em média. Não é preciso ordenar para achar o $k$-ésimo menor.
\end{chave}

% ----------------------------------------------------------------
\section{Exercícios}
\excurso
\begin{exercicio}{se-quick}[Slides]
Escreva um algoritmo de tempo \emph{esperado} $O(n)$ para determinar a mediana de uma lista com $n$ elementos, mais simples que o método da mediana das medianas.
\end{exercicio}
\begin{exercicio}{se-qs}[Slides]
Argumente como o quicksort pode ser implementado com complexidade $O(n\log n)$ no pior caso. Qual o problema dessa abordagem na prática?
\end{exercicio}

\excompl
\begin{exercicio}{se-trace}[Complementar]
Para $S=12,3,25,7,18,\ 1,22,9,15,6,\ 20,2,11,24,5,\ 17,8,23,4,14,\ 10,21,13,19,16$ (grupos de 5 na ordem dada), calcule as medianas dos grupos, o pivô $m$ e os tamanhos de $S_1$ e $S_3$. Confira que ambos são pelo menos $3\floor{n/10}$.
\end{exercicio}
\begin{exercicio}{se-grupos}[Complementar]
Refaça a análise do algoritmo com grupos de 7 e com grupos de 3. Em qual dos casos ele continua linear?
\end{exercicio}
\begin{exercicio}{se-kmenores}[Complementar]
Dê um algoritmo $O(n)$ que devolve os $k$ menores elementos de um conjunto, em qualquer ordem. E se eles forem pedidos em ordem crescente?
\end{exercicio}
\begin{exercicio}{se-oleo}[Complementar]
Um oleoduto horizontal (leste-oeste) deve ligar $n$ poços, cada um por um duto vertical. Dadas as coordenadas $y_1,\dots,y_n$ dos poços, em que altura colocar o oleoduto para minimizar o comprimento total dos dutos verticais? Mostre que a resposta pode ser encontrada em $O(n)$.
\end{exercicio}
\begin{exercicio}{se-piorqs}[Complementar]
Mostre que o quickselect pode custar $\Theta(n^2)$ no pior caso, mesmo para encontrar o mínimo.
\end{exercicio}
\begin{exercicio}{se-qstrace}[Complementar]
Execute o quickselect sobre $[7,2,9,4,3,8,6,1]$ com $k=4$, usando sempre o \emph{último} elemento do conjunto (na ordem dada) como pivô. Mostre $S_1$, $S_2$, $S_3$ e o novo $k$ em cada chamada.
\end{exercicio}
\begin{exercicio}{se-n49}[Complementar]
Calcule $n-3\floor{n/10}$ e $3n/4$ para $n=49$ e para $n=50$. O que isso mostra sobre o caso base do algoritmo?
\end{exercicio}
\begin{exercicio}{se-ordmed}[Complementar]
Suponha que, em vez de calcular a mediana das medianas recursivamente, o algoritmo ordene as $n/5$ medianas com o mergesort para achar a do meio. Qual passa a ser a complexidade?
\end{exercicio}
\begin{exercicio}{se-lista}[Complementar]
O algoritmo da mediana das medianas continua linear se os elementos estiverem numa lista encadeada, em vez de um vetor? Justifique.
\end{exercicio}
\begin{exercicio}{se-quartis}[Complementar]
Mostre como calcular os três quartis de um conjunto de $n$ números em $O(n)$.
\end{exercicio}
\begin{exercicio}{se-proximos}[Complementar]
Dados $n$ números distintos e um inteiro $k$, encontre em $O(n)$ os $k$ números mais próximos da mediana (em valor absoluto da diferença).
\end{exercicio}
\begin{exercicio}{se-ponderada}[Complementar \estrela]
Cada elemento $x_i$ tem um peso $w_i>0$, com $\sum w_i=1$. A \emph{mediana ponderada} é o elemento $x_k$ tal que a soma dos pesos dos elementos menores que $x_k$ é menor que $1/2$ e a soma dos pesos dos maiores é no máximo $1/2$. Encontre-a em $O(n)$.
\end{exercicio}
\begin{exercicio}{se-dois}[Complementar \estrela]
Dados dois vetores ordenados $X$ e $Y$, cada um com $n$ elementos, encontre a mediana dos $2n$ elementos em $O(\log n)$.
\end{exercicio}
\begin{exercicio}{se-segundo}[Complementar \estrela]
Mostre como encontrar o segundo menor de $n$ elementos com $n+\ceil{\log_2n}-2$ comparações, organizando as comparações como um torneio.
\end{exercicio}
\begin{exercicio}{se-cinco}[Complementar \estrela]
Mostre como encontrar a mediana de 5 elementos com apenas 6 comparações.
\end{exercicio}
\begin{exercicio}{se-inferior}[Complementar]
Mostre que todo algoritmo baseado em comparações que encontra a mediana de $n$ elementos distintos faz pelo menos $n-1$ comparações. (Pense no grafo cujas arestas são as comparações feitas.)
\end{exercicio}
\begin{exercicio}{se-esperado}[Complementar \estrela]
Mostre que o número esperado de comparações do quickselect satisfaz \[E[T(n)]\le n+\frac2n\sum_{i=\floor{n/2}}^{n-1}E[T(i)]\] e conclua, por substituição, que $E[T(n)]\le4n$.
\end{exercicio}
\begin{exercicio}{se-maior}[Complementar]
Como usar um algoritmo que seleciona o $k$-ésimo \emph{menor} para encontrar o $k$-ésimo \emph{maior}? E para encontrar o percentil 90 de um conjunto de $n$ notas?
\end{exercicio}

% ================================================================
\chapter{Divisão e conquista}\label{cap:dc}\label{sec:dc}
% ================================================================
\begin{objetivos}
\begin{itemize}
  \item O paradigma de divisão e conquista e um roteiro para projetar e analisar algoritmos com ele.
  \item A contagem de inversões e o subvetor contíguo de soma máxima.
  \item A multiplicação de matrizes de Strassen.
  \item Outros exemplos clássicos: busca binária, exponenciação rápida, máximo e mínimo e elemento majoritário.
\end{itemize}
\end{objetivos}

O mergesort, o quicksort e a seleção pela mediana das medianas já usaram a mesma ideia: quebrar o problema em pedaços menores do mesmo tipo, resolvê-los recursivamente e combinar as respostas. Este capítulo trata essa ideia como um método de projeto, com um roteiro para escrever o pseudocódigo e analisar a recorrência, e a aplica a vários problemas.

\section{O paradigma}
Um algoritmo de \textbf{divisão e conquista} tem três passos: (1) \emph{divisão}: quebrar o problema $\Pi$ em subproblemas menores \textbf{do mesmo tipo}, $\Pi_1,\dots,\Pi_k$; (2) \emph{conquista}: resolvê-los recursivamente, resolvendo diretamente os casos pequenos; (3) \emph{combinação}: juntar as soluções dos subproblemas, por um procedimento relativamente simples, numa solução de $\Pi$. Busca-se um \textbf{bom balanceamento}, com partes de tamanhos parecidos, que leva a poucas camadas de recursão. A técnica funciona bem quando os subproblemas são \textbf{disjuntos}; quando eles se repetem muito, a técnica adequada é a programação dinâmica do Capítulo~\ref{cap:pd}.

\begin{roteiro}[Como escrever um algoritmo de divisão e conquista]
\begin{enumerate}
  \item \textbf{Assinatura}: uma função \textsc{Resolve}$(V,ini,fim)$ que trabalha num intervalo do vetor, para não copiar vetores. Diga o que a função \textbf{devolve}; às vezes é mais de um valor, como a soma \emph{e} os índices, ou a contagem \emph{e} o vetor ordenado.
  \item \textbf{Caso base}: $ini=fim$ (ou $n$ menor que uma constante), com resposta direta.
  \item \textbf{Divisão}: $meio:=\floor{(ini+fim)/2}$.
  \item \textbf{Conquista}: duas chamadas recursivas, sobre $[ini,meio]$ e $[meio+1,fim]$.
  \item \textbf{Combinação}: trata o que ``atravessa'' o meio; é a parte criativa e costuma custar $O(n)$.
  \item \textbf{Complexidade}: escreva a recorrência e resolva-a (§\ref{sec:recorrencias}); por exemplo, $2T(n/2)+\Theta(n)=\Theta(n\log n)$.
  \item \textbf{Correção}, em uma frase: toda solução está inteira à esquerda, inteira à direita ou cruza o meio, e cada caso é tratado.
\end{enumerate}
\end{roteiro}

% ------------------------------------------------------------------
\section{Contagem de inversões}
Dada uma sequência $a_1,\dots,a_n$, um par $i<j$ é uma \textbf{inversão} se $a_i>a_j$. O número de inversões mede quão distante uma ordem está da ordem crescente, e serve como medida de semelhança entre dois rankings: se uma pessoa ordena $n$ itens como $1,\dots,n$ e outra pessoa os ordena de outra forma, as inversões da segunda ordem contam os pares em que as duas discordam. A força bruta testa todos os pares, em $\Theta(n^2)$.

\textbf{Divisão e conquista.} Divida a sequência em metades $A$ e $B$; conte recursivamente as inversões dentro de cada uma; e \textbf{combine} contando os pares $(a,b)$ com $a\in A$, $b\in B$ e $a>b$. Esse último passo é fácil se $A$ e $B$ estiverem \emph{ordenadas}: durante a intercalação, ao comparar $a_i$ com $b_j$, se $a_i\le b_j$, então $a_i$ não forma inversão com nenhum elemento restante de $B$; se $a_i>b_j$, então $b_j$ forma inversão com \textbf{todos os elementos restantes de $A$}. Por isso a função também devolve o vetor ordenado: ela é um mergesort ``instrumentado''.
\begin{pseudo}{Contagem de inversões}
\begin{algorithmic}
\Function{CalculaInversões}{$V,ini,fim$} \Comment{conta e deixa $V[ini..fim]$ ordenado}
  \If{$ini=fim$} \State \Return $0$ \EndIf
  \State $meio:=\floor{(ini+fim)/2}$
  \State $a:=$ \Call{CalculaInversões}{$V,ini,meio$}
  \State $b:=$ \Call{CalculaInversões}{$V,meio+1,fim$}
  \State $c:=$ \Call{IntercalaContando}{$V,ini,meio,fim$}
  \State \Return $a+b+c$
\EndFunction
\Function{IntercalaContando}{$V,ini,meio,fim$}
  \State $i:=ini$;\ $j:=meio+1$;\ $k:=1$;\ $cont:=0$
  \While{$i\le meio$ \textbf{e} $j\le fim$}
    \If{$V[i]\le V[j]$} \State $W[k]:=V[i]$;\ $i:=i+1$
    \Else
      \State $W[k]:=V[j]$;\ $j:=j+1$
      \State $cont:=cont+(meio-i+1)$ \Comment{$V[j]$ é menor que todos os restantes da esquerda}
    \EndIf
    \State $k:=k+1$
  \EndWhile
  \State copie os restantes de $V[i..meio]$ e de $V[j..fim]$ para $W$;\ copie $W$ para $V[ini..fim]$
  \State \Return $cont$
\EndFunction
\end{algorithmic}
\end{pseudo}
O custo é $T(n)=2T(n/2)+\Theta(n)=\Theta(n\log n)$. Dois detalhes importam: o teste é ``$\le$'', porque elementos iguais não formam inversão, e soma-se $meio-i+1$, a quantidade de elementos que restam em $A$, e não 1. Por exemplo, a sequência $1,5,4,8,10,2,6,9,12,11,3,7$ tem \textbf{22} inversões.

% ------------------------------------------------------------------
\section{Subvetor contíguo de soma máxima}
Dado um vetor $A[1..n]$ de inteiros, que podem ser negativos, queremos $i\le j$ que maximizem $A[i]+\dots+A[j]$, isto é, o subvetor contíguo \textbf{não vazio} de maior soma. Pense, por exemplo, nos resultados diários de uma empresa, positivos ou negativos, e na busca do período de dias consecutivos de maior resultado acumulado.

\textbf{Ideia.} O melhor subvetor de $A[ini..fim]$ está (i) inteiro na metade esquerda, (ii) inteiro na direita, ou (iii) \textbf{cruza o meio}, sendo da forma $A[i..meio]\cup A[meio+1..j]$. O caso (iii) se resolve em tempo linear: é o melhor sufixo da esquerda que termina em $meio$ somado ao melhor prefixo da direita que começa em $meio+1$.
\begin{pseudo}{Subvetor de soma máxima por divisão e conquista (devolve soma, início e fim)}
\begin{algorithmic}
\Function{MaxSub}{$A,ini,fim$}
  \If{$ini=fim$} \State \Return $(A[ini],\,ini,\,ini)$ \Comment{não vazio: um elemento}
  \EndIf
  \State $meio:=\floor{(ini+fim)/2}$
  \State $(s_E,i_E,j_E):=$ \Call{MaxSub}{$A,ini,meio$}
  \State $(s_D,i_D,j_D):=$ \Call{MaxSub}{$A,meio+1,fim$}
  \State $(s_C,i_C,j_C):=$ \Call{MaxCruzado}{$A,ini,meio,fim$}
  \State \Return a tripla de maior soma entre $(s_E,i_E,j_E)$, $(s_D,i_D,j_D)$ e $(s_C,i_C,j_C)$
\EndFunction
\Function{MaxCruzado}{$A,ini,meio,fim$}
  \State $soma:=0$;\ $melhorE:=-\infty$
  \For{$i:=meio$ \textbf{decrescendo até} $ini$}
    \State $soma:=soma+A[i]$
    \If{$soma>melhorE$} \State $melhorE:=soma$;\ $iMax:=i$ \EndIf
  \EndFor
  \State $soma:=0$;\ $melhorD:=-\infty$
  \For{$j:=meio+1$ \textbf{até} $fim$}
    \State $soma:=soma+A[j]$
    \If{$soma>melhorD$} \State $melhorD:=soma$;\ $jMax:=j$ \EndIf
  \EndFor
  \State \Return $(melhorE+melhorD,\ iMax,\ jMax)$
\EndFunction
\end{algorithmic}
\end{pseudo}
\begin{itemize}
  \item \textbf{Complexidade}: $T(n)=2T(n/2)+\Theta(n)=\Theta(n\log n)$.
  \item \textbf{Todos os elementos negativos}: o algoritmo funciona, porque o caso base devolve um elemento e o subvetor cruzado usa pelo menos $A[meio]$ e $A[meio+1]$; o máximo das três opções acaba sendo o maior elemento. Por exemplo, $[-8,-3,-6,-2,-5,-4]$ dá $-2$.
  \item Em $[-2,1,-3,4,-1,2,1,-5,4]$, a resposta é a soma $6$, com o subvetor $[4,-1,2,1]$, das posições 4 a 7 (com índices a partir de 1).
  \item Existem soluções lineares: a programação dinâmica de Kadane (Capítulo~\ref{cap:pd}) e uma divisão e conquista em que cada metade devolve quatro valores (Exercício~\ref{ex:dc-linear}).
\end{itemize}

% ------------------------------------------------------------------
\section{Multiplicação de matrizes: o algoritmo de Strassen}
Multiplicando matrizes $n\times n$ por blocos de tamanho $n/2$, o produto ingênuo faz 8 multiplicações de blocos, e $T(n)=8T(n/2)+\Theta(n^2)=\Theta(n^3)$, sem ganho sobre o método escolar. Strassen mostrou como fazer apenas \textbf{7} multiplicações de blocos, à custa de mais somas e subtrações: $T(n)=7T(n/2)+\Theta(n^2)=\Theta(n^{\log_27})\approx\Theta(n^{2{,}807})$, pelo caso 1 do Teorema Mestre, pois $n^2=O(n^{\log_27-\varepsilon})$.
\begin{pseudo}{Strassen$(A,B,n)$, com $n$ potência de 2 e soma e subtração de matrizes dadas}
\begin{algorithmic}
\Function{Strassen}{$A,B,n$}
  \If{$n=1$} \State \Return $[\,A[1,1]\cdot B[1,1]\,]$ \EndIf
  \State particione $A=\begin{bmatrix}A_{11}&A_{12}\\A_{21}&A_{22}\end{bmatrix}$ e $B=\begin{bmatrix}B_{11}&B_{12}\\B_{21}&B_{22}\end{bmatrix}$ em blocos $\frac n2\times\frac n2$
  \State $M_1:=$ \Call{Strassen}{$A_{11}+A_{22},\ B_{11}+B_{22},\ n/2$}
  \State $M_2:=$ \Call{Strassen}{$A_{21}+A_{22},\ B_{11},\ n/2$}
  \State $M_3:=$ \Call{Strassen}{$A_{11},\ B_{12}-B_{22},\ n/2$}
  \State $M_4:=$ \Call{Strassen}{$A_{22},\ B_{21}-B_{11},\ n/2$}
  \State $M_5:=$ \Call{Strassen}{$A_{11}+A_{12},\ B_{22},\ n/2$}
  \State $M_6:=$ \Call{Strassen}{$A_{21}-A_{11},\ B_{11}+B_{12},\ n/2$}
  \State $M_7:=$ \Call{Strassen}{$A_{12}-A_{22},\ B_{21}+B_{22},\ n/2$}
  \State $C_{11}:=M_1+M_4-M_5+M_7$;\quad $C_{12}:=M_3+M_5$
  \State $C_{21}:=M_2+M_4$;\quad $C_{22}:=M_1-M_2+M_3+M_6$
  \State \Return $\begin{bmatrix}C_{11}&C_{12}\\C_{21}&C_{22}\end{bmatrix}$
\EndFunction
\end{algorithmic}
\end{pseudo}
Cada chamada faz 18 somas ou subtrações de matrizes $\frac n2\times\frac n2$, de custo $\Theta(n^2)$. Na prática, usa-se o método clássico abaixo de um certo tamanho, pois as constantes de Strassen são maiores e o algoritmo é numericamente menos estável.

% ------------------------------------------------------------------
\section{Outros exemplos clássicos}
\begin{pseudo}{Busca binária, potência rápida e máximo e mínimo simultâneos}
\begin{algorithmic}
\Function{BuscaBin}{$V,ini,fim,x$} \Comment{$V$ ordenado; $T(n)=T(n/2)+O(1)=O(\log n)$}
  \If{$ini>fim$} \State \Return $-1$ \EndIf
  \State $meio:=\floor{(ini+fim)/2}$
  \If{$V[meio]=x$} \State \Return $meio$
  \ElsIf{$x<V[meio]$} \State \Return \Call{BuscaBin}{$V,ini,meio-1,x$}
  \Else \State \Return \Call{BuscaBin}{$V,meio+1,fim,x$}
  \EndIf
\EndFunction
\Function{Pot}{$a,n$} \Comment{$a^n$ com $O(\log n)$ multiplicações}
  \If{$n=0$} \State \Return $1$ \EndIf
  \State $t:=$ \Call{Pot}{$a,\floor{n/2}$}
  \If{$n$ par} \State \Return $t\cdot t$ \Else \State \Return $t\cdot t\cdot a$ \EndIf
\EndFunction
\Function{MaxMin}{$V,ini,fim$} \Comment{cerca de $3n/2$ comparações, contra $2n-2$ do método ingênuo}
  \If{$ini=fim$} \State \Return $(V[ini],V[ini])$ \EndIf
  \If{$fim=ini+1$} \State \Return $(\max,\min)$ de $V[ini]$ e $V[fim]$, com 1 comparação \EndIf
  \State $meio:=\floor{(ini+fim)/2}$
  \State $(M_1,m_1):=$ \Call{MaxMin}{$V,ini,meio$};\ $(M_2,m_2):=$ \Call{MaxMin}{$V,meio+1,fim$}
  \State \Return $(\max\{M_1,M_2\},\ \min\{m_1,m_2\})$
\EndFunction
\end{algorithmic}
\end{pseudo}

\begin{pseudo}{Elemento majoritário (aparece mais de $n/2$ vezes): $T(n)=2T(n/2)+\Theta(n)=\Theta(n\log n)$}
\begin{algorithmic}
\Function{Maj}{$V,ini,fim$} \Comment{devolve o majoritário de $V[ini..fim]$ ou \textsc{nenhum}}
  \If{$ini=fim$} \State \Return $V[ini]$ \EndIf
  \State $meio:=\floor{(ini+fim)/2}$
  \State $a:=$ \Call{Maj}{$V,ini,meio$};\ $b:=$ \Call{Maj}{$V,meio+1,fim$}
  \For{$c\in\{a,b\}$, $c\ne$ \textsc{nenhum}}
    \If{$c$ ocorre mais de $(fim-ini+1)/2$ vezes em $V[ini..fim]$} \State \Return $c$ \EndIf \Comment{contagem em $O(n)$}
  \EndFor
  \State \Return \textsc{nenhum}
\EndFunction
\end{algorithmic}
\end{pseudo}
A correção do último algoritmo vem de uma observação simples: se $x$ é majoritário no vetor todo, ele é majoritário em pelo menos uma das metades, pois do contrário apareceria no máximo $n/2$ vezes.

\paragraph{Mais ideias.} O algoritmo de \emph{Karatsuba} multiplica inteiros de $n$ dígitos com 3 produtos de inteiros de $n/2$ dígitos, em $\Theta(n^{\log_23})\approx\Theta(n^{1{,}585})$. O \emph{par de pontos mais próximos} no plano se resolve em $\Theta(n\log n)$, dividindo os pontos pela mediana das abscissas e combinando com uma varredura da faixa central, em que cada ponto precisa ser comparado com no máximo 7 vizinhos. O \emph{pico de um vetor unimodal} e a \emph{raiz de uma função monótona} se encontram por busca binária, em $O(\log n)$; e o número de ocorrências de $x$ num vetor ordenado sai de duas buscas binárias.

\begin{cuidado}[Erros frequentes em pseudocódigo de divisão e conquista]
\begin{itemize}
  \item Esquecer o caso base, ou escrevê-lo errado (por exemplo, permitir subvetor vazio quando o problema pede um não vazio).
  \item Dividir em $[ini,meio-1]$ e $[meio,fim]$ com $meio=\floor{(ini+fim)/2}$: com dois elementos, a recursão não termina. Use $[ini,meio]$ e $[meio+1,fim]$.
  \item Fazer uma combinação que olha um só elemento de cada lado; o subvetor cruzado precisa do \emph{melhor sufixo} e do \emph{melhor prefixo}.
  \item Não devolver a informação extra necessária, como os índices no subvetor máximo ou o vetor ordenado na contagem de inversões.
  \item Não analisar a complexidade.
\end{itemize}
\end{cuidado}

\begin{chave}
Divisão e conquista $=$ dividir ao meio, conquistar com duas chamadas e combinar o que cruza o meio. $2T(n/2)+\Theta(n)=\Theta(n\log n)$; $T(n/2)+O(1)=\Theta(\log n)$; $7T(n/2)+\Theta(n^2)=\Theta(n^{2{,}807})$; $T(n/5)+T(3n/4)+\Theta(n)=\Theta(n)$.
\end{chave}

% ----------------------------------------------------------------
\section{Exercícios}
\excurso
\begin{exercicio}{dc-t1}[Avaliação]
Dado um vetor de inteiros, determine por divisão e conquista a maior soma de um subvetor contíguo \textbf{não vazio} e o subvetor responsável, com suas posições inicial e final. Determine a complexidade. Exemplos: $[-2,1,-3,4,-1,2,1,-5,4]$ (soma 6, subvetor $[4,-1,2,1]$) e $[-8,-3,-6,-2,-5,-4]$ (soma $-2$). Compare sua solução com a da §\ref{sec:dc}.
\end{exercicio}
\begin{exercicio}{dc-t2}[Avaliação]
Escreva um algoritmo recursivo que multiplica duas matrizes $n\times n$, com $n$ potência de 2, pela estratégia de Strassen, com as sete multiplicações $M_1,\dots,M_7$, supondo dados procedimentos de soma e subtração de matrizes.
\end{exercicio}
\begin{exercicio}{dc-balanc}[Slides]
Prove: um algoritmo de divisão e conquista $c$-balanceado ($c>1$; cada uma das $k$ partes tem tamanho no máximo $n/c$, e os tamanhos diferem por uma constante), que gasta tempo linear para dividir e combinar, roda em $O(n\log n)$.
\end{exercicio}
\begin{exercicio}{dc-inv}[Slides]
Escreva o pseudocódigo completo do algoritmo de contagem de inversões por divisão e conquista.
\end{exercicio}

\excompl
\begin{exercicio}{dc-contar}[Complementar]
Conte as inversões de (a) $2,4,1,3,5$; (b) $5,4,3,2,1$; (c) $3,1,4,1,5,9,2,6$ (pares iguais não são inversão). Qual o máximo possível para $n$ elementos?
\end{exercicio}
\begin{exercicio}{dc-pico}[Complementar]
Um vetor $V[1..n]$ é \emph{unimodal}: estritamente crescente até uma posição $p$ e estritamente decrescente depois dela. Encontre $p$ em $O(\log n)$.
\end{exercicio}
\begin{exercicio}{dc-ocorr}[Complementar]
Conte quantas vezes $x$ aparece num vetor ordenado em $O(\log n)$.
\end{exercicio}
\begin{exercicio}{dc-kara}[Complementar]
Multiplique dois inteiros de $n$ dígitos com três produtos de inteiros de $n/2$ dígitos (Karatsuba): escreva $x=x_1B^{n/2}+x_0$ e $y=y_1B^{n/2}+y_0$ e use $(x_1+x_0)(y_1+y_0)$. Qual a recorrência e a sua solução?
\end{exercicio}
\begin{exercicio}{dc-maxmin}[Complementar]
Mostre que o \textsc{MaxMin} faz $\frac{3n}2-2$ comparações quando $n$ é potência de 2.
\end{exercicio}
\begin{exercicio}{dc-maj}[Complementar \estrela]
Dê um algoritmo $O(n)$, com memória $O(1)$, para encontrar o elemento majoritário, se existir.
\end{exercicio}
\begin{exercicio}{dc-bbin}[Complementar]
Escreva a busca binária de forma iterativa e prove que ela faz no máximo $\floor{\log_2n}+1$ comparações com elementos do vetor.
\end{exercicio}
\begin{exercicio}{dc-pot}[Complementar]
Mostre que \textsc{Pot}$(a,n)$ faz no máximo $2\floor{\log_2n}+2$ multiplicações. Quantas faz para $n=13$?
\end{exercicio}
\begin{exercicio}{dc-linear}[Complementar]
Mostre como resolver o subvetor de soma máxima por divisão e conquista em $O(n)$, fazendo cada chamada devolver quatro valores: a soma total do intervalo, o melhor prefixo, o melhor sufixo e o melhor subvetor.
\end{exercicio}
\begin{exercicio}{dc-blocos}[Complementar]
Escreva a multiplicação de matrizes por blocos com 8 multiplicações recursivas e mostre que ela custa $\Theta(n^3)$. Quantas somas de blocos ela faz por chamada?
\end{exercicio}
\begin{exercicio}{dc-faltante}[Complementar]
Um vetor ordenado contém os inteiros de 1 a $n+1$, exceto um. Encontre o que falta em $O(\log n)$.
\end{exercicio}
\begin{exercicio}{dc-fixo}[Complementar]
Dado um vetor ordenado $A[1..n]$ de inteiros \emph{distintos}, decida em $O(\log n)$ se existe $i$ com $A[i]=i$.
\end{exercicio}
\begin{exercicio}{dc-uns}[Complementar]
Um vetor ordenado contém apenas zeros seguidos de uns. Conte os uns em $O(\log n)$. E se o tamanho do vetor for desconhecido (pode-se consultar qualquer posição, e posições além do fim devolvem um valor especial)?
\end{exercicio}
\begin{exercicio}{dc-signif}[Complementar \estrela]
Um par $i<j$ é uma \emph{inversão significativa} se $a_i>2a_j$. Conte as inversões significativas em $O(n\log n)$.
\end{exercicio}
\begin{exercicio}{dc-par}[Complementar \estrela]
Descreva o algoritmo de divisão e conquista para o par de pontos mais próximos no plano e justifique por que cada ponto da faixa central precisa ser comparado com um número constante de vizinhos.
\end{exercicio}
\begin{exercicio}{dc-skyline}[Complementar \estrela]
Prédios retangulares são dados por triplas $(e_i,d_i,h_i)$: o prédio ocupa o intervalo $[e_i,d_i]$ e tem altura $h_i$. O \emph{contorno} (\emph{skyline}) é a curva formada pelo ponto mais alto em cada abscissa. Calcule o contorno de $n$ prédios em $O(n\log n)$.
\end{exercicio}
\begin{exercicio}{dc-recs}[Complementar]
Escreva e resolva as recorrências da busca binária, da potência rápida, do \textsc{MaxMin} e do elemento majoritário por divisão e conquista.
\end{exercicio}

% ================================================================
\chapter{Escolhendo estruturas de dados e algoritmos}\label{cap:revisao}
% ================================================================
\begin{objetivos}
\begin{itemize}
  \item Um método para escolher a estrutura de dados e o algoritmo adequados a uma situação.
  \item Um quadro comparativo das estruturas e algoritmos da Parte~I.
  \item Como reconhecer, num trecho de código, oportunidades de melhorar a complexidade.
\end{itemize}
\end{objetivos}

Conhecer cada estrutura e cada algoritmo isoladamente não basta: na prática, o problema chega descrito em palavras, ou na forma de um programa lento, e é preciso decidir qual ferramenta usar. Este capítulo reúne os critérios espalhados pelos capítulos anteriores e os organiza num método.

\section{Um método de escolha}
\begin{roteiro}[Como escolher]
\begin{enumerate}
  \item \textbf{Liste as operações} que o problema (ou o código) faz repetidamente: buscar por chave, achar o máximo, inserir, remover, unir grupos, ordenar, selecionar o $k$-ésimo.
  \item \textbf{Estime quantas vezes} cada operação acontece. A operação dominante é a que deve ficar barata.
  \item \textbf{Escolha a estrutura} que torna barata a operação dominante sem tornar as outras proibitivas.
  \item \textbf{Olhe o tipo das chaves} para escolher a ordenação: inteiros pequenos, inteiros com poucos dígitos, reais uniformes ou chaves arbitrárias.
  \item \textbf{Verifique as restrições}: memória disponível, necessidade de tempo previsível por operação (pior caso) ou apenas de custo total (amortizado), estabilidade, necessidade de manter ordem.
  \item \textbf{Justifique} com as complexidades das operações e com uma estimativa do custo total.
\end{enumerate}
\end{roteiro}

\section{Quadro comparativo}\label{sec:escolhageral}
\begin{center}
\small
\begin{tabularx}{\linewidth}{@{}XX@{}}
\toprule
O problema faz repetidamente\dots & Estrutura ou algoritmo\\
\midrule
pega o maior (ou o menor), insere e às vezes altera prioridades & \textbf{heap binário} (Capítulo~\ref{cap:heap})\\
\dots\ e une filas com frequência, exigindo tempo previsível & \textbf{heap binomial} (Capítulo~\ref{cap:heapav})\\
\dots\ com um número enorme de aumentos de prioridade, importando o custo total & \textbf{heap de Fibonacci} (Capítulo~\ref{cap:heapav})\\
``$x$ e $y$ estão no mesmo grupo?'' e une grupos & \textbf{Union-Find} com rank e compressão (Capítulo~\ref{cap:uf})\\
\dots\ e precisa listar os membros de um grupo & \textbf{Union-Find por listas} (Capítulo~\ref{cap:uf})\\
busca por chave, inserção, remoção, percurso em ordem, sucessor, mínimo e máximo, intervalos & \textbf{árvore AVL} ou rubro-negra (Capítulo~\ref{cap:avl})\\
só busca, inserção e remoção por chave exata, sem ordem & \textbf{tabela hash}, $O(1)$ esperado (Apêndice~\ref{cap:hash})\\
muitas buscas e quase nenhuma atualização & \textbf{vetor ordenado} e busca binária\\
o $k$-ésimo menor ou a mediana, uma vez & \textbf{seleção linear} (Capítulo~\ref{cap:selecao}), sem ordenar\\
os $k$ maiores de um fluxo de dados & \textbf{heap de mínimo de tamanho $k$}, $O(n\log k)$\\
ordenar inteiros num intervalo pequeno ou com poucos dígitos & \textbf{counting} ou \textbf{radix sort} (Capítulo~\ref{cap:sortlin})\\
ordenar chaves arbitrárias com garantia e estabilidade & \textbf{mergesort} (Capítulo~\ref{cap:sortcomp})\\
\bottomrule
\end{tabularx}
\end{center}

\paragraph{Padrões a reconhecer em código.} Alguns padrões indicam, quase sempre, uma estrutura inadequada:
\begin{itemize}
  \item um laço que procura o máximo de uma lista e o remove, a cada iteração ($\Theta(n)$ por iteração): use um \textbf{heap};
  \item ordenar a lista inteira depois de cada inserção: use um \textbf{heap} ou uma \textbf{AVL}, conforme seja preciso só o extremo ou a ordem toda;
  \item percorrer todos os elementos para saber se dois estão no mesmo grupo: use \textbf{Union-Find};
  \item ordenar só para depois pegar o $k$-ésimo elemento: use \textbf{seleção} em $O(n)$;
  \item ordenar inteiros pequenos com um \texttt{sort} genérico: use \textbf{counting} ou \textbf{radix sort};
  \item procurar um elemento por varredura a cada consulta: use \textbf{hash} (sem ordem) ou \textbf{AVL} (com ordem).
\end{itemize}

\section{Resumo das complexidades}
\begin{center}
\small
\textbf{Listas de prioridades} ($n$ elementos; \am\ indica custo amortizado)\\[2pt]
\begin{tabular}{@{}lcccc@{}}
\toprule
Operação & Lista não ordenada & Heap binário & Heap binomial & Heap de Fibonacci\\
\midrule
Seleção do máximo & $\Theta(n)$ & $O(1)$ & $O(\log n)$ ou $O(1)$ & $O(1)$\\
Inserção & $\Theta(1)$ & $O(\log n)$ & $O(\log n)$; $O(1)$\am & $O(1)$\am\\
Remoção do máximo & $\Theta(n)$ & $O(\log n)$ & $O(\log n)$ & $O(\log n)$\am\\
Aumentar prioridade & $\Theta(n)$ & $O(\log n)$ & $O(\log n)$ & $O(1)$\am\\
Fusão & $O(1)$ & $O(n)$ & $O(\log n)$ & $O(1)$\am\\
\bottomrule
\end{tabular}

\medskip
\textbf{Ordenação}\\[2pt]
\begin{tabular}{@{}lcccc@{}}
\toprule
Algoritmo & Pior caso & Caso médio & Memória extra & Estável?\\
\midrule
Inserção & $\Theta(n^2)$ & $\Theta(n^2)$ & $O(1)$ & sim\\
Mergesort & $\Theta(n\log n)$ & $\Theta(n\log n)$ & $\Theta(n)$ & sim\\
Heapsort & $\Theta(n\log n)$ & $\Theta(n\log n)$ & $O(1)$ & não\\
Quicksort & $\Theta(n^2)$ & $\Theta(n\log n)$ & $O(\log n)$ & não\\
Counting sort & $\Theta(n+k)$ & $\Theta(n+k)$ & $\Theta(n+k)$ & sim\\
Radix sort & $\Theta(d(n+b))$ & $\Theta(d(n+b))$ & $\Theta(n+b)$ & sim\\
Bucket sort & $\Theta(n^2)$ & $\Theta(n)$ (uniforme) & $\Theta(n)$ & depende\\
\bottomrule
\end{tabular}
\end{center}
\textbf{Outras estruturas.} Union-Find com rank e compressão: $O(\lgs n)$ amortizado por operação (ou $O(\alpha)$). AVL: busca, inserção e remoção em $O(\log n)$ no pior caso. Seleção: $O(n)$ no pior caso pela mediana das medianas. Divisão e conquista: contagem de inversões e subvetor máximo em $O(n\log n)$; Strassen em $O(n^{2{,}807})$.

\section{Exemplos}
\begin{exemplo}[Uma agenda de compromissos]
Um aplicativo guarda compromissos com data e hora, insere e remove compromissos com frequência, mostra o próximo compromisso e lista todos os compromissos de uma semana. As operações dominantes são inserção, remoção de um compromisso específico, ``próximo compromisso'' e consulta por intervalo. Um heap daria o próximo compromisso, mas não remove um elemento qualquer pela chave nem lista um intervalo; uma tabela hash não mantém ordem. Uma \textbf{AVL} indexada por data e hora faz tudo: inserção e remoção em $O(\log n)$, mínimo em $O(\log n)$ e a semana em $O(\log n+k)$, com $k$ compromissos listados.
\end{exemplo}
\begin{exemplo}[Ranking de uma competição]
São $10^6$ participantes com pontuações inteiras de 0 a 1000, e é preciso imprimir o ranking uma única vez, mantendo, entre empatados, a ordem de inscrição. As chaves são inteiros num intervalo pequeno e a estabilidade importa: \textbf{counting sort}, em $\Theta(n+k)$ com $k=1001$, estável.
\end{exemplo}

\begin{chave}[Resumo da Parte I]
Identifique a operação dominante e o tipo das chaves. Heap para extremos; binomial para fusões previsíveis; Fibonacci para muitos aumentos de prioridade. Union-Find para grupos. AVL para ordem e busca; hash para busca sem ordem. Seleção sem ordenar. Counting e radix para inteiros; mergesort para chaves arbitrárias com estabilidade; heapsort com pouca memória; quicksort para desempenho médio. Justifique sempre com complexidade, memória e o tipo de garantia (pior caso ou amortizada).
\end{chave}

% ----------------------------------------------------------------
\section{Exercícios}
Para cada situação, diga qual estrutura de dados ou algoritmo você usaria e justifique com as complexidades envolvidas.
\excompl
\begin{exercicio}{rv-maior}[Complementar]
\textbf{Enquanto} houver tarefas na lista $L$ \textbf{faça}: percorra $L$ para achar a de maior prioridade, execute-a e remova-a de $L$; novas tarefas podem ser acrescentadas a $L$ durante a execução.
\end{exercicio}
\begin{exercicio}{rv-kesimo}[Complementar]
Para cada uma de $q$ consultas: copie o vetor $V$ (que não muda) de tamanho $n$, ordene a cópia com um \texttt{sort} genérico e devolva o elemento da posição $k_i$.
\end{exercicio}
\begin{exercicio}{rv-grupos}[Complementar]
Uma rede social recebe $10^6$ pedidos de amizade, que unem comunidades, e $10^8$ perguntas ``$u$ e $v$ estão na mesma comunidade?''. O código atual responde cada pergunta com uma busca em largura no grafo de amizades.
\end{exercicio}
\begin{exercicio}{rv-dic}[Complementar]
Um sistema guarda produtos por código, com inserções e remoções frequentes, e precisa responder ``quais códigos estão entre $a$ e $b$?''. Hoje os códigos ficam num vetor ordenado e cada inserção desloca os elementos.
\end{exercicio}
\begin{exercicio}{rv-idades}[Complementar]
Ordenar as idades (inteiros entre 0 e 120) de $10^8$ pessoas com o \texttt{sort} da biblioteca.
\end{exercicio}
\begin{exercicio}{rv-filas}[Complementar]
Um sistema de atendimento tem uma fila de prioridade por guichê; quando um guichê fecha, sua fila é unida à de outro guichê. Fechamentos são frequentes e o tempo de cada operação precisa ser previsível. Hoje cada fila é um heap binário.
\end{exercicio}
\begin{exercicio}{rv-emerg}[Complementar]
Um pronto-socorro mantém os pacientes em espera ordenados por gravidade. Chegam e são atendidos milhares de pacientes por dia; a gravidade de um paciente raramente muda; o sistema roda num dispositivo com pouca memória.
\end{exercicio}
\begin{exercicio}{rv-cancel}[Complementar]
Um simulador mantém eventos futuros ordenados por horário; a todo momento processa o próximo evento, cria novos eventos e às vezes cancela um evento já agendado, identificado por um número.
\end{exercicio}
\begin{exercicio}{rv-palavras}[Complementar]
Um corretor ortográfico precisa verificar, para cada palavra de um texto, se ela está num dicionário fixo de $10^6$ palavras. Não há atualizações nem consultas por ordem alfabética.
\end{exercicio}
\begin{exercicio}{rv-posicao}[Complementar \estrela]
Um placar online recebe pontuações e precisa responder ``qual a posição de $x$ no ranking?'' e ``quem está na posição $r$?'', com inserções e remoções frequentes.
\end{exercicio}
\begin{exercicio}{rv-mediana}[Complementar]
Um sensor envia medições continuamente, e após cada medição é preciso exibir a mediana de todas as medições recebidas até então.
\end{exercicio}
\begin{exercicio}{rv-externo}[Complementar]
Ordenar um arquivo de 500 GB de registros num computador com 16 GB de memória.
\end{exercicio}
\begin{exercicio}{rv-quase}[Complementar]
Um vetor de $10^6$ elementos está ordenado, exceto por uns poucos elementos fora do lugar, cada um a no máximo 10 posições da posição correta.
\end{exercicio}
\begin{exercicio}{rv-dup}[Complementar]
Decidir se há strings repetidas numa lista de $10^7$ strings.
\end{exercicio}
\begin{exercicio}{rv-remocoes}[Complementar]
Uma malha de estradas recebe uma lista de interdições, uma por vez, e após cada interdição é preciso dizer quantas regiões desconectadas existem. Toda a lista de interdições é conhecida de antemão.
\end{exercicio}
\begin{exercicio}{rv-topk}[Complementar]
Um servidor processa um fluxo de $10^9$ requisições e, ao final, deve informar as 100 mais demoradas, usando pouca memória.
\end{exercicio}
\begin{exercicio}{rv-reais}[Complementar]
Ordenar $10^6$ números reais em $[0,1)$ cuja distribuição é desconhecida e possivelmente muito concentrada perto de 0.
\end{exercicio}
\begin{exercicio}{rv-prefixo}[Complementar]
Um programa responde a $q$ perguntas ``qual a soma de $V[i..j]$?'' sobre um vetor fixo, somando os elementos de $i$ a $j$ a cada pergunta.
\end{exercicio}
\begin{exercicio}{rv-busca}[Complementar]
A cada consulta, um programa percorre uma lista de $n$ registros para descobrir se existe um registro com uma dada chave; há muitas consultas, inserções e remoções, e também é preciso, de vez em quando, listar os registros em ordem de chave.
\end{exercicio}
\begin{exercicio}{rv-desfazer}[Complementar]
Um editor de redes permite ligar computadores, consultar se dois computadores se comunicam e desfazer as últimas ligações feitas.
\end{exercicio}


\part{Paradigmas de projeto e algoritmos em grafos}
% ================================================================
\chapter{Programação dinâmica}\label{cap:pd}\label{sec:pd}
% ================================================================
\begin{objetivos}
\begin{itemize}
  \item A ideia da programação dinâmica: subestrutura ótima, sobreposição de subproblemas, memorização e cálculo de baixo para cima.
  \item Um roteiro para projetar algoritmos de programação dinâmica.
  \item Três problemas clássicos: subsequência comum máxima, ordem de multiplicação de matrizes e mochila 0-1.
  \item Por que o algoritmo da mochila é pseudopolinomial.
\end{itemize}
\end{objetivos}

Na divisão e conquista, os subproblemas são disjuntos. Em muitos problemas, porém, a decomposição natural gera subproblemas que \emph{se repetem}, e a recursão ingênua refaz o mesmo trabalho um número exponencial de vezes. A programação dinâmica resolve cada subproblema distinto uma única vez e guarda a resposta numa tabela. Este capítulo apresenta a técnica e três aplicações; uma quarta, o algoritmo de Floyd-Warshall, é o assunto do próximo capítulo.

\section{A técnica}
No método geral de decomposição, subproblemas idênticos podem aparecer muitas vezes ao longo do processo. Se cada aparição for resolvida de novo, o algoritmo pode ser exponencial mesmo que o número de subproblemas \emph{distintos} seja polinomial. A \textbf{programação dinâmica} mantém uma tabela com uma entrada para cada subproblema distinto: na primeira vez em que um subproblema é considerado, ele é resolvido normalmente; nas vezes seguintes, sua solução é simplesmente lida da tabela.
\begin{itemize}
  \item \textbf{Divisão e conquista ou programação dinâmica?} A divisão e conquista funciona bem quando os subproblemas são \textbf{disjuntos}, como no mergesort. Quando eles \textbf{compartilham} subproblemas, é melhor trocar a abordagem de cima para baixo por uma de baixo para cima: resolver dos subproblemas menores para os maiores, de modo que, ao chegar a um subproblema, todos aqueles de que ele depende já estejam resolvidos.
  \item \textbf{Subestrutura ótima}: uma solução ótima é composta de soluções ótimas de subproblemas. A prova típica é o argumento de ``recortar e colar'': se uma parte não fosse ótima, trocá-la por uma melhor daria uma solução global melhor, o que é uma contradição.
  \item \textbf{Sobreposição de subproblemas}: a recursão ingênua revisita os mesmos subproblemas; é isso que faz a tabela valer a pena.
  \item \textbf{Memorização} (de cima para baixo): a recursão original, acrescida de ``se já calculei este subproblema, devolvo o valor guardado''. \textbf{Cálculo de baixo para cima}: laços que preenchem a tabela na ordem certa. As duas têm a mesma complexidade; a segunda evita a pilha de recursão e, às vezes, permite economizar memória.
\end{itemize}

\begin{pseudo}{Fibonacci: recursão ingênua $\Theta(\varphi^n)$; memorização $O(n)$; de baixo para cima, $O(n)$ com memória $O(1)$}
\begin{algorithmic}
\Function{FibMemo}{$n$} \Comment{$F[\,]$ começa preenchido com $-1$}
  \If{$F[n]=-1$}
    \If{$n\le1$} \State $F[n]:=n$ \Else \State $F[n]:=$ \Call{FibMemo}{$n-1$} $+$ \Call{FibMemo}{$n-2$} \EndIf
  \EndIf
  \State \Return $F[n]$
\EndFunction
\Function{FibIterativo}{$n$}
  \State $a:=0$;\ $b:=1$ \Comment{$a=F_{i}$ e $b=F_{i+1}$}
  \For{$i:=1$ \textbf{até} $n$} \State $(a,b):=(b,\ a+b)$ \EndFor
  \State \Return $a$
\EndFunction
\end{algorithmic}
\end{pseudo}

\begin{roteiro}[Como projetar um algoritmo de programação dinâmica]
\begin{enumerate}
  \item \textbf{Defina o subproblema em palavras}, por exemplo ``$C(i,j)=$ tamanho da subsequência comum máxima de $X_i$ e $Y_j$''. Sem essa frase, a recorrência não tem sentido.
  \item \textbf{Decomponha}: o que uma solução ótima faz com o \emph{último} elemento, ou qual é a \emph{primeira} decisão? Cada possibilidade vira um termo da recorrência.
  \item \textbf{Escreva a recorrência e os casos base}.
  \item \textbf{Calcule a complexidade} como o número de subproblemas vezes o custo de cada um.
  \item \textbf{Escolha a ordem de preenchimento}: cada célula só pode depender de células já calculadas (linha a linha, por diagonais, por tamanho crescente\dots).
  \item \textbf{Indique onde está a resposta}, como $C[n,m]$ ou $m[1,n]$.
  \item \textbf{Reconstrua a solução}: guarde a escolha feita em cada célula, ou recalcule-a percorrendo a tabela de volta.
\end{enumerate}
\end{roteiro}

% ------------------------------------------------------------------
\section{Subsequência comum máxima}\label{sec:lcs}
Uma sequência $Z$ é \textbf{subsequência} de $X=x_1\dots x_n$ se os seus símbolos aparecem em $X$ na mesma ordem, \textbf{não necessariamente contíguos}. Por exemplo, $6,10,12$ é subsequência de $3,6,9,10,12,14,6$, mas não de $3,9,10,12,14,6$. Dadas $X$, com $|X|=n$, e $Y$, com $|Y|=m$, o problema da \textbf{subsequência comum máxima} (LCS, do inglês \emph{longest common subsequence}) pede uma subsequência comum a $X$ e a $Y$ de tamanho máximo. A força bruta testaria as $2^n$ subsequências de $X$.

\textbf{Decomposição.} Seja $X_i=x_1\dots x_i$ e olhe para os últimos símbolos. Se $x_n=y_m$, esse símbolo termina alguma subsequência comum máxima, e o restante dela é uma subsequência comum máxima de $X_{n-1}$ e $Y_{m-1}$. Se $x_n\ne y_m$, a resposta não usa $x_n$ ou não usa $y_m$. Com $C(i,j)=$ tamanho da subsequência comum máxima de $X_i$ e $Y_j$:
\[
C(i,j)=\begin{cases}0 & i=0\text{ ou }j=0\\ 1+C(i-1,j-1) & x_i=y_j\\ \max\{C(i-1,j),\,C(i,j-1)\} & x_i\ne y_j\end{cases}\qquad\text{resposta: }C(n,m).
\]
A recursão direta é exponencial, mas há apenas $(n+1)(m+1)$ subproblemas distintos.

\begin{pseudo}{LCS de baixo para cima, com reconstrução: $\Theta(nm)$ de tempo e memória}
\begin{algorithmic}
\Function{LCS}{$X,Y$}
  \For{$j:=0$ \textbf{até} $m$} \State $C[0,j]:=0$ \EndFor
  \For{$i:=1$ \textbf{até} $n$} \State $C[i,0]:=0$ \EndFor
  \For{$i:=1$ \textbf{até} $n$}
    \For{$j:=1$ \textbf{até} $m$}
      \If{$x_i=y_j$} \State $C[i,j]:=1+C[i-1,j-1]$
      \Else \State $C[i,j]:=\max\{C[i-1,j],\,C[i,j-1]\}$
      \EndIf
    \EndFor
  \EndFor
  \State $i:=n$;\ $j:=m$;\ $Z:=$ sequência vazia \Comment{reconstrução em $O(n+m)$}
  \While{$i>0$ \textbf{e} $j>0$}
    \If{$x_i=y_j$} \State coloque $x_i$ no \textbf{início} de $Z$;\ $i:=i-1$;\ $j:=j-1$
    \ElsIf{$C[i-1,j]\ge C[i,j-1]$} \State $i:=i-1$
    \Else \State $j:=j-1$
    \EndIf
  \EndWhile
  \State \Return $Z$ \Comment{$|Z|=C[n,m]$}
\EndFunction
\end{algorithmic}
\end{pseudo}

\begin{exemplo}[$X=\mathrm{ABCBDAB}$ e $Y=\mathrm{BDCABA}$]\leavevmode
\begin{center}\small
\begin{tabular}{cc|ccccccc}
 & & $j{=}0$ & 1 & 2 & 3 & 4 & 5 & 6\\
 & & & B & D & C & A & B & A\\ \hline
$i{=}0$ & & 0&0&0&0&0&0&0\\
1 & A & 0&0&0&0&1&1&1\\
2 & B & 0&\textbf{1}&1&1&1&2&2\\
3 & C & 0&1&1&\textbf{2}&2&2&2\\
4 & B & 0&1&1&2&2&\textbf{3}&3\\
5 & D & 0&1&2&2&2&3&3\\
6 & A & 0&1&2&2&3&3&\textbf{4}\\
7 & B & 0&1&2&2&3&4&4\\
\end{tabular}
\end{center}
$C[7,6]=4$. A reconstrução, subindo nos empates, passa pelas células em negrito: de $(7,6)$ sobe para $(6,6)$, onde casa \textbf{A}; vai para $(5,5)$, sobe para $(4,5)$, onde casa \textbf{B}; vai para $(3,4)$, anda para a esquerda até $(3,3)$, onde casa \textbf{C}; vai para $(2,2)$ e anda até $(2,1)$, onde casa \textbf{B}. Como cada símbolo é colocado no início, $Z=\mathrm{BCBA}$. As subsequências comuns de tamanho 4 são exatamente BCAB, BCBA e BDAB.
\end{exemplo}
Para $X=5,2,1,3,9,10,14$ e $Y=2,3,1,10$, o tamanho é 3, com as soluções $2,3,10$ e $2,1,10$.
\begin{itemize}
  \item Se só o \textbf{tamanho} interessa, basta memória $O(\min\{n,m\})$, pois cada linha só usa a anterior. Mas então não é possível reconstruir a subsequência percorrendo a tabela.
  \item Para obter \textbf{todas} as subsequências comuns máximas, siga, na reconstrução, os dois caminhos quando $x_i\ne y_j$ e $C[i-1,j]=C[i,j-1]$, eliminando repetições; pode haver exponencialmente muitas.
\end{itemize}

% ------------------------------------------------------------------
\section{Multiplicação de cadeia de matrizes}\label{sec:cadeia}
Queremos calcular $M=M_1\times M_2\times\dots\times M_n$, onde $M_i$ tem dimensões $r_{i-1}\times r_i$. Multiplicar uma matriz $p\times q$ por uma $q\times s$ custa $p\,q\,s$ multiplicações escalares. O produto de matrizes é \textbf{associativo}: a ordem em que as multiplicações são feitas não muda o resultado, mas pode mudar muito o custo. O problema é encontrar a parentização de custo mínimo.

\begin{exemplo}\label{ex:cadeia}
Sejam $M_1:20\times50$, $M_2:50\times10$, $M_3:10\times1$ e $M_4:1\times100$, isto é, $r=(20,50,10,1,100)$. A ordem $M_1(M_2(M_3M_4))$ custa $1000+50\,000+100\,000=\mathbf{151\,000}$ multiplicações; a ordem $((M_1M_2)M_3)M_4$ custa $10\,000+200+2000=\mathbf{12\,200}$. Nenhuma das duas é a melhor: a programação dinâmica abaixo encontra $(M_1(M_2M_3))M_4$, que custa $500+1000+2000=\mathbf{3500}$.
\end{exemplo}
A força bruta é exponencial: o número de parentizações de $n$ matrizes é o número de Catalan $C_{n-1}$ (que vale $1,1,2,5,14,42,132,\dots$), e cresce como $\Omega(4^n/n^{3/2})$.

\textbf{Decomposição.} A \emph{última} multiplicação de $M_i\cdots M_j$ é $(M_i\cdots M_k)\times(M_{k+1}\cdots M_j)$, para algum $i\le k<j$. Numa solução ótima, as duas partes também estão parentizadas de forma ótima, pelo argumento de recortar e colar. Com $m[i,j]=$ número mínimo de multiplicações escalares para calcular $M_i\cdots M_j$:
\[
m[i,j]=\begin{cases}0 & i=j\\[2pt] \min\limits_{i\le k<j}\big\{\,m[i,k]+m[k+1,j]+r_{i-1}\,r_k\,r_j\,\big\} & i<j\end{cases}\qquad\text{resposta: } m[1,n].
\]
Há $O(n^2)$ subproblemas e cada um testa $O(n)$ valores de $k$: o algoritmo custa \textbf{$\Theta(n^3)$ de tempo e $\Theta(n^2)$ de memória}. Como $m[i,j]$ depende de intervalos \emph{mais curtos}, a tabela é preenchida \textbf{por diagonais}, com $d=j-i=1,2,\dots,n-1$; a diagonal principal é toda zero.

\begin{pseudo}{Ordem ótima de multiplicação, com a tabela $s$ das escolhas}
\begin{algorithmic}
\Function{OrdemÓtima}{$r_0,\dots,r_n$}
  \For{$i:=1$ \textbf{até} $n$} \State $m[i,i]:=0$ \EndFor
  \For{$d:=1$ \textbf{até} $n-1$} \Comment{para cada diagonal}
    \For{$i:=1$ \textbf{até} $n-d$} \Comment{para cada linha da diagonal}
      \State $j:=i+d$;\ $m[i,j]:=\infty$
      \For{$k:=i$ \textbf{até} $j-1$}
        \State $q:=m[i,k]+m[k+1,j]+r_{i-1}\,r_k\,r_j$
        \If{$q<m[i,j]$} \State $m[i,j]:=q$;\ $s[i,j]:=k$ \EndIf
      \EndFor
    \EndFor
  \EndFor
  \State \Return $m[1,n]$ e $s$
\EndFunction
\Procedure{Imprime}{$s,i,j$}
  \If{$i=j$} \State imprima ``$M_i$''
  \Else \State imprima ``(''; \Call{Imprime}{$s,i,s[i,j]$}; \Call{Imprime}{$s,s[i,j]+1,j$}; imprima ``)''
  \EndIf
\EndProcedure
\Function{Multiplica}{$A,s,i,j$} \Comment{calcula $A_i\cdots A_j$ na ordem ótima}
  \If{$i=j$} \State \Return $A_i$ \EndIf
  \State \Return \Call{Multiplica}{$A,s,i,s[i,j]$} $\times$ \Call{Multiplica}{$A,s,s[i,j]+1,j$}
\EndFunction
\end{algorithmic}
\end{pseudo}

\begin{center}\small
\begin{tabular}{c|cccc}
$m[i,j]$ ($s[i,j]$) & $j=1$ & 2 & 3 & 4\\ \hline
$i=1$ & 0 & 10\,000 (1) & 1500 (1) & \textbf{3500 (3)}\\
2 & & 0 & 500 (2) & 5500 (3)\\
3 & & & 0 & 1000 (3)\\
4 & & & & 0
\end{tabular}
\end{center}
Na tabela do Exemplo~\ref{ex:cadeia}, $s[1,4]=3$ indica que a última multiplicação é $(M_1M_2M_3)\cdot M_4$, e $s[1,3]=1$ indica que a primeira parte é $M_1(M_2M_3)$. Logo, a ordem ótima é $((M_1(M_2M_3))M_4)$, com custo 3500. Outros exemplos:
\begin{center}\small
\begin{tabularx}{\linewidth}{@{}>{\raggedright\arraybackslash}p{4.2cm}X@{}}
\toprule
Dimensões $r$ & Custo ótimo e ordem\\
\midrule
$(10,100,5,50)$ & $\mathbf{7500}$: $((M_1M_2)M_3)$, com $m[1,2]=5000$ e $m[2,3]=25\,000$\\
$(5,10,3,12,5,50,6)$ & $\mathbf{2010}$: $((M_1M_2)((M_3M_4)(M_5M_6)))$\\
$(30,35,15,5,10,20,25)$ & $15\,125$: $((M_1(M_2M_3))((M_4M_5)M_6))$\\
\bottomrule
\end{tabularx}
\end{center}

% ------------------------------------------------------------------
\section{Mochila 0-1}\label{sec:mochila}
Há $n$ itens; o item $i$ tem peso $p_i$ e valor $v_i$, inteiros; a mochila tem capacidade $W$. Queremos $S\subseteq\{1,\dots,n\}$ com $\sum_{i\in S}p_i\le W$ que maximize $\sum_{i\in S}v_i$, usando cada item no máximo uma vez. Seja $M(w,i)$ o maior valor obtido com os itens $1,\dots,i$ e capacidade $w$. O item $i$ fica de fora, ou entra, se couber:
\[
M(w,i)=\begin{cases}0 & i=0\\ M(w,i-1) & p_i>w\\ \max\{M(w,i-1),\ v_i+M(w-p_i,\,i-1)\} & p_i\le w\end{cases}\qquad\text{resposta: }M(W,n).
\]
\begin{pseudo}{Mochila 0-1: $\Theta(nW)$ (pseudopolinomial)}
\begin{algorithmic}
\Function{Mochila}{$p,v,n,W$}
  \For{$w:=0$ \textbf{até} $W$} \State $M[w,0]:=0$ \EndFor
  \For{$i:=1$ \textbf{até} $n$}
    \For{$w:=0$ \textbf{até} $W$}
      \State $M[w,i]:=M[w,i-1]$
      \If{$p_i\le w$ \textbf{e} $v_i+M[w-p_i,i-1]>M[w,i]$} \State $M[w,i]:=v_i+M[w-p_i,i-1]$ \EndIf
    \EndFor
  \EndFor
  \State $w:=W$ \Comment{reconstrução}
  \For{$i:=n$ \textbf{decrescendo até} $1$}
    \If{$M[w,i]\ne M[w,i-1]$} \State o item $i$ está na solução;\ $w:=w-p_i$ \EndIf
  \EndFor
  \State \Return $M[W,n]$
\EndFunction
\end{algorithmic}
\end{pseudo}
\begin{exemplo}[$W=4$ e itens $(p,v)=(1,2),(3,7),(2,1)$]\leavevmode
\begin{center}\small
\begin{tabular}{c|cccc}
$w\backslash i$ & 0 & 1 & 2 & 3\\ \hline
0 & 0 & 0 & 0 & 0\\
1 & 0 & 2 & 2 & 2\\
2 & 0 & 2 & 2 & 2\\
3 & 0 & 2 & 7 & 7\\
4 & 0 & 2 & \textbf{9} & \textbf{9}
\end{tabular}
\end{center}
A solução ótima leva os itens 1 e 2, com peso 4 e valor \textbf{9}. Repare que $M(4,2)=\max\{M(4,1),\,7+M(1,1)\}=\max\{2,\,9\}=9$: é fácil errar essa célula esquecendo que, depois de colocar o item 2, ainda sobra capacidade 1 para o item 1.
\end{exemplo}
\begin{itemize}
  \item Com $W=50$ e itens $(p,v)=(10,60),(20,100),(30,120)$, o ótimo é $\mathbf{220}$, com os itens 2 e 3. O guloso que escolhe pela razão $v/p$ (que vale $6,5,4$) pega os itens 1 e 2 e obtém só 160: o algoritmo guloso \textbf{não} resolve a mochila 0-1, embora resolva a versão fracionária, em que frações de itens são permitidas (Capítulo~\ref{cap:gulosos}).
  \item \textbf{Por que pseudopolinomial?} O tempo $\Theta(nW)$ depende do \emph{valor} de $W$, mas a entrada guarda $W$ em apenas $\log_2W$ bits. Com $W=10^9$ (30 bits), o algoritmo faz $10^9$ iterações por item: é exponencial no tamanho da entrada. O problema da mochila é NP-difícil, e não se conhece algoritmo polinomial para ele.
  \item \textbf{Memória $O(W)$.} Um único vetor $K[0..W]$ basta: para cada item, percorra $w$ \textbf{decrescendo} de $W$ até $p_i$ e faça $K[w]:=\max\{K[w],K[w-p_i]+v_i\}$. Com $w$ \emph{crescendo}, o mesmo item poderia ser reaproveitado, e o algoritmo passaria a resolver a mochila \emph{ilimitada}.
\end{itemize}

\begin{cuidado}[Armadilhas da programação dinâmica]
\begin{itemize}
  \item \textbf{Subsequência não é subvetor}: a subsequência comum máxima não exige contiguidade; o subvetor de soma máxima (Capítulo~\ref{cap:dc}) exige.
  \item Escrever a recorrência sem casos base ou sem dizer o que $C(i,j)$ significa.
  \item Preencher a tabela na ordem errada: na cadeia de matrizes, preencher por \emph{linhas} $i=1,2,\dots$ usaria $m[i+1,j]$ ainda não calculado. Use diagonais, ou $i$ decrescente.
  \item Na cadeia de matrizes, o termo de custo é $r_{i-1}\,r_k\,r_j$ (e não $r_i r_k r_j$), e $k$ vai de $i$ a $j-1$.
  \item Afirmar que $O(nW)$ é polinomial: é pseudopolinomial.
\end{itemize}
\end{cuidado}

\begin{chave}
Programação dinâmica $=$ recorrência com subproblemas que se repetem $+$ tabela. Subsequência comum máxima em $\Theta(nm)$; cadeia de matrizes em $\Theta(n^3)$, por diagonais; mochila 0-1 em $\Theta(nW)$, pseudopolinomial; Floyd-Warshall em $\Theta(n^3)$ (Capítulo~\ref{cap:fw}). Sempre: significado do subproblema, recorrência, base, ordem de preenchimento, complexidade e reconstrução.
\end{chave}

% ----------------------------------------------------------------
\section{Exercícios}
\excurso
\begin{exercicio}{pd-max}[Lista]
Na variante da multiplicação de cadeia de matrizes em que se quer \emph{maximizar} o número de multiplicações escalares, vale a subestrutura ótima?
\end{exercicio}
\begin{exercicio}{pd-lis}[Lista]
Dê um algoritmo $O(n^2)$ para a maior subsequência estritamente crescente de $n$ números. Se puder, dê um $O(n\log n)$.
\end{exercicio}
\begin{exercicio}{pd-pal}[Lista]
Dê um algoritmo eficiente para o maior palíndromo que é \emph{subsequência} de uma string; para ``character'', a resposta é ``carac''. Qual o tempo?
\end{exercicio}
\begin{exercicio}{pd-dag}[Lista]
Num grafo acíclico dirigido com pesos reais nas arestas e dois vértices $s$ e $t$, descreva uma programação dinâmica para o caminho simples de maior peso de $s$ a $t$.
\end{exercicio}
\begin{exercicio}{pd-ilim}[Lista]
Resolva por programação dinâmica a mochila \emph{inteira} (ilimitada): cada item é um tipo e pode ser usado quantas vezes se quiser.
\end{exercicio}
\begin{exercicio}{pd-corte}[Lista]
Cortar uma string de $n$ caracteres em duas partes custa $n$. Dados os pontos de corte $L[1..m]$, determine a ordem de cortes de menor custo total e o custo. Para $n=20$ e cortes após os caracteres 2, 8 e 10, cortar da esquerda para a direita custa 50, e da direita para a esquerda, 38.
\end{exercicio}
\begin{exercicio}{pd-rec}[Slides]
Modifique o algoritmo da subsequência comum máxima para devolver uma subsequência, e não só o tamanho. Aplique a $X=\mathrm{ABCBDAB}$ e $Y=\mathrm{BDCABA}$.
\end{exercicio}
\begin{exercicio}{pd-mon}[Monitoria]
(a) Monte a tabela $m[i,j]$ para $r=(10,100,5,50)$ e para $r=(20,50,10,1,100)$, com a ordem ótima. (b) Monte a tabela da mochila com $W=4$ e itens $(p,v)=(1,2),(3,7),(2,1)$. (c) Explique por que $\Theta(nW)$ é pseudopolinomial.
\end{exercicio}
\begin{exercicio}{pd-apa}[Material de apoio]
(a) Encontre a maior subsequência comum de $X=\mathrm{dacdba}$ e $Y=\mathrm{cadcda}$. (b) Resolva a mochila 0-1 com capacidade 100 e itens (valor, peso) $=(50,20),(70,30),(90,40),(120,50),(140,60)$. (c) Dê um algoritmo linear para o subvetor contíguo de soma máxima e aplique-o a $5,15,-30,10,-5,40,10$.
\end{exercicio}
\begin{exercicio}{pd-l6a}[Material de apoio]
(a) Encontre a parentização ótima para as dimensões $\langle5,10,3,12,5,50,6\rangle$. (b) Adapte o algoritmo para calcular a tabela $s$ e escreva um algoritmo recursivo que multiplica $A_i\cdots A_j$ na ordem ótima. (c) Adapte o algoritmo da subsequência comum máxima para imprimir \emph{todas} as subsequências comuns mais longas.
\end{exercicio}
\begin{exercicio}{pd-l6b}[Material de apoio]
(a) Quantas árvores binárias de busca guardam 6 chaves diferentes? (b) Construa uma árvore binária de busca ótima para as palavras com números de acessos \emph{case} (15), \emph{else} (5), \emph{end} (5), \emph{if} (35), \emph{of} (5) e \emph{then} (35), em que o custo é a soma, sobre as palavras, de (acessos)$\times$(profundidade $+1$).
\end{exercicio}
\begin{exercicio}{pd-l6c}[Material de apoio]
(a) Subsequência crescente de \emph{soma} máxima em $O(n^2)$. (b) Festa da empresa: numa árvore hierárquica com notas de sociabilidade, escolha convidados sem convidar ninguém junto com o seu chefe imediato, maximizando a soma das notas. (c) Torres gêmeas: dadas duas sequências de raios de lajotas, encontre a maior torre comum obtida removendo lajotas e mantendo a ordem. (d) Bandeiras: quantas sequências de $n$ faixas brancas, azuis e vermelhas existem sem duas faixas de mesma cor vizinhas e com toda faixa azul entre uma branca e uma vermelha?
\end{exercicio}
\begin{exercicio}{pd-mc458}[Material de apoio]
(a) Mochila exata (\emph{subset sum}) com pesos $2,7,5,2$: complete a tabela (``O'' $=$ há solução sem o item; ``I'' $=$ só usando o item; ``--'' $=$ não há solução) para os itens 3 e 4 e capacidades $0..14$. (b) Como reduzir a memória da mochila para $O(K)$? (c) Encontre a submatriz consecutiva de soma máxima numa matriz $m\times n$. (d) Corte uma laje de mármore $n\times m$, com cortes horizontais ou verticais de ponta a ponta, para maximizar o lucro, dados os preços $p_{ab}$ de cada retângulo $a\times b$, em $O(n^2m^2)$.
\end{exercicio}

\excompl
\begin{exercicio}{pd-edit}[Complementar]
A \emph{distância de edição} entre duas strings é o menor número de inserções, remoções e substituições de caracteres que transforma uma na outra. Dê a recorrência e a complexidade, e calcule a distância entre ``kitten'' e ``sitting''.
\end{exercicio}
\begin{exercicio}{pd-troco}[Complementar]
Com moedas de valores $\{1,3,4\}$, qual o menor número de moedas para dar troco de 6? O que faz o algoritmo guloso? Escreva a programação dinâmica e preencha $T(0..10)$.
\end{exercicio}
\begin{exercicio}{pd-haste}[Complementar]
Uma haste de comprimento $n$ pode ser cortada em pedaços de comprimentos inteiros; um pedaço de comprimento $i$ vale $p_i$. Com $p=(1,5,8,9,10,17,17,20)$ para $i=1,\dots,8$, qual o maior valor obtido de uma haste de comprimento 8? Dê a recorrência.
\end{exercicio}
\begin{exercicio}{pd-grade}[Complementar]
Num tabuleiro de $4\times4$ casas, quantos caminhos levam uma peça da casa do canto superior esquerdo à do canto inferior direito, andando uma casa por vez para a direita ou para baixo? E se a casa da linha 2, coluna 2 estiver bloqueada?
\end{exercicio}
\begin{exercicio}{pd-part}[Complementar]
Decida por programação dinâmica se um multiconjunto de inteiros positivos pode ser dividido em duas partes de mesma soma. Teste com $\{3,1,1,2,2,1\}$ e $\{2,3,7\}$. Qual a complexidade?
\end{exercicio}
\begin{exercicio}{pd-catalan}[Complementar]
Seja $P(n)$ o número de parentizações de um produto de $n$ matrizes. Mostre que $P(1)=1$ e $P(n)=\sum_{k=1}^{n-1}P(k)P(n-k)$, e calcule $P(n)$ para $n\le6$.
\end{exercicio}
\begin{exercicio}{pd-lcsmem}[Complementar]
Escreva uma versão do algoritmo da subsequência comum máxima que calcula apenas o tamanho usando memória $O(\min\{n,m\})$.
\end{exercicio}
\begin{exercicio}{pd-formas}[Complementar]
De quantas maneiras é possível dar troco de 5 com moedas de valores $\{1,2,5\}$, sem considerar a ordem das moedas? E de 10? Escreva a programação dinâmica que conta as maneiras.
\end{exercicio}
\begin{exercicio}{pd-substring}[Complementar]
O maior \emph{substring} comum de duas strings exige que os símbolos sejam contíguos nas duas. Dê uma programação dinâmica em $O(nm)$ e aplique-a a ``ABABC'' e ``BABCA''.
\end{exercicio}

% ================================================================
\chapter{Caminhos mínimos entre todos os pares: Floyd-Warshall}\label{cap:fw}\label{sec:fw}
% ================================================================
\begin{objetivos}
\begin{itemize}
  \item Notação e representações de grafos usadas no restante do livro.
  \item O problema dos caminhos mínimos entre todos os pares de vértices.
  \item O algoritmo de Floyd-Warshall como programação dinâmica, com reconstrução dos caminhos.
  \item Pesos negativos, ciclos negativos e a comparação com o algoritmo de Dijkstra.
\end{itemize}
\end{objetivos}

O algoritmo de Floyd-Warshall é mais uma aplicação da programação dinâmica, agora sobre um grafo: queremos o custo do caminho mais barato entre \emph{todos} os pares de vértices. Por ser o primeiro algoritmo em grafos do livro, o capítulo começa fixando a notação e as representações de grafos usadas daqui em diante.

\section{Grafos: notação e representação}\label{sec:buscas}
Um grafo é um par $G=(V,E)$, em que $V$ é o conjunto de vértices e $E$ o de arestas; escrevemos $n=|V|$ e $m=|E|$. A \textbf{vizinhança} de um vértice $v$, $N(v)$, é o conjunto dos vértices ligados a $v$ por uma aresta. Num \textbf{digrafo} (grafo direcionado), cada aresta $v\to w$ tem uma direção, e $N_{out}(v)$ é o conjunto dos vértices $w$ com uma aresta $v\to w$. Há duas formas usuais de representar um grafo na memória:
\begin{center}\small
\begin{tabular}{@{}lccc@{}}
\toprule
Representação & Memória & ``$vw\in E$?'' & Percorrer $N(v)$\\
\midrule
Matriz de adjacência (ou de pesos $W$) & $\Theta(n^2)$ & $O(1)$ & $\Theta(n)$\\
Listas de adjacência & $\Theta(n+m)$ & $O(\text{grau}(v))$ & $\Theta(\text{grau}(v))$\\
\bottomrule
\end{tabular}
\end{center}
A matriz é adequada a grafos densos, com $m$ próximo de $n^2$, e a testes frequentes de adjacência; as listas, a grafos esparsos e a algoritmos que percorrem vizinhanças. Uma \textbf{busca}, que visita todos os vértices e arestas do grafo examinando cada vizinhança uma vez, custa $\Theta(n+m)$ com listas e $\Theta(n^2)$ com matriz; as buscas em largura e em profundidade são estudadas nos Capítulos~\ref{cap:bfs} e~\ref{cap:dfs}.

\section{Matriz de pesos}\label{sec:fwmatriz}
Em problemas de caminhos mínimos, cada aresta tem um peso (um custo, uma distância, um tempo), e o grafo é dado por uma \textbf{matriz de pesos} $W$: $W(x,x)=0$; $W(x,y)$ é o peso da aresta $xy$, se ela existe; e $W(x,y)=\infty$ se não existe. O custo de um caminho é a soma dos pesos das suas arestas, e a \textbf{distância} $\operatorname{dist}(x,y)$ é o menor custo de um caminho de $x$ a $y$. Um grafo não direcionado é tratado como um digrafo em que cada aresta $xy$ de peso $p$ dá origem às arestas $x\to y$ e $y\to x$, ambas de peso $p$.

Usaremos como exemplo, neste capítulo e no Capítulo~\ref{cap:dijkstra}, o digrafo de cinco vértices dado pela matriz abaixo; os vértices são chamados $1,\dots,5$ aqui e $a,\dots,e$ no capítulo do algoritmo de Dijkstra.
\begin{center}\small
\begin{tabular}{c|ccccc}
$W$ & 1\,($a$) & 2\,($b$) & 3\,($c$) & 4\,($d$) & 5\,($e$)\\ \hline
1\,($a$) & 0 & 1 & 3 & $\infty$ & 6\\
2\,($b$) & $\infty$ & 0 & 1 & 3 & $\infty$\\
3\,($c$) & 1 & 2 & 0 & 1 & $\infty$\\
4\,($d$) & 3 & $\infty$ & $\infty$ & 0 & 2\\
5\,($e$) & $\infty$ & $\infty$ & $\infty$ & 1 & 0
\end{tabular}
\end{center}

\section{O algoritmo de Floyd-Warshall}
\textbf{Problema:} dado um digrafo com vértices $1,\dots,n$ e matriz de pesos $W$, calcular $\operatorname{dist}(i,j)$ para todos os pares ordenados $(i,j)$. A ideia é limitar, passo a passo, quais vértices podem aparecer \emph{no meio} dos caminhos:
\[
D^k(i,j)=\text{custo do caminho mínimo de }i\text{ a }j\text{ cujos vértices \emph{intermediários} estão em }\{1,\dots,k\}.
\]
Então $D^0=W$, pois sem intermediários só valem as arestas diretas, e $D^n$ é a matriz das distâncias.

\textbf{Decomposição.} Um caminho mínimo de $i$ a $j$ com intermediários em $\{1,\dots,k\}$ ou \textbf{não passa} por $k$, e custa $D^{k-1}(i,j)$, ou passa \textbf{exatamente uma vez} por $k$ --- voltar a $k$ formaria um ciclo, que não ajuda --- e se divide num caminho mínimo de $i$ a $k$ e num de $k$ a $j$, ambos com intermediários em $\{1,\dots,k-1\}$:
\[
D^k(i,j)=\begin{cases}W(i,j) & k=0\\ \min\{\,D^{k-1}(i,j),\ D^{k-1}(i,k)+D^{k-1}(k,j)\,\} & k\ge1.\end{cases}
\]

\textbf{Por que só ``prefixos'' $\{1,\dots,k\}$?} Uma formulação mais livre usaria, para cada conjunto $S\subseteq V$ de intermediários permitidos, $D(i,j,S)=\min_{v\in S}\{D(i,v,S\setminus\{v\})+D(v,j,S\setminus\{v\})\}$. Ela também é correta, mas tem $n^2\,2^n$ subproblemas. Fixando uma ordem dos vértices e usando só os $n+1$ prefixos, sobram $O(n^2)$ subproblemas por valor de $k$, e $O(n^3)$ no total.

\begin{pseudo}{Floyd-Warshall com matriz de predecessores $P$: $\Theta(n^3)$ de tempo e $\Theta(n^2)$ de memória}
\begin{algorithmic}
\Function{FloydWarshall}{$W,n$}
  \For{$i,j:=1$ \textbf{até} $n$}
    \State $D[i,j]:=W[i,j]$
    \If{$i\ne j$ \textbf{e} $W[i,j]<\infty$} \State $P[i,j]:=i$ \Else \State $P[i,j]:=\nil$ \EndIf
  \EndFor
  \For{$k:=1$ \textbf{até} $n$} \Comment{$k$ é o laço \textbf{externo}}
    \For{$i:=1$ \textbf{até} $n$}
      \For{$j:=1$ \textbf{até} $n$}
        \If{$D[i,k]+D[k,j]<D[i,j]$}
          \State $D[i,j]:=D[i,k]+D[k,j]$;\ $P[i,j]:=P[k,j]$ \Comment{o penúltimo vértice vem do trecho de $k$ a $j$}
        \EndIf
      \EndFor
    \EndFor
  \EndFor
  \State \Return $D,P$
\EndFunction
\Procedure{Caminho}{$P,i,j$}
  \If{$i=j$} \State imprima $i$
  \ElsIf{$P[i,j]=\nil$} \State imprima ``não há caminho''
  \Else \State \Call{Caminho}{$P,i,P[i,j]$};\ imprima $j$
  \EndIf
\EndProcedure
\end{algorithmic}
\end{pseudo}
A matriz $P$ guarda, em $P^k(i,j)$, o \textbf{penúltimo} vértice do caminho mínimo de $i$ a $j$ com intermediários em $\{1,\dots,k\}$. Se o caminho não melhora ao permitir $k$, $P^k(i,j)=P^{k-1}(i,j)$; se melhora, o novo caminho termina com o trecho de $k$ a $j$, e $P^k(i,j)=P^{k-1}(k,j)$. Uma única matriz $D$ basta, com atualizações ``no lugar'': na iteração $k$, a linha e a coluna $k$ não mudam, pois $D(k,k)=0$.

\begin{exemplo}[O digrafo da §\ref{sec:fwmatriz}]\leavevmode
\begin{center}\small
\begin{tabular}{c|l}
$k$ & entradas que diminuem ao permitir $k$ como intermediário\\ \hline
1 & $D(3,5)\!:\infty\to7$;\ $D(4,2)\!:\infty\to4$;\ $D(4,3)\!:\infty\to6$\\
2 & $D(1,3)\!:3\to2$;\ $D(1,4)\!:\infty\to4$;\ $D(4,3)\!:6\to5$\\
3 & $D(1,4)\!:4\to3$;\ $D(2,1)\!:\infty\to2$;\ $D(2,4)\!:3\to2$;\ $D(2,5)\!:\infty\to8$\\
4 & $D(1,5)\!:6\to5$;\ $D(2,5)\!:8\to4$;\ $D(3,5)\!:7\to3$;\ $D(5,1)\!:\infty\to4$;\ $D(5,2)\!:\infty\to5$;\ $D(5,3)\!:\infty\to6$\\
5 & nenhuma
\end{tabular}
\end{center}
\begin{center}\small
\begin{tabular}{c|ccccc}
$D^5$ & 1 & 2 & 3 & 4 & 5\\ \hline
1 & 0 & 1 & 2 & 3 & 5\\
2 & 2 & 0 & 1 & 2 & 4\\
3 & 1 & 2 & 0 & 1 & 3\\
4 & 3 & 4 & 5 & 0 & 2\\
5 & 4 & 5 & 6 & 1 & 0
\end{tabular}
\qquad
\begin{tabular}{c|ccccc}
$P^5$ & 1 & 2 & 3 & 4 & 5\\ \hline
1 & -- & 1 & 2 & 3 & 4\\
2 & 3 & -- & 2 & 3 & 4\\
3 & 3 & 3 & -- & 3 & 4\\
4 & 4 & 1 & 2 & -- & 4\\
5 & 4 & 1 & 2 & 5 & --
\end{tabular}
\end{center}
O caminho de 5 a 3 sai da matriz $P$ de trás para frente: $P(5,3)=2$, $P(5,2)=1$, $P(5,1)=4$ e $P(5,4)=5$. Logo, o caminho é $5\to4\to1\to2\to3$, com custo $1+3+1+1=6=D^5(5,3)$.
\end{exemplo}
\begin{itemize}
  \item \textbf{Pesos negativos.} O algoritmo funciona com arestas de peso negativo desde que \textbf{não haja ciclo negativo}, pois então o caminho mínimo nunca repete vértices. Há um ciclo negativo se e somente se, ao final, algum $D(i,i)$ é negativo.
  \item \textbf{Floyd-Warshall ou $n$ execuções do Dijkstra?} O Floyd-Warshall custa $\Theta(n^3)$ e é muito simples. Executar o algoritmo de Dijkstra (Capítulo~\ref{cap:dijkstra}) a partir de cada vértice custa $O(n\,(n+m)\log n)$ com heap binário, ou $O(n^3)$ com vetor. Em grafos esparsos e sem pesos negativos, as $n$ execuções do Dijkstra são mais rápidas; em grafos densos, ou com pesos negativos, o Floyd-Warshall é preferível.
\end{itemize}

\begin{cuidado}[Armadilhas do Floyd-Warshall]
\begin{itemize}
  \item Colocar o laço de $k$ \textbf{por dentro} ($i$ e $j$ por fora) dá resposta errada: $k$ precisa ser o laço externo, pois $D^k$ só pode ser calculada depois de toda a $D^{k-1}$.
  \item $D^k(i,j)$ restringe os vértices \emph{intermediários} a $\{1,\dots,k\}$; os extremos $i$ e $j$ podem ser quaisquer.
  \item $P(i,j)$ é o \textbf{penúltimo} vértice do caminho, e não o segundo; a reconstrução vai de trás para frente.
  \item Com ciclo negativo, ``caminho mínimo'' deixa de fazer sentido; o algoritmo detecta a situação por $D(i,i)<0$.
\end{itemize}
\end{cuidado}

\begin{chave}
$D^k(i,j)=\min\{D^{k-1}(i,j),\ D^{k-1}(i,k)+D^{k-1}(k,j)\}$, com $D^0=W$; $k$ é o laço externo; $\Theta(n^3)$ de tempo e $\Theta(n^2)$ de memória, com uma única matriz. Quando o caminho por $k$ melhora, $P(i,j)$ recebe $P(k,j)$. Usar só os prefixos $\{1,\dots,k\}$, e não subconjuntos quaisquer, é o que torna o algoritmo polinomial.
\end{chave}

% ----------------------------------------------------------------
\section{Exercícios}
\excurso
\begin{exercicio}{fw-sub}[Slides]
Uma primeira formulação usa $D(i,j,S)=$ custo do caminho mínimo de $i$ a $j$ com intermediários em $S\subseteq V$, com $D(i,j,\emptyset)=W(i,j)$. Escreva a recorrência. Quantos subproblemas ela tem? O que o algoritmo de Floyd-Warshall faz para evitar isso?
\end{exercicio}
\begin{exercicio}{fw-slide}[Slides]
Para a matriz de pesos $W$ da §\ref{sec:fwmatriz}, calcule $D^1,\dots,D^5$ e $P^5$ e reconstrua o caminho mínimo de 5 a 3 e o de 2 a 1. Compare com o Exemplo da §\ref{sec:fw}.
\end{exercicio}

\excompl
\begin{exercicio}{fw-neg}[Complementar]
Execute o Floyd-Warshall no digrafo com arestas $1\to2$ (3), $1\to3$ (8), $1\to5$ ($-4$), $2\to4$ (1), $2\to5$ (7), $3\to2$ (4), $4\to1$ (2), $4\to3$ ($-5$) e $5\to4$ (6). Dê a matriz final. Há ciclo negativo?
\end{exercicio}
\begin{exercicio}{fw-fecho}[Complementar]
O \emph{fecho transitivo} de um digrafo diz, para cada par $(i,j)$, se existe caminho de $i$ a $j$. Adapte o Floyd-Warshall, usando ``ou'' e ``e'' no lugar de $\min$ e $+$. Calcule o fecho para os vértices $1,\dots,5$ com as arestas $1\to2$, $2\to3$, $3\to1$ e $3\to4$.
\end{exercicio}
\begin{exercicio}{fw-ciclo}[Complementar]
Como encontrar o peso do \emph{ciclo} direcionado de menor peso de um digrafo com pesos positivos usando o Floyd-Warshall? Aplique ao digrafo com as arestas $1\to2$ (2), $2\to3$ (3), $3\to1$ (1), $3\to4$ (5) e $4\to3$ (2).
\end{exercicio}
\begin{exercicio}{fw-negciclo}[Complementar]
Prove que, se o digrafo tem um ciclo de peso negativo, ao final do Floyd-Warshall vale $D(i,i)<0$ para algum $i$ do ciclo; e que, se não tem, $D(i,i)=0$ para todo $i$.
\end{exercicio}
\begin{exercicio}{fw-dij}[Complementar]
Um grafo tem $n=10^4$ vértices, $m=5\cdot10^4$ arestas e pesos não negativos. Para obter as distâncias entre todos os pares, compare o custo do Floyd-Warshall com o de $n$ execuções do Dijkstra com heap binário. E se $m\approx n^2/2$?
\end{exercicio}
\begin{exercicio}{fw-4v}[Complementar]
Execute o Floyd-Warshall no digrafo com as arestas $1\to2$ (4), $1\to4$ (10), $2\to3$ (3), $3\to4$ (1), $4\to1$ (2) e $3\to1$ (7), mostrando $D^1,\dots,D^4$.
\end{exercicio}
\begin{exercicio}{fw-4vcam}[Complementar]
No exercício anterior, calcule a matriz $P^4$ e reconstrua os caminhos mínimos de 2 a 1 e de 4 a 3.
\end{exercicio}
\begin{exercicio}{fw-negex}[Complementar]
Dê um digrafo com dois vértices e um ciclo negativo e mostre o que o Floyd-Warshall calcula. Por que os valores de $D$ deixam de ser custos de caminhos?
\end{exercicio}
\begin{exercicio}{fw-nd}[Complementar]
Como aplicar o Floyd-Warshall a um grafo não direcionado? O que se pode dizer da matriz $D$ final? Aresta de peso negativo num grafo não direcionado é um problema? Por quê?
\end{exercicio}
\begin{exercicio}{fw-conta}[Complementar]
Quantas vezes, exatamente, o teste ``$D[i,k]+D[k,j]<D[i,j]$'' é executado? Quanta memória o algoritmo usa com e sem a matriz $P$?
\end{exercicio}
\begin{exercicio}{fw-inplace}[Complementar]
Prove que usar uma única matriz $D$, atualizada no lugar, dá o mesmo resultado que calcular $D^0,D^1,\dots,D^n$ separadamente (suponha que não há ciclos negativos).
\end{exercicio}
\begin{exercicio}{fw-largura}[Complementar]
Cada aresta tem uma \emph{capacidade}, e a capacidade de um caminho é a menor capacidade das suas arestas. Adapte o Floyd-Warshall para calcular, para todo par $(i,j)$, a maior capacidade de um caminho de $i$ a $j$.
\end{exercicio}
\begin{exercicio}{fw-diam}[Complementar]
A \emph{excentricidade} de um vértice $v$ é a maior distância de $v$ a outro vértice, e o \emph{diâmetro} do digrafo é a maior excentricidade. Use a matriz $D^5$ do Exemplo da §\ref{sec:fw} para calcular as excentricidades, o diâmetro e o \emph{centro} (o vértice de menor excentricidade).
\end{exercicio}
\begin{exercicio}{fw-aresta}[Complementar]
Dada a matriz $D$ final, como decidir se a aresta $uv$ pertence a algum caminho mínimo de $i$ a $j$? E a algum caminho mínimo qualquer do grafo?
\end{exercicio}
\begin{exercicio}{fw-atualiza}[Complementar]
Depois de calcular $D$, o peso da aresta $uv$ diminui para $w'$. Mostre como atualizar todas as distâncias em $O(n^2)$.
\end{exercicio}
\begin{exercicio}{fw-novo}[Complementar]
Depois de calcular $D$ para os vértices $1,\dots,n$, acrescenta-se um vértice $n+1$ com algumas arestas de e para ele. Mostre como atualizar a matriz em $O(n^2)$.
\end{exercicio}
\begin{exercicio}{fw-contar}[Complementar]
Com pesos positivos, adapte o Floyd-Warshall para calcular, além da distância, o \emph{número} de caminhos mínimos distintos de $i$ a $j$, para todo par.
\end{exercicio}
\begin{exercicio}{fw-bits}[Complementar \estrela]
Mostre que o fecho transitivo pode ser calculado em $O(n^3/w)$ operações, representando cada linha da matriz como um vetor de bits e operando sobre palavras de $w$ bits.
\end{exercicio}

% ================================================================
\chapter{Algoritmos gulosos e matroides}\label{cap:gulosos}
% ================================================================
\begin{objetivos}
\begin{itemize}
  \item O método guloso e como provar que ele dá a solução ótima, pelo argumento de troca.
  \item Matroides: definição, exemplos e formas equivalentes do axioma de aumento.
  \item O teorema que caracteriza as matroides como as estruturas em que o guloso é sempre ótimo.
  \item Limites do método guloso e uma visão da interseção de matroides.
\end{itemize}
\end{objetivos}

Um algoritmo guloso constrói a solução escolhendo, a cada passo, o que parece melhor naquele momento, sem nunca voltar atrás. É rápido e simples, mas nem sempre está correto. A teoria das matroides explica exatamente quando o guloso acerta: as matroides são as estruturas em que ele é ótimo para quaisquer pesos. O exemplo mais importante, a matroide gráfica, leva diretamente ao algoritmo de Kruskal do próximo capítulo.

\section{O método guloso}\label{sec:guloso}
No esquema geral, queremos um subconjunto $S'\subseteq S$ que satisfaça uma propriedade $P$ e seja máximo (ou mínimo) segundo algum critério. O algoritmo guloso começa com $S'_0=\emptyset$ e, a cada passo, acrescenta o elemento $s$ que mantém a propriedade $P$ e é o \textbf{melhor segundo o critério naquele momento}, sem jamais desfazer uma escolha. Ele para quando nenhum elemento pode mais ser acrescentado.
\begin{itemize}
  \item \textbf{Quando o guloso é ótimo?} Quando vale a \textbf{propriedade gulosa}: existe uma solução ótima que contém o elemento escolhido, e, fixado esse elemento, o que resta é um subproblema bem definido do mesmo tipo (subestrutura ótima).
  \item \textbf{Como provar?} Pelo \emph{argumento de troca}: tome uma solução ótima $O$ que concorde com a solução gulosa o máximo possível; no primeiro ponto em que elas discordam, troque um elemento de $O$ pelo escolhido pelo guloso, sem piorar a solução. Isso contradiz a escolha de $O$, ou mostra que a solução gulosa também é ótima.
  \item \textbf{O que os problemas resolvidos por gulosos têm em comum?} Em grande parte deles, a família dos conjuntos viáveis é uma \textbf{matroide} (§\ref{sec:matroides}).
  \item \textbf{Programação dinâmica ou guloso?} A programação dinâmica examina todas as escolhas possíveis de cada subproblema, guardando os resultados numa tabela; o guloso faz uma única escolha e nunca volta atrás. O guloso é mais rápido, mas só funciona quando vale a propriedade gulosa.
\end{itemize}

\begin{intuicao}[O argumento de troca]
Para mostrar que a escolha gulosa não estraga nada, imagine alguém que conhece uma solução ótima e discorda do guloso no primeiro passo. Se for possível ``trocar'' um elemento da solução ótima pelo elemento guloso, sem perder qualidade e sem violar as restrições, então existe uma solução ótima que começa como o guloso. Repetindo o argumento passo a passo (ou por indução), a solução gulosa inteira é ótima. Quando a troca falha --- quando o elemento guloso ``ocupa espaço'' que a solução ótima usaria melhor ---, o guloso costuma estar errado, e é aí que se procura um contraexemplo.
\end{intuicao}

\begin{pseudo}{Troco guloso}
\begin{algorithmic}
\Function{TrocoGuloso}{$S,\ v_1>v_2>\dots>v_n$} \Comment{moedas em ordem decrescente}
  \State $total:=0$
  \For{$i:=1$ \textbf{até} $n$}
    \State $q_i:=\floor{S/v_i}$;\ $S:=S-q_i\,v_i$;\ $total:=total+q_i$ \Comment{o máximo possível da maior moeda}
  \EndFor
  \State \Return $total$ \Comment{se $S\ne0$ ao final, não há troco exato}
\EndFunction
\end{algorithmic}
\end{pseudo}

\begin{exemplo}[Troco de 38 centavos]
Com as moedas $\{20,10,5,1\}$, o guloso dá $\floor{38/20}=1$ moeda de 20 (sobram 18), $\floor{18/10}=1$ moeda de 10 (sobram 8), $\floor{8/5}=1$ moeda de 5 (sobram 3) e três moedas de 1: seis moedas no total (Figura~\ref{fig:troco}). Cada passo é uma decisão local --- ``use a maior moeda que cabe'' ---, tomada sem olhar o que vem depois.
\begin{figura}[O troco guloso de 38 com as moedas $\{20,10,5,1\}$: em cada passo, a maior moeda que ainda cabe no valor restante.\label{fig:troco}]
\begin{tikzpicture}[moeda/.style={circle, draw, thick, fill=Goldenrod!35, minimum size=8mm, font=\small\bfseries}]
\foreach \pos/\val/\resto in {0/20/18,1.2/10/8,2.4/5/3,3.6/1/2,4.8/1/1,6/1/0} {
  \node[moeda] at (\pos,0) {\val};
  \node[rot] at (\pos,-0.75) {resta \resto};}
\node[passo, anchor=west] at (6.8,0) {6 moedas};
\end{tikzpicture}
\end{figura}
\end{exemplo}

\begin{exemplo}[O troco guloso é ótimo para as moedas $\{20,10,5,1\}$]
Seja $(a,b,c,d)$ o número de moedas de $20,10,5,1$ numa solução ótima. Então $d\le4$, pois cinco moedas de 1 podem ser trocadas por uma de 5; $c\le1$, pois duas de 5 viram uma de 10; e $b\le1$, pois duas de 10 viram uma de 20. Logo as moedas menores que 20 somam $10b+5c+d\le10+5+4=19<20$, e então $a=\floor{S/20}$, como no guloso. Do mesmo modo, $5c+d\le9<10$ força o número de moedas de 10 a ser o do guloso, e $d\le4<5$ força o de moedas de 5. A solução ótima \emph{é} a gulosa.
\end{exemplo}
O guloso, porém, \textbf{não funciona para qualquer sistema de moedas}. Com $\{6,4,1\}$ e $S=8$, ele dá $6+1+1$ (três moedas), mas $4+4$ usa duas. Com $\{25,10,1\}$ e $S=30$, ele usa 6 moedas, contra $10+10+10$. O caso geral se resolve por programação dinâmica, como a mochila ilimitada: $T(s)=1+\min_{v_i\le s}T(s-v_i)$, em $\Theta(nS)$.

% ------------------------------------------------------------------
\section{Matroides}\label{sec:matroides}
As matroides foram introduzidas por Whitney em 1935, com o objetivo de capturar as propriedades comuns à independência linear de vetores e à ausência de ciclos em grafos.
\begin{defi}[Matroide]
Uma \textbf{matroide} é um par $M=(E,\mathcal I)$, em que $E$ é um conjunto finito e $\mathcal I\subseteq2^E$ é uma família de subconjuntos, chamados \textbf{independentes}, tal que:
(I1) $\emptyset\in\mathcal I$;\quad
(I2) se $I\in\mathcal I$ e $I'\subseteq I$, então $I'\in\mathcal I$ (a família é fechada para subconjuntos);\quad
(I3) se $I_1,I_2\in\mathcal I$ e $|I_1|<|I_2|$, existe $e\in I_2\setminus I_1$ com $I_1\cup\{e\}\in\mathcal I$ (\textbf{axioma de aumento}).
\end{defi}
Em palavras: o vazio é independente; tirar elementos de um independente não o torna dependente; e um independente pequeno sempre pode ``pegar emprestado'' um elemento de qualquer independente maior.

\begin{exemplo}[Uma matroide pequena, de dois jeitos]\label{ex:matpequena}
Considere o grafo da Figura~\ref{fig:matpequena}: um triângulo com as arestas $a=12$, $b=23$ e $c=13$, e uma aresta pendente $d=34$. Na \textbf{matroide gráfica}, os independentes são os conjuntos de arestas sem ciclo. O único ciclo é $\{a,b,c\}$, então são independentes todos os $2^4=16$ subconjuntos, exceto $\{a,b,c\}$ e $\{a,b,c,d\}$: 14 independentes. Os independentes maximais, as \emph{bases}, são $\{a,b,d\}$, $\{a,c,d\}$ e $\{b,c,d\}$ --- exatamente as árvores geradoras do grafo, todas com 3 arestas.

Agora tome os vetores de $\R^3$ $a=(1,0,0)$, $b=(0,1,0)$, $c=(1,1,0)$ e $d=(0,0,1)$. Na \textbf{matroide linear}, os independentes são os conjuntos linearmente independentes. Como $c=a+b$, o conjunto $\{a,b,c\}$ é dependente, e qualquer conjunto que o contenha também; todos os outros são independentes. É \emph{a mesma} família de independentes: duas situações muito diferentes, ciclos num grafo e dependência linear, têm a mesma estrutura combinatória. É isso que a definição de matroide captura.
\begin{figura}[À esquerda, o grafo do Exemplo~\ref{ex:matpequena}; as arestas são os elementos da matroide. À direita, as suas três bases, que são as árvores geradoras.\label{fig:matpequena}]
\begin{tikzpicture}[x=1.1cm,y=1.1cm, baseline={(0,0.5)}]
\node[vtxs] (u1) at (0,1) {1}; \node[vtxs] (u2) at (1.2,1) {2}; \node[vtxs] (u3) at (0.6,0) {3}; \node[vtxs] (u4) at (1.8,0) {4};
\draw[aresta] (u1) -- node[above, font=\scriptsize] {$a$} (u2);
\draw[aresta] (u2) -- node[right, font=\scriptsize] {$b$} (u3);
\draw[aresta] (u1) -- node[left, font=\scriptsize] {$c$} (u3);
\draw[aresta] (u3) -- node[below, font=\scriptsize] {$d$} (u4);
\end{tikzpicture}
\hspace{1.2cm}
\begin{tikzpicture}[x=0.8cm,y=0.8cm, baseline={(0,0.5)}]
\node[vtxs] (u1) at (0,1) {1}; \node[vtxs] (u2) at (1.2,1) {2}; \node[vtxs] (u3) at (0.6,0) {3}; \node[vtxs] (u4) at (1.8,0) {4};
\draw[gray!30, thick] (u1)--(u2) (u2)--(u3) (u1)--(u3);
\draw[aresta, MidnightBlue] (u1)--(u2) (u2)--(u3) (u3)--(u4);
\node[font=\scriptsize] at (0.9,-0.6) {$\{a,b,d\}$};
\end{tikzpicture}\hspace{4mm}
\begin{tikzpicture}[x=0.8cm,y=0.8cm, baseline={(0,0.5)}]
\node[vtxs] (u1) at (0,1) {1}; \node[vtxs] (u2) at (1.2,1) {2}; \node[vtxs] (u3) at (0.6,0) {3}; \node[vtxs] (u4) at (1.8,0) {4};
\draw[gray!30, thick] (u1)--(u2) (u2)--(u3) (u1)--(u3);
\draw[aresta, MidnightBlue] (u1)--(u2) (u1)--(u3) (u3)--(u4);
\node[font=\scriptsize] at (0.9,-0.6) {$\{a,c,d\}$};
\end{tikzpicture}\hspace{4mm}
\begin{tikzpicture}[x=0.8cm,y=0.8cm, baseline={(0,0.5)}]
\node[vtxs] (u1) at (0,1) {1}; \node[vtxs] (u2) at (1.2,1) {2}; \node[vtxs] (u3) at (0.6,0) {3}; \node[vtxs] (u4) at (1.8,0) {4};
\draw[gray!30, thick] (u1)--(u2) (u2)--(u3) (u1)--(u3);
\draw[aresta, MidnightBlue] (u2)--(u3) (u1)--(u3) (u3)--(u4);
\node[font=\scriptsize] at (0.9,-0.6) {$\{b,c,d\}$};
\end{tikzpicture}
\end{figura}
\end{exemplo}

\begin{center}\small
\begin{tabularx}{\linewidth}{@{}lX@{}}
\toprule
Matroide & Independentes\\
\midrule
\textbf{Linear} & conjuntos de colunas (vetores) de uma matriz que são linearmente independentes.\\
\textbf{Gráfica} (de um grafo $G$) & conjuntos de arestas \textbf{acíclicos} (florestas); se $G$ é conexo, as bases são as árvores geradoras.\\
\textbf{Uniforme} $U_{n,k}$ & todos os subconjuntos $I\subseteq E$ com $|I|\le k$ (\emph{livre}, se $k=n$).\\
\textbf{De partição} & $E=E_1\cup\dots\cup E_t$ (uma partição) e os $I$ com $|I\cap E_i|\le k_i$ para todo $i$.\\
\textbf{Transversal} & num grafo bipartido com lados $X$ e $Y$: os $I\subseteq X$ \emph{emparelháveis}, isto é, cobertos por algum emparelhamento.\\
\textbf{De emparelhamento} & num grafo qualquer: os conjuntos de vértices cobertos por algum emparelhamento.\\
\textbf{Dual} $M^*$ & as bases são os complementos $E\setminus B$ das bases $B$ de $M$; o dual da gráfica é a \emph{cográfica}.\\
\bottomrule
\end{tabularx}
\end{center}

\begin{exemplo}[Uma matroide de partição]
Um restaurante monta pratos executivos com no máximo 1 entrada, 2 acompanhamentos e 1 sobremesa. O conjunto $E$ de opções é particionado em entradas, acompanhamentos e sobremesas, e um pedido é ``independente'' se respeita os três limites. O pedido \{salada, arroz, feijão, pudim\} é independente; \{arroz, feijão, purê\} não é, pois tem 3 acompanhamentos. Os axiomas valem: o pedido vazio respeita os limites; tirar itens de um pedido válido mantém-no válido; e, se um pedido válido tem menos itens que outro, há alguma categoria em que o maior tem mais itens que o menor, e um desses itens pode ser acrescentado ao menor sem estourar o limite dela.
\end{exemplo}

\textbf{O axioma (I3) para a matroide gráfica.} Uma floresta com $k$ arestas num grafo de $n$ vértices tem $n-k$ componentes. Como $|I_1|<|I_2|$, a floresta $I_2$ tem menos componentes que $I_1$, e por isso alguma aresta $e\in I_2$ liga duas componentes diferentes de $I_1$: se todas as arestas de $I_2$ ficassem dentro de componentes de $I_1$, cada componente de $I_2$ estaria contida numa de $I_1$, e $I_2$ teria pelo menos tantas componentes quanto $I_1$. Acrescentar essa aresta $e$ a $I_1$ não cria ciclo (Figura~\ref{fig:aumento}).
\textbf{O axioma (I3) para a matroide linear.} O espaço gerado por $I_1$ tem dimensão $|I_1|<|I_2|$, então algum vetor de $I_2$ está fora dele, e acrescentá-lo a $I_1$ mantém a independência.

\begin{figura}[O axioma de aumento na matroide gráfica. A floresta $I_1=\{12,34\}$ (azul) tem 3 componentes; a floresta $I_2=\{13,35,24\}$ (vermelho tracejado), com uma aresta a mais, tem só 2. A aresta $35$ de $I_2$ liga duas componentes de $I_1$ e pode ser acrescentada a ela sem formar ciclo (à direita).\label{fig:aumento}]
\begin{tikzpicture}[x=1.1cm,y=1.1cm]
\begin{scope}
\node[vtxs] (p1) at (0,1) {1}; \node[vtxs] (p2) at (1.2,1) {2}; \node[vtxs] (p3) at (0,0) {3}; \node[vtxs] (p4) at (1.2,0) {4}; \node[vtxs] (p5) at (2.4,0.5) {5};
\draw[aresta, MidnightBlue] (p1)--(p2) (p3)--(p4);
\draw[thick, BrickRed, dashed] (p1)--(p3) (p2)--(p4);
\draw[thick, BrickRed, dashed] (p3) to[bend right=35] (p5);
\node[rot] at (1.2,-0.8) {$I_1$ (azul) e $I_2$ (vermelho)};
\end{scope}
\begin{scope}[xshift=5cm]
\node[vtxs] (q1) at (0,1) {1}; \node[vtxs] (q2) at (1.2,1) {2}; \node[vtxs] (q3) at (0,0) {3}; \node[vtxs] (q4) at (1.2,0) {4}; \node[vtxs] (q5) at (2.4,0.5) {5};
\draw[aresta, MidnightBlue] (q1)--(q2) (q3)--(q4);
\draw[aresta, BrickRed] (q3) to[bend right=35] (q5);
\node[rot] at (1.2,-0.8) {$I_1\cup\{35\}$: ainda sem ciclo};
\end{scope}
\end{tikzpicture}
\end{figura}

\subsection{Formas equivalentes do axioma de aumento e bases}
\begin{itemize}
  \item \textbf{(I3$'$)} Para todo $X\subseteq E$, todos os independentes \emph{maximais} contidos em $X$ têm o mesmo tamanho, chamado \textbf{posto} de $X$, $r(X)$. Prova: se $I_1$ e $I_2$ fossem maximais em $X$ com $|I_1|<|I_2|$, (I3) aumentaria $I_1$ dentro de $X$. Reciprocamente, dados $|I_1|<|I_2|$, estenda $I_1$ a um independente maximal $B_1$ dentro de $X=I_1\cup I_2$; como $|B_1|\ge|I_2|>|I_1|$, existe $e\in B_1\setminus I_1\subseteq I_2\setminus I_1$, e $I_1\cup\{e\}\subseteq B_1$ é independente por (I2).
  \item \textbf{(I3$''$)} Todos os independentes maximais de $\mathcal I$ têm o mesmo tamanho $r$, e, se $|I|=r$ e $|I'|=r-1$, existe $x\in I\setminus I'$ com $I'\cup\{x\}\in\mathcal I$.
  \item Uma \textbf{base} é um independente maximal; todas as bases têm o mesmo tamanho. A coleção $\mathcal B$ das bases satisfaz (B1) $\mathcal B\ne\emptyset$ e (B2) (\textbf{troca de bases}): se $B_1,B_2\in\mathcal B$ e $x\in B_1\setminus B_2$, existe $y\in B_2\setminus B_1$ com $(B_1\setminus\{x\})\cup\{y\}\in\mathcal B$. Reciprocamente, toda coleção $\mathcal B$ que satisfaz (B1) e (B2) é a coleção das bases da matroide cujos independentes são os subconjuntos de elementos de $\mathcal B$.
\end{itemize}

\begin{exemplo}[Posto e troca de bases]
Na matroide do Exemplo~\ref{ex:matpequena}, o posto do triângulo $X=\{a,b,c\}$ é $r(X)=2$: qualquer par de arestas do triângulo é independente maximal dentro de $X$, e não há um trio. O posto da matroide toda é 3. Para ver a troca de bases, tome $B_1=\{a,b,d\}$, $B_2=\{b,c,d\}$ e $x=a\in B_1\setminus B_2$: retirando $a$ de $B_1$ sobra $\{b,d\}$, e o elemento $y=c\in B_2\setminus B_1$ completa uma nova base, $\{b,c,d\}$. No grafo, isso é trocar uma aresta de uma árvore geradora por outra que reconecta as duas partes --- a mesma troca que aparece na prova da propriedade do corte do Capítulo~\ref{cap:agm}.
\end{exemplo}

\subsection{O algoritmo guloso para matroides}
\textbf{Problema da base de peso máximo:} dada uma matroide $M$ e uma função de peso $c:E\to\R$, encontrar uma base $B$ com $c(B)=\sum_{e\in B}c(e)$ máximo. Em geral, os independentes são dados de forma \emph{implícita}, por um \textbf{oráculo} que responde ``$X$ é independente?'', pois listá-los todos poderia exigir tempo exponencial em $|E|$. Na matroide gráfica, por exemplo, o oráculo é ``as arestas de $X$ formam ciclo?'', respondido com Union-Find; na linear, é um escalonamento de Gauss.
\begin{pseudo}{Algoritmo guloso para matroides: $O(n\log n+n\,h)$, sendo $h$ o custo do oráculo}
\begin{algorithmic}
\Function{Guloso}{$M,c$}
  \State $B:=\emptyset$
  \State ordene $E$ de modo que $c(e_1)\ge c(e_2)\ge\dots\ge c(e_n)$ \Comment{ordem não crescente}
  \For{$i:=1$ \textbf{até} $n$}
    \If{$B\cup\{e_i\}$ é independente em $M$} \State $B:=B\cup\{e_i\}$ \EndIf
  \EndFor
  \State \Return $B$
\EndFunction
\end{algorithmic}
\end{pseudo}

\begin{exemplo}[O guloso na matroide gráfica]\label{ex:gulosograf}
No grafo da Figura~\ref{fig:gulosograf}, com as arestas $AB$ (7), $AC$ (5), $BC$ (8), $BD$ (6), $CD$ (9), $CE$ (4) e $DE$ (3), queremos uma árvore geradora de peso \emph{máximo}. Em ordem decrescente de peso, o guloso aceita $CD$ (9), $BC$ (8) e $AB$ (7); rejeita $BD$ (6), que fecharia o ciclo $B,C,D$, e $AC$ (5), que fecharia $A,B,C$; aceita $CE$ (4); e rejeita $DE$ (3). A árvore tem peso $9+8+7+4=28$, e, pelo teorema a seguir, nenhuma árvore geradora pesa mais (o que se confirma testando as poucas árvores possíveis).
\begin{figura}[O guloso para a árvore geradora de peso máximo. As arestas aceitas estão em azul; as rejeitadas, que fechariam ciclos, em cinza tracejado.\label{fig:gulosograf}]
\begin{tikzpicture}[x=1.4cm,y=1.2cm]
\node[vtx] (A) at (0,1) {$A$}; \node[vtx] (B) at (1.5,1.8) {$B$}; \node[vtx] (C) at (1.5,0) {$C$};
\node[vtx] (D) at (3,1) {$D$}; \node[vtx] (E) at (3,-0.6) {$E$};
\draw[aresta, MidnightBlue, very thick] (A) -- node[peso] {7} (B);
\draw[gray, dashed] (A) -- node[peso] {5} (C);
\draw[aresta, MidnightBlue, very thick] (B) -- node[peso] {8} (C);
\draw[gray, dashed] (B) -- node[peso] {6} (D);
\draw[aresta, MidnightBlue, very thick] (C) -- node[peso] {9} (D);
\draw[aresta, MidnightBlue, very thick] (C) -- node[peso] {4} (E);
\draw[gray, dashed] (D) -- node[peso] {3} (E);
\end{tikzpicture}
\end{figura}
\end{exemplo}

\begin{teo}\label{teo:gulosomat}
Se $M$ é uma matroide, o algoritmo \textup{\textsc{Guloso}} devolve uma base de peso máximo, qualquer que seja a função de peso $c$.
\end{teo}
\begin{proof}
Seja $B$ a saída, que é uma base, pois nada mais pode ser acrescentado a ela, e seja $B^*$ uma base de peso máximo com $|B\cap B^*|$ o maior possível. Suponha $B\ne B^*$ e seja $f$ o primeiro elemento, na ordem do algoritmo, de $B\setminus B^*$; os elementos escolhidos antes de $f$ formam um conjunto $B'\subseteq B^*$. Pela troca de bases, existe $e\in B^*\setminus B$ tal que $(B^*\setminus\{e\})\cup\{f\}$ é base. Pela otimalidade de $B^*$, $c(f)\le c(e)$. Se $c(f)=c(e)$, essa nova base também seria ótima e teria mais elementos em comum com $B$, uma contradição. Logo $c(f)<c(e)$, e $e$ vem antes de $f$ na ordem. Mas, quando $e$ foi examinado, o conjunto escolhido estava contido em $B'$, e $B'\cup\{e\}\subseteq B^*$ é independente: o algoritmo teria escolhido $e$. Contradição.
\end{proof}

\textbf{A recíproca também vale.} Se $\mathcal I$ satisfaz (I1) e (I2) e o guloso encontra um independente maximal de peso máximo para \emph{toda} função de peso, então vale (I3). De fato, suponha $|I_1|=k<|I_2|$ e que nenhum $e\in I_2\setminus I_1$ aumente $I_1$. Dê peso $k+2$ aos elementos de $I_1$, $k+1$ aos de $I_2\setminus I_1$ e 0 aos demais. O guloso escolhe $I_1$ primeiro, não consegue acrescentar nada de $I_2\setminus I_1$ e termina com peso $k(k+2)$; mas $I_2$ pesa pelo menos $(k+1)(k+1)>k(k+2)$, uma contradição. Ou seja, \textbf{as matroides são exatamente os sistemas em que o guloso é sempre ótimo}.
\begin{itemize}
  \item \textbf{Base de peso mínimo}: use a ordem não decrescente, ou os pesos $-c$. O algoritmo de Kruskal é exatamente isso na matroide gráfica (Capítulo~\ref{cap:agm}).
  \item \textbf{Testar se $k$ colunas de uma matriz $m\times k$ são linearmente independentes}: escalonamento de Gauss, em $O(m\,k\min\{m,k\})$; se $k>m$, elas já são dependentes.
  \item \textbf{Seleção de intervalos} com pesos unitários: o guloso que escolhe ``o que termina primeiro'' é ótimo, mas a família dos conjuntos de intervalos disjuntos \textbf{não} é uma matroide. Com $I_1=\{[1,10]\}$ e $I_2=\{[1,2],[3,4]\}$, nenhum intervalo de $I_2$ é compatível com $[1,10]$, e (I3) falha (Figura~\ref{fig:naomat}). A otimalidade vem de um argumento de troca específico; com pesos arbitrários, o guloso falha e é preciso programação dinâmica.
  \item \textbf{Emparelhamentos} (conjuntos de arestas sem extremos em comum) num grafo qualquer \textbf{não} formam uma matroide: no caminho $a\!-\!b\!-\!c\!-\!d$, o emparelhamento $\{bc\}$ não pode ser aumentado por $ab$ nem por $cd$, que formam o emparelhamento $\{ab,cd\}$. Um guloso que escolha $bc$ primeiro fica com uma aresta, e o máximo é 2.
\end{itemize}

\begin{figura}[Duas famílias que não são matroides. À esquerda, intervalos: $I_1=\{[1,10]\}$ tem menos elementos que $I_2=\{[1,2],[3,4]\}$, mas nenhum intervalo de $I_2$ pode ser acrescentado a $I_1$. À direita, emparelhamentos no caminho $a\!-\!b\!-\!c\!-\!d$: $\{bc\}$ (vermelho) é maximal, mas menor que $\{ab,cd\}$ (azul).\label{fig:naomat}]
\begin{tikzpicture}[x=4.5mm,y=5mm, baseline=(base)]
\draw[gray] (0.5,0) -- (11,0); \foreach \xx in {1,...,10} {\draw[gray] (\xx,0.1) -- (\xx,-0.1); \node[font=\tiny, text=gray] at (\xx,-0.45) {\xx};}
\draw[line width=3pt, BrickRed!70] (1,1) -- (10,1); \node[font=\scriptsize, anchor=west] at (10.3,1) {$I_1$};
\draw[line width=3pt, MidnightBlue!70] (1,2) -- (2,2); \draw[line width=3pt, MidnightBlue!70] (3,2) -- (4,2); \node[font=\scriptsize, anchor=west] at (10.3,2) {$I_2$};
\coordinate (base) at (0,1);
\end{tikzpicture}
\hspace{1.5cm}
\begin{tikzpicture}[x=1.2cm,y=1cm, baseline=(base)]
\foreach \nm [count=\k] in {a,b,c,d} {\node[vtxs] (m\nm) at (\k,0.5) {$\nm$};}
\draw[line width=2.5pt, MidnightBlue!70] (ma) -- (mb); \draw[line width=2.5pt, BrickRed!70] (mb) -- (mc); \draw[line width=2.5pt, MidnightBlue!70] (mc) -- (md);
\coordinate (base) at (0,0.5);
\end{tikzpicture}
\end{figura}

\begin{complemento}[Para saber mais: interseção de matroides]
Dadas duas matroides $M_1=(E,\mathcal I_1)$ e $M_2=(E,\mathcal I_2)$, o problema da \textbf{interseção de matroides} pede um conjunto $I\in\mathcal I_1\cap\mathcal I_2$ de tamanho (ou peso) máximo. Vários problemas são desse tipo: o \emph{emparelhamento máximo em grafos bipartidos} é a interseção de duas matroides de partição (no máximo uma aresta por vértice de cada lado); a \emph{floresta colorida} --- o maior conjunto de arestas que é uma floresta e tem no máximo uma aresta de cada cor --- é a interseção da matroide gráfica com uma de partição pelas cores.

A interseção de duas matroides, em geral, não é uma matroide, e o guloso pode falhar, como no caminho $a\!-\!b\!-\!c\!-\!d$ acima, que é bipartido. Mesmo assim, o problema é \textbf{polinomial}, por um algoritmo de \emph{caminhos aumentantes}. Dado um conjunto comum $I$, constrói-se um \emph{grafo de troca}: para $y\in I$ e $x\notin I$, há um arco $y\to x$ se $I-y+x\in\mathcal I_1$ e um arco $x\to y$ se $I-y+x\in\mathcal I_2$; as fontes são os $x$ com $I+x\in\mathcal I_1$ e os destinos, os $x$ com $I+x\in\mathcal I_2$. Um caminho \emph{mais curto} de uma fonte a um destino dá, por diferença simétrica com $I$, um conjunto comum com um elemento a mais; se não há caminho, $I$ é máximo. São no máximo $r$ aumentos, cada um com $O(nr)$ consultas ao oráculo. Vale ainda um teorema de dualidade: $\max|I|=\min_{U\subseteq E}\{r_1(U)+r_2(E\setminus U)\}$.

Com \textbf{três} matroides, o problema já é NP-difícil: um \emph{caminho hamiltoniano} num digrafo é um conjunto comum de tamanho $n-1$ a três matroides --- a gráfica do grafo subjacente (sem ciclos), a de partição ``grau de entrada no máximo 1'' e a de partição ``grau de saída no máximo 1''. Também é polinomial o \textbf{particionamento} de matroides, que cobre $E$ com o menor número de independentes (por exemplo, cobrir as arestas de um grafo com o menor número de florestas). Reconhecer que um problema é uma interseção ou um particionamento de matroides prova, portanto, que ele é polinomial.
\end{complemento}

\begin{chave}
Guloso: escolha local, sem retrocesso; correto quando vale a propriedade gulosa, provada por troca. Matroide: (I1) $\emptyset$ independente, (I2) fechada para subconjuntos, (I3) aumento. Nas matroides, o guloso (ordem decrescente de peso e oráculo de independência) encontra a base de peso máximo, e apenas nelas ele acerta para todo peso. Exemplos: linear, gráfica, uniforme, de partição, transversal. Não exemplos: emparelhamentos de arestas e intervalos disjuntos. A interseção de duas matroides é polinomial; a de três, NP-difícil.
\end{chave}

% ----------------------------------------------------------------
\section{Exercícios}
\excurso
\begin{exercicio}{gm-troco}[Lista]
Escreva um algoritmo guloso para o problema do troco: dar $S$ centavos com o menor número de moedas, de valores $v_1,\dots,v_n$. Mostre que ele é ótimo para as moedas $\{10,5,2,1\}$ e que não é ótimo para qualquer sistema monetário.
\end{exercicio}
\begin{exercicio}{gm-i3}[Slides]
Prove o axioma de aumento (I3) para a matroide linear, para a gráfica, para a transversal e para a de emparelhamento.
\end{exercicio}
\begin{exercicio}{gm-g}[Slides]
Mostre que, se $\mathcal I$ satisfaz (I1) e (I2) e o algoritmo guloso devolve um independente maximal de peso máximo para \emph{toda} função de peso, então vale (I3).
\end{exercicio}
\begin{exercicio}{gm-pratica}[Slides]
(a) É possível usar matroides no problema de seleção de intervalos com pesos unitários? (b) Faça um algoritmo para verificar se um conjunto $S$ de vetores-coluna de uma matriz é linearmente independente. Qual a complexidade?
\end{exercicio}
\begin{exercicio}{gm-pesq}[Pesquisa]
Investigue: os axiomas das matroides e sua interpretação; a relação entre a propriedade de troca e a correção do guloso; o algoritmo de Kruskal como guloso numa matroide; quando a maximização de pesos pelo guloso é ótima; o que muda quando há duas famílias de restrições de independência; como funciona a interseção de duas matroides por caminhos aumentantes; e por que duas matroides levam a problemas polinomiais enquanto três já podem levar a problemas NP-difíceis.
\end{exercicio}

\excompl
\begin{exercicio}{gm-ativ}[Complementar]
Seleção de atividades: 11 atividades têm inícios $s=(1,3,0,5,3,5,6,8,8,2,12)$ e términos $f=(4,5,6,7,9,9,10,11,12,14,16)$. Aplique o guloso ``a que termina primeiro'' e prove que ele é ótimo.
\end{exercicio}
\begin{exercicio}{gm-ruins}[Complementar]
Mostre com contraexemplos que, para a seleção de atividades, os gulosos ``a de menor duração primeiro'' e ``a que começa primeiro'' não são ótimos.
\end{exercicio}
\begin{exercicio}{gm-frac}[Complementar]
Na mochila \emph{fracionária}, em que se podem levar frações dos itens, prove que o guloso pela razão valor/peso é ótimo. Aplique a $W=50$ e itens $(p,v)=(10,60),(20,100),(30,120)$ e compare com a mochila 0-1.
\end{exercicio}
\begin{exercicio}{gm-unif}[Complementar]
Prove que a matroide uniforme $U_{n,k}$ e a matroide de partição satisfazem (I1), (I2) e (I3).
\end{exercicio}
\begin{exercicio}{gm-tarefas}[Complementar \estrela]
Tarefas unitárias com prazos e multas: cada tarefa leva uma unidade de tempo, tem prazo $d_i$ e multa $w_i$ se terminar depois do prazo. Um conjunto de tarefas é \emph{independente} se pode ser executado todo dentro dos prazos, o que vale se e somente se, para todo $t$, no máximo $t$ tarefas do conjunto têm prazo $\le t$. Esses conjuntos formam uma matroide; use o guloso para minimizar a multa total com $d=(4,2,4,3,1,4,6)$ e $w=(70,60,50,40,30,20,10)$.
\end{exercicio}
\begin{exercicio}{gm-huffman}[Complementar \estrela]
Os \emph{códigos de Huffman} são construídos por um guloso que junta, repetidamente, as duas árvores de menor frequência. Construa o código para as frequências $a{:}45$, $b{:}13$, $c{:}12$, $d{:}16$, $e{:}9$ e $f{:}5$ e calcule o custo $\sum(\text{frequência})\times(\text{tamanho do código})$.
\end{exercicio}
\begin{exercicio}{gm-spt}[Complementar]
Tarefas com durações $t_1,\dots,t_n$ serão executadas uma após a outra numa única máquina. Mostre que executá-las em ordem crescente de duração minimiza a soma dos instantes de término. Aplique a durações $3,1,2$.
\end{exercicio}
\begin{exercicio}{gm-naomat}[Complementar]
Seja $E=\{1,2,3\}$ e $\mathcal I=\{\emptyset,\{1\},\{2\},\{3\},\{1,2\}\}$. Verifique se $(E,\mathcal I)$ é uma matroide. Em caso negativo, mostre uma função de peso para a qual o guloso falha.
\end{exercicio}
\begin{exercicio}{gm-paralelos}[Complementar]
Seja $E=\{1,2,3,4\}$ e $\mathcal I$ a família dos conjuntos com no máximo 2 elementos, exceto $\{1,2\}$. Verifique se é uma matroide e dê uma matriz cujas colunas realizam essa matroide.
\end{exercicio}
\begin{exercicio}{gm-posto}[Complementar]
Mostre que, na matroide gráfica de um grafo com $n$ vértices e $c$ componentes conexas, todas as bases têm $n-c$ arestas. Qual o posto da matroide gráfica de um grafo com 6 vértices formado por um triângulo e um caminho com 3 vértices?
\end{exercicio}
\begin{exercicio}{gm-dual}[Complementar]
Mostre que o dual da matroide uniforme $U_{n,k}$ é $U_{n,n-k}$.
\end{exercicio}
\begin{exercicio}{gm-unifgul}[Complementar]
O que faz o algoritmo guloso na matroide uniforme $U_{n,k}$? E numa matroide de partição? Aplique ao problema de escolher, entre itens das categorias $A=\{9,7,4\}$, $B=\{8,3\}$ e $C=\{6,5,2\}$ (os números são os pesos), no máximo 2 itens por categoria com peso total máximo.
\end{exercicio}
\begin{exercicio}{gm-moedas}[Complementar]
Mostre que o troco guloso não é ótimo para as moedas $\{1,7,10\}$, dando um valor $S$ em que ele falha. O sistema $\{1,5,10,25\}$ tem esse problema?
\end{exercicio}
\begin{exercicio}{gm-pesado}[Complementar]
Dê um exemplo de seleção de intervalos \emph{com pesos} em que o guloso ``o que termina primeiro'' não encontra o conjunto compatível de peso máximo. Como resolver esse problema?
\end{exercicio}
\begin{exercicio}{gm-fracn}[Complementar \estrela]
Mostre que a mochila fracionária pode ser resolvida em $O(n)$, sem ordenar os itens, usando seleção em tempo linear.
\end{exercicio}

% ================================================================
\chapter{Árvores geradoras mínimas: Kruskal e Prim}\label{cap:agm}\label{sec:agm}
% ================================================================
\begin{objetivos}
\begin{itemize}
  \item O problema da árvore geradora mínima e a propriedade do corte.
  \item O algoritmo de Kruskal, com Union-Find, e a sua correção.
  \item O algoritmo de Prim e as suas implementações com vetor, heap binário e heap de Fibonacci.
  \item Como escolher a implementação de acordo com a densidade do grafo.
\end{itemize}
\end{objetivos}

Ligar $n$ cidades por uma rede de custo mínimo, usando apenas as ligações possíveis, é encontrar uma \textbf{árvore geradora mínima} (AGM). Os dois algoritmos clássicos são gulosos: o de Kruskal percorre as arestas em ordem crescente de peso e aceita as que não fecham ciclo, usando a estrutura Union-Find do Capítulo~\ref{cap:uf}; o de Prim faz crescer uma única árvore a partir de uma raiz, usando uma fila de prioridades (Capítulos~\ref{cap:heap} e~\ref{cap:heapav}). A correção dos dois decorre de um único lema, a propriedade do corte.

\section{O problema e a propriedade do corte}
Dado um grafo conexo $G$ com pesos $c$ nas arestas, queremos uma \textbf{árvore geradora} --- um subgrafo conexo e acíclico que contém todos os $n$ vértices e, portanto, tem $n-1$ arestas --- de peso total mínimo. Um \textbf{corte} $(S,V\setminus S)$ é uma partição dos vértices em dois lados, e uma aresta \textbf{atravessa} o corte se tem um extremo em cada lado.
\begin{lema}[Propriedade do corte]\label{lema:corte}
Seja $A$ um conjunto de arestas contido em alguma AGM e $(S,V\setminus S)$ um corte que nenhuma aresta de $A$ atravessa. Se $e$ é uma aresta de peso mínimo entre as que atravessam o corte, então $A\cup\{e\}$ está contido em alguma AGM.
\end{lema}
\begin{proof}
Seja $T$ uma AGM com $A\subseteq T$. Se $e\in T$, não há o que provar. Senão, $T+e$ tem um ciclo que passa por $e$; como $e$ atravessa o corte, o ciclo precisa atravessá-lo de novo, por outra aresta $f\in T$, e $f\notin A$, pois nenhuma aresta de $A$ atravessa o corte. Então $T'=T+e-f$ é uma árvore geradora e $c(T')=c(T)+c(e)-c(f)\le c(T)$, pois $c(e)\le c(f)$. Logo $T'$ é uma AGM e contém $A\cup\{e\}$.
\end{proof}

\section{O algoritmo de Kruskal}\label{sec:kruskal}
No algoritmo de Kruskal, a solução parcial é uma \textbf{floresta}. As arestas são percorridas em ordem crescente de peso, e cada uma é aceita se \textbf{não forma ciclo} com as já aceitas, isto é, se liga duas árvores diferentes da floresta. O teste é feito com conjuntos disjuntos: cada árvore da floresta é um conjunto, e uma aresta $uv$ fecha ciclo se e somente se $u$ e $v$ estão no mesmo conjunto.
\begin{pseudo}{Kruskal com Union-Find}
\begin{algorithmic}
\Function{Kruskal}{$G,c$}
  \State $E(T):=\emptyset$
  \For{\textbf{cada} $v\in V$} \State \MakeSet$(v)$ \Comment{$n$ árvores com um vértice} \EndFor
  \State ordene as arestas de modo que $c(e_1)\le c(e_2)\le\dots\le c(e_m)$ \Comment{$O(m\log m)=O(m\log n)$}
  \For{$j:=1$ \textbf{até} $m$}
    \State seja $e_j=uv$
    \If{\Find$(u)\ne$ \Find$(v)$} \Comment{une duas árvores distintas: não fecha ciclo}
      \State $E(T):=E(T)\cup\{e_j\}$;\ \Union$(u,v)$
    \EndIf
  \EndFor
  \State \Return $T$ \Comment{pode parar quando $|E(T)|=n-1$}
\EndFunction
\end{algorithmic}
\end{pseudo}
\textbf{Complexidade.} A ordenação domina: $O(m\log m)=O(m\log n)$, pois $m\le n^2$ implica $\log m\le2\log n$. O laço custa $O(m+n\log n)$ com listas encadeadas e união ``menor na maior'', já que cada vértice muda de conjunto no máximo $\log n$ vezes, ou $O(m\,\alpha(n))$ com floresta, rank e compressão de caminhos. O total é $O(m\log n)$.

\textbf{Correção.} Ao aceitar $e=uv$, tome como $S$ a árvore atual de $u$. Nenhuma aresta aceita atravessa esse corte, e $e$ é a mais leve que o atravessa: uma aresta mais leve que o atravessasse teria sido examinada antes, quando os seus extremos já estavam em árvores diferentes (as árvores só crescem), e teria sido aceita. Pelo Lema~\ref{lema:corte}, por indução, o conjunto de arestas aceitas está sempre contido numa AGM. O algoritmo de Kruskal também é o Teorema~\ref{teo:gulosomat} aplicado à matroide gráfica.

\begin{center}
\begin{tikzpicture}[v/.style={circle,draw,fill=yellow!40,inner sep=1.2pt,minimum size=5mm,font=\small},
  agm/.style={very thick,MidnightBlue}, nao/.style={thin,gray}, w/.style={font=\scriptsize,fill=white,inner sep=1pt}]
\node[v] (a) at (0,1.2) {$a$}; \node[v] (b) at (1.6,2.4) {$b$}; \node[v] (c) at (3.4,2.4) {$c$}; \node[v] (d) at (5.2,2.4) {$d$};
\node[v] (e) at (6.8,1.2) {$e$}; \node[v] (f) at (5.2,0) {$f$}; \node[v] (g) at (3.4,0) {$g$}; \node[v] (h) at (1.6,0) {$h$};
\node[v] (i) at (2.5,1.2) {$i$};
\draw[agm] (a)--node[w]{4}(b); \draw[agm] (a)--node[w]{8}(h); \draw[nao] (b)--node[w]{11}(h); \draw[nao] (b)--node[w]{8}(c);
\draw[nao] (h)--node[w]{7}(i); \draw[agm] (h)--node[w]{1}(g); \draw[agm] (i)--node[w]{2}(c); \draw[nao] (i)--node[w]{6}(g);
\draw[agm] (c)--node[w]{7}(d); \draw[agm] (c)--node[w]{4}(f); \draw[agm] (g)--node[w]{2}(f); \draw[nao] (d)--node[w]{14}(f);
\draw[agm] (d)--node[w]{9}(e); \draw[nao] (e)--node[w]{10}(f);
\end{tikzpicture}
\end{center}
\begin{exemplo}[Kruskal e Prim no grafo acima]\label{ex:agm}
O algoritmo de Kruskal aceita $hg\,(1)$, $ic\,(2)$, $gf\,(2)$, $ab\,(4)$, $cf\,(4)$, $cd\,(7)$, $ah\,(8)$ e $de\,(9)$, e rejeita $ig\,(6)$, $hi\,(7)$, $bc\,(8)$, $ef\,(10)$, $bh\,(11)$ e $df\,(14)$, que fechariam ciclos. A AGM tem \textbf{custo 37} (arestas em azul). O algoritmo de Prim, com raiz $a$, acrescenta $ab\,(4)$, $ah\,(8)$, $hg\,(1)$, $gf\,(2)$, $fc\,(4)$, $ci\,(2)$, $cd\,(7)$ e $de\,(9)$, nessa ordem, e chega à mesma árvore. A AGM não é única: $ah$ e $bc$ pesam 8, e trocar uma pela outra dá outra AGM de custo 37. Quando todos os pesos são distintos, a AGM é única.
\end{exemplo}

\section{O algoritmo de Prim}\label{sec:prim}
No algoritmo de Prim, a solução parcial é uma \textbf{árvore} que cresce a partir de uma raiz: a cada passo, entra a aresta mais leve com um extremo na árvore e o outro fora dela. A correção vem do Lema~\ref{lema:corte}, com $S=V(T)$.
\begin{pseudo}{Prim em $O(n^2)$ com vetores}
\begin{algorithmic}
\Function{Prim}{$G,W,r$}
  \State para todo $v$: $chave(v):=\infty$;\ $pai(v):=\nil$;\ $dentro(v):=$ falso
  \State $chave(r):=0$
  \For{$it:=1$ \textbf{até} $n$}
    \State $u:=$ vértice com $dentro(u)=$ falso e $chave(u)$ mínima \Comment{varredura em $O(n)$}
    \State $dentro(u):=$ verdadeiro \Comment{se $u\ne r$, a aresta $pai(u)\,u$ entra na árvore}
    \For{\textbf{cada} $v\in N(u)$ com $dentro(v)=$ falso} \Comment{matriz: linha $u$, $O(n)$; listas: $O(\text{grau}(u))$}
      \If{$W(u,v)<chave(v)$} \State $chave(v):=W(u,v)$;\ $pai(v):=u$ \EndIf
    \EndFor
  \EndFor
  \State \Return $\{\,pai(v)\,v : v\ne r\,\}$
\EndFunction
\end{algorithmic}
\end{pseudo}
O valor $chave(v)$ é o peso da aresta mais leve que liga $v$ à árvore. O algoritmo tem o mesmo esqueleto do algoritmo de Dijkstra (Capítulo~\ref{cap:dijkstra}), trocando ``$L(u)+W(u,v)$'' por ``$W(u,v)$''.
\begin{center}\small
\begin{tabularx}{\linewidth}{@{}>{\raggedright\arraybackslash}p{5.2cm}>{\raggedright\arraybackslash}p{3.6cm}X@{}}
\toprule
Implementação & Custo & Comentário\\
\midrule
Matriz de adjacência e vetor $chave$ & $O(n^2)$ & $n$ varreduras de $O(n)$; ótimo para matriz\\
Listas de adjacência e vetor $chave$ & $O(n^2+m)=O(n^2)$ & a varredura domina\\
Listas e heap binário (com redução de chave) & $O(m\log n)$ & $n$ remoções e até $m$ reduções, cada uma $O(\log n)$\\
Listas e heap de Fibonacci & $O(m+n\log n)$ & redução de chave em $O(1)$\am\\
\bottomrule
\end{tabularx}
\end{center}
A versão com heap \textbf{não é necessariamente melhor} que a de $O(n^2)$: num grafo denso, $m=\Theta(n^2)$ e $m\log n=\Theta(n^2\log n)$. O heap compensa quando $m=o(n^2/\log n)$.

\begin{cuidado}[Armadilhas de árvores geradoras mínimas]
\begin{itemize}
  \item ``O guloso funciona porque escolhe o melhor a cada passo'' não é uma prova. A prova é a propriedade do corte, ou a estrutura de matroide.
  \item O algoritmo de Kruskal aceita uma aresta se ela \emph{não fecha ciclo}, e não ``se o vértice ainda não foi visitado'' --- isso seria o algoritmo de Prim. No Kruskal, a solução parcial pode ser desconexa.
  \item A AGM minimiza a soma dos pesos, e não as distâncias até uma raiz; isso é o que faz o algoritmo de Dijkstra.
  \item A complexidade do Kruskal é $O(m\log n)$ por causa da \emph{ordenação}; as operações de Union-Find, sozinhas, são quase lineares.
  \item O Prim com heap não é sempre melhor; em grafos densos, perde para a versão com vetor.
\end{itemize}
\end{cuidado}

\begin{chave}
Propriedade do corte: a aresta mais leve que atravessa um corte que $A$ não atravessa pode ser acrescentada a $A$. Kruskal: arestas em ordem crescente e Union-Find, $O(m\log n)$. Prim: uma árvore que cresce pela aresta mais leve do corte, $O(n^2)$ com vetor ou matriz, $O(m\log n)$ com heap binário ou $O(m+n\log n)$ com heap de Fibonacci. Com pesos distintos, a AGM é única.
\end{chave}

% ----------------------------------------------------------------
\section{Exercícios}
\excurso
\begin{exercicio}{agm-l1}[Lista]
Prove a correção do algoritmo de Kruskal.
\end{exercicio}
\begin{exercicio}{agm-l2}[Lista]
Prove a correção do algoritmo de Prim.
\end{exercicio}
\begin{exercicio}{agm-l3}[Lista]
Mostre como implementar o algoritmo de Prim em tempo $O(m\log n)$ com o grafo representado por listas de adjacência.
\end{exercicio}
\begin{exercicio}{agm-l4}[Lista]
Mostre como implementar o algoritmo de Prim em tempo $O(n^2)$ com o grafo representado (a) por listas de adjacência; (b) por matriz de adjacência.
\end{exercicio}
\begin{exercicio}{agm-l6}[Lista]
É possível implementar o algoritmo de Prim em tempo $O(n^2)$ com matriz de adjacência? É possível fazer melhor que isso nessa representação? Justifique.
\end{exercicio}
\begin{exercicio}{agm-l7}[Lista]
Quais seriam as formas de resolver o problema da árvore geradora de peso \textbf{máximo}? Justifique.
\end{exercicio}

\excompl
\begin{exercicio}{agm-sim}[Complementar]
No grafo com as arestas $AB\,(6)$, $AC\,(1)$, $AD\,(5)$, $BC\,(5)$, $BE\,(3)$, $CD\,(5)$, $CE\,(6)$, $CF\,(4)$, $DF\,(2)$ e $EF\,(6)$, execute o Kruskal (indicando as arestas rejeitadas) e o Prim a partir de $A$ (indicando a ordem de entrada das arestas). Qual o custo da AGM? Ela é única?
\end{exercicio}
\begin{exercicio}{agm-unica}[Complementar]
Prove que, se todos os pesos são distintos, a AGM é única.
\end{exercicio}
\begin{exercicio}{agm-ciclo}[Complementar]
Prove a \emph{propriedade do ciclo}: se $e$ é a aresta de peso estritamente maior de algum ciclo de $G$, então $e$ não pertence a nenhuma AGM.
\end{exercicio}
\begin{exercicio}{agm-gargalo}[Complementar]
Mostre que, numa AGM $T$, o caminho entre dois vértices $u$ e $v$ minimiza o peso da \emph{maior} aresta, entre todos os caminhos de $u$ a $v$ em $G$.
\end{exercicio}
\begin{exercicio}{agm-nova}[Complementar]
Dada uma AGM $T$ de $G$, acrescenta-se a $G$ uma nova aresta $uv$ de peso $c$. Como obter uma AGM do novo grafo em $O(n)$?
\end{exercicio}
\begin{exercicio}{agm-prim7}[Complementar]
No grafo com as arestas $AB\,(7)$, $AD\,(5)$, $BC\,(8)$, $BD\,(9)$, $BE\,(7)$, $CE\,(5)$, $DE\,(15)$, $DF\,(6)$, $EF\,(8)$, $EG\,(9)$ e $FG\,(11)$, execute o Prim a partir de $A$, mostrando a ordem em que as arestas entram, e o Kruskal, mostrando as arestas rejeitadas. Qual o custo da AGM?
\end{exercicio}
\begin{exercicio}{agm-boruvka}[Complementar \estrela]
No algoritmo de Borůvka, cada componente da floresta atual escolhe, em paralelo, a aresta mais leve que sai dela, e todas essas arestas são acrescentadas de uma vez. Mostre que o algoritmo está correto (supondo pesos distintos) e que faz $O(\log n)$ rodadas, cada uma em $O(m)$.
\end{exercicio}
\begin{exercicio}{agm-iguais}[Complementar]
Se todas as arestas têm o mesmo peso, como encontrar uma AGM em $O(n+m)$?
\end{exercicio}
\begin{exercicio}{agm-12}[Complementar]
Se todos os pesos são 1 ou 2, mostre como encontrar uma AGM em $O(m\,\alpha(n))$, sem ordenar por comparações.
\end{exercicio}
\begin{exercicio}{agm-leve}[Complementar]
Prove que, se a aresta mais leve do grafo é única, ela pertence a toda AGM. Vale o mesmo para a segunda mais leve? E para a terceira?
\end{exercicio}
\begin{exercicio}{agm-dijk}[Complementar]
Dê um grafo conexo com pesos positivos em que a árvore de caminhos mínimos a partir de algum vértice (Capítulo~\ref{cap:dijkstra}) não é uma AGM.
\end{exercicio}
\begin{exercicio}{agm-parar}[Complementar]
No algoritmo de Kruskal, quantas arestas são examinadas no melhor e no pior caso, se o laço para ao aceitar $n-1$ arestas? Isso muda a complexidade assintótica?
\end{exercicio}
\begin{exercicio}{agm-listas}[Complementar]
Explique por que, no algoritmo de Kruskal com Union-Find por listas encadeadas e união ``menor na maior'', o laço principal custa $O(m+n\log n)$.
\end{exercicio}
\begin{exercicio}{agm-floresta}[Complementar]
Se o grafo é desconexo, o que o algoritmo de Kruskal devolve? E o de Prim? Como obter uma \emph{floresta geradora mínima}?
\end{exercicio}
\begin{exercicio}{agm-segunda}[Complementar \estrela]
A \emph{segunda melhor} árvore geradora é a de menor peso entre as diferentes de uma AGM fixada $T$. Mostre que ela pode ser obtida de $T$ trocando uma única aresta, e dê um algoritmo $O(n^2)$ para encontrá-la.
\end{exercicio}

% ================================================================
\chapter{Busca em largura}\label{cap:bfs}\label{sec:bfs}
% ================================================================
\begin{objetivos}
\begin{itemize}
  \item O que é uma busca em grafos e como a busca em largura percorre o grafo em camadas.
  \item A classificação das arestas em relação à floresta de largura.
  \item A propriedade fundamental dos níveis e suas consequências: distâncias e bipartição.
  \item A relação entre a estrutura de dados usada e o tipo de busca.
\end{itemize}
\end{objetivos}

Uma \textbf{busca} é um procedimento sistemático para visitar todos os vértices e todas as arestas de um grafo, com o objetivo de resolver um problema ou extrair informação da sua estrutura. Na busca em largura, os vértices são visitados em camadas: primeiro a raiz, depois os vizinhos dela, depois os vizinhos dos vizinhos, e assim por diante. Por isso ela calcula distâncias em número de arestas e reconhece grafos bipartidos. A notação de grafos está na §\ref{sec:buscas}.

\section{O algoritmo}
A regra da busca em largura é ``\emph{esgotar o pai antes de processar os filhos}'': a próxima aresta examinada parte sempre do vértice \textbf{mais antigo} na busca. A estrutura que realiza essa regra é uma \textbf{fila} $F$. Cada vértice recebe um rótulo $L(v)$, com a ordem em que foi alcançado, e um nível $nivel(v)$, a sua profundidade na floresta de largura; os ponteiros $pai(v)$ definem a floresta.
\begin{pseudo}{Busca em largura com classificação das arestas}
\begin{algorithmic}
\Procedure{BuscaLargura}{$G$}
  \State $t:=0$;\ $F:=\emptyset$;\ para todo $v$: $L(v):=0$;\ $pai(v):=\nil$
  \While{existe $v$ com $L(v)=0$} \Comment{nova raiz, se o grafo for desconexo}
    \State $nivel(v):=0$;\ $t:=t+1$;\ $L(v):=t$;\ insira $v$ no fim de $F$
    \While{$F\ne\emptyset$}
      \State $v:=$ primeiro elemento de $F$;\ retire $v$ de $F$
      \For{\textbf{cada} $w\in N(v)$}
        \If{$L(w)=0$} \Comment{aresta-\textbf{pai}}
          \State $pai(w):=v$;\ $nivel(w):=nivel(v)+1$;\ $t:=t+1$;\ $L(w):=t$;\ insira $w$ no fim de $F$
        \ElsIf{$nivel(w)=nivel(v)$ \textbf{e} $w\in F$} \Comment{$w\in F$: para não classificar duas vezes}
          \State se $pai(w)=pai(v)$, $vw$ é uma aresta-\textbf{irmão}; senão, uma aresta-\textbf{primo}
        \ElsIf{$nivel(w)=nivel(v)+1$} \State $vw$ é uma aresta-\textbf{tio}
        \EndIf
      \EndFor
    \EndWhile
  \EndWhile
\EndProcedure
\end{algorithmic}
\end{pseudo}

\begin{center}
\begin{tikzpicture}[v/.style={circle,draw,fill=yellow!40,inner sep=1.2pt,minimum size=5mm,font=\small},
  pai/.style={very thick,MidnightBlue}, irm/.style={thick,BrickRed}, pri/.style={thick,dotted,orange!85!black}, tio/.style={thick,dashed,ForestGreen},
  lab/.style={font=\scriptsize,text=gray!70!black}]
\node[v] (a) at (0.9,3.6) {$a$};
\node[v] (b) at (0,2.4) {$b$}; \node[v] (c) at (1.8,2.4) {$c$};
\node[v] (d) at (0,1.2) {$d$}; \node[v] (e) at (1.8,1.2) {$e$};
\node[v] (f) at (0,0) {$f$}; \node[v] (g) at (1.4,0) {$g$}; \node[v] (h) at (2.8,0) {$h$};
\draw[pai] (a)--(b); \draw[pai] (a)--(c); \draw[pai] (b)--(d); \draw[pai] (c)--(e); \draw[pai] (d)--(f); \draw[pai] (d)--(g); \draw[pai] (e)--(h);
\draw[irm] (b)--(c); \draw[irm] (f)--(g); \draw[pri] (d)--(e); \draw[pri] (g)--(h); \draw[tio] (c)--(d); \draw[tio] (e)--(f);
\node[lab,right=1pt of a] {1/0}; \node[lab,left=1pt of b] {2/1}; \node[lab,right=1pt of c] {3/1}; \node[lab,left=1pt of d] {4/2};
\node[lab,right=1pt of e] {5/2}; \node[lab,below=1pt of f] {6/3}; \node[lab,below=1pt of g] {7/3}; \node[lab,below=1pt of h] {8/3};
\end{tikzpicture}
\qquad
\begin{minipage}[b]{0.5\linewidth}\small
\textbf{Exemplo.} Busca a partir de $a$, com os vizinhos em ordem alfabética; os rótulos são $L/nivel$. Em azul grosso, as arestas-pai; em vermelho, as arestas-irmão ($bc$ e $fg$); em amarelo pontilhado, as arestas-primo ($de$ e $gh$); em verde tracejado, as arestas-tio ($cd$ e $ef$).
As distâncias a partir de $a$ são 1 para $b,c$; 2 para $d,e$; e 3 para $f,g,h$. O grafo não é bipartido: o triângulo $abc$ é um ciclo ímpar, e isso aparece nas arestas-irmão e primo.
\end{minipage}
\end{center}

\section{Propriedades e aplicações}
A floresta de largura alcança todos os vértices, e somente as arestas-pai pertencem a ela. As demais arestas são caracterizadas pelos níveis:
\begin{itemize}
  \item $vw$ é aresta-\textbf{pai} se e somente se $nivel(w)=nivel(v)+1$ e $L(w)=0$ no momento da visita;
  \item $vw$ é aresta-\textbf{tio} se e somente se $nivel(w)=nivel(v)+1$ e $L(w)\ne0$ no momento da visita;
  \item $vw$ é aresta-\textbf{irmão} se e somente se $nivel(w)=nivel(v)$ e $pai(w)=pai(v)$;
  \item $vw$ é aresta-\textbf{primo} se e somente se $nivel(w)=nivel(v)$ e $pai(w)\ne pai(v)$.
\end{itemize}

\begin{teo}[Propriedade fundamental]
Para toda aresta $vw$ de um grafo e toda floresta de largura, $|nivel(v)-nivel(w)|\le1$.
\end{teo}
A demonstração se reduz ao caso das arestas-tio e usa uma observação sobre a fila: em qualquer momento da busca, os vértices na fila pertencem a no máximo dois níveis consecutivos, em ordem não decrescente de nível. Desse teorema decorrem as principais aplicações.
\begin{itemize}
  \item \textbf{Distâncias.} Numa busca com raiz $x$, $\operatorname{dist}(x,v)=nivel(v)$ para todo $v$, e o caminho de $x$ a $v$ na árvore é um caminho \textbf{mínimo} em número de arestas. Por exemplo, o menor caminho da entrada à saída de um labirinto, dado por uma matriz de casas livres e paredes, é obtido por uma busca em largura no grafo das casas livres.
  \item \textbf{Bipartição.} Um grafo é bipartido se e somente se a sua floresta de largura não tem arestas-irmão nem arestas-primo, que fecham ciclos ímpares. Uma 2-coloração é dada pela paridade do nível.
  \item \textbf{Complexidade.} $O(n+m)$ com listas de adjacência, pois cada vértice entra na fila uma vez e cada lista é percorrida uma vez; $O(n^2)$ com matriz de adjacência.
\end{itemize}

\paragraph{Pilha ou fila.} Considere o seguinte esquema genérico: marque a raiz e insira-a numa estrutura $D$; enquanto $D$ não estiver vazia, seja $v$ o primeiro elemento de $D$; se $v$ tem um vizinho desmarcado $w$, marque $w$, faça $pai(w):=v$ e insira $w$ em $D$; senão, remova $v$ de $D$. Se $D$ é uma \textbf{pilha}, esse esquema é a busca em profundidade (Capítulo~\ref{cap:dfs}); se $D$ é uma \textbf{fila}, é a busca em largura; e, com uma \textbf{fila de prioridades}, chega-se ao algoritmo de Dijkstra (Capítulo~\ref{cap:dijkstra}).

\begin{cuidado}[Armadilhas da busca em largura]
\begin{itemize}
  \item Esquecer o laço externo: num grafo desconexo, a busca a partir de um vértice só alcança a sua componente.
  \item O caminho da busca em largura é mínimo apenas em \emph{número de arestas}; com pesos, é preciso o algoritmo de Dijkstra.
  \item Uma aresta entre vértices do mesmo nível deve ser classificada uma única vez; daí o teste ``$w\in F$'' no pseudocódigo.
\end{itemize}
\end{cuidado}

\begin{chave}
Fila; ``esgotar o pai antes de processar os filhos''. Arestas-pai (árvore), irmão, primo e tio. Para toda aresta, $|nivel(v)-nivel(w)|\le1$. $nivel(v)=\operatorname{dist}(raiz,v)$ em arestas, e o caminho na árvore é mínimo. Bipartido $\iff$ sem arestas-irmão e primo. $O(n+m)$ com listas.
\end{chave}

% ----------------------------------------------------------------
\section{Exercícios}
\excurso
\begin{exercicio}{bl-arv}[Slides]
Dado um grafo conexo $G$ e uma árvore geradora $T$ de $G$, decida se $T$ é uma árvore de largura de $G$, isto é, se existe uma busca em largura em $G$ que produz $T$.
\end{exercicio}
\begin{exercicio}{bl-pilha}[Slides]
Mostre que o esquema genérico da §\ref{sec:bfs} funciona como uma busca em profundidade se $D$ for uma pilha e como uma busca em largura se $D$ for uma fila.
\end{exercicio}
\begin{exercicio}{bl-apl}[Slides]
Descreva, usando a busca em largura: (a) o cálculo das distâncias de um vértice $x$ a todos os demais; (b) o menor caminho da entrada à saída de um labirinto dado por uma matriz; (c) o teste de bipartição e a 2-coloração.
\end{exercicio}

\excompl
\begin{exercicio}{bl-sim}[Complementar]
Execute a busca em largura a partir de $s$ no grafo com as arestas $rs$, $rv$, $sw$, $wt$, $wx$, $tx$, $tu$, $xu$, $xy$ e $uy$ (vizinhos em ordem alfabética). Dê a ordem de visita, os níveis, os pais e a classificação das arestas que não são de árvore. O grafo é bipartido?
\end{exercicio}
\begin{exercicio}{bl-diam}[Complementar]
O \emph{diâmetro} de uma árvore é o maior número de arestas de um caminho entre dois vértices. Mostre que duas buscas em largura bastam: a partir de qualquer vértice, encontre o vértice $a$ mais distante; a partir de $a$, o mais distante, $b$; o diâmetro é $\operatorname{dist}(a,b)$. Aplique à árvore com as arestas $12$, $13$, $24$, $25$, $36$ e $57$.
\end{exercicio}
\begin{exercicio}{bl-num}[Complementar]
Adapte a busca em largura para contar, para cada vértice $v$, \emph{quantos} caminhos mínimos (em número de arestas) existem da raiz até $v$.
\end{exercicio}
\begin{exercicio}{bl-prova}[Complementar]
Prove a propriedade fundamental: para toda aresta $vw$, $|nivel(v)-nivel(w)|\le1$. Conclua que $nivel(v)=\operatorname{dist}(raiz,v)$.
\end{exercicio}
\begin{exercicio}{bl-grade}[Complementar]
Num labirinto $4\times4$ com as linhas \texttt{S..\#}, \texttt{\#\#.\#}, \texttt{....} e \texttt{\#.\#E} (\texttt{\#} é parede; anda-se na horizontal ou na vertical), qual o menor número de passos de \texttt{S} a \texttt{E}? Dê o caminho e a distância de \texttt{S} a cada casa livre.
\end{exercicio}
\begin{exercicio}{bl-dig}[Complementar]
Como fica a busca em largura num digrafo, percorrendo $N_{out}(v)$? Ela ainda calcula distâncias? Que tipos de arestas podem aparecer entre níveis?
\end{exercicio}
\begin{exercicio}{bl-cam}[Complementar]
Escreva um procedimento que, depois de uma busca em largura com raiz $x$, imprime um caminho mínimo de $x$ até um vértice $y$, ou informa que não há caminho.
\end{exercicio}
\begin{exercicio}{bl-custo}[Complementar]
Justifique os custos $O(n+m)$ da busca em largura com listas de adjacência e $O(n^2)$ com matriz de adjacência. Em que tipo de grafo as duas representações dão o mesmo custo assintótico?
\end{exercicio}
\begin{exercicio}{bl-01}[Complementar \estrela]
Num grafo cujas arestas têm peso 0 ou 1, mostre como calcular as distâncias a partir de uma raiz em $O(n+m)$, usando uma fila de duas pontas (inserindo no início os vértices alcançados por arestas de peso 0 e no fim os alcançados por arestas de peso 1).
\end{exercicio}
\begin{exercicio}{bl-multi}[Complementar]
Num mapa, há várias saídas de emergência. Calcule, para cada local, a distância (em número de passos) até a saída mais próxima, com uma única busca em largura.
\end{exercicio}
\begin{exercicio}{bl-bip}[Complementar]
Use a busca em largura para decidir se o ciclo com 5 vértices e o ciclo com 6 vértices são bipartidos, mostrando as arestas-irmão ou primo encontradas.
\end{exercicio}
\begin{exercicio}{bl-comp}[Complementar]
Mostre como contar as componentes conexas de um grafo e rotular cada vértice com o número da sua componente usando a busca em largura, em $O(n+m)$.
\end{exercicio}
\begin{exercicio}{bl-cavalo}[Complementar]
Num tabuleiro de xadrez $8\times8$, qual o menor número de movimentos de um cavalo para ir da casa a1 à casa h8? Modele como um grafo e use a busca em largura.
\end{exercicio}
\begin{exercicio}{bl-jarros}[Complementar]
Com um jarro de 3 litros e outro de 5 litros, sem marcações, pode-se encher um jarro, esvaziá-lo ou despejar o conteúdo de um no outro até esvaziar o primeiro ou encher o segundo. Qual o menor número de operações para obter exatamente 4 litros num dos jarros? Modele os estados como vértices de um grafo.
\end{exercicio}
\begin{exercicio}{bl-impar}[Complementar]
Mostre que toda aresta-irmão ou aresta-primo fecha um ciclo de comprimento ímpar com os caminhos da floresta de largura. Por que uma aresta-tio não fecha, necessariamente, um ciclo ímpar?
\end{exercicio}
\begin{exercicio}{bl-girth}[Complementar \estrela]
A \emph{cintura} de um grafo não direcionado é o comprimento do menor ciclo. Mostre como calculá-la em $O(n(n+m))$ com buscas em largura.
\end{exercicio}
\begin{exercicio}{bl-k33}[Complementar]
Execute a busca em largura no grafo bipartido completo $K_{3,3}$ (lados $\{a,b,c\}$ e $\{x,y,z\}$) a partir de $a$. Quais são os níveis? Que tipos de arestas aparecem?
\end{exercicio}

% ================================================================
\chapter{Caminhos mínimos de uma origem: Dijkstra}\label{cap:dijkstra}\label{sec:dijkstra}
% ================================================================
\begin{objetivos}
\begin{itemize}
  \item O problema dos caminhos mínimos a partir de uma origem, com pesos não negativos.
  \item O algoritmo de Dijkstra como uma busca guiada por uma fila de prioridades.
  \item A prova de correção e o papel dos pesos não negativos.
  \item As implementações com vetor, heap binário e heap de Fibonacci, e quando usar cada uma.
\end{itemize}
\end{objetivos}

Quando as arestas têm pesos, a busca em largura deixa de fornecer distâncias. O algoritmo de Dijkstra generaliza a busca: em vez de uma fila, usa uma \textbf{fila de prioridades} e processa sempre o vértice mais próximo da origem que ainda não teve a distância fixada. É o primeiro algoritmo deste livro em que a escolha da implementação da fila de prioridades --- vetor, heap binário ou heap de Fibonacci --- muda a complexidade de forma decisiva. O exemplo do capítulo usa o digrafo da §\ref{sec:fwmatriz}.

\section{O algoritmo}
\textbf{Problema:} dado um digrafo com pesos \textbf{não negativos} e uma raiz $s$, calcular os caminhos mínimos de $s$ a todos os vértices. O algoritmo mantém, para cada vértice $w$, um rótulo $L(w)$, que é o custo do melhor caminho de $s$ a $w$ encontrado até o momento, e um ponteiro $pai(w)$. A cada passo, retira da fila de prioridades o vértice de menor rótulo, fixa esse rótulo como definitivo e tenta melhorar os rótulos dos seus vizinhos.
\begin{pseudo}{Algoritmo de Dijkstra}
\begin{algorithmic}
\Procedure{Dijkstra}{$G,W,s$}
  \State para todo $w$: $L(w):=\infty$;\ $pai(w):=\nil$;\ $w$ desmarcado
  \State $L(s):=0$;\ insira $s$ em $D$
  \While{$D\ne\emptyset$}
    \State $v:=$ elemento de $D$ com menor $L(v)$;\ remova $v$ de $D$;\ marque $v$ \Comment{$L(v)$ agora é definitivo}
    \For{\textbf{cada} $w\in N_{out}(v)$ desmarcado}
      \If{$w\notin D$} \State insira $w$ em $D$ com prioridade $L(w)$ \EndIf
      \If{$L(v)+W(v,w)<L(w)$} \Comment{``relaxamento'' da aresta $vw$}
        \State $L(w):=L(v)+W(v,w)$;\ reposicione $w$ em $D$;\ $pai(w):=v$
      \EndIf
    \EndFor
  \EndWhile
\EndProcedure
\end{algorithmic}
\end{pseudo}
\begin{exemplo}[Raiz $a$ no digrafo da §\ref{sec:fwmatriz}]\leavevmode
\begin{center}\small
\begin{tabular}{l|ccccc|l}
Removido & $L(a)$ & $L(b)$ & $L(c)$ & $L(d)$ & $L(e)$ & mudanças em $pai$\\ \hline
$a$ (0) & \textbf{0} & 1 & 3 & $\infty$ & 6 & $b,c,e\leftarrow a$\\
$b$ (1) & & \textbf{1} & 2 & 4 & 6 & $c,d\leftarrow b$\\
$c$ (2) & & & \textbf{2} & 3 & 6 & $d\leftarrow c$\\
$d$ (3) & & & & \textbf{3} & 5 & $e\leftarrow d$\\
$e$ (5) & & & & & \textbf{5} & ---
\end{tabular}
\end{center}
Os ponteiros $pai$ formam a \textbf{árvore de caminhos mínimos}: $a\to b\to c\to d\to e$. Por exemplo, o caminho mínimo de $a$ a $e$ é $a,b,c,d,e$, com custo $L(e)=5$, embora a aresta direta custe 6. A primeira linha da matriz de distâncias calculada pelo Floyd-Warshall (Capítulo~\ref{cap:fw}) confirma os valores $0,1,2,3,5$.
\end{exemplo}
Para um grafo não direcionado, basta trocar cada aresta $xy$ de peso $p$ pelas arestas direcionadas $x\to y$ e $y\to x$, ambas de peso $p$. E, se todos os pesos são iguais a 1, o algoritmo se comporta como a busca em largura: a fila de prioridades passa a liberar os vértices em ordem de nível, e $L(w)$ é a distância em número de arestas.

\section{Correção}
\begin{teo}
Com pesos não negativos, quando um vértice $v$ é removido de $D$, vale $L(v)=\operatorname{dist}(s,v)$.
\end{teo}
\begin{proof}
Sempre vale $L(v)\ge\operatorname{dist}(s,v)$, pois $L(v)$ é o custo de algum caminho. Suponha que $u$ seja o primeiro vértice removido com $L(u)>\operatorname{dist}(s,u)$. Num caminho mínimo de $s$ a $u$, seja $y$ o primeiro vértice ainda não marcado e $x$ o vértice anterior a ele, que já foi marcado e, portanto, tem $L(x)=\operatorname{dist}(s,x)$. Quando $x$ foi marcado, a aresta $xy$ foi relaxada, e então $L(y)\le\operatorname{dist}(s,x)+W(x,y)=\operatorname{dist}(s,y)\le\operatorname{dist}(s,u)<L(u)$. A segunda desigualdade usa que os \textbf{pesos não são negativos}: o trecho do caminho de $y$ até $u$ custa pelo menos 0. Mas então $u$ não tinha o menor rótulo em $D$, uma contradição.
\end{proof}
A hipótese sobre os pesos é essencial. Com as arestas $x\to y$ (1), $x\to z$ (2) e $z\to y$ ($-2$), o algoritmo fixa $L(y)=1$, mas o caminho $x,z,y$ custa 0. Para pesos negativos sem ciclos negativos, existem outros algoritmos, como o de Bellman-Ford, em $O(nm)$; para todos os pares, o Floyd-Warshall.

\section{Implementações e complexidade}
\begin{center}\small
\begin{tabularx}{\linewidth}{@{}>{\raggedright\arraybackslash}p{5.2cm}>{\raggedright\arraybackslash}p{3.6cm}X@{}}
\toprule
Estrutura $D$ & Custo & Quando usar\\
\midrule
Vetor $L$ com varredura (matriz de adjacência) & $O(n^2)$ & grafo \textbf{denso}, com $m=\Theta(n^2)$\\
Heap binário de mínimo com a posição de cada vértice & $O((n+m)\log n)$ & grafo esparso\\
Heap de Fibonacci & $O(m+n\log n)$ & teórico: redução de chave em $O(1)$\am\\
\bottomrule
\end{tabularx}
\end{center}
Os custos vêm de $n$ inserções, $n$ remoções do mínimo e até $m$ reduções de prioridade (Capítulos~\ref{cap:heap} e~\ref{cap:heapav}). No heap binário, ``reposicionar'' é uma operação \emph{subir}, e exige um vetor $pos[w]$ com a posição de cada vértice no heap.

\begin{cuidado}[Armadilhas do algoritmo de Dijkstra]
\begin{itemize}
  \item Usar o algoritmo com uma aresta negativa: o rótulo de um vértice já fixado pode estar errado.
  \item Esquecer o vetor de posições no heap binário: sem ele, encontrar o vértice a reposicionar custaria $O(n)$.
  \item Confundir com a árvore geradora mínima: a árvore de caminhos mínimos minimiza a distância de cada vértice à raiz; a AGM (Capítulo~\ref{cap:agm}) minimiza a \emph{soma} dos pesos da árvore. Em geral, são árvores diferentes.
  \item Num grafo não direcionado, cada aresta deve ser considerada nos dois sentidos.
\end{itemize}
\end{cuidado}

\begin{chave}
Fila de prioridades na chave $L$; remove-se o menor $L$ ainda não fixado; relaxa-se $L(w):=\min\{L(w),L(v)+W(v,w)\}$. Exige pesos não negativos. Custos: $O(n^2)$ com vetor, $O((n+m)\log n)$ com heap binário, $O(m+n\log n)$ com heap de Fibonacci. Os ponteiros $pai$ formam a árvore de caminhos mínimos. Com todos os pesos iguais a 1, é a busca em largura.
\end{chave}

% ----------------------------------------------------------------
\section{Exercícios}
\excurso
\begin{exercicio}{dj-bfs}[Slides]
Mostre que o algoritmo de Dijkstra se comporta como uma busca em largura se todos os pesos são iguais a 1: a árvore de caminhos mínimos é uma árvore de largura e cada $L(w)$ é a distância em arestas.
\end{exercicio}
\begin{exercicio}{dj-nd}[Slides]
Como encontrar caminhos mínimos de uma raiz a todos os vértices num grafo \emph{não direcionado} com pesos positivos?
\end{exercicio}
\begin{exercicio}{dj-slide}[Slides]
Execute o algoritmo de Dijkstra a partir de $a$ no digrafo da §\ref{sec:fwmatriz}, mostrando os rótulos $L$ e os pais a cada vértice removido.
\end{exercicio}

\excompl
\begin{exercicio}{dj-sim}[Complementar]
Execute o algoritmo de Dijkstra a partir de $s$ no digrafo com as arestas $s\to t$ (10), $s\to y$ (5), $t\to x$ (1), $t\to y$ (2), $x\to z$ (4), $y\to t$ (3), $y\to x$ (9), $y\to z$ (2), $z\to s$ (7) e $z\to x$ (6). Dê a ordem de remoção, as distâncias e a árvore de caminhos mínimos.
\end{exercicio}
\begin{exercicio}{dj-neg}[Complementar]
Dê um digrafo com três vértices e uma aresta negativa, sem ciclo negativo, em que o algoritmo de Dijkstra erra. Somar uma constante a todos os pesos, para torná-los positivos, resolve? Por quê?
\end{exercicio}
\begin{exercicio}{dj-conf}[Complementar]
Cada aresta $uv$ de uma rede funciona com probabilidade $p(u,v)\in(0,1]$, independentemente das outras. Como encontrar o caminho \emph{mais confiável} (de maior probabilidade de funcionar) de $s$ a $t$ com o algoritmo de Dijkstra?
\end{exercicio}
\begin{exercicio}{dj-fontes}[Complementar]
Há várias fontes $s_1,\dots,s_k$ (por exemplo, quartéis de bombeiros). Como calcular, para cada vértice, a distância até a fonte mais próxima com uma única execução do algoritmo de Dijkstra?
\end{exercicio}
\begin{exercicio}{dj-escolha}[Complementar]
Para um grafo com $n=10^5$ vértices e $m=3\cdot10^5$ arestas, e para outro com $n=10^4$ e $m=5\cdot10^7$, qual implementação do algoritmo de Dijkstra você usaria, vetor ou heap binário? Estime o número de operações.
\end{exercicio}
\begin{exercicio}{dj-und}[Complementar]
Execute o algoritmo de Dijkstra a partir do vértice 1 no grafo não direcionado com as arestas $12$ (7), $13$ (9), $16$ (14), $23$ (10), $24$ (15), $34$ (11), $36$ (2), $45$ (6) e $56$ (9). Dê a ordem de remoção, as distâncias e os pais.
\end{exercicio}
\begin{exercicio}{dj-zero}[Complementar]
O algoritmo de Dijkstra continua correto se algumas arestas têm peso zero? Justifique com a prova de correção.
\end{exercicio}
\begin{exercicio}{dj-alvo}[Complementar]
Se só interessa a distância de $s$ a um único vértice $t$, em que momento o algoritmo pode parar? Por que isso está correto?
\end{exercicio}
\begin{exercicio}{dj-vert}[Complementar]
Cada \emph{vértice} $v$ tem um custo $c(v)\ge0$ de passagem, além dos pesos das arestas. Como transformar o problema para aplicar o algoritmo de Dijkstra?
\end{exercicio}
\begin{exercicio}{dj-k}[Complementar]
Queremos o caminho mínimo de $s$ a $t$ que usa no máximo $k$ arestas. Por que o algoritmo de Dijkstra não resolve isso diretamente? Mostre como resolver com um grafo de estados $(v,j)$, em que $j$ é o número de arestas usadas.
\end{exercicio}
\begin{exercicio}{dj-gargalo}[Complementar]
A capacidade de um caminho é a menor capacidade das suas arestas. Adapte o algoritmo de Dijkstra para encontrar, a partir de $s$, o caminho de maior capacidade até cada vértice.
\end{exercicio}
\begin{exercicio}{dj-conta}[Complementar]
Com pesos positivos, adapte o algoritmo de Dijkstra para contar quantos caminhos mínimos distintos existem de $s$ até cada vértice.
\end{exercicio}
\begin{exercicio}{dj-dag}[Complementar]
Num digrafo \emph{acíclico} com pesos quaisquer, inclusive negativos, mostre como calcular os caminhos mínimos a partir de $s$ em $O(n+m)$, sem fila de prioridades.
\end{exercicio}
\begin{exercicio}{dj-longo}[Complementar]
Por que não basta trocar ``menor'' por ``maior'' no algoritmo de Dijkstra para encontrar caminhos simples \emph{mais longos}?
\end{exercicio}
\begin{exercicio}{dj-dial}[Complementar \estrela]
Se os pesos são inteiros entre 0 e $C$, mostre como implementar o algoritmo de Dijkstra em $O(m+nC)$, usando um vetor de baldes indexado pelo valor do rótulo.
\end{exercicio}
\begin{exercicio}{dj-inv}[Complementar]
Prove que, durante toda a execução, $L(v)\ge\operatorname{dist}(s,v)$ para todo vértice $v$, e que $L(v)$ nunca aumenta.
\end{exercicio}
\begin{exercicio}{dj-fib}[Complementar]
Quantas operações de cada tipo (inserção, remoção do mínimo, redução de chave) o algoritmo de Dijkstra faz no pior caso? Use isso para justificar o custo $O(m+n\log n)$ com heap de Fibonacci.
\end{exercicio}

% ================================================================
\chapter{Busca em profundidade}\label{cap:dfs}\label{sec:dfs}\index{busca em profundidade}
% ================================================================
\begin{objetivos}
\begin{itemize}
  \item A busca em profundidade, com as profundidades de entrada e de saída de cada vértice.
  \item A árvore de profundidade, as arestas de retorno e o teorema dos intervalos aninhados.
  \item Aplicações: conexidade, caminhos, ciclos e bipartição.
  \item Articulações, pontes e blocos, com o valor $back(v)$.
\end{itemize}
\end{objetivos}

Na busca em profundidade, a próxima aresta examinada parte sempre do vértice visitado mais recentemente: a busca avança o quanto puder e só volta quando não há mais para onde ir. Os instantes de entrada e de saída de cada vértice formam intervalos aninhados que revelam a estrutura do grafo: componentes, ciclos e bipartição e, com um valor auxiliar chamado $back(v)$, as articulações e os blocos --- tudo em tempo linear. Este capítulo trata de grafos não direcionados; o próximo, de digrafos.

\begin{intuicao}
É a estratégia de quem explora um labirinto desenrolando um novelo de lã: em cada encruzilhada, segue por um corredor ainda não explorado; quando chega a um beco sem saída, ou a um lugar onde já esteve, volta pelo fio até a última encruzilhada com corredores novos. O fio é a pilha de recursão: ele contém exatamente o caminho da entrada até a posição atual. A busca em largura, ao contrário, avança em ondas, como um grupo que se divide em todas as direções ao mesmo tempo.
\end{intuicao}

\section{O algoritmo}\label{sec:dfsalg}
A regra da busca em profundidade é ``\emph{esgotar o filho antes de esgotar o pai}'': a busca é naturalmente recursiva, ou, de forma equivalente, usa uma pilha. Um relógio global $t$ marca, para cada vértice $v$, a sua profundidade de entrada $PE(v)$, quando ele é alcançado, e a sua profundidade de saída $PS(v)$, quando todos os seus vizinhos já foram examinados.
\begin{pseudo}{Busca em profundidade, com o cálculo de $back(v)$ e dos blocos}
\begin{algorithmic}
\Procedure{BuscaProfundidade}{$G$}
  \State $t:=0$;\ para todo $v$: $PE(v):=0$;\ $PS(v):=0$;\ $pai(v):=\nil$
  \While{existe $v$ com $PE(v)=0$} \Call{P}{$v$} \Comment{nova raiz; uma por componente}
  \EndWhile
\EndProcedure
\Procedure{P}{$v$}
  \State $t:=t+1$;\ $PE(v):=t$;\ $back(v):=PE(v)$
  \For{\textbf{cada} $w\in N(v)$}
    \If{$PE(w)=0$} \Comment{aresta de \textbf{árvore}}
      \State $pai(w):=v$;\ empilhe $vw$;\ \Call{P}{$w$}
      \If{$back(w)\ge PE(v)$} \State desempilhe e imprima as arestas até $vw$ \Comment{um bloco} \EndIf
      \State $back(v):=\min\{back(v),\,back(w)\}$
    \ElsIf{$PS(w)=0$ \textbf{e} $w\ne pai(v)$} \Comment{aresta de \textbf{retorno}\index{aresta!de retorno}}
      \State empilhe $vw$;\ $back(v):=\min\{back(v),\,PE(w)\}$
    \EndIf
  \EndFor
  \State $t:=t+1$;\ $PS(v):=t$
\EndProcedure
\end{algorithmic}
\end{pseudo}
As linhas com $back$ e com a pilha de arestas só são necessárias para encontrar articulações e blocos (§\ref{sec:articulacoes}); sem elas, o procedimento é a busca básica.

\begin{figura}[Busca em profundidade a partir de $e$, com os vizinhos em ordem alfabética; os rótulos são $PE/PS$. Em azul grosso, a árvore de profundidade; em vermelho tracejado, as arestas de retorno.\label{fig:dfs}]
\begin{tikzpicture}[arv/.style={very thick,MidnightBlue}, ret/.style={thick,dashed,BrickRed}, lab/.style={font=\scriptsize,text=gray!70!black}]
\node[vtx] (a) at (0,3.6) {$a$}; \node[vtx] (b) at (1.9,3.6) {$b$};
\node[vtx] (c) at (0,2.4) {$c$}; \node[vtx] (d) at (1.9,2.4) {$d$};
\node[vtx] (e) at (0,1.2) {$e$}; \node[vtx] (f) at (1.9,1.2) {$f$};
\node[vtx] (g) at (0.95,0) {$g$};
\draw[arv] (e)--(c); \draw[arv] (c)--(a); \draw[arv] (a)--(b); \draw[arv] (b)--(d); \draw[arv] (e)--(f); \draw[arv] (f)--(g);
\draw[ret] (b)--(c); \draw[ret] (d)--(c); \draw[ret] (d)--(e); \draw[ret] (g)--(e);
\node[lab,left=1pt of a] {3/8}; \node[lab,right=1pt of b] {4/7}; \node[lab,left=1pt of c] {2/9}; \node[lab,right=1pt of d] {5/6};
\node[lab,left=1pt of e] {1/14}; \node[lab,right=1pt of f] {10/13}; \node[lab,below=1pt of g] {11/12};
\end{tikzpicture}
\end{figura}

\begin{exemplo}[A busca passo a passo]\label{ex:dfs}
Na Figura~\ref{fig:dfs}, a busca começa em $e$ (instante 1) e examina os vizinhos em ordem alfabética. O primeiro é $c$, que entra (2); de $c$, o primeiro vizinho não visitado é $a$ (3), e de $a$, o $b$ (4). Em $b$, o vizinho $c$ já está ativo --- a sua chamada ainda não terminou ---, e $bc$ é uma aresta de retorno; o vizinho $d$ é novo e entra (5). Em $d$, os vizinhos $b$ (o pai), $c$ e $e$ já estão ativos: $dc$ e $de$ são arestas de retorno. Sem vizinhos novos, $d$ sai (6); depois saem $b$ (7), $a$ (8) e $c$ (9). De volta a $e$, o vizinho $d$ já terminou e não é examinado de novo; $f$ é novo (10), e dele se chega a $g$ (11), cujo vizinho $e$ está ativo ($ge$ é de retorno). Saem $g$ (12), $f$ (13) e $e$ (14).

No instante 5, quando $d$ entra, a pilha de recursão é $e,c,a,b,d$: exatamente o caminho da raiz até $d$ na árvore. É por isso que uma aresta para um vértice ativo sempre vai a um ancestral.
\end{exemplo}

\section{Propriedades}\label{sec:dfsprop}
\begin{itemize}
  \item Se $G$ é conexo, as arestas de árvore formam uma \textbf{árvore geradora}, a árvore de profundidade. Se não é, o laço externo produz uma floresta, com uma árvore por componente.
  \item Num grafo \textbf{não direcionado}, só existem arestas de \textbf{árvore} e de \textbf{retorno}. Toda aresta de retorno liga um vértice $v$ a um \textbf{ancestral} de $v$ na árvore e fecha um ciclo.
  \item \textbf{Intervalos de vida}\index{intervalo de vida}: seja $I(v)=[PE(v),PS(v)]$. Dois intervalos são disjuntos ou um contém o outro (o \emph{teorema dos parênteses}\index{teorema dos parenteses@teorema dos parênteses}), e $v$ é descendente de $w$ se e somente se $PE(w)<PE(v)$ e $PS(v)<PS(w)$: $v$ ``entra depois e sai antes'' de $w$.
  \item \textbf{Complexidade}: $O(n+m)$ com listas de adjacência. Cada vértice é visitado uma vez, e cada aresta é examinada duas vezes, uma a partir de cada extremo.
\end{itemize}

\begin{figura}[Os intervalos de vida $[PE,PS]$ da busca da Figura~\ref{fig:dfs}, desenhados sobre a linha do tempo. Escrevendo ``('' na entrada e ``)'' na saída de cada vértice, obtém-se $(e\,(c\,(a\,(b\,(d\ d)\,b)\,a)\,c)\,(f\,(g\ g)\,f)\,e)$: os parênteses se aninham corretamente.\label{fig:parenteses}]
\begin{tikzpicture}[x=7.5mm,y=5.5mm]
\draw[gray] (0.5,0.7) -- (14.5,0.7);
\foreach \tt in {1,...,14} {\draw[gray] (\tt,0.6) -- (\tt,0.8); \node[idx] at (\tt,1.15) {\tt};}
\foreach \nome/\ini/\fim/\prof in {e/1/14/0,c/2/9/1,a/3/8/2,b/4/7/3,d/5/6/4,f/10/13/1,g/11/12/2} {
  \fill[azul, draw=MidnightBlue] (\ini-0.15,-\prof-0.3) rectangle (\fim+0.15,-\prof+0.3);
  \node[font=\small] at ({(\ini+\fim)/2},-\prof) {$\nome$};}
\end{tikzpicture}
\end{figura}

\section{Aplicações}\label{sec:dfsapl}
\begin{enumerate}
  \item \textbf{Conexidade.} O número de vezes que o laço externo chama \textsc{P} é o número de componentes conexas; $G$ é conexo se e somente se esse número é 1.
  \item \textbf{Caminho de $v$ a $w$.} Execute \textsc{P}$(v)$. Se ao final $PE(w)=0$, não há caminho; senão, suba de $w$ até $v$ pelos ponteiros $pai$, colocando cada vértice à esquerda do caminho. Esse caminho \textbf{não} é necessariamente o mais curto; para isso, use a busca em largura. Um labirinto pode ser tratado da mesma forma, modelado como um grafo.
  \item \textbf{Ciclo.} Se a busca não encontra nenhuma aresta de retorno, $G$ é acíclico, e as arestas de árvore formam uma floresta geradora. Se encontra uma aresta de retorno $vw$, então $w$ é ancestral de $v$, e o ciclo é obtido subindo de $v$ até $w$ pelos ponteiros $pai$ e fechando com a aresta $vw$.
  \item \textbf{Bipartição.} $G$ é bipartido (2-colorível) se e somente se não contém ciclo ímpar. Durante a busca, dê à raiz a cor 0 e, a cada aresta de árvore, faça $cor(w):=1-cor(v)$. Uma aresta de retorno com $cor(w)=cor(v)$ fecha um ciclo ímpar, e então o grafo não é bipartido.
\end{enumerate}

\begin{exemplo}[As aplicações no grafo da Figura~\ref{fig:dfs}]\label{ex:dfsapl}
\textbf{Conexidade}: a chamada \textsc{P}$(e)$ alcança os sete vértices, e o laço externo não faz mais nenhuma chamada: o grafo é conexo. Num grafo com as componentes $\{1,2\}$ e $\{3,4,5\}$, o laço faria duas chamadas, uma para cada componente.
\textbf{Caminho}: para ir de $e$ até $a$, sobe-se de $a$ pelos pais, $pai(a)=c$ e $pai(c)=e$, e o caminho é $e,c,a$.
\textbf{Ciclo}: a primeira aresta de retorno encontrada é $bc$; subindo de $b$ até $c$ pelos pais ($b\to a\to c$) e fechando com $bc$, obtém-se o ciclo $c,a,b,c$.
\textbf{Bipartição}: com $cor(e)=0$, a árvore dá $cor(c)=1$, $cor(a)=0$, $cor(b)=1$ e $cor(d)=0$. A aresta de retorno $bc$ liga dois vértices de cor 1: o ciclo $c,a,b$ é ímpar, e o grafo não é bipartido.
\end{exemplo}

\section{Articulações, pontes e blocos}\label{sec:articulacoes}
Uma \textbf{articulação}\index{articulacao@articulação} é um vértice cuja remoção aumenta o número de componentes conexas. Um \textbf{bloco}\index{bloco} é um subgrafo maximal conexo e sem articulações; um bloco com uma única aresta é uma \textbf{ponte}\index{ponte}, e os demais blocos são subgrafos biconexos maximais, em que quaisquer dois vértices são ligados por dois caminhos internamente disjuntos. Os blocos particionam as \emph{arestas} do grafo; já um vértice pode estar em vários blocos, e isso acontece exatamente quando ele é uma articulação. Pensando no grafo como uma rede, articulações são nós críticos e pontes são ligações críticas.

Para encontrá-los, a busca calcula, para cada vértice $v$, o valor $back(v)$: a menor profundidade de entrada que se pode alcançar a partir de $v$ descendo por arestas da árvore e usando, no máximo, uma aresta de retorno.
\begin{teo}[Articulações]
Seja $T$ uma árvore de profundidade de $G$. A raiz de $T$ é articulação se e somente se tem pelo menos dois filhos em $T$. Um vértice $v$ que não é a raiz é articulação se e somente se tem um filho $w$ com $back(w)\ge PE(v)$; nesse caso, nenhum descendente de $w$ ``escapa'' para cima de $v$, e $w$ é chamado de \emph{demarcador} de $v$.
\end{teo}
Os blocos são obtidos com uma pilha de arestas, como no pseudocódigo: ao voltar de um filho $w$ com $back(w)\ge PE(v)$, desempilham-se as arestas até $vw$, que formam um bloco. As \textbf{pontes} são as arestas de árvore $vw$ com $back(w)>PE(v)$. Tudo isso custa $O(n+m)$.

\begin{intuicao}
Se $v$ não é a raiz e for removido, a subárvore de um filho $w$ só continua ligada ao resto do grafo se algum vértice dela tiver uma aresta de retorno que pule \emph{por cima} de $v$, para um ancestral próprio de $v$. O valor $back(w)$ mede exatamente quão alto a subárvore de $w$ consegue ``alcançar'': se $back(w)<PE(v)$, há uma rota de fuga; se $back(w)\ge PE(v)$, não há, e $v$ é uma articulação. Para uma aresta $vw$ ser ponte, a exigência é mais forte: nem o próprio $v$ pode ser alcançado sem a aresta, isto é, $back(w)>PE(v)$.
\end{intuicao}

No exemplo da Figura~\ref{fig:dfs}, todos os vértices têm $back(v)=1$. Com a raiz $e$, a regra da raiz identifica $e$ como articulação, pois ele tem dois filhos, $c$ e $f$; os blocos são $\{a,b,c,d,e\}$ e $\{e,f,g\}$. Se a busca começar em $a$, o vértice $e$ deixa de ser a raiz, e passa a ser identificado pelo outro caso do teorema: ele tem o filho $f$ com $back(f)=5\ge PE(e)=5$.

\begin{exemplo}[Articulações, pontes e blocos]\label{ex:articulacoes}
O grafo da Figura~\ref{fig:articulacoes} tem um quadrado $a,b,c,d$, uma aresta $de$, um triângulo $e,f,g$ e uma aresta pendente $bh$. A busca a partir de $a$, com os vizinhos em ordem alfabética, desce $a,b,c,d,e,f,g$ em linha ($PE=1,\dots,7$), encontra as arestas de retorno $ge$ e $da$, volta até $b$ e só então entra em $h$ ($PE(h)=13$). Os valores de $back$ são 1 para $a,b,c,d$ (por causa da aresta $da$), 5 para $e,f,g$ (por causa de $ge$) e 13 para $h$.
\begin{itemize}
  \item $e$ é articulação: o filho $f$ tem $back(f)=5\ge PE(e)=5$; o triângulo não escapa de $e$.
  \item $d$ é articulação: o filho $e$ tem $back(e)=5\ge PE(d)=4$; e, como $5>4$, a aresta $de$ é uma \textbf{ponte}.
  \item $b$ é articulação: o filho $h$ tem $back(h)=13\ge PE(b)=2$, e $bh$ também é ponte.
  \item $c$ não é articulação: o filho $d$ tem $back(d)=1<PE(c)=3$, pela aresta de retorno $da$. E a raiz $a$ tem um único filho, $b$.
\end{itemize}
Os blocos, na ordem em que a pilha de arestas os imprime, são $\{e,f,g\}$, $\{d,e\}$, $\{b,h\}$ e $\{a,b,c,d\}$.
\begin{figura}[À esquerda, o grafo do Exemplo~\ref{ex:articulacoes}, com as articulações em amarelo e as pontes em vermelho. À direita, a árvore de profundidade a partir de $a$, com as arestas de retorno tracejadas e os rótulos $PE/back$.\label{fig:articulacoes}]
\begin{tikzpicture}[x=1cm,y=1cm, baseline=(base)]
\node[vtx] (a) at (0,1) {$a$}; \node[vtx, fill=destaque] (b) at (1.2,1) {$b$}; \node[vtx] (c) at (1.2,0) {$c$}; \node[vtx, fill=destaque] (d) at (0,0) {$d$};
\node[vtx, fill=destaque] (e) at (0,-1.3) {$e$}; \node[vtx] (f) at (-0.7,-2.4) {$f$}; \node[vtx] (g) at (0.7,-2.4) {$g$};
\node[vtx] (h) at (2.4,1) {$h$};
\draw[aresta] (a)--(b) (b)--(c) (c)--(d) (d)--(a) (e)--(f) (f)--(g) (g)--(e);
\draw[aresta, BrickRed, very thick] (d)--(e) (b)--(h);
\coordinate (base) at (0,-0.6);
\end{tikzpicture}
\hspace{2cm}
\begin{tikzpicture}[x=1cm,y=0.85cm, baseline=(base), lab/.style={font=\small,text=gray!70!black}]
\foreach \nome/\yy/\rot in {a/0/{1/1},b/-1/{2/1},c/-2/{3/1},d/-3/{4/1},e/-4/{5/5},f/-5/{6/5},g/-6/{7/5}} {
  \node[vtx] (t\nome) at (0,\yy) {$\nome$}; \node[lab, right=2pt of t\nome] {\rot};}
\node[vtx] (th) at (1.8,-2) {$h$}; \node[lab, right=2pt of th] {13/13};
\draw[aresta, MidnightBlue] (ta)--(tb)--(tc)--(td)--(te)--(tf)--(tg); \draw[aresta, MidnightBlue] (tb)--(th);
\draw[thick, dashed, BrickRed] (td) to[bend left=50] (ta);
\draw[thick, dashed, BrickRed] (tg) to[bend left=50] (te);
\coordinate (base) at (0,-3);
\end{tikzpicture}
\end{figura}
\end{exemplo}

\begin{cuidado}[Armadilhas da busca em profundidade]
\begin{itemize}
  \item Esquecer o laço externo: num grafo desconexo, \textsc{P}$(v)$ só alcança a componente de $v$.
  \item Num grafo não direcionado, a aresta $vw$ com $w=pai(v)$ \emph{não} é de retorno: é a própria aresta de árvore, vista do outro extremo.
  \item O caminho encontrado pela busca em profundidade não é mínimo.
  \item Para a raiz, a regra das articulações é ``pelo menos dois filhos'', e não o teste com $back$: na raiz, $back(w)\ge PE(\text{raiz})=1$ é sempre verdadeiro.
\end{itemize}
\end{cuidado}

\begin{chave}
Recursão ou pilha; $PE$ e $PS$; intervalos de vida aninhados (descendente $\iff$ intervalo contido). Em grafos não direcionados, só há arestas de árvore e de retorno, e há ciclo $\iff$ há aresta de retorno. Aplicações: componentes, caminhos, ciclos, bipartição, articulações (raiz com pelo menos dois filhos, ou filho $w$ com $back(w)\ge PE(v)$), pontes ($back(w)>PE(v)$) e blocos (pilha de arestas). Tudo em $O(n+m)$.
\end{chave}

% ----------------------------------------------------------------
\section{Exercícios}
\excurso
\begin{exercicio}{bp-arv}[Slides]<sec:dfsprop>
Dado um grafo conexo $G$ e uma árvore geradora $T$ de $G$, decida se $T$ é uma árvore de profundidade de $G$, isto é, se existe uma busca em profundidade em $G$ que produz $T$.
\end{exercicio}
\begin{exercicio}{bp-pilha}[Slides]<fig:dfs,sec:articulacoes>
Simule o cálculo dos blocos com a pilha de arestas no grafo da figura (raiz $e$, vizinhos em ordem alfabética), mostrando o conteúdo da pilha quando cada bloco é impresso.
\end{exercicio}
\begin{exercicio}{bp-pontes}[Slides]<sec:articulacoes>
Faça pequenos acréscimos ao procedimento com $back(v)$ para listar as articulações e as pontes do grafo.
\end{exercicio}
\begin{exercicio}{bp-demarc}[Slides]<sec:articulacoes>
Verifique se cada bloco tem o seu próprio demarcador, isto é, se o número de demarcadores é igual ao número de blocos.
\end{exercicio}
\begin{exercicio}{bp-apl}[Slides]<sec:dfsapl>
Descreva como a busca em profundidade resolve: (a) contar as componentes conexas; (b) encontrar um caminho de $v$ a $w$; (c) encontrar um ciclo ou concluir que o grafo é acíclico; (d) decidir se o grafo é 2-colorível.
\end{exercicio}

\excompl
\begin{exercicio}{bp-sim}[Complementar]<ex:articulacoes>
Execute a busca em profundidade, com o cálculo de $back$, a partir do vértice 1 no grafo com as arestas $12$, $13$, $23$, $34$, $45$, $46$, $56$ e $67$ (vizinhos em ordem crescente). Dê $PE$, $PS$, $back$, as arestas de retorno, as articulações, as pontes e os blocos.
\end{exercicio}
\begin{exercicio}{bp-cruz}[Complementar]<sec:dfsprop>
Prove que, numa busca em profundidade num grafo \emph{não direcionado}, não existem arestas de cruzamento: toda aresta liga um vértice a um ancestral ou a um descendente.
\end{exercicio}
\begin{exercicio}{bp-m}[Complementar]<sec:dfsapl>
Mostre que um grafo não direcionado com $m\ge n$ arestas tem um ciclo, e use isso para decidir se há ciclo em tempo $O(n)$, independentemente de $m$.
\end{exercicio}
\begin{exercicio}{bp-iter}[Complementar]<sec:dfsalg>
Escreva a busca em profundidade sem recursão, com uma pilha explícita, mantendo os mesmos valores de $PE$ e $PS$ da versão recursiva.
\end{exercicio}
\begin{exercicio}{bp-trace}[Complementar]<ex:dfs>
Execute a busca em profundidade a partir do vértice 1 no grafo com as arestas $12$, $13$, $24$, $25$, $36$, $37$, $48$, $58$, $67$ e $78$ (vizinhos em ordem crescente). Dê $PE$, $PS$, as arestas de árvore e as de retorno. O grafo tem articulações?
\end{exercicio}
\begin{exercicio}{bp-parenteses}[Complementar]<sec:dfsprop>
Prove o teorema dos parênteses: para quaisquer vértices $u$ e $v$, os intervalos $[PE(u),PS(u)]$ e $[PE(v),PS(v)]$ são disjuntos ou um está contido no outro.
\end{exercicio}
\begin{exercicio}{bp-branco}[Complementar \estrela]<sec:dfsprop>
Prove o \emph{teorema do caminho livre}: $v$ é descendente de $u$ na floresta de profundidade se e somente se, no instante $PE(u)$, existe um caminho de $u$ a $v$ formado apenas por vértices ainda não alcançados pela busca.
\end{exercicio}
\begin{exercicio}{bp-comp}[Complementar]<sec:dfsapl>
Quantas componentes conexas tem o grafo com os vértices $1,\dots,9$ e as arestas $12$, $23$, $31$, $45$, $67$, $78$, $89$ e $96$? Quais vértices são raízes da floresta de profundidade, se a busca percorre os vértices em ordem crescente?
\end{exercicio}
\begin{exercicio}{bp-cicloex}[Complementar]<ex:dfsapl>
No grafo do Exercício~\ref{ex:bp-trace}, use a primeira aresta de retorno encontrada para exibir um ciclo.
\end{exercicio}
\begin{exercicio}{bp-bipex}[Complementar]<ex:dfsapl>
Use a 2-coloração por busca em profundidade para decidir se é bipartido o grafo com as arestas $12$, $23$, $34$, $41$ e $13$. Indique a aresta que revela o ciclo ímpar, se houver.
\end{exercicio}
\begin{exercicio}{bp-arvore}[Complementar]<sec:articulacoes>
Quais são as pontes e as articulações de uma árvore? E de um ciclo? E de um caminho com $n\ge3$ vértices?
\end{exercicio}
\begin{exercicio}{bp-blocos}[Complementar]<ex:articulacoes>
Encontre as articulações, as pontes e os blocos do grafo com as arestas $12$, $23$, $31$, $34$, $45$, $56$, $64$ e $47$.
\end{exercicio}
\begin{exercicio}{bp-biconexo}[Complementar]<sec:articulacoes>
Um grafo é \emph{biconexo} se é conexo e não tem articulações. Como decidir isso em $O(n+m)$? Mostre que, num grafo biconexo com pelo menos 3 vértices, toda aresta está em algum ciclo.
\end{exercicio}
\begin{exercicio}{bp-labirinto}[Complementar]<sec:dfsapl>
Dê um exemplo de grafo e de vértices $v$ e $w$ em que o caminho de $v$ a $w$ encontrado pela busca em profundidade é mais longo que o caminho mínimo.
\end{exercicio}
\begin{exercicio}{bp-euler}[Complementar \estrela]<sec:dfsalg>
Um \emph{circuito euleriano} percorre cada aresta exatamente uma vez e volta ao início. Sabendo que um grafo conexo tem circuito euleriano se e somente se todos os vértices têm grau par, mostre como encontrar o circuito em $O(n+m)$ com uma busca que avança por arestas ainda não usadas.
\end{exercicio}

% ================================================================
\chapter{Busca em profundidade em digrafos}\label{cap:dfsdig}\label{sec:dfsdig}
% ================================================================
\begin{objetivos}
\begin{itemize}
  \item A busca em profundidade em digrafos e os quatro tipos de aresta: árvore, retorno, avanço e cruzamento.
  \item A caracterização dos tipos de aresta pelos intervalos de vida e o teste de aciclicidade.
  \item As componentes fortemente conexas, calculadas numa única busca com o valor $old(v)$.
  \item A ordenação topológica de digrafos acíclicos.
\end{itemize}
\end{objetivos}

Em grafos direcionados, a mesma busca em profundidade do capítulo anterior produz quatro tipos de aresta, e não apenas dois; os intervalos de vida dos vértices continuam a caracterizá-los. Com essa classificação, decide-se se um digrafo tem ciclo, obtém-se uma ordenação topológica de um digrafo acíclico e calculam-se as componentes fortemente conexas --- tudo em tempo linear.

\section{Classificação das arestas}
A busca é a mesma do Capítulo~\ref{cap:dfs}, percorrendo a vizinhança de saída $N_{out}(v)$ e com o laço externo, de modo que a floresta de profundidade alcança todos os vértices. Quando a busca, a partir de $v$, examina uma aresta $vw$, há quatro possibilidades.
\begin{pseudo}{Classificação das arestas na busca em profundidade em digrafos}
\begin{algorithmic}
\Procedure{P}{$v$}
  \State $t:=t+1$;\ $PE(v):=t$
  \For{\textbf{cada} $w\in N_{out}(v)$}
    \If{$PE(w)=0$} \State $vw$ é de \textbf{árvore};\ $pai(w):=v$;\ \Call{P}{$w$}
    \ElsIf{$PS(w)=0$} \State $vw$ é de \textbf{retorno} \Comment{$w$ ainda está ativo: é ancestral de $v$}
    \ElsIf{$PE(v)<PE(w)$} \State $vw$ é de \textbf{avanço} \Comment{$w$ é um descendente que já terminou}
    \Else \State $vw$ é de \textbf{cruzamento} \Comment{$w$ terminou antes de $v$ começar}
    \EndIf
  \EndFor
  \State $t:=t+1$;\ $PS(v):=t$
\EndProcedure
\end{algorithmic}
\end{pseudo}
\begin{center}\small
\begin{tabularx}{\linewidth}{@{}l>{\raggedright\arraybackslash}p{6.2cm}X@{}}
\toprule
Tipo de $vw$ & Caracterização pelos intervalos & Observação\\
\midrule
Árvore & $I(v)\supset I(w)$ e $PE(w)=0$ na hora da visita & liga pai e filho\\
Avanço & $I(v)\supset I(w)$ e $PE(w)\ne0$ na hora da visita & vai a um descendente que não é filho\\
Retorno & $I(v)\subset I(w)$ & vai a um ancestral e \textbf{fecha um ciclo direcionado}\\
Cruzamento & $I(v)$ inteiro à \textbf{direita} de $I(w)$ & vai ``para trás'' no tempo\\
\bottomrule
\end{tabularx}
\end{center}
Nunca há uma aresta $vw$ com $I(v)$ inteiro à esquerda de $I(w)$: quando $v$ examina a aresta $vw$, o vértice $w$ ou já foi alcançado ou será visitado como descendente de $v$. Como toda aresta de retorno fecha um ciclo direcionado, e todo ciclo contém uma aresta de retorno (a que chega ao primeiro vértice do ciclo alcançado pela busca), vale: \textbf{um digrafo é acíclico se e somente se uma busca em profundidade --- qualquer uma --- não produz arestas de retorno}.

\section{Componentes fortemente conexas}
Uma \textbf{componente fortemente conexa} (CFC) de um digrafo é um subdigrafo maximal em que, para quaisquer dois vértices $v$ e $w$, existe um caminho direcionado de $v$ para $w$ \emph{e} um de $w$ para $v$. As CFC particionam os \emph{vértices}; algumas arestas não pertencem a nenhuma CFC, e essas são, numa busca, arestas de árvore, de avanço ou de cruzamento. As componentes não dependem da ordem da busca, embora o tipo de cada aresta dependa.

Para calculá-las numa única busca, define-se $old(v)$ como a profundidade de entrada do vértice mais antigo, \textbf{na mesma CFC} que $v$, que pode ser alcançado a partir de $v$ descendo por arestas de árvore e usando, no máximo, uma aresta de retorno ou de cruzamento. Um vértice $v$ é \textbf{forte} --- é a ``raiz'' da sua componente na floresta --- se $old(v)=PE(v)$ no momento em que $v$ sai da busca. Os vértices são empilhados numa pilha $Q$ ao serem alcançados; ao encontrar um vértice forte, desempilham-se todos os vértices até ele, que formam uma componente.
\begin{pseudo}{Componentes fortemente conexas: $O(n+m)$}
\begin{algorithmic}
\Procedure{P}{$v$}
  \State $t:=t+1$;\ $PE(v):=t$;\ $old(v):=PE(v)$;\ empilhe $v$ em $Q$
  \For{\textbf{cada} $w\in N_{out}(v)$}
    \If{$PE(w)=0$} \State $pai(w):=v$;\ \Call{P}{$w$};\ $old(v):=\min\{old(v),\,old(w)\}$ \Comment{árvore}
    \ElsIf{$PS(w)=0$} \State $old(v):=\min\{old(v),\,PE(w)\}$ \Comment{retorno}
    \ElsIf{$PE(v)<PE(w)$} \State nada a fazer \Comment{avanço}
    \ElsIf{$w$ está em $Q$} \State $old(v):=\min\{old(v),\,PE(w)\}$ \Comment{cruzamento dentro de uma componente aberta}
    \EndIf
  \EndFor
  \State $t:=t+1$;\ $PS(v):=t$
  \If{$old(v)=PE(v)$} \State desempilhe de $Q$ todos os vértices até $v$, inclusive: \textbf{eles formam uma CFC} \EndIf
\EndProcedure
\end{algorithmic}
\end{pseudo}

\begin{center}
\begin{tikzpicture}[v/.style={circle,draw,fill=yellow!40,inner sep=1.2pt,minimum size=5mm,font=\small},
  >={Stealth[length=2.2mm]}, arv/.style={->,very thick,MidnightBlue}, ret/.style={->,thick,dashed,BrickRed},
  ava/.style={->,thick,dotted,orange!85!black}, cru/.style={->,thick,ForestGreen}, lab/.style={font=\scriptsize,text=gray!70!black}]
\node[v] (a) at (0,1.5) {$a$}; \node[v] (b) at (1.8,1.5) {$b$}; \node[v] (c) at (3.6,1.5) {$c$}; \node[v] (d) at (5.4,1.5) {$d$};
\node[v] (e) at (0,0) {$e$}; \node[v] (f) at (1.8,0) {$f$}; \node[v] (g) at (3.6,0) {$g$}; \node[v] (h) at (5.4,0) {$h$};
\draw[arv] (a)--(b); \draw[arv] (b)--(c); \draw[arv] (b)--(e); \draw[arv] (c) to[bend left=18] (d); \draw[arv] (c)--(g);
\draw[arv] (d)--(h); \draw[arv] (g) to[bend left=18] (f);
\draw[ret] (d) to[bend left=18] (c); \draw[ret] (e)--(a); \draw[ret] (f) to[bend left=18] (g);
\draw[ava] (b)--(f); \draw[cru] (e)--(f); \draw[cru] (g)--(h);
\node[lab,above=1pt of a] {1/16}; \node[lab,above=1pt of b] {2/15}; \node[lab,above=1pt of c] {3/12}; \node[lab,above=1pt of d] {4/7};
\node[lab,below=1pt of e] {13/14}; \node[lab,below=1pt of f] {9/10}; \node[lab,below=1pt of g] {8/11}; \node[lab,below=1pt of h] {5/6};
\end{tikzpicture}
\end{center}
\begin{exemplo}[Componentes fortemente conexas]\label{ex:cfc}
Busca a partir de $a$, com os vizinhos em ordem alfabética; os rótulos são $PE/PS$. Em azul, as arestas de árvore; em vermelho tracejado, as de retorno ($dc$, $ea$ e $fg$); em laranja pontilhado, a de avanço ($bf$); em verde, as de cruzamento ($ef$ e $gh$). Os valores de $old$ são $1$ para $a$, $b$ e $e$; $3$ para $c$ e $d$; $8$ para $g$ e $f$; e $5$ para $h$. Os vértices fortes, na ordem em que saem da busca, são $h$, $g$, $c$ e $a$, e as componentes são $\{h\}$, $\{f,g\}$, $\{c,d\}$ e $\{a,b,e\}$. A aresta $ef$ é de cruzamento, mas $f$ já saiu de $Q$ --- a sua componente já tinha sido fechada ---, e por isso ela não altera $old(e)$.
\end{exemplo}

\section{Ordenação topológica}\label{sec:topo}
Uma \textbf{ordenação topológica} de um digrafo acíclico é uma ordem $v_1,\dots,v_n$ dos vértices tal que toda aresta $v_iv_j$ tem $i<j$. Ela aparece sempre que há tarefas com dependências: se cada aresta $a\to b$ indica que $b$ só pode ser feita depois de $a$, uma ordenação topológica é uma sequência viável de execução.

\textbf{Algoritmo.} Execute uma busca em profundidade; se houver uma aresta de retorno, o digrafo tem ciclo e não há ordenação. Senão, ordene os vértices por \textbf{$PS$ decrescente}. Para obter $O(n+m)$, basta colocar cada vértice na frente de uma lista no momento em que ele recebe o seu $PS$, o que dispensa a ordenação.
\begin{teo}
Num digrafo acíclico, toda aresta $vw$ satisfaz $PS(w)<PS(v)$.
\end{teo}
\begin{proof}
Não há arestas de retorno. Se $vw$ é de árvore ou de avanço, $w$ é descendente de $v$ e sai antes dele. Se é de cruzamento, $I(w)$ está inteiramente à esquerda de $I(v)$.
\end{proof}
Uma outra forma de ver a correção: o vértice de menor $PS$ é um \emph{sumidouro} (não tem arestas de saída) e pode ser o último da ordem; retirando-o, o de menor $PS$ entre os restantes é um sumidouro do digrafo que sobrou, e assim por diante.

\begin{exemplo}
No digrafo acíclico com as arestas $a{\to}c$, $a{\to}d$, $d{\to}c$, $b{\to}c$, $e{\to}b$, $e{\to}f$, $e{\to}g$, $f{\to}g$, $h{\to}f$, $i{\to}h$, $i{\to}j$, $j{\to}h$ e $j{\to}g$, uma busca com as raízes em ordem alfabética dá, por $PS$ decrescente, a ordem $i,j,h,e,f,g,b,a,d,c$. Com as raízes na ordem $i,c,b,e,d,a$, obtém-se $a,d,e,b,c,i,j,h,f,g$. As duas são válidas: a ordenação topológica não é única.
\end{exemplo}

\paragraph{Uma alternativa sem busca.} Todo digrafo acíclico tem um sumidouro --- o último vértice de um caminho de comprimento máximo --- e uma fonte. Retirar repetidamente uma fonte e colocá-la no fim da ordem (ou um sumidouro, no início) também produz uma ordenação topológica. Com contadores de grau de entrada e uma fila dos vértices cujo contador chegou a zero, isso custa $O(n+m)$; é o \emph{algoritmo de Kahn}.

\begin{cuidado}[Armadilhas em digrafos]
\begin{itemize}
  \item ``Acíclico'' refere-se a ciclos \emph{direcionados}: arestas de cruzamento e de avanço não indicam ciclo; só as de retorno indicam.
  \item O tipo de cada aresta depende da ordem da busca; as componentes fortemente conexas, não.
  \item A ordenação topológica é por $PS$ \emph{decrescente}, e não por $PE$; e ela não é única.
  \item No cálculo das componentes, uma aresta de cruzamento só atualiza $old(v)$ se o outro extremo ainda está na pilha $Q$, isto é, numa componente que ainda não foi fechada.
\end{itemize}
\end{cuidado}

\begin{chave}
Quatro tipos de aresta: árvore e avanço ($I(v)\supset I(w)$), retorno ($I(v)\subset I(w)$, fecha ciclo) e cruzamento ($I(v)$ à direita de $I(w)$). Acíclico $\iff$ sem arestas de retorno. Ordenação topológica: $PS$ decrescente, em $O(n+m)$. Componentes fortemente conexas: quando $old(v)=PE(v)$ na saída, desempilha-se uma componente. Tudo em $O(n+m)$.
\end{chave}

% ----------------------------------------------------------------
\section{Exercícios}
\excurso
\begin{exercicio}{dd-sorv}[Slides]
Prove que todo digrafo acíclico tem um sumidouro. Use isso para um algoritmo alternativo de ordenação topológica: repetidamente, encontre um sumidouro, retire-o e coloque-o no início da ordenação. É possível implementá-lo em tempo $O(n+m)$?
\end{exercicio}
\begin{exercicio}{dd-porque}[Slides]
Por que ordenar os vértices por $PS$ decrescente dá uma ordenação topológica? Aplique ao digrafo acíclico do exemplo da §\ref{sec:topo}.
\end{exercicio}
\begin{exercicio}{dd-papel}[Slides]
Uma aresta $vw$ pode ser de cruzamento numa busca e de avanço em outra? Dê um exemplo. Por que, mesmo assim, as componentes fortemente conexas encontradas são sempre as mesmas?
\end{exercicio}

\excompl
\begin{exercicio}{dd-sim}[Complementar]
No digrafo com as arestas $1\to2$, $1\to3$, $2\to3$, $2\to4$, $3\to1$, $4\to5$, $5\to6$, $6\to4$, $7\to3$ e $7\to6$, execute a busca em profundidade a partir de 1, com raízes e vizinhos em ordem crescente. Dê $PE$, $PS$, o tipo de cada aresta, os valores $old$ e as componentes fortemente conexas.
\end{exercicio}
\begin{exercicio}{dd-roupa}[Complementar]
Para se vestir: cueca antes de calça e de sapatos; calça antes de cinto e de sapatos; camisa antes de cinto e de gravata; cinto e gravata antes do paletó; meias antes dos sapatos; o relógio é independente. Dê uma ordenação topológica pela busca em profundidade, com raízes e vizinhos em ordem alfabética.
\end{exercicio}
\begin{exercicio}{dd-cond}[Complementar]
O \emph{grafo das componentes} (condensação) tem um vértice por componente fortemente conexa e uma aresta $C\to C'$ se há uma aresta de algum vértice de $C$ para algum de $C'$. Prove que ele é acíclico.
\end{exercicio}
\begin{exercicio}{dd-unica}[Complementar]
Mostre que um digrafo acíclico tem uma \emph{única} ordenação topológica se e somente se tem um caminho hamiltoniano, isto é, um caminho que passa por todos os vértices. Como testar isso em $O(n+m)$?
\end{exercicio}
\begin{exercicio}{dd-prereq}[Complementar]
Um curso tem $n$ disciplinas e $m$ pré-requisitos. Como decidir se é possível cursar todas e, se for, qual o menor número de semestres, se cada disciplina leva um semestre e não há limite de disciplinas por semestre?
\end{exercicio}
\begin{exercicio}{dd-classif}[Complementar]
No digrafo com as arestas $1\to2$, $1\to4$, $2\to3$, $3\to4$, $4\to2$, $5\to3$, $5\to6$ e $6\to5$, execute a busca em profundidade com raízes e vizinhos em ordem crescente. Classifique todas as arestas e dê $PE$ e $PS$.
\end{exercicio}
\begin{exercicio}{dd-cfcex}[Complementar]
Calcule, com o procedimento da §\ref{sec:dfsdig}, os valores $old$ e as componentes fortemente conexas do digrafo do exercício anterior, na ordem em que são encontradas.
\end{exercicio}
\begin{exercicio}{dd-kahn}[Complementar]
Execute o algoritmo de Kahn no digrafo acíclico com as arestas $a\to b$, $a\to c$, $b\to d$, $c\to d$, $d\to e$ e $f\to c$, escolhendo sempre, entre as fontes disponíveis, a de menor nome. Qual a ordem obtida?
\end{exercicio}
\begin{exercicio}{dd-ciclo}[Complementar]
Mostre como, ao encontrar uma aresta de retorno $vw$ numa busca em profundidade num digrafo, imprimir um ciclo direcionado.
\end{exercicio}
\begin{exercicio}{dd-forte}[Complementar]
Mostre que um digrafo é fortemente conexo se e somente se, a partir de um vértice qualquer $s$, todos os vértices são alcançáveis no digrafo e também no digrafo com as arestas invertidas. Conclua que é possível testar a conexidade forte com duas buscas.
\end{exercicio}
\begin{exercicio}{dd-kosaraju}[Complementar \estrela]
O algoritmo de Kosaraju faz uma busca em profundidade no digrafo, registrando a ordem de saída dos vértices; depois, faz buscas no digrafo invertido, iniciando-as em ordem decrescente de $PS$. Mostre que cada árvore da segunda busca é uma componente fortemente conexa.
\end{exercicio}
\begin{exercicio}{dd-conta}[Complementar]
Quantas ordenações topológicas tem o digrafo acíclico com as arestas $a\to c$, $b\to c$, $c\to d$ e $c\to e$? Liste-as.
\end{exercicio}
\begin{exercicio}{dd-longo}[Complementar]
No digrafo acíclico com as arestas $s\to a$ (3), $s\to b$ (2), $a\to c$ (4), $a\to t$ (2), $b\to c$ (6), $b\to t$ (3) e $c\to t$ (1), encontre o caminho de maior peso de $s$ a $t$ processando os vértices em ordem topológica.
\end{exercicio}
\begin{exercicio}{dd-semestres}[Complementar]
As disciplinas $C_1,C_2,C_3,A_1,P_1,P_2$ têm os pré-requisitos $C_1\to C_2$, $C_2\to C_3$, $C_1\to A_1$, $A_1\to C_3$ e $P_1\to P_2$. Qual o menor número de semestres para cursar todas, e em que semestre cada uma pode ser cursada?
\end{exercicio}
\begin{exercicio}{dd-minarestas}[Complementar \estrela]
Num digrafo cujo grafo das componentes tem $f$ fontes e $s$ sumidouros (e mais de uma componente), mostre que são necessárias pelo menos $\max\{f,s\}$ novas arestas para torná-lo fortemente conexo.
\end{exercicio}
\begin{exercicio}{dd-lex}[Complementar]
Como obter, entre todas as ordenações topológicas, a lexicograficamente menor (comparando os nomes dos vértices)? Qual o custo?
\end{exercicio}
\begin{exercicio}{dd-nd}[Complementar]
Por que, num grafo não direcionado, não aparecem arestas de avanço nem de cruzamento, enquanto num digrafo elas aparecem?
\end{exercicio}


\appendix
% ================================================================
\chapter{Tabelas hash}\label{cap:hash}\label{sec:hash}
% ================================================================
\begin{objetivos}
\begin{itemize}
  \item Dicionários, endereçamento direto e a ideia de uma função hash.
  \item Funções hash clássicas e o problema inevitável das colisões.
  \item Tratamento de colisões por encadeamento e por endereçamento aberto.
  \item A complexidade em função do fator de carga, e quando preferir uma tabela hash a uma árvore de busca.
\end{itemize}
\end{objetivos}

Um \textbf{dicionário} guarda pares (chave, valor) e oferece inserção, busca e remoção por chave. As árvores AVL do Capítulo~\ref{cap:avl} implementam dicionários com todas as operações em $O(\log n)$ e ainda mantêm as chaves em ordem. Quando a ordem não importa, é possível fazer melhor: as \textbf{tabelas hash} transformam a chave diretamente num endereço e realizam as operações em tempo $O(1)$ esperado.

\section{Endereçamento direto e funções hash}
Se as chaves são inteiros de um universo pequeno $U=\{0,\dots,M-1\}$, basta um vetor de tamanho $M$ indexado pela própria chave: é o \textbf{endereçamento direto}, com todas as operações em $O(1)$. Mas o universo de chaves costuma ser enorme ou infinito (todos os inteiros de 64 bits, todas as strings), e a memória é finita. Uma \textbf{função hash} $h:U\to\{0,\dots,m-1\}$ transforma cada chave num endereço de uma tabela de $m$ posições, com $m$ da ordem do número de chaves efetivamente guardadas.

\begin{exemplo}[Do endereçamento direto à função hash]\label{ex:enddireto}
Uma turma tem alunos com matrículas de 0 a 9, e estão presentes os de matrícula 2, 5 e 7. Com endereçamento direto, um vetor de 10 posições guarda cada aluno na posição da sua matrícula, e ``o aluno 5 está presente?'' é uma única consulta (Figura~\ref{fig:hashideia}). Se as matrículas fossem números de 9 dígitos, o vetor teria $10^9$ posições para guardar 3 alunos. Uma função hash, como $h(k)=k\bmod 7$, comprime o universo: as matrículas $201\,345\,002$, $731\,000\,006$ e $555\,777\,016$ vão para os endereços 5, 2 e 4 de uma tabela com 7 posições.
\begin{figura}[À esquerda, endereçamento direto: a chave é o próprio endereço. À direita, uma função hash leva chaves de um universo enorme a uma tabela pequena.\label{fig:hashideia}]
\small
\begin{tabular}{@{}c@{\hspace{8mm}}c@{}}
\vetorz[2,5,7]{-,-,2,-,-,5,-,7,-,-} &
\begin{tikzpicture}[baseline=-0.6ex, x=10mm]
\foreach \pos/\val in {0/{},1/{},2/{\tiny\ldots006},3/{},4/{\tiny\ldots016},5/{\tiny\ldots002},6/{}} {
  \node[cel, minimum width=10mm] (t\pos) at (\pos,0) {\val}; \node[idx] at (\pos,0.48) {\pos};}
\end{tikzpicture}\\[2pt]
\scriptsize vetor de 10 posições & \scriptsize tabela de 7 posições, $h(k)=k\bmod7$
\end{tabular}
\end{figura}
\end{exemplo}

\paragraph{Funções clássicas.}
\begin{itemize}
  \item \emph{Divisão}: $h(k)=k\bmod m$, com $m$ primo e longe de potências de 2, para que o resultado dependa de todos os bits da chave.
  \item \emph{Multiplicação}: $h(k)=\floor{m\,(kA\bmod1)}$, com uma constante $0<A<1$; um valor muito usado é $A\approx(\sqrt5-1)/2$. Aqui o valor de $m$ é pouco crítico, e pode ser uma potência de 2.
  \item \emph{Dobra} e \emph{meio do quadrado}: combinam partes da chave, ou extraem os dígitos centrais de $k^2$.
  \item Para \emph{strings}, a função polinomial $h(s)=\big(\sum_is_i\,b^{i}\big)\bmod m$, calculada pela regra de Horner.
  \item \emph{Hashing universal}: a função é sorteada, no início, de uma família, como $h_{a,b}(k)=((ak+b)\bmod p)\bmod m$ com $p$ primo. Isso garante bom desempenho esperado para \emph{qualquer} conjunto de chaves, sem depender de a entrada ser ``aleatória''.
\end{itemize}

\begin{exemplo}[Calculando endereços]
\textbf{Divisão}, com $m=13$: $h(100)=100\bmod13=9$ e $h(2026)=2026\bmod13=11$.
\textbf{Multiplicação}, com $m=16$ e $A=0{,}6180339887$: $100A\approx61{,}8034$, cuja parte fracionária é $0{,}8034$, e $h(100)=\floor{16\cdot0{,}8034}=12$; já $101A\approx62{,}4214$, e $h(101)=\floor{16\cdot0{,}4214}=6$. Chaves vizinhas caem em endereços distantes.
\textbf{String}, com as letras valendo $\mathtt a=1,\dots,\mathtt z=26$, $b=31$ e $m=97$: para ``gato'' (7, 1, 20, 15), a regra de Horner faz $h:=7$; $h:=(7\cdot31+1)\bmod97=24$; $h:=(24\cdot31+20)\bmod97=85$; $h:=(85\cdot31+15)\bmod97=31$. Em cada passo, $h$ fica menor que $m$, e nunca é preciso calcular $31^3$.
\textbf{Universal}, com $p=17$ e $m=6$: com $(a,b)=(3,4)$, as chaves 5 e 9 colidem, pois $((3\cdot5+4)\bmod17)\bmod6=2$ e $((3\cdot9+4)\bmod17)\bmod6=2$; com $(a,b)=(5,1)$, elas vão para os endereços 3 e 0. Um adversário que não sabe qual função foi sorteada não consegue escolher chaves que colidam sempre.
\end{exemplo}

\section{Colisões}
Uma \textbf{colisão} ocorre quando duas chaves diferentes recebem o mesmo endereço: $h(k_1)=h(k_2)$ com $k_1\ne k_2$. Colisões são inevitáveis, pelo princípio da casa dos pombos, pois $|U|>m$. E elas aparecem muito antes de a tabela ficar cheia: é o ``paradoxo do aniversário'', pelo qual, numa turma de apenas 23 pessoas, a chance de duas fazerem aniversário no mesmo dia (365 ``endereços'') já passa de 50\%. Há duas famílias de soluções.

\paragraph{Encadeamento.} No \textbf{encadeamento exterior}, cada posição da tabela aponta para uma lista com as chaves que receberam aquele endereço. Inserir é colocar a chave no início da lista, em $O(1)$; buscar e remover percorrem a lista. No \textbf{encadeamento interior}, as listas são formadas dentro da própria tabela, usando uma área de colisões ou ligando posições livres por ponteiros.

\begin{exemplo}[Encadeamento exterior]
Inserimos $12,44,13,88,23,94,11,39,20,16,5$, nessa ordem, numa tabela de $m=11$ posições com $h(k)=k\bmod11$ (Figura~\ref{fig:encadeamento}). As chaves 44, 88 e 11 vão todas para o endereço 0, e como cada uma entra no início da lista, ela fica $11\to88\to44$. Buscar 44 percorre a lista inteira (3 comparações); buscar 33, que também iria para o endereço 0, faz 3 comparações e conclui que a chave não está lá. Buscar uma chave de endereço 3 não faz nenhuma comparação: a lista está vazia.
\begin{figura}[Encadeamento exterior com $m=11$. Cada posição aponta para a lista das chaves que receberam aquele endereço; as inserções são feitas no início da lista.\label{fig:encadeamento}]
\begin{tikzpicture}[x=1cm,y=-0.62cm, no/.style={draw, thick, minimum width=8mm, minimum height=4.6mm, font=\scriptsize, inner sep=1pt, fill=white}]
\foreach \pos in {0,...,10} {\node[cel, minimum width=8mm, minimum height=5.2mm, fill=azul] (s\pos) at (0,\pos) {\pos};}
\foreach \pos/\lista in {0/{11,88,44},1/{23,12},2/{13},5/{5,16},6/{39,94},9/{20}} {
  \foreach \val [count=\k] in \lista {\node[no] (l\pos-\k) at (\k*1.3+0.3,\pos) {\val};}
  \draw[seta] (s\pos.east) -- (l\pos-1.west);
  \foreach \val [count=\k, evaluate=\k as \prox using int(\k+1)] in \lista {\ifnum\k>1 \pgfmathtruncatemacro{\ant}{\k-1}\draw[seta] (l\pos-\ant.east) -- (l\pos-\k.west); \fi}}
\end{tikzpicture}
\end{figura}
\end{exemplo}

\paragraph{Endereçamento aberto.} Todas as chaves ficam na própria tabela. Em caso de colisão, examina-se uma sequência de posições, chamada sequência de \emph{sondagem}, até encontrar uma livre:
\begin{itemize}
  \item \emph{sondagem linear}: $h(k,i)=(h(k)+i)\bmod m$, simples, mas sujeita a \textbf{agrupamento primário}: longas sequências de posições ocupadas, que tendem a crescer;
  \item \emph{sondagem quadrática}: $h(k,i)=(h(k)+c_1i+c_2i^2)\bmod m$, que reduz o agrupamento;
  \item \emph{duplo hashing}: $h(k,i)=(h_1(k)+i\,h_2(k))\bmod m$, com $h_2(k)$ primo com $m$, para que a sequência visite todas as posições.
\end{itemize}
No endereçamento aberto, uma remoção não pode simplesmente esvaziar a posição, pois isso interromperia as sequências de sondagem de outras chaves; a posição é marcada como ``removida''.

\begin{exemplo}[Sondagem linear e agrupamento]\label{ex:linear}
Inserimos $18,41,22,44,59,32,31,73$, nessa ordem, numa tabela de $m=13$ posições com $h(k)=k\bmod13$ e sondagem linear (Figura~\ref{fig:sondagem}). As chaves 18, 41, 22 e 59 vão direto para os endereços 5, 2, 9 e 7. O 44 tem $h=5$, ocupado pelo 18, e vai para 6 (2 sondagens). O 32 tem $h=6$: as posições 6 e 7 estão ocupadas, e ele vai para 8 (3 sondagens). O 31 tem $h=5$ e só encontra lugar na posição 10, depois de 6 sondagens; o 73, de $h=8$, vai para 11, com 4. As chaves formam um \textbf{agrupamento} contínuo das posições 5 a 11, e qualquer chave cujo endereço caia nesse trecho vai parar no seu fim, aumentando-o ainda mais.

Se outras funções de sondagem fossem usadas, a chave 31 seguiria outro caminho: com a sondagem quadrática $h(k,i)=(h(k)+i^2)\bmod13$, a sequência seria $5,6,9,1,\dots$; com o duplo hashing de $h_2(k)=1+(k\bmod11)=10$, seria $5,2,12,9,6,3,\dots$ --- saltos que espalham as tentativas e evitam o agrupamento.
\begin{figura}[Sondagem linear com $m=13$ depois das oito inserções do Exemplo~\ref{ex:linear}. O número abaixo de cada chave é a quantidade de sondagens que a sua inserção fez; as posições 5 a 11 formam um agrupamento.\label{fig:sondagem}]
\begin{tikzpicture}[x=8mm,y=7mm]
\fill[destaque] (4.5,-0.3) rectangle (11.5,0.3);
\foreach \val [count=\k from 0] in {{},{},41,{},{},18,44,59,32,22,31,73,{}} {\node[cel, minimum width=8mm] at (\k,0) {\val}; \node[idx] at (\k,0.55) {\k};}
\foreach \pos/\sd in {2/1,5/1,6/2,7/1,8/3,9/1,10/6,11/4} {\node[font=\tiny, text=BrickRed] at (\pos,-0.6) {\sd};}
\draw[seta, BrickRed!70, bend left=40] (5,0.85) to node[above, font=\tiny] {31: 5, 6, 7, 8, 9, 10} (10,0.85);
\end{tikzpicture}
\end{figura}
\end{exemplo}

\begin{exemplo}[Por que a remoção só marca a posição]
Na tabela da Figura~\ref{fig:sondagem}, remova o 44, que está na posição 6. Se a posição simplesmente ficasse vazia, uma busca por 32 começaria em $h(32)=6$, encontraria a posição vazia e concluiria que 32 não está na tabela --- embora ele esteja na posição 8, porque na sua inserção as posições 6 e 7 estavam ocupadas. Marcando a posição 6 como ``removida'', a busca passa por ela e continua até encontrar o 32; uma inserção futura pode reaproveitar a posição.
\end{exemplo}

\section{Complexidade}
Seja $\alpha=n/m$ o \textbf{fator de carga}. Supondo que a função hash espalha as chaves de modo uniforme:
\begin{itemize}
  \item com encadeamento, a busca custa $\Theta(1+\alpha)$ em média;
  \item com endereçamento aberto ($\alpha<1$), uma busca sem sucesso faz no máximo $\frac1{1-\alpha}$ sondagens em média.
\end{itemize}
Mantendo $\alpha=O(1)$ --- por exemplo, dobrando o tamanho da tabela e reinserindo as chaves sempre que $\alpha$ passa de um limite, o que custa $O(1)$ amortizado por inserção (Capítulo~\ref{cap:heapav}) ---, todas as operações custam $O(1)$ \textbf{esperado}. O \textbf{pior caso}, porém, é $\Theta(n)$: todas as chaves podem cair no mesmo endereço.

\begin{intuicao}
No encadeamento, as $n$ chaves se espalham por $m$ listas, e uma lista típica tem $\alpha$ chaves: a busca custa o cálculo do hash mais uma lista. No endereçamento aberto, cada sondagem encontra uma posição ocupada com probabilidade de cerca de $\alpha$; o número esperado de tentativas até achar uma livre é o de uma ``moeda'' que dá ``livre'' com probabilidade $1-\alpha$, isto é, $\frac1{1-\alpha}$. Enquanto a tabela tem folga, isso é quase 1; perto de cheia, explode (Figura~\ref{fig:carga}).
\end{intuicao}

\begin{figura}[O custo médio de uma busca sem sucesso em função do fator de carga $\alpha$: com encadeamento, $1+\alpha$ (cresce devagar e admite $\alpha>1$); com endereçamento aberto, $\frac1{1-\alpha}$ (explode quando a tabela fica cheia). Com $\alpha=0{,}75$, o endereçamento aberto faz em média até 4 sondagens.\label{fig:carga}]
\begin{tikzpicture}
\begin{axis}[width=9cm, height=5.4cm, xmin=0, xmax=1, ymin=0, ymax=12, xlabel={fator de carga $\alpha$}, ylabel={sondagens}, tick label style={font=\scriptsize}, label style={font=\small},
  legend style={font=\scriptsize, at={(0.03,0.97)}, anchor=north west}]
\addplot[very thick, MidnightBlue, domain=0:0.93, samples=80] {1/(1-x)};
\addplot[very thick, ForestGreen, domain=0:1] {1+x};
\addplot[only marks, mark=*, BrickRed] coordinates {(0.75,4)};
\legend{endereçamento aberto: $\frac1{1-\alpha}$, encadeamento: $1+\alpha$}
\end{axis}
\end{tikzpicture}
\end{figura}

\begin{exemplo}[Dimensionando a tabela]
Uma tabela com encadeamento e $m=800$ posições guarda $n=600$ chaves: $\alpha=0{,}75$, e cada lista tem, em média, menos de uma chave; uma busca sem sucesso examina, em média, 0,75 chave depois de calcular o hash. Se as mesmas 600 chaves estivessem numa tabela com endereçamento aberto de 800 posições, uma busca sem sucesso faria em média até $\frac1{1-0{,}75}=4$ sondagens; com só 650 posições ($\alpha\approx0{,}92$), seriam até cerca de 13. Por isso, no endereçamento aberto, a tabela é redimensionada bem antes de encher.
\end{exemplo}

\begin{roteiro}[Tabela hash ou árvore de busca?]
Use uma \textbf{tabela hash} quando só forem necessárias inserção, busca e remoção por chave exata e o desempenho médio importar mais que a garantia de pior caso. Use uma \textbf{árvore balanceada} (AVL ou rubro-negra) quando for preciso manter a ordem das chaves --- mínimo, máximo, sucessor, percurso em ordem, consultas por intervalo --- ou quando for exigida uma garantia de $O(\log n)$ no pior caso.
\end{roteiro}

\begin{cuidado}[Armadilhas das tabelas hash]
\begin{itemize}
  \item Esperar $O(1)$ no pior caso: o tempo constante é \emph{esperado}, sob a hipótese de uma boa função hash; com chaves escolhidas por um adversário e uma função fixa, tudo pode cair no mesmo endereço.
  \item Usar uma função que ignora parte da chave (por exemplo, só os primeiros caracteres de strings longas e parecidas): as colisões se multiplicam.
  \item Deixar a tabela de endereçamento aberto encher: o custo cresce como $\frac1{1-\alpha}$, e com $\alpha=1$ não há onde inserir.
  \item Remover no endereçamento aberto esvaziando a posição, em vez de marcá-la como ``removida''.
\end{itemize}
\end{cuidado}

\begin{chave}
Função hash: chave $\to$ endereço. Colisões são inevitáveis; tratam-se por encadeamento (listas) ou por endereçamento aberto (sondagem linear, quadrática ou duplo hashing). Com fator de carga $\alpha=O(1)$, as operações custam $O(1)$ esperado; o pior caso é $\Theta(n)$. Tabelas hash não mantêm ordem.
\end{chave}

% ----------------------------------------------------------------
\section{Exercícios}
\excurso
\begin{exercicio}{hs-pesq}[Pesquisa]
Disserte sobre: a definição e a motivação das tabelas hash; definição e exemplos de funções de dispersão clássicas; o problema das colisões; o tratamento de colisões por encadeamento exterior e interior e por endereçamento aberto; e a complexidade de pior caso da busca por uma chave.
\end{exercicio}

\excompl
\begin{exercicio}{hs-abertos}[Complementar]
Insira as chaves $10,22,31,4,15,28,17,88,59$, nessa ordem, numa tabela de tamanho $m=11$ com $h(k)=k\bmod11$, usando (a) sondagem linear; (b) sondagem quadrática com $h(k,i)=(h(k)+i+3i^2)\bmod11$; (c) duplo hashing com $h_2(k)=1+(k\bmod10)$.
\end{exercicio}
\begin{exercicio}{hs-encad}[Complementar]
Insira $5,28,19,15,20,33,12,17,10$ numa tabela com encadeamento exterior, $m=9$ e $h(k)=k\bmod9$. Qual o tamanho da maior lista?
\end{exercicio}
\begin{exercicio}{hs-primo}[Complementar]
Por que, no método da divisão, convém escolher $m$ primo e longe de potências de 2? Dê um exemplo do que dá errado com $m=2^r$.
\end{exercicio}
\begin{exercicio}{hs-remove}[Complementar]
No endereçamento aberto, por que a remoção não pode simplesmente esvaziar a posição? Como contornar o problema?
\end{exercicio}
\begin{exercicio}{hs-carga}[Complementar]
Uma tabela com encadeamento tem $m=500$ posições e guarda $n=1000$ chaves. Qual o fator de carga e o comprimento médio das listas? Quantas comparações, em média, faz uma busca sem sucesso?
\end{exercicio}
\begin{exercicio}{hs-sondagens}[Complementar]
Com endereçamento aberto e hashing uniforme, quantas sondagens, em média, faz uma busca sem sucesso quando a tabela está com fator de carga $0{,}5$? E com $0{,}9$? O que isso sugere sobre quando redimensionar a tabela?
\end{exercicio}
\begin{exercicio}{hs-mult}[Complementar]
Calcule $h(k)=\floor{m\,(kA\bmod1)}$ com $m=1000$, $A=(\sqrt5-1)/2\approx0{,}6180339887$, para $k=123456$ e $k=123457$.
\end{exercicio}
\begin{exercicio}{hs-universal}[Complementar]
O que é uma família universal de funções hash? Por que sortear a função de uma família universal protege contra entradas escolhidas para causar muitas colisões?
\end{exercicio}
\begin{exercicio}{hs-string}[Complementar]
Escreva o cálculo de $h(s)=\big(\sum_{i=0}^{L-1}s_i\,b^{L-1-i}\big)\bmod m$ pela regra de Horner, sem calcular potências grandes. Por que $b$ e $m$ não devem ter fatores em comum?
\end{exercicio}
\begin{exercicio}{hs-doissoma}[Complementar]
Dado um vetor de $n$ inteiros e um valor $x$, decida em $O(n)$ esperado se há dois elementos com soma $x$.
\end{exercicio}
\begin{exercicio}{hs-repetido}[Complementar]
Encontre, em $O(n)$ esperado, o primeiro elemento que se repete numa sequência de $n$ valores.
\end{exercicio}
\begin{exercicio}{hs-ordem}[Complementar]
Mostre que encontrar o sucessor de uma chave numa tabela hash pode custar $\Theta(n)$, enquanto numa AVL custa $O(\log n)$. Que outras operações têm a mesma limitação?
\end{exercicio}
\begin{exercicio}{hs-rehash}[Complementar]
Uma tabela dobra de tamanho e reinsere todas as chaves sempre que o fator de carga passa de $1/2$. Mostre que cada inserção custa $O(1)$ amortizado.
\end{exercicio}
\begin{exercicio}{hs-agrup}[Complementar]
Explique o agrupamento primário da sondagem linear e por que, quando uma posição fica no fim de um grupo de $k$ posições ocupadas, ela tem probabilidade maior de ser ocupada na próxima inserção.
\end{exercicio}
\begin{exercicio}{hs-duplo}[Complementar]
No duplo hashing, por que $h_2(k)$ precisa ser primo com $m$? O que acontece se $m=10$ e $h_2(k)=5$?
\end{exercicio}
\begin{exercicio}{hs-remencad}[Complementar]
Com encadeamento exterior, como remover uma chave em $O(1)$, dado um ponteiro para o seu nó? Que tipo de lista é necessário?
\end{exercicio}
\begin{exercicio}{hs-perfeito}[Complementar \estrela]
Para um conjunto \emph{fixo} de $n$ chaves, descreva a ideia do \emph{hashing perfeito} em dois níveis, que garante buscas em $O(1)$ no pior caso com memória $O(n)$ esperada.
\end{exercicio}
\begin{exercicio}{hs-distintos}[Complementar]
Conte quantos valores distintos há numa sequência de $n$ inteiros arbitrários, em $O(n)$ esperado. Como fazer em $O(n\log n)$ no pior caso?
\end{exercicio}
\begin{exercicio}{hs-tamanho}[Complementar]
Uma tabela com endereçamento aberto vai guardar $10^6$ chaves, e o fator de carga não deve passar de $0{,}75$. Qual o tamanho mínimo da tabela? Que valor de $m$ você escolheria para usar o método da divisão?
\end{exercicio}
\begin{exercicio}{hs-busca}[Complementar]
Na tabela da Figura~\ref{fig:sondagem} (sondagem linear, $m=13$), quantas sondagens faz uma busca por 31? E uma busca sem sucesso por 57? E por 26? O que esses números mostram sobre o agrupamento?
\end{exercicio}
\begin{exercicio}{hs-reinsere}[Complementar]
Na mesma tabela, remova o 44, marcando a posição como ``removida'', e depois insira o 70, que não está na tabela. Em que posição o 70 fica? Quantas sondagens faz, em seguida, uma busca pelo 32?
\end{exercicio}
\begin{exercicio}{hs-univsim}[Complementar]
Com $p=17$, $m=6$ e a família $h_{a,b}(k)=((ak+b)\bmod17)\bmod6$, calcule os endereços das chaves 3, 8 e 14 para $(a,b)=(2,7)$ e para $(a,b)=(11,3)$. Em qual dos casos há colisão?
\end{exercicio}
\begin{exercicio}{hs-horner}[Complementar]
Com as letras valendo $\mathtt a=1,\dots,\mathtt z=26$, $b=31$ e $m=97$, calcule pela regra de Horner o endereço de ``casa'' e o de ``saca''. As duas palavras têm as mesmas letras: elas colidem? O que aconteceria com uma função que só somasse os valores das letras?
\end{exercicio}

% ================================================================
\chapter{Respostas e sugestões}\label{ap:respostas}
% ================================================================
Este apêndice traz respostas, sugestões e esboços de demonstração para todos os exercícios do livro, organizados por capítulo. O número de cada resposta é um link que leva de volta ao enunciado. Quando a solução completa já aparece no texto de um capítulo, a resposta indica a seção correspondente. Recomenda-se tentar cada exercício antes de consultar a resposta.

\small
\respcap{cap:complex}
\resp{cx-slide} (a) Não: $(2n^2+1)/n\to\infty$. (b) Sim: $2n^2+1\le3n^2$ para $n\ge1$. (c) Sim: $2n^2+1\le3n^3$ para $n\ge1$. (d) Sim: $2n^2+1\ge n$ para $n\ge1$. (e) Não: $(2n^2+1)/n^3\to0$. (f) Sim, por (b) e porque $2n^2+1\ge n^2$.

\resp{cx-tec} Se o tempo disponível permite $T$ passos na máquina antiga, a nova faz $10T$. Para $n$: o tamanho é multiplicado por 10. Para $n^2$: por $\sqrt{10}\approx3{,}16$. Para $2^n$: $2^{n'}=10\cdot2^n$, logo $n'=n+\log_210\approx n+3{,}3$ --- só se \emph{soma} uma constante. Algoritmos polinomiais aproveitam o avanço tecnológico multiplicando o tamanho tratável; os exponenciais quase não ganham.

\resp{cx-k} O tamanho da entrada é o número de bits: $n$ números de até $\log_2k$ bits cada, ou seja, $\Theta(n\log k)$. Com $k=2^n$, a entrada tem $\Theta(n^2)$ bits, mas o algoritmo gasta $\Theta(2^n)$ (vetor de contagem de tamanho $k$): exponencial no tamanho da entrada. Com $k=O(n)$, o custo é $\Theta(n)$. Algoritmos assim, polinomiais no \emph{valor} dos números mas não no seu tamanho em bits, são \textbf{pseudopolinomiais}.

\resp{cx-ordem} $\log^2n\ \prec\ \sqrt n\ \prec\ 2^{\log_2n}=n\ \prec\ n\log n\ \prec\ n^{\log_23}\approx n^{1{,}585}\ \prec\ n^{100}\ \prec\ 1{,}01^n\ \prec\ n!\ \prec\ 2^{2^n}$. Toda exponencial $a^n$ com $a>1$ supera qualquer polinômio, por menor que seja $a$.

\resp{cx-vf} (a) V: $2^{n+1}=2\cdot2^n$. (b) F: $2^{2n}/2^n=2^n\to\infty$. (c) V: $\log(n^2)=2\log n$. (d) V: polilogaritmos são $o(n^\varepsilon)$ para todo $\varepsilon>0$. (e) V: $n^{\log n}=2^{\log^2n}$ e $\log^2n=o(n)$. (f) F: $f=2n$ e $g=n$ dão $2^{2n}\ne O(2^n)$.

\resp{cx-fat} Cota superior: $n!\le n^n$, logo $\log n!\le n\log n$. Cota inferior: os $n/2$ maiores fatores de $n!$ são $\ge n/2$, então $n!\ge(n/2)^{n/2}$ e $\log n!\ge\frac n2\log\frac n2=\frac n2\log n-\frac n2=\Omega(n\log n)$.

\resp{cx-incomp} $f(n)=n$ se $n$ é par e $1$ se $n$ é ímpar; $g(n)=1$ se $n$ é par e $n$ se $n$ é ímpar. Nos pares, $f/g=n\to\infty$; nos ímpares, $g/f\to\infty$. Outro exemplo: $f=n$ e $g=n^{1+\sin n}$.

\resp{cx-lacos} (a) O laço interno roda $i$ vezes para $i=1,2,4,\dots,\le n$: o total é $1+2+4+\cdots+2^{\floor{\log n}}<2n$, isto é, $\Theta(n)$ --- e não $n\log n$. (b) Para cada $i$, o laço interno roda $\floor{\log_2n}+1$ vezes: $\Theta(n\log n)$.

\resp{cx-primo} Não. A entrada é o número $N$, que ocupa $b=\Theta(\log N)$ bits. O algoritmo faz $\Theta(\sqrt N)=\Theta(2^{b/2})$ divisões: é exponencial em $b$ (pseudopolinomial). Existe algoritmo polinomial para testar primalidade, mas não é esse.

\resp{cx-def3} Cota superior: $3n^3-100n^2+7\le3n^3+7\le4n^3$ para $n\ge2$. Cota inferior: $3n^3-100n^2+7\ge2n^3$ equivale a $n^3\ge100n^2-7$, que vale para $n\ge100$. Logo $2n^3\le3n^3-100n^2+7\le4n^3$ para $n\ge100$: constantes $c_1=2$, $c_2=4$, $n_0=100$.

\resp{cx-nlogn} Se fosse, existiriam $c,n_0$ com $n^2\le c\,n\log n$, isto é, $n/\log n\le c$, para todo $n\ge n_0$. Mas $n/\log n\to\infty$. Contradição.

\resp{cx-base} $\log_an=\log_bn/\log_ba$, e $1/\log_ba$ é uma constante positiva; logo cada um é constante vezes o outro. Já $a^n/b^n=(a/b)^n$ tende a 0 ou a $\infty$ quando $a\ne b$: a base de uma exponencial importa.

\resp{cx-prop} (a) $f\le c_1g$ para $n\ge n_1$ e $g\le c_2h$ para $n\ge n_2$ dão $f\le c_1c_2h$ para $n\ge\max\{n_1,n_2\}$. (b) $\max\{f,g\}\le f+g\le2\max\{f,g\}$. (c) Se $c_1g\le f\le c_2g$, então $\frac1{c_2}f\le g\le\frac1{c_1}f$.

\resp{cx-somak} Cota superior: cada termo é $\le n^k$, logo a soma é $\le n^{k+1}$. Cota inferior: os $\ceil{n/2}$ últimos termos são $\ge(n/2)^k$, logo a soma é $\ge\frac n2(n/2)^k=n^{k+1}/2^{k+1}$.

\resp{cx-harm} Como $1/x$ é decrescente, $\int_i^{i+1}\frac{dx}x\le\frac1i\le\int_{i-1}^i\frac{dx}x$. Somando, $\ln(n+1)\le H_n\le1+\ln n$. Logo $H_n=\Theta(\log n)$.

\resp{cx-busca} Pior caso: $n$ comparações ($x$ na última posição, ou ausente). Melhor caso: 1. Caso médio: $\sum_{i=1}^n\frac1n\cdot i=\frac{n+1}2$.

\resp{cx-triplo} O corpo roda $j-i+1$ vezes para cada par $i\le j$: $\sum_{1\le i\le j\le n}(j-i+1)=\sum_{d=0}^{n-1}(n-d)(d+1)=\frac{n(n+1)(n+2)}6$ (para $n=1$ e $n=2$ dá 1 e 4). Logo $\Theta(n^3)$.

\resp{cx-loglog} Após $t$ iterações, $j=2^{2^t}$. O laço para quando $2^{2^t}>n$, isto é, quando $2^t>\log n$, ou $t>\log\log n$: $\Theta(\log\log n)$.

\resp{cx-fatexp} $\frac{2^n}{n!}=\prod_{i=1}^n\frac2i\le2\cdot1\cdot\left(\frac23\right)^{n-2}\to0$. $\frac{n!}{n^n}=\prod_{i=1}^n\frac in\le\frac1n\to0$ (cada fator é $\le1$ e o primeiro é $1/n$).

\resp{cx-binaria} No pior caso, $\floor{\log_2n}+1$ comparações (cada passo descarta pelo menos metade do intervalo). Para $n=1000$: $\floor{9{,}97}+1=10$. Para $n=10^6$: $\floor{19{,}93}+1=20$.

\resp{cx-inf} Argumento do adversário: o adversário responde ``diferente'' a toda comparação. Se o algoritmo parar depois de menos de $n$ comparações, alguma posição $p$ nunca foi comparada com $x$. Se ele responder ``não está'', o adversário coloca $x$ na posição $p$; se responder ``está'', o adversário põe em $p$ um valor diferente de $x$. Em ambos os casos, as respostas dadas continuam consistentes e o algoritmo erra. Logo são necessárias $n$ comparações, e a busca sequencial, que faz $n$, é ótima.

\resp{cx-max} (a) V: $\max\{f,g\}\le f+g\le2\max\{f,g\}$. (b) V: de $f\le cg$ vem $\log f\le\log c+\log g\le(\abs{\log c}+1)\log g$, pois $\log g\ge1$. (c) V: $f^2\le c^2g^2$. (d) F: para $f=2^n$, $f(n)/f(n/2)=2^{n/2}\to\infty$.

\resp{cx-insercao4} Em $[2,4,1,3]$: $1+2+2=5$ comparações (o 4 compara só com o 2; o 1 compara com o 4 e o 2 e chega ao início; o 3 compara com o 4 e para no 2). Em $[4,3,1,2]$: $1+2+3=6$, o máximo para $n=4$. Cada elemento faz uma comparação com cada anterior maior que ele e, se não chegar ao início, mais uma com o primeiro anterior menor. O 2, último a ser inserido, encontra $[1,3,4]$: compara com o 4 e o 3, que são maiores, e ainda com o 1, que o faz parar. São 3 comparações, as mesmas de quando se chega ao início; só os deslocamentos (2, contra 3 no vetor decrescente) diminuem.

\resp{cx-tercos} Para $n=100$, $i$ passa por $100,33,11,3,1$: 4 iterações. Para $n=1000$, por $1000,333,111,37,12,4,1$: 6 iterações. Depois de $k$ iterações, $i=\floor{n/3^k}$, e o laço para quando $\floor{n/3^k}\le1$, isto é, quando $3^k>n/2$. O número de iterações é o menor $k$ com $3^k>n/2$, ou seja, $\floor{\log_3(n/2)}+1$, que é $\Theta(\log n)$.

\respcap{cap:heap}
\resp{hp-l1} (a) Não: a posição 6 ($29$) tem pai na posição 3 ($28$) e $29>28$. (b) Sim: as 8 relações pai--filho valem ($33\ge32,28$; $32\ge31,29$; $28\ge26,25$; $31\ge30,27$).

\resp{hp-l2} Falso. Contraexemplo: $S=[10,9,1,8,7,0,0,6]$, $i=3$ ($1$), $j=4$ ($8$). Após a troca, $[10,9,8,1,7,0,0,6]$: a posição 8 ($6$) tem pai na posição 4 ($1$). O valor pequeno foi para uma posição cujos filhos podem ser maiores que ele.

\resp{hp-l3} Falso. $S=[10,9]$, $i=1$, $j=2$: após a troca, $[9,10]$ tem o filho maior que o pai.

\resp{hp-l4} $H(0)=H(1)=1$ e $H(n)=\binom{n-1}{L}H(L)H(n-1-L)$, onde $L$ é o tamanho da subárvore esquerda da raiz, determinado pelo formato da árvore completa: a raiz é obrigatoriamente o extremo, escolhem-se os $L$ valores que vão para a esquerda e cada lado é um heap. Valores: $1,1,2,3,8,20,80,210,896,3360$ para $n=1,\dots,10$. A demonstração está na §\ref{sec:heapprops}.

\resp{hp-l5} Processam-se $i=4,3,2,1$: $34\leftrightarrow50$; $41\leftrightarrow52$; $25\leftrightarrow50$ e $25\leftrightarrow48$; $18\leftrightarrow52$ e $18\leftrightarrow41$. Resultado: $[52,50,41,48,14,10,18,34,25]$.

\resp{hp-l6} É o Lema~\ref{lema:alturas}: o nó $i$ tem altura $\ge j$ se e somente se $i\,2^{j-1}\le n$; assim, o número de nós com altura exatamente $j$ é $\floor{x}-\floor{x/2}=\ceil{\floor{x}/2}\le\ceil{n/2^j}$, com $x=n/2^{j-1}$.

\resp{hp-l16} Basta tornar \emph{subir} e \emph{descer} iterativos:
\begin{pseudo}{Subir e descer sem recursão}
\begin{algorithmic}
\Procedure{subir}{$i$}
  \While{$i>1$ \textbf{e} $T[i].chave>T[\floor{i/2}].chave$} \State $T[i]\Leftrightarrow T[\floor{i/2}]$;\ $i:=\floor{i/2}$ \EndWhile
\EndProcedure
\Procedure{descer}{$i,n$}
  \While{$2i\le n$}
    \State $j:=2i$
    \If{$j<n$ \textbf{e} $T[j+1].chave>T[j].chave$} \State $j:=j+1$ \EndIf
    \If{$T[i].chave<T[j].chave$} \State $T[i]\Leftrightarrow T[j]$;\ $i:=j$ \Else \State \textbf{pare} \EndIf
  \EndWhile
\EndProcedure
\end{algorithmic}
\end{pseudo}
Inserir: $T[n+1]:=x$; $n:=n+1$; subir$(n)$. Remover: $T[1]:=T[n]$; $n:=n-1$; descer$(1,n)$. Aumentar a chave de $i$: subir$(i)$; diminuir: descer$(i,n)$. Construir: \textbf{para} $i:=\floor{n/2}$ \textbf{decrescendo até} 1, descer$(i,n)$.

\resp{hp-r1} (a) $s_i\le s_{\floor{i/2}}$ para $1<i\le n$: cada nó é maior ou igual aos filhos. (b) Pela transitividade ao longo do caminho até a raiz (Proposição~\ref{prop:ancestral}), a raiz é o máximo; basta ler $T[1]$. (c) $O(1)$, $O(\log n)$, $O(\log n)$, $O(\log n)$ e $O(n)$.

\resp{hp-r2} (a) O 98 sobe trocando com 78 (posição 3) e com 95 (posição 1): $[98,60,95,39,28,78,70,33]$. (b) Com 33 na raiz, ele desce trocando com 78 e depois com 70: $[78,60,70,39,28,66,33]$. (c) Em $i=3$ nada muda; em $i=2$, $15\leftrightarrow40$; em $i=1$, $20\leftrightarrow40$: $[40,20,30,15,10,25,5]$. (d) Em $i=3$, $15\leftrightarrow20$; em $i=2$, $2\leftrightarrow9$; em $i=1$, $7\leftrightarrow20$ e $7\leftrightarrow15$: $[20,9,15,2,1,7,13]$.

\resp{hp-r5} As operações estão invertidas. Se o novo valor ficou \emph{maior que o pai}, a violação é com o pai e o elemento precisa \textbf{subir}; caso contrário, a única violação possível é com os filhos e ele precisa \textbf{descer}. Versão correta: \textbf{se} $i>1$ \textbf{e} $H[i]>H[\floor{i/2}]$ \textbf{então} Subir$(i)$ \textbf{senão} Descer$(i,n)$. Equivalentemente, comparando com o valor antigo: se aumentou, sobe; se diminuiu, desce.

\resp{hp-tele} Use um heap de mínimo de pares $(t_{\text{livre}},\ \text{número do vendedor})$, comparados lexicograficamente, começando com $(0,1),\dots,(0,n)$. Para cada ligação de duração $d$, remova o mínimo $(t,v)$ --- o vendedor que fica livre primeiro, com desempate pelo menor número ---, conte uma ligação para $v$ e reinsira $(t+d,v)$. Cada ligação custa $O(\log n)$, e o total é $O(n+L\log n)$.

\resp{hp-dario} (a) Pai de $i$: $\floor{(i-2)/d}+1$; filhos: $d(i-1)+2,\dots,di+1$. (b) $\Theta(\log_dn)$. (c) Subir custa $O(\log_dn)$, com uma comparação por nível; descer custa $O(d\log_dn)$, pois em cada nível acha o maior de $d$ filhos. (d) Aumentos de prioridade usam \emph{subir}, que fica mais barato com $d$ maior, enquanto as remoções (descer) ficam mais caras; quando há muito mais aumentos, compensa aumentar $d$.

\resp{hp-kmerge} Mantenha um heap de mínimo com o primeiro elemento de cada lista (e a lista de onde veio). Repita $n$ vezes: remova o mínimo, coloque-o na saída e insira o próximo elemento da mesma lista. O heap tem sempre no máximo $k$ elementos, e o custo total é $O(n\log k)$, mais $O(k)$ da construção.

\resp{hp-kmaiores} Use um heap de \textbf{mínimo} com até $k$ elementos. Para cada número $x$: se o heap tem menos de $k$ elementos, insira $x$; senão, se $x$ é maior que o mínimo do heap, substitua o mínimo por $x$ e faça-o descer. Ao final, o heap contém os $k$ maiores. Tempo $O(n\log k)$ e memória $O(k)$.

\resp{hp-verif} \textbf{Para} $i:=2$ \textbf{até} $n$: \textbf{se} $s_i>s_{\floor{i/2}}$, \textbf{retorna} falso. Ao final, \textbf{retorna} verdadeiro. São $n-1$ comparações, uma por aresta da árvore, e isso basta porque a propriedade de heap é exatamente essa condição local.

\resp{hp-estavel} Entrada $[1,\ 2_a,\ 2_b]$. O \emph{arranjar} troca 1 com $2_a$ (o primeiro filho, já que $2_b$ não é maior): $[2_a,1,2_b]$. O primeiro passo do heapsort troca a raiz com a última posição: $[2_b,1,2_a]$; descer não muda nada. O segundo passo dá $[1,2_b,2_a]$. Na saída, $2_b$ aparece antes de $2_a$: a ordem relativa dos iguais foi invertida.

\resp{hp-crescente} Ao ser inserido, $i$ é o maior elemento até então e sobe até a raiz, com $\floor{\log_2i}$ trocas. O total, $\sum_{i=2}^n\floor{\log_2i}$, é $\Theta(n\log n)$, pois metade dos termos é $\ge\log_2(n/2)-1$. O \emph{arranjar} sobre o vetor $[1,2,\dots,n]$ continua $O(n)$, pela análise que soma as alturas (Lema~\ref{lema:alturas}).

\resp{hp-ins} $[15]$; $[40,15]$; $[40,15,25]$; $[50,40,25,15]$ (50 sobe duas vezes); $[50,40,25,15,10]$; $[50,40,45,15,10,25]$ (45 troca com 25).

\resp{hp-rem2} Primeira remoção: 20 vai para a raiz e desce trocando com 80 e depois com 50, resultando em $[80,70,50,30,\allowbreak60,20,40,10]$. Segunda: 10 vai para a raiz e desce trocando com 70 e depois com 60, resultando em $[70,60,50,30,\allowbreak10,20,40]$.

\resp{hp-sort5} Após o \emph{arranjar}: $[10,5,3,4,1]$. Passos do laço, para $m=5,4,3,2$:\\ $[5,4,3,1\,|\,10]$,\ \ $[4,1,3\,|\,5,10]$,\ \ $[3,1\,|\,4,5,10]$,\ \ $[1\,|\,3,4,5,10]$.

\resp{hp-remi} Copie $T[n]$ para a posição $i$ e faça $n:=n-1$. O valor que chegou a $i$ pode ser maior que o pai de $i$ (ele vinha de outro ramo da árvore) ou menor que algum filho de $i$; então, \textbf{se} $i>1$ \textbf{e} $T[i]>T[\floor{i/2}]$, suba; \textbf{senão}, desça. Apenas uma das duas violações pode ocorrer, e o custo é $O(\log n)$.

\resp{hp-mediana} Guarde a metade menor dos elementos num heap de \textbf{máximo} $A$ e a metade maior num heap de \textbf{mínimo} $B$, com $|A|=|B|$ ou $|A|=|B|+1$. Para inserir $x$: se $x\le$ topo de $A$, insira em $A$, senão em $B$; depois, se os tamanhos ficarem desequilibrados, mova o topo de um heap para o outro. A mediana (inferior) é o topo de $A$, lido em $O(1)$; cada inserção custa $O(\log n)$.

\resp{hp-comp} Cada passo do \emph{descer} faz no máximo duas comparações (uma entre os dois filhos e outra entre o maior filho e o pai). Um nó de altura $j$ faz no máximo $j-1$ passos. Pelo Lema~\ref{lema:alturas}, o total é no máximo $2\sum_{j\ge2}(j-1)\ceil{n/2^j}\le2n+O(\log^2n)$.

\respcap{cap:heapav}
\resp{hb-l7} Faça $n$ inserções sucessivas. Inserir é somar 1 em binário, e cada \emph{link} corresponde a um ``vai um''. Contando de 0 até $n$, o ``vai um'' do bit $k$ para o bit $k+1$ ocorre no máximo $\floor{n/2^{k+1}}$ vezes, de modo que o número total de links é $\sum_{k\ge1}\floor{n/2^k}<n$, e o custo total é $O(n)$. Alternativa equivalente: agrupe os elementos em pares formando árvores $B_1$, depois os pares de $B_1$ formando $B_2$, e assim por diante, como numa soma binária de $n$ uns.

\resp{hb-l8} A inserção é o pior caso quando o novo $B_0$ provoca ``vai um'' em todas as posições: o heap contém $B_0,B_1,\dots,B_{k-1}$, isto é, $n=2^k-1$, com todos os bits iguais a 1. As alturas são $0,1,\dots,k-1$ em arestas (ou $1,\dots,k$ em níveis). A inserção faz $k$ links e produz uma única árvore $B_k$.

\resp{hb-l10} A remoção percorre as $k$ raízes, transforma os $h$ filhos da raiz removida no heap $H''=\{B_0,\dots,B_{h-1}\}$ e funde $H''$ com o restante $H'$. O número de links é máximo quando $H'$ também contém $B_0,\dots,B_{h-1}$ e $h$ é o maior possível: $H=\{B_0,B_1,\dots,B_{k-1}\}$ com o máximo na raiz de $B_{k-1}$, isto é, $h=k-1$ (em arestas).

\resp{hb-l11} Falso. $H_1$ com $3=11_2$ nós ($k_1=2$) e $H_2$ com $5=101_2$ nós ($k_2=2$). A fusão tem $8=1000_2$ nós, numa única árvore: $k=1<2$.

\resp{hb-l12} Verdadeiro. A união das listas de raízes tem $k_1+k_2$ árvores e cada link transforma duas árvores em uma; logo $k=k_1+k_2-\#\text{links}\le k_1+k_2$.

\resp{hb-l13} Suponha árvores de posto $k$ formadas por subordinações de árvores de mesmo posto, com cortes controlados por marcas em nós que não são raízes, e meça alturas em arestas. \textbf{Altura mínima}: $\ceil{k/2}$. Pelo lema $\operatorname{grau}(y_i)\ge i-2$, a menor altura possível satisfaz $A(k)=1+A(k-2)$, com $A(0)=0$ e $A(1)=1$. \textbf{Altura máxima}: $k$, atingida pela própria $B_k$; por indução, subordinar duas árvores de posto $h$ e altura $\le h$ dá posto $h+1$ e altura $\le h+1$, e cortes não aumentam a altura. Se as raízes puderem perder filhos livremente, só o \emph{grau} continua sendo $O(\log n)$; a altura pode crescer.

\resp{hb-l14} Ver §\ref{sec:binomial} e §\ref{sec:fib}. Heap binomial: a seleção percorre as raízes (ou mantém um ponteiro); a inserção é a fusão com $\{B_0\}$; a fusão intercala as listas de raízes por ordem e faz links como numa soma binária; a remoção do máximo retira a árvore, transforma os filhos num heap e funde; o aumento sobe trocando com o pai; a diminuição remove o nó (aumentando-o para $+\infty$ e removendo o máximo) e o reinsere; a construção faz $n$ inserções. Heap de Fibonacci: inserção e fusão concatenam listas circulares; a remoção do máximo promove os filhos e consolida por grau; o aumento corta o nó e faz cortes em cascata; a diminuição remove e reinsere.

\resp{hb-r6} (a) $B_0$ é um nó; $B_k$ é obtida ligando duas $B_{k-1}$. Tem $2^k$ nós e a raiz tem grau $k$, que é a ordem da árvore. (b) Como uma lista de raízes de árvores com a condição de heap e \textbf{ordens distintas}; $B_k$ aparece se e somente se o bit $k$ de $n$ é 1. (c) A raiz 10 fica por cima e a raiz 15 vira sua filha; o resultado é uma $B_3$, com 8 nós e raiz de grau 3. (d) O extremo está numa das no máximo $\floor{\log n}+1$ raízes: $O(\log n)$, ou $O(1)$ com ponteiro. (e) $O(\log n)$: a união percorre só as listas de raízes, que têm $O(\log n)$ árvores, fazendo links como numa soma binária; os nós internos não são visitados.

\resp{hb-r9} (a) O heap de Fibonacci é preguiçoso: permite várias árvores de mesmo grau na lista de raízes, consolidando só na remoção, e usa listas circulares duplas e marcas nos nós. (b) Inserção, união, seleção e redução de chave (\emph{decrease-key}) custam $O(1)$ amortizado; \emph{Extract-Min} custa $O(\log n)$ amortizado. (c) A nova chave pode ficar menor que a do pai. Corta-se $x$ com sua subárvore, e ele vai para a lista de raízes, desmarcado. (d) Se o pai já tinha perdido um filho (estava marcado), ele também é cortado, e assim por diante. Isso garante que cada nó perde no máximo um filho antes de ser cortado, o que mantém o grau $O(\log n)$. Com o potencial $\Phi=t+2m$, cada corte em cascata desmarca um nó: o potencial cai 2 e sobe 1, e o corte paga a si mesmo. Por isso o custo amortizado é $O(1)$.

\resp{hb-m1} São $21=10101_2$ e $3=11_2$ elementos. Os links, na ordem, são: $B_0+B_0\to B_1$; esse $B_1$ com o $B_1$ de $H_2$ forma $B_2$; essa $B_2$ com a $B_2$ de $H_1$ forma $B_3$. O resultado é $\{B_3,B_4\}$, isto é, $24=11000_2$.

\resp{hb-m2} (1) $X$ é cortado e vai para a lista de raízes, desmarcado. (2) $Y$, que estava marcado, perdeu o segundo filho: é cortado, desmarcado e vai para as raízes. (3) O mesmo acontece com $Z$. (4) Se o pai $W$ de $Z$ é raiz, nada muda; se não é raiz e está desmarcado, é marcado e a cascata termina; se estava marcado, a cascata continua. Vão para a lista de raízes $X$, $Y$ e $Z$ (e possivelmente outros); ao final, $X,Y,Z$ estão desmarcados e $W$ marcado, se não for raiz.

\resp{hb-m3} Sem as marcas, um nó poderia perder quase todos os filhos sem que nada acontecesse: cortando todos os netos de uma raiz, obtém-se uma árvore de posto $k$ com apenas $k+1$ nós. O tamanho deixaria de crescer exponencialmente com o posto (como $F_{k+2}$), o grau máximo deixaria de ser $O(\log n)$ e a consolidação, que depende do número de graus distintos, ficaria cara.

\resp{hb-t1} \textbf{Heap binomial}. A fusão custa $O(\log n)$ \emph{no pior caso de cada execução} e cresce só logaritmicamente com o número de aeronaves: é previsível. No heap binário, a fusão exige reconstruir o vetor, $\Theta(n)$ por fusão; milhares de fusões com $10^5$ elementos estouram as janelas de tempo. No heap de Fibonacci, a fusão é $O(1)$ apenas \emph{amortizado}: o trabalho adiado aparece de uma vez numa remoção, cuja consolidação pode custar $\Theta(n)$, o que é imprevisível; além disso, a grande vantagem dessa estrutura, o aumento de prioridade em $O(1)$ amortizado, não seria usada, já que as alterações são raras.

\resp{hb-t2} \textbf{Heap binário em vetor}. Os dados ocupam $2\cdot10^6\times16=32\,000\,000$ bytes, cerca de 30,5 MB, abaixo do limite de $41\,943\,040$ bytes. Um heap binomial precisaria de pelo menos 3 ponteiros por nó (pai, filho, irmão), isto é, 24 bytes, mais o grau: pelo menos 40 bytes por nó e $80\cdot10^6$ bytes no total, o que não cabe; o de Fibonacci, com 4 ponteiros, grau e marca, ainda menos. Como não há fusões e as alterações são poucas, inserção e remoção em $O(\log n)$ bastam; o heap binário é o único que cabe e o mais simples.

\resp{hb-t3} \textbf{Heap de Fibonacci}. O aumento de prioridade custa $O(1)$ amortizado, sem retirar e reinserir a tarefa. Uma estimativa do custo total é $5\cdot10^6+8\cdot10^7+5\cdot10^6\log_2(5\cdot10^6)\approx2\cdot10^8$ passos. Nos heaps binário e binomial, cada um dos $8\cdot10^7$ aumentos custa $O(\log n)\approx22$ passos, cerca de $2\cdot10^9$ no total. Como importa o custo total da execução, e não a previsibilidade de cada operação, a garantia amortizada é adequada; as poucas uniões não pesam.

\resp{hb-vetor} Com $k$ elementos e capacidade $c$, vale $c\ge k\ge c/2$ depois da primeira duplicação, logo $\Phi=2k-c\ge0$. Uma inserção que não duplica tem custo 1 e $\Delta\Phi=2$: custo amortizado 3. Uma inserção que duplica (com $k=c$) custa $c+1$, pois copia $c$ elementos e insere um; o potencial passa de $2c-c=c$ para $2(c+1)-2c=2$, ou seja, $\Delta\Phi=2-c$, e o custo amortizado é $(c+1)+(2-c)=3$. Logo $n$ inserções custam no máximo $3n$.

\resp{hb-contador} Comece em $2^{k-1}-1=0111\ldots1$. Um incremento leva a $1000\ldots0$, trocando $k$ bits, e um decremento volta ao estado anterior, trocando $k$ bits. Alternando as duas operações, cada uma custa $k$, e o total é $\Theta(mk)$. Com apenas incrementos, o potencial ``número de bits iguais a 1'' pagava os ``vai um'' porque os uns, uma vez zerados, só voltam um a um; com decrementos, o mesmo trabalho pode ser repetido indefinidamente sem que nada o pague.

\resp{hb-fila} \emph{Enfileirar}: empilhe na pilha de entrada. \emph{Desenfileirar}: se a pilha de saída estiver vazia, transfira todos os elementos da entrada para ela, invertendo a ordem; depois desempilhe da saída. Cada elemento é empilhado e desempilhado no máximo uma vez em cada pilha, de modo que $m$ operações custam $O(m)$. Com o potencial $\Phi=2\cdot|\text{entrada}|$, enfileirar custa 3 amortizado e desenfileirar custa 1.

\resp{hb-forma} $100=1100100_2$: $\{B_2,B_5,B_6\}$. $101=1100101_2$: $\{B_0,B_2,B_5,B_6\}$, numa inserção sem links. De 100 a $104=1101000_2$: $\{B_3,B_5,B_6\}$. As quatro inserções ($100\to101\to102\to103\to104$) fazem $0,1,0,3$ links, respectivamente: o número de uns no final de $100$, $101$, $102$ e $103$.

\resp{hb-fibbin} Por indução. As árvores nascem com um único nó ($B_0$). Sem cortes, uma árvore só muda por \emph{link}, que na consolidação liga duas raízes de mesmo grau $k$; se ambas são árvores binomiais $B_k$ (não ordenadas), o resultado é $B_{k+1}$. Na remoção do máximo, os filhos da raiz de uma $B_k$ são raízes de $B_0,\dots,B_{k-1}$, também binomiais. Logo todas as árvores são binomiais; a diferença em relação ao heap binomial é que, antes da consolidação, pode haver várias árvores de mesma ordem.

\resp{hb-fibcaro} Insira $n$ elementos --- cada inserção só acrescenta uma raiz, e ficam $n$ raízes --- e depois remova o máximo: a consolidação percorre as $n-1$ raízes restantes e faz $\Theta(n)$ links, com custo real $\Theta(n)$. Não há contradição, porque o custo amortizado $O(\log n)$ vale para a \emph{sequência}: as $n$ inserções pagaram, cada uma com um crédito (o aumento de potencial $\Delta t=+1$), pelo trabalho da consolidação.

\resp{hb-binsim} Depois do 4: $\{4\}$. Depois do 9: as duas $B_0$ se juntam, com o 9 por cima: $\{9\to4\}$ (1 link). Depois do 2: $\{2,\ 9\to4\}$ (nenhum). Depois do 7: o 7 se junta ao 2, formando $7\to2$, que se junta a $9\to4$, com o 9 por cima: uma $B_2$ de raiz 9, cujos filhos são $7\to2$ e 4 (2 links). Depois do 5: $\{5,\ B_2\}$ (nenhum). Os links por inserção são $0,1,0,2,0$, exatamente os ``vai um'' das somas $0+1$, $1+1$, $10_2+1$, $11_2+1$ e $100_2+1$; o total é 3, e o heap final tem $5=101_2$ elementos.

\resp{hb-fibsim} As sete inserções criam sete raízes, e $\Phi=7$. O 9 sai, e a consolidação percorre $6,2,4,8,1,7$: o 2 colide com o 6, formando $6\to2$, de grau 1; o 4 fica em $A[0]$; o 8 colide com o 4, formando $8\to4$, que colide com $6\to2$: o 8 fica por cima, numa árvore de grau 2, com os filhos $6\to2$ e 4; o 1 vai para $A[0]$, e o 7 colide com ele, formando $7\to1$, de grau 1. Restam duas raízes, 8 (grau 2) e 7 (grau 1), depois de 4 links; o novo máximo é 8, e $\Phi$ cai de 7 para 2.

\respcap{cap:uf}
\resp{uf-r1} (a) Manter uma partição de $n$ elementos em conjuntos disjuntos, respondendo se dois elementos estão no mesmo conjunto e unindo conjuntos. (b) \MakeSet$(x)$ cria $\{x\}$ ($pai[x]:=x$); \Find$(x)$ sobe pelos pais até a raiz, que é o representante; \Union$(x,y)$ encontra as duas raízes e faz uma apontar para a outra. (c) O \textbf{pai} de cada elemento na floresta; a raiz guarda também o rank (ou o tamanho). (d) A árvore de menor rank é ligada sob a de maior; o rank só aumenta quando os dois são iguais, e então a árvore pelo menos dobra de tamanho. Por isso uma árvore de rank $h$ tem pelo menos $2^h$ nós, e a altura é no máximo $\log n$ (Teorema~\ref{teo:ufaltura}).

\resp{uf-r2} Ver §\ref{sec:ufexemplos}. (a) Representantes finais $A$ (de $\{A,B,C,D\}$, com rank 2) e $E$ (de $\{E,F\}$): 2 conjuntos. (b) Devolve $A$ e deixa $B,C,D,E$ apontando diretamente para $A$. (c) A raiz é $A$, com rank 2: só no empate a altura cresce. (d) Só com rank, \Find{} custa $\Theta(\log n)$; só com compressão, o total é $O(m\log n)$, com \Find{} individual possivelmente $\Theta(n)$; juntas, $O(m\lgs n)$ (ou $O(m\,\alpha)$).

\resp{uf-l1} A prova é o Teorema~\ref{teo:ufaltura}: por indução, uma árvore de rank (igual à altura, sem compressão) $h$ tem pelo menos $2^h$ nós, pois a altura só aumenta ao unir duas árvores de mesmo rank $h$, cada uma com pelo menos $2^h$ nós. Logo $2^h\le n$. Para atingir a cota com $n=2^k$, una os elementos aos pares (altura 1), depois as árvores resultantes aos pares (altura 2), e assim por diante. Após $k$ rodadas há uma única árvore de altura $k=\log_2n$, que é uma árvore binomial $B_k$.

\resp{uf-l2} Faça \Union$(1,2)$, \Union$(1,3)$, \dots, \Union$(1,n)$, sempre pendurando a raiz da árvore grande sob o novo elemento unitário ($pai[\text{raiz}]:=i$). Forma-se uma cadeia de altura $n-1$ sem nenhum \Find{} caro no caminho; um \Find{} no fundo custa $\Theta(n)$ e achata tudo. Com compressão e união arbitrária, $m$ operações custam $O(m\log n)$. A ideia da prova é classificar cada passo $x\to pai[x]$ da compressão como ``leve'' quando a subárvore de $x$ tem menos da metade dos nós da subárvore do pai: cada nó só passa por $O(\log n)$ passos leves, e cada \Find{} tem $O(\log n)$ passos pesados.

\resp{uf-t4} (a) Union-Find. Como as consultas dominam ($10^8$ contra $10^5$ uniões) e precisam ser rápidas, a \textbf{lista encadeada com ponteiro direto ao representante} responde \Find{} em $O(1)$ \emph{no pior caso}, deixando o custo para as raras uniões: $O(\log n)$ amortizado com ``menor na maior'', no máximo $n\log_2n\approx10^7\cdot23$ atualizações de ponteiro no total. A floresta com rank e compressão também é defensável ($O(\alpha)$ amortizado por operação e menos memória), desde que se diga que a garantia é amortizada; floresta sem balanceamento não serve. (b) \textbf{Anexar o menor ao maior}. Na lista, cada elemento muda de representante no máximo $\log_2n$ vezes, pois o tamanho do seu grupo pelo menos dobra a cada mudança; na árvore, a altura fica no máximo $\log_2n$. Sem isso, listas e árvores degeneram e cada operação pode custar $\Theta(n)$.

\resp{uf-doc} Ver os cenários da §\ref{sec:ufescolha}. (a) Não há como desfazer uniões, e a compressão destrói a forma das árvores. Resolva \emph{offline}, de trás para frente, transformando cada destruição numa união a partir do grafo final; guarde o número de componentes a cada passo e apresente as respostas na ordem original. (b) Floresta sem compressão, só com união por rank ou tamanho (\Find{} em $O(\log n)$), empilhando cada ponteiro e rank alterado por \Union; desfazer é desempilhar e restaurar, em $O(1)$. (c) A floresta só tem ponteiros para cima, e listar membros custaria $O(n)$; use listas encadeadas com ponteiro ao representante e união ``menor na maior'': \emph{MesmoGrupo} em $O(1)$, \emph{ListarMembros} em $O(k)$ e \emph{Juntar} em $O(\log n)$ amortizado.

\resp{uf-cf} Cada bloco de células contíguas é um conjunto, representado por uma \textbf{lista encadeada} dos gatinhos na ordem das células, com cabeça e cauda guardadas no representante. No dia $i$, encontre os representantes de $x_i$ e $y_i$ e concatene uma lista na outra (cauda de uma $\to$ cabeça da outra), unindo também os conjuntos. Como os blocos de $x_i$ e $y_i$ eram vizinhos, colocá-los lado a lado é consistente com alguma ordem inicial. Ao final, percorra a única lista restante. Custo $O(n\,\alpha(n))$.

\resp{uf-sim} \Union$(1,2)$ torna 1 raiz, com rank 1; o mesmo ocorre com 3 (unido a 4), 5 (a 6) e 7 (a 8). \Union$(1,3)$: empate, $pai[3]=1$ e $\rank(1)=2$. \Union$(5,7)$: $pai[7]=5$ e $\rank(5)=2$. \Union$(1,5)$: empate, $pai[5]=1$ e $\rank(1)=3$. Antes do \Find, o vetor $pai$ (posições $1..8$) é $1,1,1,3,1,5,5,7$. \Find$(8)$ percorre $8\to7\to5\to1$, devolve 1 e comprime: $pai=1,1,1,3,1,5,1,1$. Os ranks não mudam: $\rank(1)=3$, $\rank(5)=2$, $\rank(3)=\rank(7)=1$ e os demais são 0.

\resp{uf-ciclo} Faça \MakeSet{} para todo vértice; para cada aresta $uv$, se \Find$(u)=$ \Find$(v)$, a aresta fecha um ciclo; senão, faça \Union$(u,v)$. Custo $O(n+m\,\alpha(n))$. Uma busca em profundidade (Capítulo~\ref{cap:dfs}) também resolve o problema, em $O(n+m)$.

\resp{uf-rankpai} Quando $x$ deixa de ser raiz, ele é ligado a uma raiz $r$ com $\rank(r)>\rank(x)$, ou com $\rank(r)=\rank(x)$, e nesse caso $\rank(r)$ é incrementado; em ambos os casos, $\rank(x)<\rank(r)$. Depois disso, o rank de $x$ nunca muda (só raízes mudam de rank), e o rank de $pai[x]$ nunca diminui. A compressão troca $pai[x]$ por um ancestral, que, pela mesma desigualdade aplicada ao longo do caminho, tem rank ainda maior. Logo a desigualdade se mantém.

\resp{uf-max} Guarde na raiz o maior elemento e o tamanho do conjunto. Na \Union{} de raízes $r$ e $s$, com $s$ indo para baixo de $r$, faça $max[r]:=\max\{max[r],max[s]\}$ e $tam[r]:=tam[r]+tam[s]$. As consultas fazem \Find$(x)$ e leem o valor na raiz: $O(\alpha(n))$ amortizado.

\resp{uf-tamanho} Indução: uma árvore de altura $h$ tem pelo menos $2^h$ nós. Ao pendurar $T_v$ (a menor) sob a raiz de $T_w$, com $|T_v|\le|T_w|$, a nova altura é $\max\{h_w,h_v+1\}$. Se é $h_w$, vale pela hipótese. Se é $h_v+1$, o novo tamanho é $|T_v|+|T_w|\ge2|T_v|\ge2\cdot2^{h_v}=2^{h_v+1}$. Logo $2^h\le n$.

\resp{uf-ingenua} Faça \Union$(1,2)$, \Union$(2,3)$, \dots, \Union$(n-1,n)$: a raiz do conjunto de $i$ é sempre $i$ (o último elemento a virar raiz é o de $y$), e cada união faz a raiz atual apontar para o novo elemento, formando a cadeia $1\to2\to\cdots\to n$, de altura $n-1$. Sem compressão, $n$ chamadas \Find$(1)$ custam $\Theta(n)$ cada: $\Theta(n^2)$ no total.

\resp{uf-listas} Na primeira rodada, 4 uniões de listas de tamanho 1: 4 atualizações. Na segunda, 2 uniões de listas de tamanho 2, atualizando 2 ponteiros cada: 4. Na terceira, 1 união de listas de tamanho 4: 4. Total: 12. Para $n=2^k$, são $k$ rodadas com $n/2$ atualizações cada: $\frac n2\log_2n$.

\resp{uf-listasprova} O ponteiro de um elemento $x$ só é atualizado quando a lista de $x$ é a menor das duas (ou igual); a nova lista tem pelo menos o dobro do tamanho da lista anterior de $x$. Como nenhuma lista tem mais de $n$ elementos, o ponteiro de $x$ é atualizado no máximo $\log_2n$ vezes. Somando sobre os $n$ elementos, no máximo $n\log_2n$ atualizações.

\resp{uf-semponderada} Una sempre o conjunto grande ao elemento seguinte, anexando a lista grande à unitária: \Union$(2,1)$, \Union$(3,\{1,2\})$, \dots, com a lista do novo elemento como base. A $i$-ésima união atualiza $i$ ponteiros, e o total é $1+2+\cdots+(n-1)=\Theta(n^2)$.

\resp{uf-logstar} Com $F=1,2,4,16,65536,2^{65536}$: $\lgs2=1$, $\lgs16=3$, $\lgs65536=4$ e $\lgs2^{65536}=5$. Como $65536<10^{80}\le2^{65536}$, $\lgs10^{80}=5$.

\resp{uf-comp} Uniões efetivas: $12$, $34$, $23$, $56$, $78$, $89$ (seis). A aresta $65$ liga vértices que já estão juntos (repete $56$) e a aresta $97$ fecha o ciclo $7,8,9$. Componentes finais: $\{1,2,3,4\}$, $\{5,6\}$, $\{7,8,9\}$ e $\{10\}$, ou seja, $10-6=4$ componentes.

\resp{uf-percol} Crie um elemento por casa e dois elementos virtuais, $\top$ e $\bot$. Ao abrir uma casa, una-a a cada vizinha já aberta; se ela está na primeira linha, una-a a $\top$; se está na última, a $\bot$. Há caminho se e somente se \Find$(\top)=$ \Find$(\bot)$. Cada abertura faz no máximo 6 operações; com rank e compressão, $n^2$ aberturas custam $O(n^2\,\alpha(n^2))$.

\resp{uf-equiv} Cada declaração $x\equiv y$ é uma \Union$(x,y)$, e cada consulta compara \Find$(x)$ com \Find$(y)$. No exemplo, $a,b,c,d$ ficam equivalentes ($a\equiv b$, $b\equiv c$, $c\equiv d$), e $e,f$ formam outra classe.

\resp{uf-kruskal} Para cada uma das $m$ arestas, duas operações \Find{} (uma por extremo): $2m$. Uniões efetivas: exatamente $n-1$, uma por aresta da árvore. Com rank e compressão, essas operações custam $O(m\,\alpha(n))$, muito menos que a ordenação das arestas, $O(m\log m)=O(m\log n)$, que domina o algoritmo.

\resp{uf-rankcomp} Atualizar os ranks exigiria recalcular alturas de subárvores depois de cada compressão, o que custaria caro. Com compressão, o rank deixa de ser a altura e passa a ser uma \emph{cota superior} para ela. Ele continua servindo para decidir as uniões porque as propriedades usadas na análise --- uma árvore de rank $k$ tem pelo menos $2^k$ nós e o rank cresce ao subir na árvore --- dependem apenas de como os ranks evoluem nas uniões, e não da altura real.

\respcap{cap:avl}
\resp{avl-r1} (a) Uma ABB é balanceada se sua altura é $O(\log n)$ e o mesmo vale para cada subárvore (com $m$ nós, altura $O(\log m)$). (b) AVL: em todo nó, as alturas das subárvores esquerda e direita diferem de no máximo 1. (c) Toda AVL é balanceada, pois $h<1{,}44\log_2(n+2)$ (Teorema~\ref{teo:avlaltura}); nem toda árvore balanceada é AVL. (d) $\bal(v)=h_D(v)-h_E(v)$.

\resp{avl-r2} (a) $\bal(p)=1\to0$: a altura não muda e o processo para. $0\to-1$: a altura aumenta e o processo continua com o pai. $-1\to-2$: rotação, e o processo para. (b) Com sinais iguais ($-2$ e $-1$), rotação \textbf{simples à direita} em $p$. Com sinais opostos ($-2$ e $+1$), rotação \textbf{dupla à direita}: à esquerda em $u$ e depois à direita em $p$. No primeiro caso o excesso está ``em linha'' (esquerda-esquerda) e uma rotação resolve; no segundo, está em zigue-zague (esquerda-direita), e é preciso primeiro alinhar.

\resp{avl-r4} (a) $h$ indica se a altura da subárvore aumentou. (b) Se a altura não mudou, nenhum ancestral muda de balanço; só quando $h$ é verdadeiro o pai precisa se atualizar. (c) A rotação muda a raiz da subárvore; como $pt$ é passado por referência, atribuir a nova raiz a $pt$ religa a subárvore ao pai (ou atualiza a raiz da árvore). (d) Os testes precisam ser exclusivos: com ``se'' em sequência, um nó com balanço 1 vira 0 e, logo em seguida, o teste ``balanço $=0$'' dispara e o deixa com $-1$, com $h$ verdadeiro, o que está errado.

\resp{avl-r5} Ver §\ref{sec:avlexemplos}. (a) Rotação simples à direita em 40 ao inserir 10, e dupla à direita em 40 ao inserir 35. Árvore final: $20^{+1}(10,\,35(30,40))$. (b) Rotação simples à esquerda em 10 (ao inserir 30) e simples à esquerda em \textbf{30} (ao inserir 50; o primeiro desregulado é 30, não 20). Árvore final: $20^{+1}(10,\,40(30,50))$. (c) Inserir 49 desregula 20, cujo filho direito 40 tem balanço $+1$: rotação simples à esquerda em 20. Árvore final: $40(20(10,30),\,50^{-1}(49,\cdot))$.

\resp{avl-t5} Ver o Exemplo~\ref{ex:teste1q5}. Inserir 5 desregula 30 (caso esquerda-esquerda, rotação simples à direita em 30); 15 não causa rotação; 35 desregula 50 (caso esquerda-direita, rotação dupla à direita: à esquerda em 20, à direita em 50, e 30 vira a raiz); 45 não causa rotação. Árvore final: $30(20(10(5,15),25),\allowbreak\,50(40(35,45),70(60,80)))$. (c) Numa ABB, inserções em ordem crescente, decrescente ou em zigue-zague geram altura $\Theta(n)$ e operações $\Theta(n)$; a AVL garante $O(\log n)$ qualquer que seja a ordem de chegada, com $O(1)$ rotações por inserção.

\resp{avl-l3} Use uma árvore de Fibonacci (uma AVL mínima) em que todos os nós internos têm balanço $-1$ e remova a folha mais à direita. Cada ancestral perde altura do lado que já era mais baixo, fica com balanço $-2$ e tem o filho esquerdo com balanço $-1$; a rotação simples diminui a altura da subárvore e propaga o problema para cima. Por exemplo, na árvore de Fibonacci de altura 7, com 33 nós (chaves $1,\dots,33$) e raiz 21, remover a folha 33 (cujos ancestrais são 32, 29 e 21) causa rotações em 32, 29 e 21.

\resp{avl-l4} Se todas as chaves de $T_1$ são menores que as de $T_2$ (junção): retire o mínimo $x$ de $T_2$; supondo $h(T_1)\ge h(T_2)$, desça pela espinha direita de $T_1$ até um nó $v$ de altura $h(T_2)$ ou $h(T_2)+1$; crie um nó $x$ com filhos $T(v)$ e $T_2$, coloque-o no lugar de $v$ e suba rebalanceando como numa inserção. Custo $O(\log n)$. No caso geral, com chaves intercaladas: faça os percursos em ordem das duas árvores ($O(n+m)$), intercale as duas listas ordenadas e construa uma árvore completa tomando sempre a mediana como raiz, em $O(n+m)$ no total. Alternativa: inserir os nós da menor na maior, em $O(m\log(n+m))$.

\resp{avl-l5} Sim: cada conjunto é uma AVL e cada nó guarda um ponteiro para o pai. \Find$(x)$ sobe até a raiz da AVL de $x$, que identifica o conjunto: $O(\log n)$. \Union{} insere os elementos da AVL menor na maior: $O(k\log n)$, sendo $k$ o tamanho da menor; com ``menor na maior'', cada elemento é reinserido no máximo $\log n$ vezes, e o total é $O(n\log^2n)$. Ganha-se ordem dentro de cada conjunto, mas a estrutura é bem pior que o Union-Find, de custo praticamente $O(1)$ amortizado.

\resp{avl-r6} (a) AVL: a ABB pode degenerar para altura $\Theta(n)$ com inserções em ordem, enquanto a AVL garante $O(\log n)$ para busca, inserção e remoção. (b) Heap: o máximo está na raiz ($O(1)$) e inserir ou remover custa $O(\log n)$, num vetor simples. A AVL também faria tudo em $O(\log n)$, mas é mais pesada, e só compensaria se também fosse preciso buscar ou remover tarefas específicas por identificador.

\resp{avl-pseudo} A inserção e as rotações estão na §\ref{sec:avl}. Remoção: remova como numa ABB (com dois filhos, troque pelo maior da subárvore esquerda), suba a partir do nó fisicamente removido usando a tabela ``a altura diminuiu?'' e, ao rotacionar, trate o caso em que o filho do lado mais pesado tem balanço 0 com rotação \emph{simples}; nesse caso a altura não muda e o processo para.

\resp{avl-juiz} 1195: insira as chaves numa ABB e imprima os três percursos. 1200 e 1201: ABB com inserção, busca, percursos e, no segundo, remoção. 1191: o primeiro símbolo da pré-ordem é a raiz; sua posição no percurso em ordem separa a subárvore esquerda (à esquerda dela) da direita; resolva recursivamente as duas partes e imprima a raiz por último. Custa $O(n^2)$ na forma simples, ou $O(n)$ com uma tabela de posições.

\resp{avl-rn} Ver §\ref{sec:rn}: definição (raiz e folhas nulas pretas, nunca dois vermelhos seguidos, mesma altura negra em todos os caminhos); $h\le2\log_2(n+1)$, porque a subárvore de $x$ tem pelo menos $2^{bh(x)}-1$ nós internos e pelo menos metade dos nós de um caminho é preta; inserção com no máximo duas rotações (casos do tio vermelho e do tio preto) e remoção com no máximo três.

\resp{avl-1a7} Há rotação simples à esquerda em 1 (ao inserir 3), em 3 (ao inserir 5), em 2 (ao inserir 6; 2 é a raiz nesse momento) e em 5 (ao inserir 7). A árvore final é perfeita: $4(2(1,3),\,6(5,7))$, com todos os balanços iguais a 0.

\resp{avl-3rot} Inserir 30: rotação simples à esquerda em 10, dando $20(10,30)$. Inserir 25 e 28: o nó 30 fica com $-2$ e o filho 25 com $+1$; rotação dupla à direita em 30, dando $20(10,\,28(25,30))$. Inserir 27 (à direita de 25): o nó 20 fica com $+2$ e o filho 28 com $-1$; rotação dupla à esquerda em 20. Árvore final: $25^{0}(20^{-1}(10,\cdot),\ 28^{0}(27,30))$.

\resp{avl-rem} Remover 1: $4(2^{+1}(\cdot,3),\,6(5,7))$, sem rotação. Remover 3: $4^{+1}(2,\,6(5,7))$, sem rotação. Remover 2: o nó 4 fica com $+2$ e o filho 6 com balanço 0, caso exclusivo da remoção; rotação simples à esquerda em 4, dando $6^{-1}(4^{+1}(\cdot,5),\,7)$, e a altura não diminuiu.

\resp{avl-altmax} $N(1),\dots,N(7)=1,2,4,7,12,20,33$. Com 20 nós cabe uma AVL mínima de altura 6 (uma árvore de Fibonacci), e a altura 7 exigiria 33 nós: a altura máxima é 6. Com 32 nós, ainda $N(7)=33>32$: a altura máxima continua 6.

\resp{avl-inordem} Na rotação simples à direita em $p$, com filho esquerdo $u$: antes, o percurso em ordem é $T_1,u,T_2,p,T_3$; depois, com $u$ na raiz, $T_1$ à esquerda e $p$ à direita com $T_2$ e $T_3$, ele é $T_1,u,T_2,p,T_3$ --- o mesmo. A rotação à esquerda é simétrica, e a dupla é a composição de duas simples. Como uma árvore binária é ABB exatamente quando o percurso em ordem é crescente, a propriedade se mantém.

\resp{avl-catalan} $C_3=5$ e $C_4=14$ (números de Catalan). Escolhendo a raiz $r$, a subárvore esquerda guarda $r-1$ chaves e a direita $n-r$, logo $C_n=\sum_{r=1}^nC_{r-1}C_{n-r}$, com $C_0=1$. Os valores são $1,1,2,5,14,42,132$.

\resp{avl-rem30} O maior elemento da subárvore esquerda é 25, que toma o lugar de 30 e é retirado da sua posição. O nó 20 fica sem filho direito e com a subárvore esquerda $10(5,15)$ de altura 2: $\bal(20)=-2$, com o filho 10 de balanço 0. Rotação simples à direita em 20 (caso exclusivo da remoção): a subárvore passa a ser $10^{+1}(5,\,20^{-1}(15,\cdot))$, sem diminuir de altura, e o processo para. Árvore final: $25(10(5,20(15,\cdot)),\,50(40(35,45),70(60,80)))$.

\resp{avl-seq} Rotações: ao inserir 1, simples à direita em 3; ao inserir 5, simples à esquerda em 3; ao inserir 6, simples à esquerda em 2 (a raiz); ao inserir 7, simples à esquerda em 5; ao inserir 15, dupla à esquerda em 7; ao inserir 14, dupla à esquerda em 6. As inserções de 3, 2, 4 e 16 não causam rotação. Árvore final: $4^{+1}(2(1,3),\ 7(6^{-1}(5,\cdot),\,15(14,16)))$.

\resp{avl-hmin} Uma árvore binária de altura $h$ tem no máximo $1+2+\cdots+2^{h-1}=2^h-1$ nós; logo $n\le2^h-1$, isto é, $h\ge\log_2(n+1)$, e, como $h$ é inteiro, $h\ge\ceil{\log_2(n+1)}$. Com o Teorema~\ref{teo:avlaltura}, a altura de uma AVL fica entre $\log_2(n+1)$ e $1{,}44\log_2(n+2)$: é $\Theta(\log n)$.

\resp{avl-verif} Uma função recursiva que, para cada nó, devolve (é válida, altura, menor chave, maior chave) da subárvore: a subárvore de $v$ é válida se as dos filhos são válidas, se a maior chave da esquerda é $\le v.chave\le$ a menor chave da direita e se as alturas dos filhos diferem de no máximo 1. Cada nó é visitado uma vez: $O(n)$.

\resp{avl-suc} Se $x$ tem filho direito, o sucessor é o mínimo da subárvore direita (desça uma vez à direita e depois sempre à esquerda). Senão, suba pelos pais enquanto $x$ for filho direito; o sucessor é o primeiro ancestral do qual $x$ está na subárvore esquerda (ou não existe, se chegar à raiz por um caminho só de filhos direitos). Em ambos os casos o caminho percorrido tem comprimento no máximo $h$.

\respcap{cap:sortcomp}
\resp{sc-merge} Ver o procedimento \textsc{Intercala} da §\ref{sec:sortcomp}: dois índices percorrem as metades, o menor elemento apontado vai para o vetor auxiliar (``$\le$'' para manter a estabilidade), o que sobra de uma das metades é copiado ao final e o auxiliar volta para o vetor. Custo $\Theta(n)$.

\resp{sc-iter} Intercale blocos de tamanho $1,2,4,\dots$:
\begin{pseudo}{Mergesort iterativo (de baixo para cima)}
\begin{algorithmic}
\State $tam:=1$
\While{$tam<n$}
  \State $ini:=1$
  \While{$ini+tam\le n$}
    \State $meio:=ini+tam-1$;\ $fim:=\min\{ini+2\,tam-1,\ n\}$
    \State \Call{Intercala}{$V,ini,meio,fim$};\ $ini:=ini+2\,tam$
  \EndWhile
  \State $tam:=2\,tam$
\EndWhile
\end{algorithmic}
\end{pseudo}
São $\ceil{\log_2n}$ passadas de custo $O(n)$: $\Theta(n\log n)$, sem pilha de recursão.

\resp{sc-semcomp} Se as chaves são inteiros em $[0,k]$, conte as ocorrências de cada valor e use as contagens para posicionar os elementos (\emph{counting sort}, $\Theta(n+k)$). Se têm $d$ dígitos na base $b$, ordene dígito a dígito, do menos significativo para o mais significativo, com um método estável (\emph{radix sort}, $\Theta(d(n+b))$). Se são reais uniformemente distribuídos em $[0,1)$, distribua-os em $n$ baldes (\emph{bucket sort}, $\Theta(n)$ esperado). Todas usam a estrutura das chaves (Capítulo~\ref{cap:sortlin}); para chaves arbitrárias, só resta comparar.

\resp{sc-estavel} (a) Sim: na intercalação, em caso de empate, copie o elemento da metade esquerda (teste ``$\le$''). Por indução, as metades estão ordenadas de forma estável, e a intercalação preserva a ordem relativa dos iguais, que vêm todos da esquerda antes dos da direita. (b) A versão in-place usual não é estável: em $[2_a,2_b,1]$, com pivô 1, a troca de $2_a$ com 1 coloca $2_a$ depois de $2_b$. É possível torná-lo estável com uma partição que copia os menores, os iguais e os maiores para vetores auxiliares na ordem original, ao custo de $\Theta(n)$ de memória extra, ou usando como chave o par (valor, posição original).

\resp{sc-mestre} (a) $E=2$ e $f=n=O(n^{2-\varepsilon})$: $\Theta(n^2)$. (b) $f=\Theta(n^2)$: $\Theta(n^2\log n)$. (c) $f=\Omega(n^{2+\varepsilon})$ e $4(n/2)^3=n^3/2$: $\Theta(n^3)$. (d) $E=1/2$ e $f=\Theta(\sqrt n)$: $\Theta(\sqrt n\log n)$. (e) $E=2$: $\Theta(n^2)$. (f) $E=0$, $f=n$ e $n/2\le cn$: $\Theta(n)$.

\resp{sc-arvore} Cada nível custa no máximo $n$, pois os subproblemas de um nível somam no máximo $n$. O ramo mais curto ($n\to n/3\to\cdots$) tem $\log_3n$ níveis completos e o mais longo ($n\to2n/3\to\cdots$) tem $\log_{3/2}n$ níveis. Logo $n\log_3n\le T(n)\le n\log_{3/2}n+O(n)$: $\Theta(n\log n)$.

\resp{sc-pior8} $W(1)=0$ e $W(n)=2W(n/2)+(n-1)$, pois a intercalação de duas listas de tamanho $n/2$ faz no máximo $n-1$ comparações. $W(2)=1$, $W(4)=5$ e $W(8)=17$. Em geral, $W(n)=n\log_2n-n+1$ para $n$ potência de 2.

\resp{sc-max} Todo elemento que não é o máximo precisa ``perder'' pelo menos uma comparação; senão, o algoritmo não teria como descartá-lo, e um adversário poderia torná-lo o maior. Cada comparação tem um único perdedor, logo são necessárias pelo menos $n-1$ comparações. A árvore de decisão dá só $\log_2n$, porque há apenas $n$ respostas possíveis --- um limite fraco.

\resp{sc-iguais} Com todos os elementos iguais, o teste $V[j]<x$ nunca é verdadeiro: o pivô fica na primeira posição e as partes têm tamanhos $0$ e $n-1$. Então $T(n)=T(n-1)+\Theta(n)=\Theta(n^2)$. Correção: particionar em três partes ($<x$, $=x$ e $>x$) e não recursar na parte dos iguais, ou usar uma partição que divide os iguais entre os dois lados.

\resp{sc-insercao} Ao inserir $V[i]$, ele passa por cima exatamente dos elementos anteriores maiores que ele, isto é, das inversões $(j,i)$ com $j<i$. Cada deslocamento desfaz uma inversão e nenhuma é criada. O total é o número de inversões: até $\binom n2$ no pior caso (vetor invertido) e $\binom n2/2$ em média, pois cada par está invertido com probabilidade $1/2$. Logo $\Theta(n^2)$ nos dois casos.

\resp{sc-intercalar} Considere a entrada $a_1<b_1<a_2<b_2<\cdots<a_n<b_n$. Cada par de elementos vizinhos na saída, $(a_i,b_i)$ ou $(b_i,a_{i+1})$, precisa ser comparado: se um deles não fosse, seria possível trocar os dois valores sem alterar nenhuma resposta das comparações feitas, e a saída ficaria errada. São $2n-1$ pares vizinhos, todos distintos.

\resp{sc-mtrace} As intercalações acontecem nesta ordem: $[27,38]$, $[3,43]$, $[3,27,38,43]$, $[9,82]$, $[9,10,82]$ e, por fim, $[3,9,10,27,38,43,82]$.

\resp{sc-lomuto} Resultado $[2,1,3,4,7,5,6,8]$, e a posição devolvida é 4 (o pivô 4 ocupa sua posição definitiva).

\resp{sc-quickpior} Num vetor ordenado, o pivô (o último elemento) é sempre o máximo do intervalo: uma partição de $k$ elementos faz $k-1$ comparações e deixa uma parte com $k-1$ elementos e outra vazia. O total é $\sum_{k=2}^n(k-1)=\frac{n(n-1)}2$.

\resp{sc-subst} Para $n=1$, $T(1)=0=1\cdot\log_21$. Supondo $T(n/2)=\frac n2\log_2\frac n2$, temos $T(n)=2\cdot\frac n2(\log_2n-1)+n=n\log_2n-n+n=n\log_2n$.

\resp{sc-selecao} A ordenação por seleção faz sempre $\sum_{i=1}^{n-1}(n-i)=\frac{n(n-1)}2$ comparações e no máximo $n-1$ trocas. Não é estável: em $[2_a,2_b,1]$, a primeira troca leva $2_a$ para o fim, depois de $2_b$. A ordenação por inserção é estável, pois um elemento só passa por cima dos estritamente maiores.

\resp{sc-cinco} A árvore de decisão precisa de pelo menos $5!=120$ folhas, logo de altura $\ge\ceil{\log_2120}=7$. Para 12 elementos, $\ceil{\log_2(12!)}=\ceil{28{,}8}=29$.

\resp{sc-kdist} Se cada um dos $k$ valores aparece $n/k$ vezes, há $\frac{n!}{((n/k)!)^k}$ saídas distintas possíveis, e a altura da árvore de decisão é pelo menos o logaritmo desse número, que é $n\log n-n\log(n/k)-O(n)=\Omega(n\log k)$ (usando $\log m!=m\log m-O(m)$). Algoritmo: insira os elementos numa árvore AVL de valores \emph{distintos}, com um contador por valor; a árvore tem no máximo $k$ nós, então cada inserção custa $O(\log k)$; depois, percorra-a em ordem, escrevendo cada valor tantas vezes quanto seu contador. Total $O(n\log k)$.

\resp{sc-nlog2} O nível $j$ tem $2^j$ subproblemas de tamanho $n/2^j$, com custo total $n\log(n/2^j)=n(\log n-j)$. Somando de $j=0$ a $\log n-1$: $n\sum_{j}(\log n-j)=n\cdot\frac{\log n(\log n+1)}2=\Theta(n\log^2n)$. O Teorema Mestre não se aplica porque $f(n)=n\log n$ é maior que $n^E=n$, mas não por um fator polinomial $n^\varepsilon$: cai no ``vão'' entre os casos 2 e 3.

\resp{sc-lineares} (a) $T(n)=1+2+\cdots+n=\frac{n(n+1)}2$. (b) $T(n)=2^n-1$: por indução, $2(2^{n-1}-1)+1=2^n-1$. (c) $T(n)=1+\sum_{i=2}^n\log i=1+\log(n!)=\Theta(n\log n)$.

\respcap{cap:sortlin}
\resp{sl-t2} (A) \textbf{Counting sort}: $k=101$ é muito menor que $n$, o custo é $\Theta(n+k)=\Theta(n)$ e bastam 101 contadores; uma ordenação por comparação faria cerca de $n\log_2n\approx2{,}3\cdot10^8$ comparações. (B) \textbf{Radix sort}: o counting sort é inviável ($k=10^{12}$); com $d=12$ e $b=10$ são 12 passadas lineares, $\Theta(d(n+b))\approx1{,}2\cdot10^7$ (ou 4 e 3 passadas nas bases $10^3$ e $10^4$). (C) \textbf{Bucket sort} com $n$ baldes: $\Theta(n)$ esperado para entrada uniforme; counting e radix não se aplicam a reais. (D) \textbf{Radix sort} depois de somar $10^9$ a todos os valores (as chaves passam a estar em $[0,2\cdot10^9]$, com cerca de 31 bits: 2 passadas na base $2^{16}$); o counting sort exigiria $2\cdot10^9$ contadores. Mergesort ou quicksort, em $\Theta(n\log n)$, seriam aceitáveis, mas o radix é linear. (E) Só resta comparar: \textbf{mergesort}, com garantia de $\Theta(n\log n)$ e estabilidade, boa para registros; ou quicksort, mais rápido em média e in-place, mas $\Theta(n^2)$ no pior caso e não estável.

\resp{sl-var} Radix: agrupe os números pelo número de dígitos (um counting sort por $d_i$). Como um número com menos dígitos é menor, basta ordenar cada grupo separadamente com um radix de $d$ passadas, sendo $d$ o tamanho dos números do grupo, e concatenar os grupos em ordem crescente de $d$. Um grupo com $n_d$ números de $d$ dígitos custa $O(d\,(n_d+b))$; somando, $O\big(\sum_id_i+b\cdot D\big)$, onde $D\le\max d_i$ é o número de tamanhos distintos. American Flag: primeiro separe por tamanho (mesma ideia); dentro de cada grupo, a recursão por dígito processa cada dígito de cada número uma única vez, com $O(b)$ de trabalho por chamada.

\resp{sl-str} Radix (do último caractere para o primeiro): complete as strings menores à direita com um símbolo ``menor que todos'' até o comprimento máximo $L$ e ordene da última posição para a primeira, com counting sort estável sobre o alfabeto, em $O(L(n+|\Sigma|))$; para não pagar o custo de completar, agrupe por comprimento como no exercício anterior. American Flag (do primeiro caractere para o último): distribua pelo primeiro caractere, com um balde extra para ``a string terminou'', que vem antes de todos, e chame recursivamente cada balde com o caractere seguinte; strings que terminam saem da recursão.

\resp{sl-count} Contagens: $C=[2,0,2,3,0,1]$. Depois das somas acumuladas: $C=[2,2,4,7,7,8]$. Percorrendo $A$ de trás para frente, cada $A[j]$ vai para a posição $C[A[j]]$, que é então decrementada. Saída: $B=[0,0,2,2,3,3,3,5]$.

\resp{sl-radix} Pelas unidades: $720,355,436,457,657,329,839$. Pelas dezenas: $720,329,436,839,355,457,657$. Pelas centenas: $329,355,436,457,657,720,839$.

\resp{sl-n3} Escreva cada número na base $n$: como $x<n^3$, ele tem no máximo 3 dígitos, $x\bmod n$, $\floor{x/n}\bmod n$ e $\floor{x/n^2}$. O radix sort com 3 passadas de counting sort, cada uma com $k=n$, custa $\Theta(3(n+n))=\Theta(n)$.

\resp{sl-frente} A saída continua ordenada, pois cada valor ainda ocupa o seu bloco, mas os elementos iguais saem em ordem \emph{invertida}: o primeiro a ser lido vai para a última posição do bloco. O counting sort deixa de ser estável, e o radix sort que o usa em cada passada \textbf{deixa de ordenar corretamente}, porque cada passada desfaz a ordem obtida pelos dígitos anteriores.

\resp{sl-bucket} Se todos os elementos caírem no mesmo balde, a inserção custa $\Theta(n^2)$. Ordene cada balde com mergesort ou heapsort: o pior caso passa a $O(n\log n)$, e o caso esperado continua $\Theta(n)$, porque um balde com $n_i$ elementos custa $O(n_i\log n_i)\le O(n_i^2)$ e $\sum_iE[n_i^2]=O(n)$.

\resp{sl-msd} No radix que começa pelo dígito menos significativo, a passada pelo dígito $i$ precisa preservar, entre elementos com o mesmo dígito $i$, a ordem estabelecida pelos dígitos $1,\dots,i-1$, que é a única informação sobre eles. No que começa pelo mais significativo, depois de separar pelo primeiro dígito, cada bloco é ordenado recursivamente e de forma independente pelos dígitos seguintes; a ordem dentro do bloco antes da recursão não importa, porque a recursão a refaz.

\resp{sl-neg} Some $m$ a todos os valores, que passam a estar em $[0,2m]$, ordene com o counting sort e subtraia $m$ na saída (ou indexe o vetor de contadores por $x+m$). O custo é $\Theta(n+m)$.

\resp{sl-n2} O counting sort direto usaria $k=n^2$ contadores e custaria $\Theta(n^2)$. Na base $n$, cada número tem 2 dígitos, e o radix sort faz 2 passadas de counting sort com $k=n$: $\Theta(n)$ de tempo e de memória.

\resp{sl-strtrace} Pela terceira letra: SEA TEA MOB TAB DOG RUG DIG BIG BAR EAR TAR COW ROW NOW BOX FOX. Pela segunda: TAB BAR EAR TAR SEA TEA DIG BIG MOB DOG COW ROW NOW BOX FOX RUG. Pela primeira: BAR BIG BOX COW DIG DOG EAR FOX MOB NOW ROW RUG SEA TAB TAR TEA.

\resp{sl-buckettrace} Balde 1: $0{,}17;\ 0{,}12$. Balde 2: $0{,}26;\ 0{,}21;\ 0{,}23$. Balde 3: $0{,}39$. Balde 6: $0{,}68$. Balde 7: $0{,}78;\ 0{,}72$. Balde 9: $0{,}94$. Ordenando cada balde e concatenando: $0{,}12;\ 0{,}17;\ 0{,}21;\ 0{,}23;\ 0{,}26;\ 0{,}39;\ 0{,}68;\ 0{,}72;\ 0{,}78;\ 0{,}94$.

\resp{sl-base} Com $b=2^4$: 8 passadas de custo $\approx n+16$. Com $b=2^8$: 4 passadas de custo $\approx n+256$. Com $b=2^{16}$: 2 passadas de custo $\approx n+65536$. Como $n=10^6$, o termo $b$ é desprezível nas três, e a base $2^{16}$ faz menos passadas, com custo total $\approx2\cdot10^6$. Na prática, a base $2^8$ também é muito usada, porque o vetor de contadores cabe na memória cache.

\resp{sl-prova} Hipótese: depois da passada $i$, a lista está ordenada pelos $i$ dígitos menos significativos. Para $i=1$ vale, pois a passada ordena pelo dígito 1. Suponha válido para $i-1$ e tome dois números $x$ e $y$. Se os seus dígitos $i$ diferem, a passada $i$ os coloca na ordem desse dígito, que decide a comparação entre os $i$ últimos dígitos. Se são iguais, a passada, por ser estável, mantém a ordem relativa que eles tinham, que, pela hipótese de indução, é a ordem pelos $i-1$ últimos dígitos. Em ambos os casos, ficam na ordem certa pelos $i$ últimos dígitos.

\resp{sl-registros} Counting sort pela nota ($k=11$), que é \textbf{estável}: alunos com a mesma nota continuam na ordem em que estavam, isto é, em ordem alfabética. O custo é $\Theta(n)$. Um algoritmo não estável, como o quicksort ou o heapsort, poderia desfazer a ordem alfabética.

\resp{sl-datas} Do menos para o mais significativo: primeiro pelo dia (base 31 ou 32), depois pelo mês (base 12) e por fim pelo ano (base igual à amplitude dos anos, ou em duas passadas de dígitos), sempre com passadas estáveis.

\resp{sl-dup} Ordene com o radix sort na base $n+1$: os números de $[0,n^2]$ têm no máximo 3 dígitos nessa base (2 bastam para valores menores que $(n+1)^2$), e a ordenação custa $O(n)$. Depois, verifique se há dois elementos vizinhos iguais.

\resp{sl-intervalo} Calcule as contagens acumuladas $C[i]=\#\{j:A[j]\le i\}$, como no counting sort, em $O(n+k)$. A resposta para $[a,b]$ é $C[b]-C[a-1]$ (com $C[-1]=0$), em $O(1)$.

\resp{sl-bucketruim} Por exemplo, $x_i=\frac{i}{n^2}$ para $i=1,\dots,n$: todos ficam no intervalo $[0,1/n)$ e caem no balde 0. Se estiverem em ordem decrescente na entrada, a inserção nesse balde faz $\Theta(n^2)$ deslocamentos.

\resp{sl-flag} Cada nível da recursão usa um vetor de $b$ contadores (e as posições de início dos blocos); a recursão tem profundidade $d$, de modo que a memória auxiliar é $O(b\cdot d)$, mais a pilha de recursão. Diz-se que ele é in-place porque essa memória não depende de $n$: os elementos são permutados dentro do próprio vetor, sem um vetor auxiliar de tamanho $n$.

\respcap{cap:selecao}
\resp{se-quick} Quickselect com pivô aleatório (§\ref{sec:selecao}): particione em $S_1=\{<m\}$, $S_2=\{=m\}$ e $S_3=\{>m\}$ e recursione apenas no lado que contém o $k$-ésimo, com $k=\ceil{n/2}$. Com probabilidade pelo menos $1/2$, o pivô fica entre os quartis e o lado mantido tem no máximo $3n/4$ elementos; agrupando as chamadas em fases, $E[T(n)]\le E[T(3n/4)]+2cn$, e portanto $E[T(n)]\le8cn=O(n)$.

\resp{se-qs} Use como pivô a mediana exata, obtida pela mediana das medianas em $O(n)$. As duas partes têm no máximo $n/2$ elementos, e $T(n)=2T(n/2)+O(n)=O(n\log n)$ no \textbf{pior} caso. Na prática, as constantes da mediana das medianas (grupos, medianas, duas recursões) tornam esse algoritmo bem mais lento que o quicksort com pivô aleatório ou mediana de três, cujo pior caso é raríssimo. Por isso, implementações reais preferem trocar para o heapsort quando a recursão fica profunda demais.

\resp{se-trace} Medianas dos grupos: $12,9,11,14,16$; pivô $m=12$, a mediana das medianas. $|S_1|=11$ (os elementos $1,\dots,11$) e $|S_3|=13$; ambos são pelo menos $3\floor{25/10}=6$. (O algoritmo, de fato, ordenaria diretamente, pois $n\le49$; aqui o exemplo serve só para praticar a escolha do pivô.)

\resp{se-grupos} Grupos de 7: pelo menos metade das $n/7$ medianas é maior ou igual a $m$, e cada uma traz 4 elementos maiores ou iguais a $m$; há pelo menos $4\cdot\frac n{14}=\frac{2n}7$ elementos de cada lado, e cada parte tem no máximo $\frac{5n}7$ (mais uma constante). Então $T(n)\le T(n/7)+T(5n/7)+O(n)$, com $\frac17+\frac57=\frac67<1$: o algoritmo é \textbf{linear}. Grupos de 3: só há $2\cdot\frac n6=\frac n3$ elementos garantidos de cada lado, e $T(n)\le T(n/3)+T(2n/3)+O(n)$, com $\frac13+\frac23=1$: $\Theta(n\log n)$, \textbf{não} linear.

\resp{se-kmenores} Encontre o $k$-ésimo menor $x$ pela mediana das medianas, em $O(n)$, e percorra o conjunto separando os elementos menores que $x$, completando com cópias de $x$ até ter $k$ elementos: $O(n)$. Para devolvê-los em ordem crescente, ordene esses $k$ elementos: $O(n+k\log k)$.

\resp{se-oleo} O comprimento total é $\sum_i|y_i-y|$, que é minimizado em qualquer \textbf{mediana} dos $y_i$: se há mais poços acima do que abaixo de $y$, subir um pouco o oleoduto diminui a soma, e vice-versa; na mediana, os dois lados se equilibram. A mediana é encontrada em $O(n)$.

\resp{se-piorqs} Se o pivô escolhido for sempre o maior elemento --- por exemplo, com o pivô determinístico ``último elemento'' num vetor crescente, ou por um azar sistemático na escolha aleatória ---, cada chamada descarta apenas o pivô, e $T(n)=T(n-1)+\Theta(n)=\Theta(n^2)$. Isso acontece mesmo procurando o mínimo ($k=1$). Com pivô aleatório, esse caso tem probabilidade desprezível.

\resp{se-qstrace} 1ª chamada: pivô 1, $S_1=\emptyset$, $S_2=\{1\}$, $S_3=\{7,2,9,4,3,8,6\}$; como $0<4$ e $1<4$, segue para $S_3$ com $k=3$. 2ª: pivô 6, $S_1=\{2,4,3\}$, $S_2=\{6\}$, $S_3=\{7,9,8\}$; como $|S_1|=3\ge3$, segue para $S_1$ com $k=3$. 3ª: pivô 3, $S_1=\{2\}$, $S_2=\{3\}$, $S_3=\{4\}$; como $1<3$ e $2<3$, segue para $S_3$ com $k=1$. 4ª: pivô 4, $S_2=\{4\}$: devolve \textbf{4}, que é de fato o quarto menor.

\resp{se-n49} Para $n=49$: $49-3\cdot4=37>36{,}75=3\cdot49/4$. Para $n=50$: $50-3\cdot5=35\le37{,}5$. A garantia ``cada lado tem no máximo $3n/4$'' só vale a partir de $n=50$; por isso os conjuntos com até 49 elementos são resolvidos diretamente.

\resp{se-ordmed} Ordenar as $n/5$ medianas custa $O(\frac n5\log\frac n5)=O(n\log n)$, e a recorrência passa a ser $T(n)=T(3n/4)+O(n\log n)$, cuja solução é $O(n\log n)$ (soma geométrica). O algoritmo continua correto, mas deixa de ser linear: a recursão sobre as medianas é essencial.

\resp{se-lista} Sim. Formar os grupos de 5, calcular suas medianas e particionar em $S_1$, $S_2$, $S_3$ são operações que percorrem a sequência uma vez, o que se faz numa lista encadeada em tempo linear; nenhum passo depende de acesso direto por posição.

\resp{se-quartis} Selecione o elemento de posição $\ceil{n/2}$ (a mediana) e particione; depois selecione, em cada metade, o elemento do meio. São três seleções sobre conjuntos de tamanho no máximo $n$, em $O(n)$ no total.

\resp{se-proximos} Ache a mediana $m$ em $O(n)$. Calcule $d_i=|x_i-m|$ para todo $i$ e selecione o $k$-ésimo menor $d$, digamos $d^*$, em $O(n)$. Os $k$ números mais próximos são os $x_i$ com $d_i<d^*$, completados com elementos de $d_i=d^*$ até somar $k$ (excluindo a própria mediana, se ela não deve contar). Total $O(n)$.

\resp{se-ponderada} Recursão parecida com a seleção: ache a mediana $m$ (por valor) do conjunto atual em $O(n)$ e some os pesos $W_<$ dos elementos menores que $m$ e $W_>$ dos maiores. Se $W_<<1/2$ e $W_>\le1/2$, $m$ é a resposta. Se $W_<\ge1/2$, a resposta está entre os menores: descarte os maiores (e $m$), somando o peso deles a um contador do lado direito. Caso contrário, faça o simétrico. Cada passo descarta metade dos elementos: $T(n)=T(n/2)+O(n)=O(n)$.

\resp{se-dois} Compare as medianas $X[\ceil{n/2}]$ e $Y[\ceil{n/2}]$. Se a de $X$ é menor, a mediana procurada não está na metade inferior de $X$ nem na metade superior de $Y$; descarte essas duas metades, de mesmo tamanho, o que não muda a mediana do conjunto restante. Se é maior, descarte as outras duas metades. Repita até sobrarem poucos elementos e resolva diretamente. Cada passo divide $n$ por 2: $O(\log n)$.

\resp{se-segundo} Organize um torneio eliminatório: compare os elementos aos pares e faça os vencedores (os menores) se enfrentarem, até restar o mínimo, com $n-1$ comparações. O segundo menor perdeu, necessariamente, para o mínimo em algum confronto. O mínimo disputou no máximo $\ceil{\log_2n}$ confrontos; o menor dos que perderam para ele, encontrado com $\ceil{\log_2n}-1$ comparações, é o segundo menor. Total: $n+\ceil{\log_2n}-2$.

\resp{se-cinco} Sejam $a,b,c,d,e$. Compare $a$ com $b$ e $c$ com $d$; suponha $a<b$ e $c<d$ (2 comparações). Compare os menores, $a$ e $c$; suponha $a<c$ (3). Então $a$ é menor que $b$, $c$ e $d$, e não pode ser a mediana; a mediana dos cinco é o segundo menor de $\{b,c,d,e\}$, com $c<d$ já sabido. Compare $b$ com $e$; suponha $b<e$ (4). Compare $b$ com $c$ (5): se $b<c$, o menor dos quatro é $b$ e a resposta é $\min\{c,e\}$; se $c<b$, o menor é $c$ e a resposta é $\min\{b,d\}$. A última comparação (6) decide.

\resp{se-inferior} Considere o grafo com os $n$ elementos como vértices e uma aresta para cada comparação feita. Se ele fosse desconexo, haveria uma componente $C$ que não contém a mediana devolvida $x$. Somando a todos os elementos de $C$ uma constante muito grande (se algum elemento de $C$ é menor que $x$) ou subtraindo (se todos são maiores), todas as comparações feitas continuariam com o mesmo resultado, mas o número de elementos menores que $x$ mudaria, e $x$ deixaria de ser a mediana. Logo o grafo é conexo e tem pelo menos $n-1$ arestas.

\resp{se-esperado} Com pivô aleatório, cada posição final $i$ do pivô tem probabilidade $1/n$, e a partição custa no máximo $n$ comparações. No pior caso, a recursão vai para o lado maior, de tamanho $\max\{i-1,n-i\}$, e cada tamanho entre $\floor{n/2}$ e $n-1$ aparece no máximo duas vezes; daí a desigualdade. Supondo $E[T(i)]\le4i$ para $i<n$: $E[T(n)]\le n+\frac2n\sum_{i=\floor{n/2}}^{n-1}4i\le n+\frac8n\cdot\frac{3n^2}8=4n$, usando $\sum_{i=\floor{n/2}}^{n-1}i\le\frac{3n^2}8$.

\resp{se-maior} O $k$-ésimo maior é o $(n-k+1)$-ésimo menor. O percentil 90 de $n$ notas é o elemento de posição $\ceil{0{,}9n}$ na ordem crescente, obtido por uma seleção em $O(n)$.

\resp{se-qstrace2} Pivô 9, com $k=7$: $S_1=\{3,7,1,5\}$, $S_2=\{9\}$ e $S_3=\{14,21,16,11\}$; como $4+1<7$, segue em $S_3$ com $k=7-5=2$. Pivô 16: $S_1=\{14,11\}$, $S_2=\{16\}$ e $S_3=\{21\}$; como $|S_1|=2\ge2$, segue em $S_1$ com $k=2$. Pivô 11: $S_1=\emptyset$, $S_2=\{11\}$ e $S_3=\{14\}$; como $0+1<2$, segue em $S_3$ com $k=1$. Pivô 14: devolve \textbf{14}. De fato, a ordem crescente é $1,3,5,7,9,11,\mathbf{14},16,21$.

\resp{se-trace2} Medianas dos grupos: $15$, $19$ e $18$; o pivô é a mediana delas, $m=18$. Há 7 elementos menores ($8,15,4,12,2,10,6$) e 7 maiores ($31,22,27,19,35,26,30$), bem mais que a garantia de $3\floor{15/10}=3$ de cada lado. Aqui o pivô coincide com a mediana verdadeira, o 8º menor dos 15 elementos; em geral, ele é só uma aproximação, como no Exemplo~\ref{ex:mmedianas}.

\respcap{cap:dc}
\resp{dc-t1} Ver §\ref{sec:dc}: o melhor subvetor está inteiro na metade esquerda, inteiro na direita ou cruza o meio; o cruzado é o melhor sufixo que termina em $meio$ mais o melhor prefixo que começa em $meio+1$, ambos calculados em $O(n)$. $T(n)=2T(n/2)+\Theta(n)=\Theta(n\log n)$. O caso base com um elemento garante subvetor não vazio e resolve o caso em que todos os elementos são negativos.

\resp{dc-t2} Ver §\ref{sec:dc}: particione $A$ e $B$ em blocos $\frac n2\times\frac n2$, calcule recursivamente $M_1,\dots,M_7$ e combine $C_{11}=M_1+M_4-M_5+M_7$, $C_{12}=M_3+M_5$, $C_{21}=M_2+M_4$ e $C_{22}=M_1-M_2+M_3+M_6$. $T(n)=7T(n/2)+\Theta(n^2)=\Theta(n^{\log_27})$.

\resp{dc-balanc} Na árvore de recursão, os subproblemas de um mesmo nível são partes disjuntas da entrada, então a soma dos seus tamanhos é no máximo $n$, e o trabalho linear de dividir e combinar, somado no nível, é $O(n)$. Como o tamanho cai por um fator de pelo menos $c$ a cada nível, há no máximo $\log_cn+1$ níveis. O total é $O(n\log_cn)=O(n\log n)$. Para partes iguais, é a recorrência $T(n)=c\,T(n/c)+an$, resolvida pelo caso 2 do Teorema Mestre.

\resp{dc-inv} Ver as funções \textsc{CalculaInversões} e \textsc{IntercalaContando} da §\ref{sec:dc}: um mergesort que, ao copiar um elemento da metade direita, soma $meio-i+1$, a quantidade de elementos que restam na esquerda, todos maiores que ele. $\Theta(n\log n)$.

\resp{dc-contar} (a) 3: $(2,1)$, $(4,1)$ e $(4,3)$. (b) $10=\binom52$. (c) 8. O máximo é $\binom n2$, atingido pelo vetor decrescente.

\resp{dc-pico} Compare $V[meio]$ com $V[meio+1]$: se $V[meio]<V[meio+1]$, o pico está em $[meio+1,fim]$; senão, está em $[ini,meio]$. $T(n)=T(n/2)+O(1)=O(\log n)$.

\resp{dc-ocorr} Faça duas buscas binárias: uma pela primeira posição com $V[i]\ge x$ e outra pela primeira posição com $V[i]>x$; a resposta é a diferença entre elas. $O(\log n)$.

\resp{dc-kara} $xy=x_1y_1B^n+\big((x_1+x_0)(y_1+y_0)-x_1y_1-x_0y_0\big)B^{n/2}+x_0y_0$: são três produtos de números com cerca de $n/2$ dígitos, mais somas e deslocamentos em $O(n)$. $T(n)=3T(n/2)+\Theta(n)=\Theta(n^{\log_23})\approx\Theta(n^{1{,}585})$, contra $\Theta(n^2)$ do método escolar.

\resp{dc-maxmin} $T(2)=1$ e $T(n)=2T(n/2)+2$: uma comparação para o máximo e outra para o mínimo na combinação. Por indução, $T(n)=\frac{3n}2-2$, pois $2\big(\frac{3n}4-2\big)+2=\frac{3n}2-2$. Os valores são $1,4,10,22$ para $n=2,4,8,16$, contra $2n-2$ do método ingênuo.

\resp{dc-maj} Algoritmo de votação de Boyer e Moore: mantenha um candidato e um contador. Para cada elemento, se o contador é 0, o elemento vira o candidato, com contador 1; senão, some 1 se ele é igual ao candidato e subtraia 1 se é diferente. Cada decremento ``cancela'' um par de elementos diferentes, e um elemento majoritário não pode ser totalmente cancelado; logo, se existir, ele é o candidato final. Uma segunda passada confere se o candidato aparece mais de $n/2$ vezes. Tempo $O(n)$ e memória $O(1)$.

\resp{dc-bbin} \textbf{Enquanto} $ini\le fim$: $meio:=\floor{(ini+fim)/2}$; se $V[meio]=x$, devolva $meio$; se $x<V[meio]$, faça $fim:=meio-1$; senão, $ini:=meio+1$. Ao final, devolva $-1$. Cada iteração reduz o tamanho $s$ do intervalo para no máximo $\floor{s/2}$, então, partindo de $n$, após $t$ iterações sobram no máximo $\floor{n/2^t}$ elementos; o intervalo fica vazio após no máximo $\floor{\log_2n}+1$ iterações, cada uma com uma comparação (contando o teste de igualdade e o de ordem como uma comparação de três resultados).

\resp{dc-pot} Cada nível da recursão faz uma ou duas multiplicações, e há $\floor{\log_2n}+1$ níveis com $n\ge1$; logo, no máximo $2\floor{\log_2n}+2$ multiplicações. Para $n=13=1101_2$: os níveis com $n=1$, $3$ e $13$ são ímpares (2 multiplicações cada, contando a do nível $n=1$, que multiplica $1\cdot1\cdot a$) e o nível com $n=6$ é par (1 multiplicação): 7 ao todo, ou 5 se não contarmos as duas multiplicações por 1 do nível $n=1$.

\resp{dc-linear} Para cada intervalo, devolva $(tot,pref,suf,melhor)$. No caso base, com um elemento $a$, todos valem $a$. Na combinação de $E$ (esquerda) com $D$ (direita): $tot=tot_E+tot_D$; $pref=\max\{pref_E,\ tot_E+pref_D\}$; $suf=\max\{suf_D,\ tot_D+suf_E\}$; $melhor=\max\{melhor_E,\ melhor_D,\ suf_E+pref_D\}$. A combinação custa $O(1)$, e $T(n)=2T(n/2)+O(1)=\Theta(n)$.

\resp{dc-blocos} $C_{11}=A_{11}B_{11}+A_{12}B_{21}$, $C_{12}=A_{11}B_{12}+A_{12}B_{22}$, $C_{21}=A_{21}B_{11}+A_{22}B_{21}$ e $C_{22}=A_{21}B_{12}+A_{22}B_{22}$: 8 produtos recursivos e 4 somas de blocos $\frac n2\times\frac n2$ por chamada. $T(n)=8T(n/2)+\Theta(n^2)=\Theta(n^{\log_28})=\Theta(n^3)$.

\resp{dc-faltante} Antes do número que falta, $A[i]=i$; a partir dele, $A[i]=i+1$. Por busca binária, encontre a primeira posição $i$ com $A[i]\ne i$: o número que falta é $i$ (ou $n+1$, se não houver tal posição). $O(\log n)$.

\resp{dc-fixo} Como os elementos são inteiros distintos e ordenados, $A[i+1]-(i+1)\ge A[i]-i$: a função $g(i)=A[i]-i$ é não decrescente. Faça uma busca binária por um zero de $g$: se $g(meio)<0$, procure à direita; se $g(meio)>0$, à esquerda. $O(\log n)$.

\resp{dc-uns} Busca binária pela primeira posição com valor 1; o número de uns é $n$ menos essa posição mais 1. Se $n$ é desconhecido, consulte as posições $1,2,4,8,\dots$ até encontrar um 1 (ou o fim): isso delimita, em $O(\log n)$ consultas, um intervalo onde fazer a busca binária (busca exponencial).

\resp{dc-signif} Como na contagem de inversões, recursione nas metades, que voltam ordenadas. Antes de intercalar, conte os pares cruzados com dois ponteiros: para cada $a_i$ da esquerda, em ordem crescente, avance um ponteiro $j$ na direita enquanto $a_i>2b_j$; o número de elementos da direita já ultrapassados é a contagem de $a_i$, e o ponteiro nunca volta, porque os $a_i$ crescem. Esse passo custa $O(n)$; depois intercale normalmente. Total $O(n\log n)$.

\resp{dc-par} Ordene os pontos pela abscissa e divida-os pela mediana. Resolva cada lado, obtendo as distâncias mínimas $\delta_E$ e $\delta_D$, e seja $\delta=\min\{\delta_E,\delta_D\}$. Um par melhor teria de cruzar a reta divisória e ter os dois pontos na faixa de largura $2\delta$ em torno dela. Ordenando os pontos da faixa pela ordenada, cada ponto só precisa ser comparado com os 7 seguintes: num retângulo $\delta\times2\delta$ acima dele cabem no máximo 8 pontos com distâncias mútuas $\ge\delta$ dentro de cada lado. Mantendo as listas ordenadas por ordenada durante a recursão (como no mergesort), a combinação custa $O(n)$ e o total é $O(n\log n)$.

\resp{dc-skyline} Represente um contorno como uma lista de pontos $(x,h)$ em que a altura muda. Divida os prédios em duas metades, calcule recursivamente os dois contornos e \emph{intercale-os} como no mergesort: percorra as duas listas em ordem de $x$, mantendo a altura atual de cada contorno, e em cada $x$ registre $\max\{h_1,h_2\}$ sempre que esse máximo mudar. A intercalação custa $O(n)$ e o total é $O(n\log n)$.

\resp{dc-recs} Busca binária: $T(n)=T(n/2)+O(1)=\Theta(\log n)$. Potência: $T(n)=T(n/2)+O(1)=\Theta(\log n)$ multiplicações. \textsc{MaxMin}: $T(n)=2T(n/2)+2$, isto é, $\frac{3n}2-2$ comparações. Elemento majoritário: $T(n)=2T(n/2)+\Theta(n)=\Theta(n\log n)$, por causa das contagens na combinação.

\respcap{cap:revisao}
\resp{rv-maior} Cada iteração custa $\Theta(n)$ só para achar o máximo, e o total é $\Theta(n^2)$. Com um \textbf{heap} de máximo, achar e remover custa $O(\log n)$, e inserir também: $O(n\log n)$ no total.

\resp{rv-kesimo} Cada consulta custa $\Theta(n\log n)$. Se $q$ é pequeno, use \textbf{seleção linear} (quickselect ou mediana das medianas) sobre a cópia: $O(n)$ por consulta. Se $q$ é grande, \textbf{ordene uma única vez}, em $O(n\log n)$, e responda cada consulta em $O(1)$: total $O(n\log n+q)$.

\resp{rv-grupos} Cada busca em largura custa $O(n+m)$, o que é inviável para $10^8$ perguntas. Use \textbf{Union-Find}: cada amizade é uma \Union{} e cada pergunta compara \Find$(u)$ com \Find$(v)$. Com rank e compressão, o custo é praticamente $O(1)$ por operação; com listas e ponteiro ao representante, \Find{} é $O(1)$ no pior caso.

\resp{rv-dic} Inserir ou remover num vetor ordenado custa $\Theta(n)$. Use uma \textbf{AVL} (ou rubro-negra): inserção e remoção em $O(\log n)$, e a consulta por intervalo em $O(\log n+k)$, sendo $k$ o número de códigos no intervalo (desça até $a$ e percorra em ordem até passar de $b$). Uma tabela hash não serve, pois não mantém a ordem.

\resp{rv-idades} Uma ordenação por comparação custa $\Theta(n\log n)$, cerca de $2{,}7\cdot10^9$ comparações. As chaves são inteiros num intervalo pequeno ($k=121$): \textbf{counting sort}, em $\Theta(n+k)$, com 121 contadores.

\resp{rv-filas} No heap binário a fusão custa $\Theta(n)$. Use o \textbf{heap binomial}: fusão em $O(\log n)$ no pior caso de \emph{cada} operação. O heap de Fibonacci faria a fusão em $O(1)$, mas só amortizado --- uma remoção pode consolidar tudo de uma vez ---, o que fere a previsibilidade exigida.

\resp{rv-emerg} \textbf{Heap binário em vetor}: inserção e remoção do mais grave em $O(\log n)$, sem fusões, com as raras alterações de gravidade em $O(\log n)$ (com um vetor de posições), e sem gastar memória com ponteiros.

\resp{rv-cancel} \textbf{Heap} com um vetor (ou tabela) de posições indexado pelo número do evento, para encontrar o evento cancelado e removê-lo de uma posição qualquer em $O(\log n)$ (Exercício~\ref{ex:hp-remi}). Uma \textbf{AVL} indexada por (horário, número) também resolve, com todas as operações em $O(\log n)$.

\resp{rv-palavras} \textbf{Tabela hash}: cada verificação custa $O(1)$ esperado (mais o custo de calcular o hash da palavra). Um vetor ordenado com busca binária também serviria, em $O(\log n)$ por consulta, já que o dicionário é fixo.

\resp{rv-posicao} Uma \textbf{AVL aumentada}: cada nó guarda o tamanho da sua subárvore, atualizado nas inserções, remoções e rotações em $O(1)$ por nó afetado. A posição de $x$ se calcula descendo da raiz e somando os tamanhos das subárvores à esquerda do caminho; o elemento da posição $r$ se encontra descendo e comparando $r$ com esses tamanhos. Tudo em $O(\log n)$.

\resp{rv-mediana} Dois heaps: um de máximo com a metade menor das medições e um de mínimo com a metade maior, com tamanhos que diferem de no máximo 1. Cada medição é inserida em $O(\log n)$, com um rebalanceamento eventual, e a mediana é lida no topo de um dos heaps em $O(1)$.

\resp{rv-externo} \textbf{Mergesort externo}: leia o arquivo em blocos que cabem na memória (por exemplo, de 12 GB), ordene cada bloco e grave-o; depois intercale os cerca de 42 blocos ordenados com um heap de mínimo de tamanho 42 (intercalação de $k$ listas), lendo e gravando sequencialmente. O custo é dominado pelas duas passadas pelo disco.

\resp{rv-quase} \textbf{Ordenação por inserção}: cada elemento se desloca no máximo 10 posições, então o custo é $O(10n)=O(n)$. Um heap de mínimo de tamanho 11 percorrendo o vetor também ordena em $O(n\log 11)$.

\resp{rv-dup} \textbf{Tabela hash}: insira as strings uma a uma e verifique se já estão presentes, em $O(1)$ esperado cada, mais o custo de calcular os hashes. Ordenar (por exemplo, com radix sort de strings ou mergesort) e comparar vizinhas também resolve.

\resp{rv-remocoes} \textbf{Union-Find offline, de trás para frente}: comece com a malha sem as estradas que serão interditadas em algum momento, una as regiões ligadas pelas estradas restantes e conte as componentes; depois percorra as interdições da última para a primeira, transformando cada uma numa \Union{} e atualizando o contador. Guarde as respostas e apresente-as na ordem original.

\resp{rv-topk} Um \textbf{heap de mínimo de tamanho 100}: para cada requisição, se ela for mais demorada que o mínimo do heap, substitua o mínimo. Memória $O(100)$ e tempo $O(n\log100)$.

\resp{rv-reais} O bucket sort perde a garantia sem distribuição uniforme: pode cair no pior caso, com quase tudo num balde. Use uma ordenação por comparação com garantia, como o \textbf{mergesort} ou o \textbf{heapsort}, em $\Theta(n\log n)$; ou converta os reais para inteiros de 64 bits preservando a ordem e use radix sort.

\resp{rv-prefixo} Cada pergunta custa $O(n)$. Pré-calcule as \textbf{somas de prefixo} $P[i]=V[1]+\cdots+V[i]$ em $O(n)$; cada pergunta passa a custar $O(1)$: $P[j]-P[i-1]$.

\resp{rv-busca} Uma \textbf{AVL} indexada pela chave: busca, inserção e remoção em $O(\log n)$ e listagem em ordem em $O(n)$. Uma tabela hash daria buscas em $O(1)$ esperado, mas a listagem em ordem exigiria ordenar a cada vez.

\resp{rv-desfazer} \textbf{Union-Find sem compressão}, com união por rank ou tamanho e uma pilha das alterações: cada ligação empilha o ponteiro e o rank alterados; desfazer é desempilhar e restaurar em $O(1)$; a consulta custa $O(\log n)$.

\respcap{cap:pd}
\resp{pd-max} \textbf{Sim.} Se a última multiplicação de uma parentização máxima de $M_i\cdots M_j$ é feita no ponto $k$, o custo é $c(i..k)+c(k+1..j)+r_{i-1}r_kr_j$, e o último termo depende só de $k$, não de como cada parte foi parentizada. Se uma das partes não fosse máxima, trocá-la por uma máxima aumentaria o total, uma contradição. A recorrência é a mesma, com $\max$ no lugar de $\min$, e o custo é $\Theta(n^3)$.

\resp{pd-lis} Seja $L[i]$ o tamanho da maior subsequência crescente que \emph{termina em} $a_i$: $L[i]=1+\max\{L[j]:j<i,\ a_j<a_i\}$, com $\max\emptyset=0$; a resposta é $\max_iL[i]$, e guardar o $j$ escolhido em $ant[i]$ permite reconstruí-la. Custo $O(n^2)$. Por exemplo, $5,2,8,6,3,6,9,7$ dá $L=1,1,2,2,2,3,4,4$ e resposta 4 ($2,3,6,9$). Em $O(n\log n)$: mantenha $d[\ell]=$ menor final possível de uma subsequência crescente de tamanho $\ell$, que é um vetor crescente; para cada $a_i$, encontre por busca binária o primeiro $d[\ell]\ge a_i$ e substitua-o por $a_i$ (ou acrescente $a_i$ no fim, se não houver). O tamanho final de $d$ é a resposta.

\resp{pd-pal} Seja $P(i,j)$ o tamanho do maior palíndromo que é subsequência de $s_i..s_j$: $P(i,i)=1$ e $P(i,i-1)=0$; se $s_i=s_j$, $P(i,j)=2+P(i+1,j-1)$; senão, $P(i,j)=\max\{P(i+1,j),P(i,j-1)\}$. Preencha por tamanho crescente do trecho: $\Theta(n^2)$ de tempo e memória. Para ``character'', o valor é 5 (``carac''). Esse valor também é o tamanho da subsequência comum máxima entre $s$ e $s$ invertida.

\resp{pd-dag} Ordene os vértices topologicamente (Capítulo~\ref{cap:dfsdig}). Faça $D(s)=0$ e $D(v)=-\infty$ para os demais; na ordem topológica, $D(v)=\max_{uv\in E}\{D(u)+w(u,v)\}$, guardando o predecessor que deu o máximo. A resposta é $D(t)$, em $\Theta(n+m)$. Num grafo acíclico todo caminho é simples, e por isso a subestrutura ótima vale; em grafos gerais, o caminho simples mais longo é um problema NP-difícil.

\resp{pd-ilim} Para $w=0..W$: \[K(w)=\max\big\{0,\ \max_{i:\,p_i\le w}\{K(w-p_i)+v_i\}\big\}.\] $\Theta(nW)$. Por exemplo, com itens $(p,v)=(3,5),(4,7),(2,3)$ e $W=10$, o máximo é 17 (por exemplo, $4+4+2$). É o mesmo que a mochila 0-1 em vetor único, com $w$ \emph{crescente}.

\resp{pd-corte} Ordene os pontos e acrescente as sentinelas $L_0=0$ e $L_{m+1}=n$. Seja $C(i,j)$ o custo mínimo de fazer todos os cortes estritamente entre $L_i$ e $L_j$: $C(i,i+1)=0$ e $C(i,j)=(L_j-L_i)+\min_{i<k<j}\{C(i,k)+C(k,j)\}$, em que $L_k$ é o primeiro corte feito no trecho. A resposta é $C(0,m+1)$; guardando o $k$ ótimo, a sequência de cortes sai em pré-ordem. Custo $\Theta(m^3)$. No exemplo, o ótimo é \textbf{38}: corte primeiro em 10 e depois em 2 e 8, em qualquer ordem (as seis ordens custam $50,46,40,40,38,38$).

\resp{pd-rec} Percorra a tabela a partir de $(n,m)$: se $x_i=y_j$, coloque $x_i$ no início da resposta e vá para $(i-1,j-1)$; senão, vá para a vizinha, $(i-1,j)$ ou $(i,j-1)$, de maior valor. Custo $O(n+m)$. No exemplo, subindo nos empates, obtém-se \textbf{BCBA}; as subsequências comuns de tamanho 4 são BCAB, BCBA e BDAB (§\ref{sec:lcs}).

\resp{pd-mon} (a) Para $r=(10,100,5,50)$: $m[1,2]=5000$, $m[2,3]=25\,000$ e $m[1,3]=\min\{75\,000,\,7500\}=7500$, com a ordem $((M_1M_2)M_3)$. Para $r=(20,50,10,1,100)$: $m[1,4]=3500$, com a ordem $((M_1(M_2M_3))M_4)$; a tabela completa está na §\ref{sec:cadeia}. (b) A última linha é $M(4,\cdot)=0,2,9,9$; o ótimo é 9, com os itens 1 e 2. (c) O tempo depende do \emph{valor} de $W$, mas $W$ ocupa só $\log_2W$ bits: o algoritmo é exponencial no tamanho da entrada.

\resp{pd-apa} (a) Tamanho 4: \textrm{acda} ou \textrm{dcda}. (b) 240, com os itens 1, 2 e 4 (pesos $20+30+50=100$). (c) Algoritmo de Kadane: $S(i)=\max\{a_i,\ S(i-1)+a_i\}$ é a melhor soma de um subvetor que \emph{termina} em $i$, e a resposta é $\max_iS(i)$, em $O(n)$. No exemplo, 55, com o subvetor $10,-5,40,10$.

\resp{pd-l6a} (a) Custo 2010, com $((A_1A_2)((A_3A_4)(A_5A_6)))$. (b) Guarde em $s[i,j]$ o $k$ que deu o mínimo e use \textsc{Multiplica}$(A,s,i,j)$ da §\ref{sec:cadeia}. (c) Na reconstrução, quando $x_i\ne y_j$ e $C[i-1,j]=C[i,j-1]$, siga os dois caminhos recursivamente e elimine as repetições; pode haver exponencialmente muitas respostas.

\resp{pd-l6b} (a) O número de Catalan $C_6=\binom{12}6/7=132$. (b) $c(i,j)=\sum_{k=i}^ja_k+\min_{i\le r\le j}\{c(i,r-1)+c(r+1,j)\}$, com $c(i,i-1)=0$, em $\Theta(n^3)$. Resultado: custo 185, com raiz \emph{if}; à esquerda, \emph{case}, com o caminho \emph{else}$\to$\emph{end} à direita; à direita da raiz, \emph{then}, com \emph{of} como filho esquerdo.

\resp{pd-l6c} (a) $S[i]=a_i+\max\{S[j]:j<i,\ a_j<a_i\}$ (0 se não houver tal $j$), e a resposta é $\max_iS[i]$; por exemplo, $1,101,2,3,100,4,5$ dá 106 ($1,2,3,100$).\\ (b) Em pós-ordem: $com(v)=s(v)+\sum_fsem(f)$ e $sem(v)=\sum_f\max\{com(f),sem(f)\}$, somando sobre os filhos $f$; a resposta é $\max\{com,sem\}$ na raiz, em $\Theta(n)$.\\ (c) É a subsequência comum máxima das duas sequências.\\ (d) $f(1)=f(2)=2$ e $f(n)=f(n-1)+f(n-2)$. A última faixa é branca ou vermelha; se a penúltima não é azul, a sequência vem de uma de tamanho $n-1$, estendida de um único jeito; se é azul, vem de uma de tamanho $n-2$, estendida pela faixa azul e pela cor que falta. Valores: $2,2,4,6,10,16,26,\dots$

\resp{pd-mc458} (a) Capacidades $0..14$:\\
\hspace*{1em}item 3 (peso 5): \texttt{O - O - - I - O - O - - I - I}\\
\hspace*{1em}item 4 (peso 2): \texttt{O - O - I O - O - O - I O - O}\\
(b) Um único vetor $K[0..W]$: para cada item, percorra $w$ \emph{decrescendo} de $W$ até $p_i$, fazendo $K[w]:=\max\{K[w],K[w-p_i]+v_i\}$. A memória cai para $O(K)$; com $w$ crescendo, o item poderia ser reaproveitado. (c) Fixe o par de linhas $(i_1,i_2)$, some cada coluna entre elas e aplique Kadane ao vetor de somas: $O(m^2n)$, usando somas de prefixo por coluna. (d) $P(a,b)=\max\{p_{ab},\ \max_{a'<a}P(a',b)+P(a-a',b),\ \max_{b'<b}P(a,b')+P(a,b-b')\}$ para todos os $a\le n$ e $b\le m$: $O(nm(n+m))\subseteq O(n^2m^2)$.

\resp{pd-edit} $D(i,0)=i$, $D(0,j)=j$ e $D(i,j)=\min\{D(i-1,j)+1,\ D(i,j-1)+1,\ D(i-1,j-1)+[x_i\ne y_j]\}$, em $\Theta(nm)$. De ``kitten'' para ``sitting'' a distância é \textbf{3}: trocar k por s, trocar e por i e inserir g.

\resp{pd-troco} São 2 moedas ($3+3$). O guloso pega $4+1+1$, com 3 moedas. Programação dinâmica: $T(0)=0$ e $T(s)=1+\min_{c\le s}T(s-c)$; os valores de $T(0)$ a $T(10)$ são $0,1,2,1,1,2,2,2,2,3,3$. Custo $\Theta(nS)$.

\resp{pd-haste} $r(0)=0$ e $r(n)=\max_{1\le i\le n}\{p_i+r(n-i)\}$. Os valores de $r(1)$ a $r(8)$ são $1,5,8,10,13,17,18,22$: o máximo para 8 é \textbf{22}, com os pedaços $2+6$.

\resp{pd-grade} São 6 passos, 3 para a direita e 3 para baixo, em qualquer ordem: $\binom63=20$ caminhos. Programação dinâmica: $P(i,j)=P(i-1,j)+P(i,j-1)$, com $P(1,1)=1$ e $P=0$ nas casas bloqueadas. Com a casa $(2,2)$ bloqueada, sobram 8 caminhos.

\resp{pd-part} Se a soma $S$ é ímpar, não é possível. Senão, decida se algum subconjunto soma $S/2$: $R(i,s)=R(i-1,s)\lor R(i-1,s-a_i)$, em $\Theta(nS)$, pseudopolinomial. $\{3,1,1,2,2,1\}$: sim ($3+2=5$ de um lado). $\{2,3,7\}$: a soma é 12, mas nenhum subconjunto soma 6: não.

\resp{pd-catalan} A última multiplicação separa os $k$ primeiros fatores dos $n-k$ últimos, e cada lado pode ser parentizado de forma independente: $P(n)=\sum_{k=1}^{n-1}P(k)P(n-k)$. Os valores para $n=1,\dots,6$ são $1,1,2,5,14,42$ --- os números de Catalan $C_{n-1}$.

\resp{pd-lcsmem} Troque $X$ e $Y$ se for preciso, para que $Y$ seja a menor. Mantenha dois vetores, $ant[0..m]$ e $atual[0..m]$, com as linhas $i-1$ e $i$ da tabela: $atual[0]:=0$ e, para $j=1..m$, $atual[j]:=1+ant[j-1]$ se $x_i=y_j$ e $\max\{ant[j],atual[j-1]\}$ caso contrário; ao fim de cada linha, troque os papéis dos vetores. A resposta é o último valor calculado. Memória $O(\min\{n,m\})$ e tempo $\Theta(nm)$.

\resp{pd-formas} Seja $F(s)$ o número de maneiras; processando as moedas uma de cada vez (para não contar ordens diferentes como maneiras diferentes): para cada moeda $c$ e para $s=c,\dots,S$, faça $F(s):=F(s)+F(s-c)$, com $F(0)=1$. Para 5 há \textbf{4} maneiras: $5$; $2+2+1$; $2+1+1+1$; $1+1+1+1+1$. Para 10 há 10 maneiras.

\resp{pd-substring} Seja $L(i,j)$ o tamanho do maior substring comum que \emph{termina} em $x_i$ e em $y_j$: $L(i,j)=L(i-1,j-1)+1$ se $x_i=y_j$ e $0$ caso contrário. A resposta é o maior $L(i,j)$, em $O(nm)$. Para ``ABABC'' e ``BABCA'', o valor é 4, com o substring ``BABC''.

\respcap{cap:fw}
\resp{fw-sub} $D(i,j,S)=\min\big\{W(i,j),\ \min_{v\in S}\{D(i,v,S\setminus\{v\})+D(v,j,S\setminus\{v\})\}\big\}$. Há $n^2\cdot2^n$ subproblemas, um por par e por subconjunto: exponencial. O Floyd-Warshall fixa uma ordem dos vértices e só usa os $n+1$ prefixos $\{1,\dots,k\}$; em cada passo, decide-se apenas se o vértice $k$ entra ou não como intermediário. Isso dá $O(n^2)$ subproblemas por valor de $k$ e $\Theta(n^3)$ no total.

\resp{fw-slide} Ver o Exemplo da §\ref{sec:fw}: as linhas de $D^5$ são $(0,1,2,3,5)$, $(2,0,1,2,4)$, $(1,2,0,1,3)$, $(3,4,5,0,2)$ e $(4,5,6,1,0)$. De 5 a 3: $P(5,3)=2$, $P(5,2)=1$, $P(5,1)=4$ e $P(5,4)=5$, logo $5\to4\to1\to2\to3$, com custo 6. De 2 a 1: $P(2,1)=3$ e $P(2,3)=2$, logo $2\to3\to1$, com custo 2.

\resp{fw-neg} Linhas da matriz final, de 1 a 5: $(0,1,-3,2,-4)$, $(3,0,-4,1,-1)$, $(7,4,0,5,3)$, $(2,-1,-5,0,-2)$ e $(8,5,1,6,0)$. A diagonal é toda nula: não há ciclo negativo. Por exemplo, $D(1,2)=1$ pelo caminho $1\to5\to4\to3\to2$, de custo $-4+6-5+4$.

\resp{fw-fecho} $T^0(i,j)$ é verdadeiro se $i=j$ ou se $ij\in E$, e $T^k(i,j)=T^{k-1}(i,j)\lor\big(T^{k-1}(i,k)\land T^{k-1}(k,j)\big)$, com $\Theta(n^3)$ operações lógicas. No exemplo, os vértices $1$, $2$ e $3$ alcançam $\{1,2,3,4\}$; o $4$ alcança apenas a si mesmo, e o $5$ também.

\resp{fw-ciclo} Inicialize $D(i,i)=\infty$, em vez de 0, e execute o algoritmo: ao final, $D(i,i)$ é o peso do menor ciclo que passa por $i$, e a resposta é $\min_iD(i,i)$. No exemplo, $D(1,1)=D(2,2)=D(3,3)=6$ (o ciclo $1\to2\to3\to1$) e $D(4,4)=7$ (o ciclo $4\to3\to4$): a resposta é 6.

\resp{fw-negciclo} Se não há ciclo negativo, todo passeio fechado tem peso maior ou igual a zero; $D(i,i)$ começa em 0 e só poderia diminuir por meio de um passeio fechado de peso negativo, logo continua 0. Se há um ciclo negativo, seja $k$ o vértice de maior índice do ciclo e $i$ outro vértice dele. Os dois trechos do ciclo, de $i$ a $k$ e de $k$ a $i$, só têm intermediários de índice menor que $k$; por isso, depois da iteração $k-1$, $D(i,k)$ e $D(k,i)$ são no máximo os pesos desses trechos. Na iteração $k$, então, $D(i,i)\le D(i,k)+D(k,i)\le$ peso do ciclo $<0$.

\resp{fw-dij} Floyd-Warshall: $n^3=10^{12}$ operações. $n$ execuções do Dijkstra com heap: $n\cdot(n+m)\log n\approx10^4\cdot6\cdot10^4\cdot14\approx8\cdot10^9$, bem melhor nesse grafo esparso. Com $m\approx n^2/2$, as $n$ execuções do Dijkstra custam cerca de $n\cdot m\log n\approx n^3\log n/2$, mais que o Floyd-Warshall; com o Dijkstra em vetor ($O(n^2)$ por execução), o total também é $\Theta(n^3)$, e o Floyd-Warshall é mais simples e aceita pesos negativos.

\resp{fw-4v} $D^1$: linhas $(0,4,\infty,10)$, $(\infty,0,3,\infty)$, $(7,11,0,1)$ e $(2,6,\infty,0)$. $D^2$: $(0,4,7,10)$, $(\infty,0,3,\infty)$, $(7,11,0,1)$ e $(2,6,9,0)$. $D^3$: $(0,4,7,8)$, $(10,0,3,4)$, $(7,11,0,1)$ e $(2,6,9,0)$. $D^4$: $(0,4,7,8)$, $(6,0,3,4)$, $(3,7,0,1)$ e $(2,6,9,0)$.

\resp{fw-4vcam} $P^4$ tem as linhas $(-,1,2,3)$, $(4,-,2,3)$, $(4,1,-,3)$ e $(4,1,2,-)$. De 2 a 1: $P(2,1)=4$, $P(2,4)=3$ e $P(2,3)=2$, logo $2\to3\to4\to1$, com custo $3+1+2=6$. De 4 a 3: $P(4,3)=2$, $P(4,2)=1$ e $P(4,1)=4$, logo $4\to1\to2\to3$, com custo $2+4+3=9$.

\resp{fw-negex} Com as arestas $1\to2$ (1) e $2\to1$ ($-3$), o ciclo $1\to2\to1$ pesa $-2$. Na iteração $k=1$, $D(2,2)$ passa a $D(2,1)+D(1,2)=-2$; na iteração $k=2$, $D(1,1)$ passa a $D(1,2)+D(2,1)=-2$. Os valores negativos na diagonal sinalizam o ciclo. Os demais valores passam a corresponder a \emph{passeios} que repetem vértices e não são mais custos de caminhos mínimos, que nem estão definidos, pois percorrer o ciclo mais vezes diminuiria o custo indefinidamente.

\resp{fw-nd} Troque cada aresta não direcionada por duas arestas direcionadas de mesmo peso; a matriz $W$ fica simétrica, e $D$ também. Uma aresta não direcionada de peso negativo $uv$ é, ela mesma, um ciclo negativo $u\to v\to u$; por isso, com o Floyd-Warshall, pesos negativos só são aceitáveis em digrafos.

\resp{fw-conta} O teste é executado exatamente $n^3$ vezes (três laços de 1 a $n$). A matriz $D$ ocupa $n^2$ posições; com a matriz $P$, o total é $2n^2$. Calcular $D^0,\dots,D^n$ separadamente usaria $\Theta(n^3)$ de memória, o que a atualização no lugar evita.

\resp{fw-inplace} Na iteração $k$, os valores usados são $D(i,k)$ e $D(k,j)$. Eles não mudam durante a iteração: $D^k(i,k)=\min\{D^{k-1}(i,k),\ D^{k-1}(i,k)+D^{k-1}(k,k)\}=D^{k-1}(i,k)$, pois $D(k,k)=0$ sem ciclos negativos, e o mesmo vale para $D(k,j)$. Assim, não importa se uma entrada da linha ou da coluna $k$ já foi atualizada; e as demais entradas só podem diminuir para valores que correspondem a caminhos reais, nunca abaixo de $D^k$.

\resp{fw-largura} Troque $\min$ por $\max$ e $+$ por $\min$: $C^k(i,j)=\max\{C^{k-1}(i,j),\ \min\{C^{k-1}(i,k),\,C^{k-1}(k,j)\}\}$, com $C^0(i,j)$ igual à capacidade da aresta $ij$ (0 se ela não existe e $\infty$ para $i=j$). O argumento de correção é o mesmo, e o custo é $\Theta(n^3)$.

\resp{fw-diam} Excentricidades (o maior valor de cada linha de $D^5$): 5, 4, 3, 5 e 6 para os vértices 1 a 5. O diâmetro é 6 (a distância de 5 a 3), e o centro é o vértice 3, de excentricidade 3.

\resp{fw-aresta} A aresta $uv$ está em algum caminho mínimo de $i$ a $j$ se e somente se $D(i,u)+W(u,v)+D(v,j)=D(i,j)$. Ela está em algum caminho mínimo do grafo se e somente se $W(u,v)=D(u,v)$, pois então ela própria é um caminho mínimo de $u$ a $v$.

\resp{fw-atualiza} Os novos caminhos mínimos ou não usam a aresta $uv$, e já estão em $D$, ou a usam uma única vez (com pesos que não criam ciclo negativo). Logo, para todo par, $D'(i,j)=\min\{D(i,j),\ D(i,u)+w'+D(v,j)\}$: $O(n^2)$.

\resp{fw-novo} Calcule primeiro as distâncias de e para o novo vértice $n+1$: $D(i,n+1)=\min_u\{D(i,u)+W(u,n+1)\}$ e $D(n+1,j)=\min_u\{W(n+1,u)+D(u,j)\}$, em $O(n^2)$ no total. Depois, para todo par $(i,j)$, faça $D(i,j):=\min\{D(i,j),\ D(i,n+1)+D(n+1,j)\}$, em $O(n^2)$. É exatamente uma iteração do Floyd-Warshall com $k=n+1$.

\resp{fw-contar} Mantenha $C(i,j)$, com $C(i,i)=1$ e $C(i,j)=1$ se há aresta $ij$ (0 se não há). Na iteração $k$, para $i\ne k\ne j$: se $D(i,k)+D(k,j)<D(i,j)$, faça $D(i,j):=D(i,k)+D(k,j)$ e $C(i,j):=C(i,k)\cdot C(k,j)$; se $D(i,k)+D(k,j)=D(i,j)$, faça $C(i,j):=C(i,j)+C(i,k)\cdot C(k,j)$. Com pesos positivos, cada caminho mínimo é contado exatamente uma vez: na iteração do seu vértice intermediário de maior índice.

\resp{fw-bits} Represente a linha $i$ da matriz como um vetor de $n$ bits, ocupando $\ceil{n/w}$ palavras. A atualização $T(i,\cdot):=T(i,\cdot)\lor T(k,\cdot)$, feita sempre que $T(i,k)$ é verdadeiro, custa $O(n/w)$ operações sobre palavras. Para cada $k$ e cada $i$, isso dá $O(n^2\cdot n/w)=O(n^3/w)$ no total.

\respcap{cap:gulosos}
\resp{gm-troco} Ver a §\ref{sec:guloso}: em ordem decrescente, use o máximo possível de cada moeda. Para $\{10,5,2,1\}$, uma solução ótima tem no máximo uma moeda de 1, uma de 5 e duas de 2, e nunca $2+2+1$; então as moedas pequenas somam no máximo 9, e o número de moedas de 10 é $\floor{S/10}$, como no guloso; o mesmo argumento fixa as demais. Contraexemplo geral: $\{4,3,1\}$ com $S=6$ (o guloso dá $4+1+1$; o ótimo é $3+3$).

\resp{gm-i3} \emph{Linear} e \emph{gráfica}: ver a §\ref{sec:matroides} (dimensão do espaço gerado; número de componentes de uma floresta). \emph{Transversal}: sejam $I_1,I_2\subseteq X$ emparelháveis, com $|I_1|<|I_2|$, por emparelhamentos $M_1$ e $M_2$ que cobrem, do lado $X$, exatamente $I_1$ e $I_2$. A diferença simétrica $M_1\triangle M_2$ é uma união de caminhos e ciclos alternantes; como $|M_2|>|M_1|$, algum caminho tem mais arestas de $M_2$, e começa e termina em vértices cobertos só por $M_2$. Esse caminho tem um número ímpar de arestas e alterna os lados, então um dos extremos está em $X$; ele pertence a $I_2\setminus I_1$. Trocando as arestas do caminho em $M_1$, obtém-se um emparelhamento que cobre $I_1$ e mais esse vértice. \emph{Emparelhamento}: o mesmo argumento, com caminhos alternantes num grafo qualquer: o caminho com uma aresta de $M_2$ a mais acrescenta à cobertura um vértice de $I_2$ sem descobrir os de $I_1$.

\resp{gm-g} Ver a §\ref{sec:matroides}. Suponha $|I_1|=k<|I_2|$ e que nenhum $e\in I_2\setminus I_1$ aumente $I_1$. Com pesos $k+2$ em $I_1$, $k+1$ em $I_2\setminus I_1$ e 0 nos demais, o guloso escolhe $I_1$ primeiro e depois só elementos de peso 0, totalizando $k(k+2)$; mas $I_2$ pesa pelo menos $(k+1)^2>k(k+2)$, uma contradição.

\resp{gm-pratica} (a) Não diretamente: os conjuntos de intervalos dois a dois disjuntos não formam uma matroide (com $I_1=\{[1,10]\}$ e $I_2=\{[1,2],[3,4]\}$, (I3) falha). Mesmo assim, o guloso ``o que termina primeiro'' é ótimo para pesos unitários, por um argumento de troca próprio (Exercício~\ref{ex:gm-ativ}); com pesos arbitrários, é preciso programação dinâmica (Exercício~\ref{ex:gm-pesado}). (b) Escalonamento de Gauss na matriz $m\times k$ formada pelas colunas de $S$: elas são linearmente independentes se e somente se o posto é $k$. O custo é $O(m\,k\min\{m,k\})$; se $k>m$, elas já são dependentes.

\resp{gm-pesq} A §\ref{sec:matroides} e o quadro ``Para saber mais'' reúnem os pontos centrais: os axiomas; o teorema do guloso, provado pela troca de bases; o algoritmo de Kruskal na matroide gráfica; a interseção de duas matroides por caminhos aumentantes no grafo de troca; o emparelhamento bipartido e a floresta colorida como interseções de duas matroides; e o caminho hamiltoniano como interseção de três. Para aprofundar, consulte as obras de Oxley e de Schrijver listadas nas Referências.

\resp{gm-ativ} Em ordem de término, o guloso escolhe as atividades $1,4,8,11$, isto é, os intervalos $[1,4)$, $[5,7)$, $[8,11)$ e $[12,16)$: quatro atividades, que é o máximo possível. Prova pelo argumento de troca: seja $O$ uma solução ótima e $g$ a atividade que termina primeiro. Se $g\notin O$, troque a primeira atividade de $O$ por $g$; como $g$ termina antes, continua compatível com as demais, e $|O|$ não muda. Logo existe uma solução ótima que contém $g$, e o restante é o mesmo problema restrito às atividades que começam depois do fim de $g$; por indução, o guloso é ótimo.

\resp{gm-ruins} Menor duração: com $[1,5)$, $[4,7)$ e $[6,10)$, o guloso escolhe $[4,7)$, de duração 3, e perde as outras duas, que são compatíveis entre si (o ótimo é 2). Começa primeiro: com $[0,10)$, $[1,2)$ e $[3,4)$, o guloso escolhe $[0,10)$ e fica com uma atividade; o ótimo é 2.

\resp{gm-frac} Ordene os itens por $v_i/p_i$ decrescente e leve o máximo possível de cada um. Pela troca: se uma solução ótima leva menos do item $i$ de maior razão e algo de um item $j$ de razão menor, trocar uma quantidade $\varepsilon$ de peso de $j$ por $i$ não diminui o valor; repetindo, chega-se à solução gulosa. No exemplo: os itens 1 e 2 inteiros e $20/30$ do item 3, valendo $60+100+80=$ \textbf{240}. Na mochila 0-1, o ótimo é 220 (itens 2 e 3), e o guloso pela razão dá só 160.

\resp{gm-unif} (I1) e (I2) são imediatos: as condições $|I|\le k$ e $|I\cap E_i|\le k_i$ valem para $\emptyset$ e passam para os subconjuntos. (I3) na uniforme: se $|I_1|<|I_2|\le k$, qualquer $e\in I_2\setminus I_1$ serve, pois $|I_1\cup\{e\}|\le k$. (I3) na de partição: como $|I_1|<|I_2|$, há uma parte $E_i$ com $|I_2\cap E_i|>|I_1\cap E_i|$; tome $e\in(I_2\setminus I_1)\cap E_i$. Então $|(I_1\cup\{e\})\cap E_i|\le|I_2\cap E_i|\le k_i$, e as outras partes não mudam.

\resp{gm-tarefas} Minimizar a multa das tarefas atrasadas é o mesmo que maximizar o peso das que cumprem o prazo, isto é, encontrar uma base de peso máximo da matroide. Em ordem decrescente de multa, o guloso aceita as tarefas 1, 2, 3 e 4, rejeita a 5 e a 6 (cada uma tornaria o conjunto impossível de escalonar) e aceita a 7. A multa total é $w_5+w_6=30+20=$ \textbf{50}. Uma ordem válida de execução: as tarefas aceitas em ordem crescente de prazo ($2,4,1,3,7$), seguidas das atrasadas.

\resp{gm-huffman} As junções são: $f+e=14$; $c+b=25$; $14+d=30$; $25+30=55$; $a+55=100$. Um código possível: $a=0$, $c=100$, $b=101$, $f=1100$, $e=1101$, $d=111$. O custo é $45\cdot1+13\cdot3+12\cdot3+16\cdot3+9\cdot4+5\cdot4=$ \textbf{224}.

\resp{gm-spt} Se, numa ordem, uma tarefa $i$ vem imediatamente antes de uma tarefa $j$ com $t_i>t_j$, trocar as duas não altera os términos das demais, e a soma dos términos diminui de $t_i-t_j$. Logo, numa ordem ótima, não há tal par, e as durações aparecem em ordem crescente. Para as durações $3,1,2$, a ordem $1,2,3$ termina as tarefas nos instantes $1,3,6$, com soma \textbf{10}; a pior ordem, $3,2,1$, soma 14.

\resp{gm-naomat} Não é matroide: com $I_1=\{3\}$ e $I_2=\{1,2\}$, nem $\{1,3\}$ nem $\{2,3\}$ são independentes, e (I3) falha. Com $c(3)=3$ e $c(1)=c(2)=2$, o guloso escolhe 3 e não consegue acrescentar mais nada, ficando com peso 3, enquanto $\{1,2\}$ pesa 4.

\resp{gm-paralelos} É uma matroide: é a matroide linear das colunas $1=(1,0)$, $2=(1,0)$, $3=(0,1)$ e $4=(1,1)$ de $\R^2$. Os pares independentes são todos, exceto $\{1,2\}$, cujas colunas são iguais (elementos ``paralelos''), e nenhum trio é independente em $\R^2$.

\resp{gm-posto} Uma base é uma floresta geradora maximal, isto é, uma árvore geradora em cada componente; uma componente com $n_i$ vértices contribui com $n_i-1$ arestas, e o total é $\sum_i(n_i-1)=n-c$. No exemplo, $n=6$ e $c=2$: o posto é 4 (duas arestas do triângulo e as duas do caminho).

\resp{gm-dual} As bases de $U_{n,k}$ são todos os subconjuntos de tamanho $k$; os seus complementos são todos os subconjuntos de tamanho $n-k$, que são exatamente as bases de $U_{n,n-k}$.

\resp{gm-unifgul} Na uniforme $U_{n,k}$, o guloso escolhe os $k$ elementos mais pesados. Na de partição, escolhe os $k_i$ mais pesados de cada parte $E_i$. No exemplo: $9$ e $7$ de $A$, $8$ e $3$ de $B$, $6$ e $5$ de $C$, com peso total 38.

\resp{gm-moedas} Com $S=14$, o guloso dá $10+1+1+1+1$ (5 moedas), e o ótimo é $7+7$ (2 moedas). O sistema $\{1,5,10,25\}$ não tem esse problema: o guloso é ótimo para todo valor, o que se prova com um argumento como o do exemplo da §\ref{sec:guloso}.

\resp{gm-pesado} Com $A=[0,3)$ de peso 1 e $B=[1,5)$ de peso 10, o guloso escolhe $A$, que termina primeiro, e não pode escolher $B$, ficando com peso 1; o ótimo é 10. O problema se resolve por programação dinâmica: ordene os intervalos por término, seja $p(j)$ o último intervalo compatível com o $j$-ésimo e $OPT(j)=\max\{OPT(j-1),\ w_j+OPT(p(j))\}$, em $O(n\log n)$.

\resp{gm-fracn} Encontre, por seleção em $O(n)$, a razão mediana $r$ dos itens e some os pesos dos itens com razão maior que $r$. Se essa soma é no máximo a capacidade restante, leve todos esses itens inteiros e continue com os itens de razão no máximo $r$ e a capacidade que sobra; senão, descarte os itens de razão no máximo $r$ e continue com os de razão maior. Cada passo descarta metade dos itens: $T(n)=T(n/2)+O(n)=O(n)$. No fim, sobra no máximo um item, levado fracionado.

\respcap{cap:agm}
\resp{agm-l1} Por indução, o conjunto $A$ das arestas aceitas está contido em alguma AGM. Ao aceitar $e=uv$, tome como $S$ a árvore atual de $u$ na floresta. Nenhuma aresta de $A$ atravessa o corte $(S,V\setminus S)$, e $e$ é a mais leve que o atravessa: uma aresta mais leve que o atravessasse teria sido examinada antes, quando os seus extremos já estavam em árvores diferentes (as árvores só crescem), e teria sido aceita. Pelo Lema~\ref{lema:corte}, $A\cup\{e\}$ está contido em alguma AGM. Ao final, a saída é acíclica por construção e conexa, pois uma aresta entre duas componentes finais teria sido aceita; logo é uma árvore geradora, e é mínima. (Também segue do Teorema~\ref{teo:gulosomat} aplicado à matroide gráfica.)

\resp{agm-l2} Invariante: as arestas de $T$ estão contidas em alguma AGM. Vale no início, com $T$ sem arestas. A cada passo, o corte $(V(T),V\setminus V(T))$ não é atravessado por nenhuma aresta de $T$, e o algoritmo de Prim escolhe a mais leve que o atravessa; pelo Lema~\ref{lema:corte}, o invariante se mantém. Depois de $n-1$ passos, $T$ é geradora e está contida numa AGM, logo é uma AGM.

\resp{agm-l3} Use um heap de mínimo com os vértices fora da árvore, com chave $chave(v)=$ peso da aresta mais leve de $v$ até a árvore, e um vetor de posições para reduzir chaves. São $n$ remoções do mínimo e, para cada aresta $uv$ examinada quando $u$ entra na árvore, no máximo uma redução de chave (uma operação \emph{subir}, em $O(\log n)$). O total é $O((n+m)\log n)=O(m\log n)$, pois um grafo conexo tem $m\ge n-1$.

\resp{agm-l4} Use o vetor $chave$, sem heap (pseudocódigo da §\ref{sec:prim}): em cada um dos $n$ passos, uma varredura de $O(n)$ encontra o vértice de menor chave fora da árvore, e depois as chaves dos vizinhos são atualizadas. (a) Com listas, as atualizações percorrem cada lista uma vez, $O(m)$ no total; o custo é $O(n^2+m)=O(n^2)$. (b) Com matriz, percorre-se a linha $u$, em $O(n)$, a cada passo: $O(n^2)$.

\resp{agm-l6} Sim, como no exercício anterior. E não é possível fazer melhor nessa representação: qualquer algoritmo correto precisa ler $\Omega(n^2)$ posições da matriz. Se alguma posição $(u,v)$ não fosse lida, um adversário poderia colocar ali uma aresta levíssima, que estaria em toda AGM, e o algoritmo erraria. Logo $O(n^2)$ é ótimo com matriz, e a versão com heap seria pior, $O(n^2\log n)$.

\resp{agm-l7} (i) Executar o Kruskal com as arestas em ordem \textbf{não crescente} de peso, o que é o guloso de matroides aplicado diretamente aos pesos $c$ (Teorema~\ref{teo:gulosomat}). (ii) Trocar os pesos por $c'=C-c$, com $C>\max c$ para mantê-los positivos, ou por $-c$, e executar qualquer algoritmo de AGM: toda árvore geradora tem exatamente $n-1$ arestas, então $c'(T)=(n-1)C-c(T)$, e minimizar $c'$ é maximizar $c$. (iii) Executar o Prim escolhendo sempre a aresta \emph{mais pesada} do corte. No grafo do Exemplo~\ref{ex:agm}, a árvore geradora máxima pesa 71.

\resp{agm-sim} Kruskal: aceita $AC\,(1)$, $DF\,(2)$, $BE\,(3)$ e $CF\,(4)$; rejeita $AD\,(5)$ e $CD\,(5)$, que fecham ciclo com $A,C,F,D$; aceita $BC\,(5)$; rejeita todas as arestas de peso 6. Prim a partir de $A$: $AC\,(1)$, $CF\,(4)$, $FD\,(2)$, $CB\,(5)$ e $BE\,(3)$. O custo é \textbf{15}. A AGM é única: cada aresta escolhida é a \emph{única} mais leve de algum corte --- $AC$ no corte $\{A\}$, $DF$ em $\{D\}$, $BE$ em $\{E\}$, $CF$ em $\{A,C\}$ e $BC$ em $\{B,E\}$ ---, e uma aresta assim pertence a toda AGM, pois uma AGM sem ela poderia ser melhorada pela troca do Lema~\ref{lema:corte}.

\resp{agm-unica} Suponha duas AGMs distintas, $T_1\ne T_2$, e seja $e$ a aresta de menor peso em $T_1\triangle T_2$; digamos que $e\in T_1$. $T_2+e$ tem um ciclo, que contém alguma aresta $f\in T_2\setminus T_1$, pois do contrário o ciclo estaria todo em $T_1$. Como $e$ é a menor aresta da diferença simétrica e os pesos são distintos, $c(e)<c(f)$; então $T_2+e-f$ é uma árvore geradora mais leve que $T_2$, uma contradição.

\resp{agm-ciclo} Suponha que $e=uv$ pertença a uma AGM $T$. Retirando $e$, $T$ se divide em dois lados; o ciclo de $G$ que contém $e$ precisa atravessar esse corte de novo, por outra aresta $f$, com $c(f)<c(e)$. Então $T-e+f$ é uma árvore geradora mais leve, uma contradição.

\resp{agm-gargalo} Suponha um caminho $P$ de $u$ a $v$ em $G$ cuja maior aresta pese menos que a maior aresta $e$ do caminho de $u$ a $v$ em $T$. Retirando $e$ de $T$, $u$ e $v$ ficam em lados diferentes; $P$ atravessa esse corte por alguma aresta $f$ com $c(f)<c(e)$. Então $T-e+f$ é mais leve que $T$, uma contradição.

\resp{agm-nova} Em $T+uv$ há um único ciclo, formado pelo caminho de $u$ a $v$ em $T$ e pela nova aresta. Encontre esse caminho com uma busca em $T$, em $O(n)$, tome a aresta mais pesada do ciclo e retire-a; se for a própria $uv$, a AGM não muda. Pela propriedade do ciclo, a aresta retirada não está em nenhuma AGM do novo grafo, e as outras arestas de $G$ fora de $T$ continuam excluídas pelos mesmos ciclos de antes.

\resp{agm-prim7} Prim a partir de $A$: $AD\,(5)$, $DF\,(6)$, $AB\,(7)$, $BE\,(7)$, $EC\,(5)$ e $EG\,(9)$. Kruskal: aceita $AD\,(5)$, $CE\,(5)$, $DF\,(6)$, $AB\,(7)$, $BE\,(7)$ e $EG\,(9)$, e rejeita $BC\,(8)$, $EF\,(8)$, $BD\,(9)$, $FG\,(11)$ e $DE\,(15)$. O custo da AGM é \textbf{39}.

\resp{agm-boruvka} Com pesos distintos, a aresta mais leve que sai de uma componente $C$ é a mais leve do corte $(C,V\setminus C)$, e pertence à AGM (que é única) pela propriedade do corte; logo todas as arestas escolhidas numa rodada pertencem à AGM e podem ser acrescentadas juntas, sem formar ciclo. Cada componente se junta a pelo menos uma outra em cada rodada, então o número de componentes pelo menos cai à metade, e há no máximo $\log_2n$ rodadas. Cada rodada examina todas as arestas uma vez e une as componentes, em $O(m)$ (ou $O(m\,\alpha(n))$ com Union-Find). Total $O(m\log n)$.

\resp{agm-iguais} Toda árvore geradora tem o mesmo peso, $(n-1)c$, e portanto é mínima. Basta encontrar uma árvore geradora qualquer, por exemplo a árvore de uma busca em largura ou em profundidade, em $O(n+m)$.

\resp{agm-12} Separe as arestas em duas listas, as de peso 1 e as de peso 2, em $O(m)$; essa é a ``ordenação''. Execute o Kruskal processando primeiro a lista de peso 1 e depois a de peso 2. O custo é o das operações de Union-Find, $O(m\,\alpha(n))$.

\resp{agm-leve} A aresta mais leve $e=uv$, única, é a mais leve do corte $(\{u\},V\setminus\{u\})$, e é a única mais leve desse corte; pela troca do Lema~\ref{lema:corte}, uma AGM sem ela poderia ser melhorada, então ela está em toda AGM. A segunda mais leve, $f=xy$, se única, também está em toda AGM (num grafo sem arestas paralelas): um dos extremos de $f$, digamos $y$, não é extremo da mais leve, e então $f$ é a única aresta mais leve do corte $(\{y\},V\setminus\{y\})$. A terceira, não necessariamente: ela pode fechar um triângulo com as duas primeiras.

\resp{agm-dijk} O triângulo com $AB\,(2)$, $BC\,(2)$ e $AC\,(3)$. A partir de $A$, a árvore de caminhos mínimos usa $AB$ e $AC$ (pois $\operatorname{dist}(A,C)=3<4$), com peso total 5; a AGM é $\{AB,BC\}$, com peso 4.

\resp{agm-parar} No melhor caso, as $n-1$ primeiras arestas da ordem formam uma árvore, e só elas são examinadas. No pior caso, a última aresta da ordem é necessária (por exemplo, é a única que liga um vértice ao resto), e todas as $m$ arestas são examinadas. A complexidade assintótica não muda, porque a ordenação, de custo $O(m\log m)$, é feita de qualquer forma.

\resp{agm-listas} São $2m$ operações \Find, cada uma em $O(1)$ (um ponteiro), e $n-1$ uniões. Com ``menor na maior'', cada vértice muda de representante no máximo $\log_2n$ vezes, pois o tamanho do seu conjunto pelo menos dobra a cada mudança; o custo total das uniões é, portanto, $O(n\log n)$. O laço custa $O(m+n\log n)$.

\resp{agm-floresta} O Kruskal, executado até o fim, devolve uma \emph{floresta geradora mínima}: uma AGM de cada componente. O Prim só alcança a componente da raiz; para obter a floresta, repita-o a partir de um vértice ainda não alcançado, até cobrir todos os vértices.

\resp{agm-segunda} Seja $T'$ a segunda melhor árvore e $e\in T'\setminus T$ a aresta de menor peso dessa diferença. Pelo argumento de troca, $T+e$ tem um ciclo com uma aresta $f\in T\setminus T'$, e $T+e-f$ não é mais pesada que $T'$; repetindo o argumento, conclui-se que alguma árvore da forma $T+e-f$ é uma segunda melhor. Algoritmo: para cada vértice $u$, faça uma busca em $T$ a partir de $u$ calculando, para todo $v$, a aresta de maior peso no caminho de $u$ a $v$ em $T$, $\max(u,v)$: $O(n^2)$ no total. Depois, para cada aresta $uv\notin T$, o custo da troca é $c(uv)-\max(u,v)$; a melhor troca dá a segunda melhor árvore, em $O(n^2+m)=O(n^2)$.

\resp{agm-sim3} Kruskal: aceita $AC\,(2)$, $CE\,(3)$, $AB\,(4)$ e $DE\,(4)$; rejeita $BC\,(5)$, que fecharia o ciclo $A,B,C$; aceita $EF\,(8)$; rejeita $BD\,(10)$ e $DF\,(11)$. Prim a partir de $F$: $FE\,(8)$, $EC\,(3)$, $CA\,(2)$, $ED\,(4)$ e $AB\,(4)$; quando $A$ entra, $B$ e $D$ empatam com chave 4, e entra $D$, cuja chave foi definida antes. O custo é \textbf{21}.

\resp{agm-corte} O primeiro corte é atravessado por $AC\,(2)$, $BC\,(5)$ e $BD\,(10)$; a mais leve, $AC$, está em alguma AGM (de fato, na do exercício anterior). O segundo é atravessado por $BD\,(10)$, $DE\,(4)$ e $DF\,(11)$, e a mais leve, $DE$, também está na AGM.

\respcap{cap:bfs}
\resp{bl-arv} Fixe uma raiz candidata $r$ (tente cada vértice). \emph{Condição 1}: $\operatorname{dist}_T(r,v)=\operatorname{dist}_G(r,v)$ para todo $v$ (calcule por busca em largura em $G$); ela é necessária, porque o nível na árvore de largura é a distância. \emph{Condição 2 (ordem)}: numa busca em largura, os vértices saem da fila nível a nível; dentro de um nível, irmãos ficam juntos, e os blocos de irmãos seguem a ordem dos pais; e cada $w$ é filho do \emph{primeiro} vizinho de $w$ do nível anterior a sair da fila. Assim, a única liberdade da busca é a ordem dos filhos de cada vértice, e a condição 2 se traduz em restrições ``$x$ antes de $y$'': para cada $w$ e cada outro vizinho $v$ de $w$ no nível anterior, $pai(w)$ deve vir antes de $v$. Se $x$ e $y$ são irmãos, é uma restrição sobre a ordem dos filhos do pai deles; se não são, ela equivale a ``$pai(x)$ antes de $pai(y)$'', e sobe-se até chegar a dois irmãos. $T$ é árvore de largura com raiz $r$ se e somente se, para cada vértice, as restrições sobre a ordem dos seus filhos não formam ciclo, o que se verifica com uma ordenação topológica (Capítulo~\ref{cap:dfsdig}).

\resp{bl-pilha} Com pilha, o vértice consultado é sempre o último inserido que ainda não foi esgotado: a busca avança a partir do vértice mais recente e só volta quando ele não tem vizinhos desmarcados, que é a busca em profundidade; os pais formam a árvore de profundidade. Com fila, o vértice consultado é o mais antigo: todos os seus vizinhos desmarcados são inseridos antes que qualquer outro vértice seja examinado, ou seja, ``esgota-se o pai antes dos filhos'', que é a busca em largura.

\resp{bl-apl} (a) Busca em largura com raiz $x$: $\operatorname{dist}(x,v)=nivel(v)$, e o caminho na árvore é mínimo. (b) Modele cada casa livre como um vértice, ligando casas vizinhas livres; faça a busca a partir da entrada; se a saída for alcançada, o caminho na árvore (subindo pelos pais) é o menor, com $nivel(\text{saída})$ passos. (c) O grafo é bipartido se e somente se a floresta de largura não tem arestas-irmão nem primo, que fecham ciclos ímpares; a 2-coloração é dada pela paridade do nível.

\resp{bl-sim} Ordem de visita: $s,r,w,v,t,x,u,y$. Níveis: $s$ no nível 0; $r,w$ no 1; $v,t,x$ no 2; $u,y$ no 3. Pais: $r,w\leftarrow s$; $v\leftarrow r$; $t,x\leftarrow w$; $u\leftarrow t$; $y\leftarrow x$. Arestas que não são de árvore: $tx$ é irmão (mesmo nível e mesmo pai $w$); $xu$ é tio ($u$ já tinha sido alcançado por $t$, um nível abaixo de $x$); $uy$ é primo (mesmo nível, pais $t$ e $x$). O grafo não é bipartido: o triângulo $w,t,x$ é um ciclo ímpar.

\resp{bl-diam} O vértice $a$ mais distante de qualquer vértice é extremo de algum caminho de comprimento máximo; se não fosse, trocando trechos de caminhos obter-se-ia um caminho mais longo que o diâmetro. Então o vértice mais distante de $a$ dá o diâmetro. No exemplo, a partir do vértice 1, o mais distante é o 7 (distância 3); a partir de 7, o mais distante é o 6, à distância \textbf{5} (caminho $7,5,2,1,3,6$).

\resp{bl-num} Faça $cnt(raiz)=1$. Ao processar uma aresta $vw$ com $nivel(w)=nivel(v)+1$ (aresta-pai ou aresta-tio), faça $cnt(w):=cnt(w)+cnt(v)$. Como todos os vértices do nível $\ell$ são processados antes dos do nível $\ell+1$, $cnt(v)$ já está completo quando $v$ é processado. Custo $O(n+m)$.

\resp{bl-prova} Os vértices na fila pertencem sempre a no máximo dois níveis consecutivos, em ordem não decrescente de nível; por indução, retira-se um vértice de nível $\ell$ e entram no fim vértices de nível $\ell+1$. Seja $vw$ uma aresta, com $v$ processado antes de $w$. Se $w$ ainda não havia sido alcançado, vira filho de $v$ e $nivel(w)=nivel(v)+1$. Se já estava na fila, seu nível é $nivel(v)$ ou $nivel(v)+1$. Ele não pode já ter saído da fila, pois então $w$ teria sido processado antes de $v$. Logo $|nivel(v)-nivel(w)|\le1$. Quanto à distância: o caminho na árvore tem $nivel(v)$ arestas, então $\operatorname{dist}\le nivel(v)$; e qualquer caminho da raiz a $v$ com $k$ arestas muda o nível em no máximo 1 por aresta, então $nivel(v)\le k$.

\resp{bl-grade} Numerando linhas e colunas a partir de 0, o caminho mínimo é $(0,0)\to(0,1)\to(0,2)\to(1,2)\to(2,2)\to(2,3)\to(3,3)$, com \textbf{6} passos. Distâncias: $(0,1)$: 1; $(0,2)$: 2; $(1,2)$: 3; $(2,2)$: 4; $(2,1)$ e $(2,3)$: 5; $(2,0)$, $(3,1)$ e $(3,3)$: 6.

\resp{bl-dig} O algoritmo é o mesmo, percorrendo $N_{out}(v)$; ele calcula a distância (em arestas) da raiz a cada vértice alcançável. Uma aresta $v\to w$ sempre satisfaz $nivel(w)\le nivel(v)+1$, mas pode ``voltar'' vários níveis ($nivel(w)$ muito menor que $nivel(v)$), pois a desigualdade só vale num sentido.

\resp{bl-cam} \textbf{Se} $L(y)=0$, não há caminho. \textbf{Senão}, suba de $y$ por $pai(\cdot)$ até $x$, empilhando os vértices, e depois desempilhe-os, imprimindo de $x$ para $y$ (ou insira cada vértice à esquerda numa lista). Custo $O(nivel(y))$.

\resp{bl-custo} Cada vértice entra na fila e é processado uma vez. Com listas, processar $v$ custa $O(1+\text{grau}(v))$, e a soma dos graus é $2m$: $O(n+m)$. Com matriz, processar $v$ exige percorrer a linha inteira, $O(n)$: $O(n^2)$. As duas coincidem assintoticamente em grafos densos, com $m=\Theta(n^2)$.

\resp{bl-01} Mantenha uma fila de duas pontas com os vértices e rótulos de distância $d$. Ao retirar $v$ do início, para cada aresta $vw$ de peso $p$: se $d(v)+p<d(w)$, atualize $d(w)$ e insira $w$ no \emph{início} se $p=0$, ou no \emph{fim} se $p=1$. A fila mantém a propriedade de conter no máximo dois valores de distância consecutivos, com os menores no início, como a fila de prioridades do algoritmo de Dijkstra; cada vértice é processado um número constante de vezes, e o custo é $O(n+m)$.

\resp{bl-multi} Coloque todas as saídas na fila no início, com nível 0, como se fossem filhas de uma raiz virtual, e execute a busca em largura normalmente. O nível final de cada local é a distância até a saída mais próxima. Custo $O(n+m)$.

\resp{bl-bip} No ciclo $1,2,3,4,5$, a partir de 1: os vértices 2 e 5 ficam no nível 1, e 3 e 4 no nível 2; a aresta $34$ liga dois vértices do nível 2 com pais diferentes (2 e 5): é uma aresta-primo, e o ciclo não é bipartido. No ciclo de 6 vértices, a partir de 1: 2 e 6 no nível 1, 3 e 5 no nível 2 e 4 no nível 3, alcançado por 3; a aresta $54$ é uma aresta-tio, e não há arestas-irmão nem primo: o ciclo é bipartido.

\resp{bl-comp} Use o laço externo da busca em largura com um contador $c$: cada vez que um vértice ainda não alcançado inicia uma nova busca, incremente $c$ e rotule com $c$ todos os vértices alcançados nessa busca. Ao final, $c$ é o número de componentes. Custo $O(n+m)$.

\resp{bl-cavalo} Cada casa é um vértice, ligado às até 8 casas alcançáveis por um movimento de cavalo. Uma busca em largura a partir de a1 dá distância \textbf{6} até h8.

\resp{bl-jarros} Os estados são pares $(a,b)$, com $0\le a\le3$ e $0\le b\le5$, ligados pelas operações possíveis. Uma busca em largura a partir de $(0,0)$ encontra 4 litros com \textbf{6} operações: encher o de 5, despejar no de 3 (sobram 2 no de 5), esvaziar o de 3, despejar os 2 litros no de 3, encher o de 5 e despejar no de 3 até enchê-lo, o que deixa 4 litros no de 5.

\resp{bl-impar} Numa aresta-irmão $vw$, $v$ e $w$ têm o mesmo pai $u$: o ciclo $u,v,w$ tem 3 arestas. Numa aresta-primo $vw$, $v$ e $w$ estão no mesmo nível $\ell$; os caminhos na árvore de $v$ e de $w$ até o seu ancestral comum mais próximo $u$, no nível $j$, têm $\ell-j$ arestas cada; com a aresta $vw$, o ciclo tem $2(\ell-j)+1$ arestas, um número ímpar. Numa aresta-tio, os níveis dos extremos diferem de 1, e o ciclo formado tem comprimento par.

\resp{bl-girth} Faça uma busca em largura a partir de cada vértice $s$. Durante a busca, toda aresta $vw$ que não é de árvore fecha, com os caminhos da árvore até $s$, um passeio fechado de comprimento $nivel(v)+nivel(w)+1$, que contém um ciclo de comprimento no máximo esse valor. O menor valor dessa expressão, sobre todas as raízes e arestas, é exatamente a cintura: quando $s$ pertence a um ciclo mínimo, a aresta desse ciclo ``oposta'' a $s$ dá o seu comprimento exato. Cada busca custa $O(n+m)$, e o total é $O(n(n+m))$.

\resp{bl-k33} A partir de $a$: $x,y,z$ ficam no nível 1 (filhos de $a$) e $b,c$ no nível 2 (filhos de $x$, o primeiro a sair da fila). As arestas $yb$, $yc$, $zb$ e $zc$ são arestas-tio. Não há arestas-irmão nem primo, coerente com o fato de $K_{3,3}$ ser bipartido.

\respcap{cap:dijkstra}
\resp{dj-bfs} Por indução nos vértices removidos: com pesos 1, os rótulos dos vértices na fila de prioridades assumem, em cada momento, no máximo dois valores consecutivos, $d$ e $d+1$, e os de valor $d$ saem antes, como na fila da busca em largura. Quando $v$ sai com $L(v)=d$, cada vizinho $w$ ainda não alcançado recebe $L(w)=d+1$ e $pai(w)=v$, exatamente como na busca em largura; vizinhos já alcançados têm $L(w)\le d+1$ e não mudam. Logo $L(w)=nivel(w)=\operatorname{dist}(raiz,w)$, e a árvore é de largura (com os empates resolvidos na mesma ordem).

\resp{dj-nd} Transforme cada aresta $xy$ de peso $p$ em duas arestas direcionadas, $x\to y$ e $y\to x$, ambas de peso $p$ (e use $\infty$ onde não há aresta), e execute o algoritmo no digrafo resultante.

\resp{dj-slide} Ver o Exemplo da §\ref{sec:dijkstra}: são removidos $a\,(0)$, $b\,(1)$, $c\,(2)$, $d\,(3)$ e $e\,(5)$, e a árvore é $a\to b\to c\to d\to e$.

\resp{dj-sim} São removidos $s\,(0)$, $y\,(5)$, $z\,(7)$, $t\,(8)$ e $x\,(9)$. Pais: $y\leftarrow s$, $z\leftarrow y$, $t\leftarrow y$ e $x\leftarrow t$. Por exemplo, o caminho mínimo até $x$ é $s\to y\to t\to x$, com custo $5+3+1$.

\resp{dj-neg} Com $x\to y$ (1), $x\to z$ (2) e $z\to y$ ($-2$), o algoritmo remove $x$, depois $y$ com $L(y)=1$, e o fixa; mas $x\to z\to y$ custa 0. Somar uma constante $c$ a todos os pesos \textbf{não} resolve: um caminho com $k$ arestas ganha $kc$, o que favorece os caminhos com menos arestas. No exemplo, com $c=2$: $x\to y$ passa a valer 3 e $x\to z\to y$ passa a valer $4+0=4$; o caminho ótimo muda.

\resp{dj-conf} A probabilidade de um caminho é o produto $\prod p(e)$; maximizá-la equivale a minimizar $\sum-\log p(e)$. Como $p(e)\in(0,1]$, os pesos $-\log p(e)$ são não negativos, e o algoritmo de Dijkstra se aplica. (Alternativa: um Dijkstra que mantém o maior produto e remove sempre o maior rótulo, o que funciona porque os fatores são no máximo 1.)

\resp{dj-fontes} Crie uma \emph{superfonte} $s^*$, com arestas de peso 0 para cada $s_i$, e execute o algoritmo a partir de $s^*$; ou, de forma equivalente, comece com $L(s_i)=0$ para todas as fontes, todas inseridas na fila de prioridades. A árvore resultante também indica qual fonte é a mais próxima de cada vértice.

\resp{dj-escolha} No primeiro grafo, esparso, use heap binário: $(n+m)\log n\approx4\cdot10^5\cdot17\approx7\cdot10^6$ operações, contra $n^2=10^{10}$ com vetor. No segundo, denso ($m\approx n^2/2$), use vetor: $n^2=10^8$, contra $(n+m)\log n\approx5\cdot10^7\cdot14=7\cdot10^8$ com heap.

\resp{dj-und} São removidos $1\,(0)$, $2\,(7)$, $3\,(9)$, $6\,(11)$, $4\,(20)$ e $5\,(20)$. Pais: $2\leftarrow1$, $3\leftarrow1$, $6\leftarrow3$, $4\leftarrow3$ e $5\leftarrow6$. Por exemplo, o caminho mínimo até 5 é $1,3,6,5$, com custo $9+2+9=20$, mais barato que $1\to6\to5$, que custa 23.

\resp{dj-zero} Sim. A prova de correção só usa que o trecho de $y$ até $u$ num caminho mínimo custa pelo menos 0, o que continua valendo com pesos nulos.

\resp{dj-alvo} Pode parar assim que $t$ for removido da fila de prioridades: pelo teorema de correção, nesse momento $L(t)=\operatorname{dist}(s,t)$, e as iterações seguintes só fixam vértices mais distantes, sem alterar $L(t)$.

\resp{dj-vert} Divida cada vértice $v$ em dois, $v_{in}$ e $v_{out}$, ligados por uma aresta $v_{in}\to v_{out}$ de peso $c(v)$; cada aresta $u\to v$ do grafo original passa a ser $u_{out}\to v_{in}$. Os caminhos mínimos no novo grafo correspondem aos do problema original. (Alternativa: some $c(v)$ ao peso de cada aresta que entra em $v$.)

\resp{dj-k} O algoritmo de Dijkstra só guarda o melhor custo até cada vértice, sem considerar quantas arestas foram usadas; o caminho mais barato até um vértice intermediário pode usar arestas demais. No grafo de estados, o vértice $(v,j)$ significa ``estou em $v$ depois de usar $j$ arestas'', com arestas $(u,j)\to(v,j+1)$ de peso $W(u,v)$ para $j<k$. A resposta é $\min_{j\le k}\operatorname{dist}((s,0),(t,j))$; como esse grafo é acíclico (o $j$ cresce), basta uma programação dinâmica em $O(k(n+m))$.

\resp{dj-gargalo} Troque $L(v)+W(v,w)$ por $\min\{L(v),\ cap(v,w)\}$, comece com $L(s)=\infty$ e $L(w)=0$ para os demais, e remova sempre o vértice de \emph{maior} rótulo. A prova de correção é análoga, porque estender um caminho nunca aumenta a sua capacidade (o papel dos pesos não negativos).

\resp{dj-conta} Mantenha $cnt(s)=1$ e, ao relaxar a aresta $vw$ ao remover $v$: se $L(v)+W(v,w)<L(w)$, faça $L(w):=L(v)+W(v,w)$ e $cnt(w):=cnt(v)$; se $L(v)+W(v,w)=L(w)$, faça $cnt(w):=cnt(w)+cnt(v)$. Com pesos positivos, quando $v$ é removido, $cnt(v)$ já está completo, pois todos os predecessores de $v$ em caminhos mínimos foram removidos antes dele.

\resp{dj-dag} Ordene os vértices topologicamente (Capítulo~\ref{cap:dfsdig}) e processe-os nessa ordem, relaxando as arestas que saem de cada um, a partir de $L(s)=0$. Quando um vértice é processado, todos os seus predecessores já foram, e o seu rótulo é definitivo. O custo é $O(n+m)$, e pesos negativos não são problema, pois não há ciclos.

\resp{dj-longo} A correção do algoritmo depende de que estender um caminho nunca o torne ``melhor''; para caminhos mais longos, estender sempre melhora, e um vértice fixado cedo pode ter um rótulo errado. Além disso, a melhor extensão pode precisar repetir vértices. O problema do caminho simples mais longo é NP-difícil em grafos gerais; só em grafos acíclicos ele é fácil.

\resp{dj-dial} Os rótulos finitos estão sempre entre $d$ e $d+C$, onde $d$ é o último rótulo removido. Use um vetor de baldes (listas) indexado pelo rótulo, de tamanho $nC+1$ (ou circular, de tamanho $C+1$). Para remover o mínimo, avance um ponteiro pelos baldes até encontrar um não vazio; o ponteiro só avança, no máximo $nC$ posições no total. Reduzir a chave é mover o vértice de um balde para outro, em $O(1)$. Total $O(m+nC)$.

\resp{dj-inv} Por indução: no início, $L(s)=0=\operatorname{dist}(s,s)$ e os demais valem $\infty$. Um relaxamento só muda $L(w)$ para $L(v)+W(v,w)\ge\operatorname{dist}(s,v)+W(v,w)\ge\operatorname{dist}(s,w)$, pela hipótese e pela desigualdade triangular; e só muda quando o novo valor é menor, de modo que $L(w)$ nunca aumenta.

\resp{dj-fib} No máximo $n$ inserções, $n$ remoções do mínimo e $m$ reduções de chave (no máximo uma por aresta, quando a aresta é relaxada). Com heap de Fibonacci, inserção e redução de chave custam $O(1)$ amortizado e remoção do mínimo $O(\log n)$ amortizado: o total é $O(n+m+n\log n)=O(m+n\log n)$.

\respcap{cap:dfs}
\resp{bp-arv} Numa árvore de profundidade de um grafo não direcionado, toda aresta fora de $T$ liga um vértice a um ancestral (Exercício~\ref{ex:bp-cruz}). Reciprocamente, se para alguma raiz $r$ toda aresta de $G$ liga ancestral e descendente em $T$ enraizada em $r$, a busca que, em cada vértice, explora primeiro os filhos que ele tem em $T$ reproduz $T$. Algoritmo: para cada raiz candidata, calcule $PE$ e $PS$ de $T$ e teste, para cada aresta $uv$ de $G$, se um intervalo contém o outro: $O(n+m)$ por raiz.

\resp{bp-pilha} Empilham-se $ec$, $ca$, $ab$, $bc$ (retorno), $bd$, $dc$ (retorno) e $de$ (retorno). Ao voltar de $c$ para $e$, $back(c)=1\ge PE(e)=1$: desempilha-se até $ec$, e o bloco tem os vértices $\{a,b,c,d,e\}$. Depois empilham-se $ef$, $fg$ e $ge$ (retorno); ao voltar de $f$, $back(f)=1\ge PE(e)$, e o bloco é $\{e,f,g\}$.

\resp{bp-pontes} Depois da chamada \textsc{P}$(w)$ pela aresta de árvore $vw$: se $v$ não é a raiz e $back(w)\ge PE(v)$, marque $v$ como articulação; se $back(w)>PE(v)$, a aresta $vw$ é ponte. Para a raiz, conte os filhos e marque-a se forem pelo menos dois.

\resp{bp-demarc} Sim. Cada bloco é impresso exatamente quando a busca volta de um demarcador $w$ para o seu pai, e cada demarcador imprime um único bloco; há, portanto, uma bijeção entre demarcadores e blocos (contando os filhos da raiz, que são sempre demarcadores).

\resp{bp-apl} (a) Conte as chamadas do laço externo. (b) Execute \textsc{P}$(v)$; se $PE(w)\ne0$, suba de $w$ até $v$ pelos pais. (c) Sem aresta de retorno, é acíclico; uma aresta de retorno $vw$ fecha o ciclo formado pelo caminho de $w$ a $v$ na árvore e pela aresta $vw$. (d) Faça $cor(w):=1-cor(v)$ nas arestas de árvore; uma aresta de retorno com cores iguais revela um ciclo ímpar.

\resp{bp-sim} A busca desce $1,2,\dots,7$ em linha: $PE=1,\dots,7$ e $PS(7),\dots,PS(1)=8,\dots,14$. Retornos: $31$ e $64$. $back$: 1 para os vértices 1, 2 e 3; 4 para 4, 5 e 6; 7 para o 7. Articulações: 3, 4 e 6. Pontes: $34$ e $67$. Blocos: $\{1,2,3\}$, $\{3,4\}$, $\{4,5,6\}$ e $\{6,7\}$.

\resp{bp-cruz} Seja $vw$ uma aresta, com $v$ alcançado antes de $w$. Enquanto $v$ está ativo, ele examina a aresta $vw$; se $w$ ainda não tiver sido alcançado nesse momento, ele será alcançado dentro da chamada de $v$. Em ambos os casos, $w$ é alcançado entre $PE(v)$ e $PS(v)$, e é descendente de $v$.

\resp{bp-m} Uma floresta com $n$ vértices tem no máximo $n-1$ arestas; logo $m\ge n$ implica ciclo. Se $m\ge n$, a resposta é ``sim''; senão, $m<n$ e a busca custa $O(n+m)=O(n)$.

\resp{bp-iter} Empilhe pares (vértice, próximo índice na lista de adjacência). Olhe o topo $(v,i)$: se $i$ passou do grau de $v$, desempilhe e faça $PS(v)$; senão, avance o índice e, se o $i$-ésimo vizinho $w$ não foi alcançado, faça $pai(w):=v$, $PE(w)$ e empilhe $(w,1)$. Os tempos coincidem com os da versão recursiva.

\resp{bp-trace} $PE$ de 1 a 8: $1,2,8,3,5,9,7,4$; $PS$: $16,15,11,14,6,10,12,13$. Arestas de árvore: $12$, $24$, $48$, $85$, $87$, $73$ e $36$. Retornos: $52$, $31$ e $67$. Não há articulações: todo filho $w$ de um vértice não raiz tem $back(w)<PE(v)$, e a raiz tem um único filho.

\resp{bp-parenteses} Suponha $PE(u)<PE(v)$. Se $PE(v)<PS(u)$, então $v$ foi alcançado enquanto $u$ estava ativo e é descendente de $u$; a chamada de $v$ termina antes da de $u$, e $I(v)\subset I(u)$. Senão, $PS(u)<PE(v)$, e os intervalos são disjuntos.

\resp{bp-branco} Se $v$ é descendente de $u$, o caminho de $u$ a $v$ na árvore só tem vértices alcançados depois de $PE(u)$. Reciprocamente, se no instante $PE(u)$ há um caminho de vértices não alcançados até $v$ e $v$ não fosse descendente, tome o primeiro vértice do caminho que não é descendente de $u$; o anterior a ele é descendente e o examina antes de terminar, tornando-o descendente também, uma contradição.

\resp{bp-comp} Três componentes: $\{1,2,3\}$, $\{4,5\}$ e $\{6,7,8,9\}$. As raízes são 1, 4 e 6.

\resp{bp-cicloex} A primeira aresta de retorno é $52$: subindo de 5 por $pai$ até 2 ($5\to8\to4\to2$) e fechando com $52$, obtém-se o ciclo $2,4,8,5,2$.

\resp{bp-bipex} A partir de 1: $cor(1)=0$, $cor(2)=1$, $cor(3)=0$, $cor(4)=1$. A aresta de retorno $13$ liga dois vértices de cor 0: o ciclo $1,2,3,1$ é ímpar, e o grafo não é bipartido.

\resp{bp-arvore} Numa árvore, todas as arestas são pontes, e as articulações são os vértices de grau pelo menos 2. Num ciclo, não há pontes nem articulações. Num caminho com $n\ge3$ vértices, todas as arestas são pontes, e as articulações são os vértices internos.

\resp{bp-blocos} Articulações: 3 e 4. Pontes: $34$ e $47$. Blocos: $\{1,2,3\}$, $\{3,4\}$, $\{4,5,6\}$ e $\{4,7\}$.

\resp{bp-biconexo} Faça uma busca em profundidade: o grafo é biconexo se ela alcança todos os vértices e não encontra articulações, em $O(n+m)$. Num grafo biconexo com pelo menos 3 vértices não há pontes, e uma aresta que não é ponte está em algum ciclo.

\resp{bp-labirinto} No ciclo $1,2,3,4$ com a aresta $14$, a busca a partir de 1 com vizinhos em ordem crescente percorre $1\to2\to3\to4$, encontrando o caminho de 1 a 4 com 3 arestas, enquanto a aresta $14$ é um caminho com uma aresta.

\resp{bp-euler} Algoritmo de Hierholzer: a partir de um vértice, siga arestas ainda não usadas, marcando-as, até não haver saída (como todos os graus são pares, isso só acontece no vértice inicial); ao voltar na recursão, insira cada vértice no início do circuito. Os subciclos encontrados ao retroceder são encaixados automaticamente. Cada aresta é usada uma vez: $O(n+m)$.

\respcap{cap:dfsdig}
\resp{dd-sorv} Tome um caminho direcionado de comprimento máximo; o seu último vértice $u$ é um sumidouro, pois uma aresta $u\to x$ com $x$ fora do caminho daria um caminho maior, e com $x$ no caminho daria um ciclo. Algoritmo linear: guarde o grau de saída de cada vértice e as listas de \emph{entrada}, e ponha numa fila todos os sumidouros. Ao retirar um sumidouro $u$, coloque-o no início da ordenação e decremente o grau de saída de cada $x$ com $x\to u$; quem chegar a zero entra na fila. Cada aresta é processada uma vez: $O(n+m)$. Se a fila esvaziar antes de todos os vértices saírem, há um ciclo.

\resp{dd-porque} Num digrafo acíclico não há arestas de retorno. Numa aresta $vw$ de árvore ou de avanço, $w$ é descendente de $v$ e sai antes dele; numa de cruzamento, $w$ já tinha saído quando $v$ começou. Logo $PS(w)<PS(v)$ para toda aresta, e a ordem por $PS$ decrescente põe $v$ antes de $w$. No exemplo, com as raízes em ordem alfabética, obtém-se $i,j,h,e,f,g,b,a,d,c$.

\resp{dd-papel} Sim. Com as arestas $a\to b$, $a\to c$ e $b\to c$: se a busca a partir de $a$ visita $b$ primeiro, $c$ vira filho de $b$ e $ac$ é de avanço; se visita $c$ primeiro, $c$ termina antes de $b$ começar, e $bc$ passa a ser de cruzamento (e $ac$, de árvore). As componentes são definidas pela existência de caminhos nos dois sentidos, uma propriedade do digrafo, e o algoritmo é correto qualquer que seja a ordem da busca.

\resp{dd-sim} $PE$ dos vértices 1 a 7: $1,2,3,5,6,7,13$; $PS$: $12,11,4,10,9,8,14$. Árvore: $1\to2$, $2\to3$, $2\to4$, $4\to5$ e $5\to6$. Retorno: $3\to1$ e $6\to4$. Avanço: $1\to3$. Cruzamento: $7\to3$ e $7\to6$. $old$: 1 para os vértices 1, 2 e 3; 5 para 4, 5 e 6; 13 para o 7. Componentes, na ordem em que são encontradas: $\{4,5,6\}$ (quando 4 sai), $\{1,2,3\}$ (quando 1 sai) e $\{7\}$.

\resp{dd-roupa} Relógio, meias, cueca, camisa, gravata, calça, sapatos, cinto, paletó. Qualquer outra ordem que respeite as dependências também é válida.

\resp{dd-cond} Se houvesse um ciclo $C_1\to C_2\to\cdots\to C_k\to C_1$ com $k\ge2$, de qualquer vértice de $C_1$ se chegaria a qualquer vértice de $C_2$ (seguindo o ciclo de componentes e andando dentro de cada uma) e vice-versa. Então $C_1\cup C_2$ seria fortemente conexo, contradizendo a maximalidade das componentes.

\resp{dd-unica} Se há um caminho hamiltoniano $v_1\to v_2\to\cdots\to v_n$, toda ordenação topológica precisa pôr $v_i$ antes de $v_{i+1}$, e portanto é única. Se não há, em alguma ordenação topológica dois vértices consecutivos $v_i$ e $v_{i+1}$ não são ligados por aresta, e trocá-los dá outra ordenação válida. Teste: obtenha uma ordenação topológica e verifique se cada par de vértices consecutivos é uma aresta, em $O(n+m)$.

\resp{dd-prereq} Monte o digrafo com uma aresta $a\to b$ se $a$ é pré-requisito de $b$. Se houver ciclo (aresta de retorno numa busca), é impossível. Senão, o menor número de semestres é o número de vértices do caminho mais longo: $sem(v)=1+\max\{sem(u):u\to v\}$ (1 se não há pré-requisitos), calculado em ordem topológica em $O(n+m)$ --- é a programação dinâmica de caminho mais longo num digrafo acíclico do Capítulo~\ref{cap:pd}.

\resp{dd-classif} $PE$ dos vértices 1 a 6: $1,2,3,4,9,10$; $PS$: $8,7,6,5,12,11$. Árvore: $1\to2$, $2\to3$, $3\to4$ e $5\to6$. Retorno: $4\to2$ e $6\to5$. Avanço: $1\to4$. Cruzamento: $5\to3$.

\resp{dd-cfcex} $old$: 1 para o vértice 1; 2 para 2, 3 e 4; 9 para 5 e 6. Componentes, na ordem em que são encontradas: $\{2,3,4\}$ (quando 2 sai), $\{1\}$ e $\{5,6\}$. A aresta $5\to3$ é de cruzamento, mas 3 já saiu da pilha, e não altera $old(5)$.

\resp{dd-kahn} As fontes iniciais são $a$ e $f$; sai $a$, e $b$ vira fonte; sai $b$; sai $f$, e $c$ vira fonte; sai $c$, e $d$ vira fonte; saem $d$ e $e$. Ordem: $a,b,f,c,d,e$.

\resp{dd-ciclo} Como a aresta é de retorno, $w$ é ancestral de $v$. Suba de $v$ até $w$ pelos ponteiros $pai$, guardando os vértices; o ciclo é $w\to\cdots\to v\to w$, impresso na ordem inversa da subida, seguido de $w$.

\resp{dd-forte} Se o digrafo é fortemente conexo, de $s$ se alcança todo vértice e todo vértice alcança $s$, o que equivale a alcançar todos a partir de $s$ no digrafo invertido. Reciprocamente, se isso vale, para quaisquer $u$ e $v$ há caminhos $u\leadsto s\leadsto v$ e $v\leadsto s\leadsto u$. O teste faz uma busca no digrafo e outra no digrafo invertido, em $O(n+m)$.

\resp{dd-kosaraju} Para duas componentes $C$ e $C'$ com uma aresta de $C$ para $C'$, o maior $PS$ de $C$ é maior que o maior $PS$ de $C'$ (se a busca entra primeiro em $C$, termina $C'$ antes; se entra primeiro em $C'$, termina $C'$ inteira antes de alcançar $C$). No digrafo invertido, as arestas vão de $C'$ para $C$. Começando pelo vértice de maior $PS$, a segunda busca está numa componente da qual, no digrafo invertido, só saem arestas para componentes já visitadas; por isso, cada árvore da segunda busca é exatamente uma componente.

\resp{dd-conta} São 4: $a,b,c,d,e$; $a,b,c,e,d$; $b,a,c,d,e$; $b,a,c,e,d$.

\resp{dd-longo} Em ordem topológica $s,a,b,c,t$: $L(s)=0$, $L(a)=3$, $L(b)=2$, $L(c)=\max\{3+4,\ 2+6\}=8$ e $L(t)=\max\{3+2,\ 2+3,\ 8+1\}=9$. O caminho de maior peso é $s\to b\to c\to t$, com peso \textbf{9}.

\resp{dd-semestres} São necessários \textbf{3} semestres: no 1º, $C_1$ e $P_1$; no 2º, $C_2$, $A_1$ e $P_2$; no 3º, $C_3$. O caminho mais longo, $C_1\to C_2\to C_3$ (ou $C_1\to A_1\to C_3$), tem 3 vértices.

\resp{dd-minarestas} Num digrafo fortemente conexo com mais de um vértice, todo vértice tem grau de entrada e de saída positivos. Cada fonte do grafo das componentes precisa receber pelo menos uma nova aresta entrando nela, e cada sumidouro, pelo menos uma saindo dele; uma nova aresta entra em no máximo uma fonte e sai de no máximo um sumidouro. Logo são necessárias pelo menos $\max\{f,s\}$ arestas. (Sabe-se que esse número também é suficiente.)

\resp{dd-lex} Use o algoritmo de Kahn com um heap de mínimo no lugar da fila: entre as fontes disponíveis, retire sempre a de menor nome. Custo $O((n+m)\log n)$.

\resp{dd-nd} Num grafo não direcionado, toda aresta $vw$ é examinada a partir dos dois extremos; quando o primeiro dos dois a ser alcançado, digamos $v$, a examina, $w$ ainda não foi alcançado ou está ativo como descendente, e a aresta vira de árvore ou de retorno. Num digrafo, a aresta $v\to w$ só é examinada a partir de $v$, e $w$ pode já ter sido completamente explorado antes, numa outra parte da floresta, o que dá as arestas de cruzamento; ou pode ser um descendente alcançado por outro caminho, o que dá as de avanço.

\respcap{cap:hash}
\resp{hs-pesq} O Apêndice~\ref{cap:hash} serve de roteiro: motivação (dicionários e endereçamento direto), funções de divisão, multiplicação, dobra e polinomial, hashing universal, encadeamento exterior e interior, sondagens linear e quadrática e duplo hashing, e a complexidade em função do fator de carga, com pior caso $\Theta(n)$. O texto da pesquisa é um trabalho individual, a ser escrito com as próprias palavras.

\resp{hs-abertos} Posições $0..10$; ``--'' indica posição vazia.\\
\hspace*{1em}(a) Linear: $22,88,-,-,4,15,28,17,59,31,10$.\\
\hspace*{1em}(b) Quadrática: $22,-,88,17,4,-,28,59,15,31,10$.\\
\hspace*{1em}(c) Duplo hashing: $22,-,59,17,4,15,28,88,-,31,10$.

\resp{hs-encad} Lista 1: $28,19,10$; lista 2: $20$; lista 3: $12$; lista 5: $5$; lista 6: $15,33$; lista 8: $17$; as demais ficam vazias. A maior lista tem 3 elementos.

\resp{hs-primo} Com $m=2^r$, $k\bmod m$ são apenas os $r$ bits menos significativos de $k$: chaves que diferem só nos bits mais altos (por exemplo, múltiplos de $2^r$, ou strings codificadas que diferem só no começo) caem todas no mesmo endereço. Um primo longe de potências de 2 faz o resultado depender de todos os bits da chave.

\resp{hs-remove} Uma busca por $k$ para ao encontrar uma posição vazia. Se uma chave no meio da sequência de sondagem de $k$ for apagada, a busca pararia ali e não encontraria $k$, que está mais adiante. A solução é marcar a posição como ``removida'': a busca continua passando por ela, e a inserção pode reutilizá-la. Com muitas remoções, o desempenho cai e convém reconstruir a tabela.

\resp{hs-carga} $\alpha=1000/500=2$: as listas têm, em média, 2 elementos. Uma busca sem sucesso percorre uma lista inteira, com 2 comparações em média, mais o cálculo do hash: $\Theta(1+\alpha)$.

\resp{hs-sondagens} No máximo $\frac1{1-\alpha}$: 2 sondagens com $\alpha=0{,}5$ e 10 com $\alpha=0{,}9$. O custo cresce rapidamente quando a tabela se aproxima de cheia; por isso é comum redimensionar quando $\alpha$ passa de algo entre $0{,}5$ e $0{,}75$.

\resp{hs-mult} $123456\cdot A\approx76300{,}0041$, cuja parte fracionária é $0{,}0041$: $h=\floor{4{,}1}=4$. $123457\cdot A\approx76300{,}6221$: $h=\floor{622{,}1}=622$. Chaves vizinhas vão para endereços distantes.

\resp{hs-universal} Uma família $\mathcal H$ de funções de $U$ em $\{0,\dots,m-1\}$ é universal se, para quaisquer chaves distintas $k_1\ne k_2$, no máximo $|\mathcal H|/m$ funções da família as fazem colidir, isto é, a probabilidade de colisão é no máximo $1/m$ quando a função é sorteada. Como o sorteio é feito depois de a entrada ser fixada, nenhuma entrada específica é ruim para a maioria das funções, e o número esperado de colisões de uma chave é no máximo $\alpha$, qualquer que seja o conjunto de chaves.

\resp{hs-string} $h:=0$; \textbf{para} $i:=0$ \textbf{até} $L-1$: $h:=(h\cdot b+s_i)\bmod m$. Cada passo mantém $h<m$, sem calcular potências grandes. Se $b$ e $m$ tivessem um fator comum $d$, os termos $s_ib^j$ com $j\ge1$ seriam múltiplos de $d$ módulo $m$, e só uma fração dos endereços seria usada, aumentando as colisões.

\resp{hs-doissoma} Percorra o vetor mantendo uma tabela hash dos valores já vistos; para cada $a_i$, verifique se $x-a_i$ está na tabela e depois insira $a_i$. São $n$ buscas e $n$ inserções, em $O(n)$ esperado.

\resp{hs-repetido} Percorra a sequência inserindo cada valor numa tabela hash; o primeiro valor que já estiver na tabela no momento da inserção é a resposta. $O(n)$ esperado.

\resp{hs-ordem} A função hash espalha as chaves sem relação com a ordem, então o sucessor de $k$ pode estar em qualquer posição, e é preciso examinar todas as chaves: $\Theta(n)$. Na AVL, ele é encontrado descendo ou subindo um caminho, em $O(\log n)$. Mínimo, máximo, percurso em ordem e consultas por intervalo têm a mesma limitação nas tabelas hash.

\resp{hs-rehash} É a análise do vetor dinâmico do Capítulo~\ref{cap:heapav}: uma reconstrução com $k$ chaves custa $O(k)$, mas só acontece depois de $\Theta(k)$ inserções baratas desde a anterior. Com o potencial $\Phi=$ constante vezes (número de chaves inseridas desde a última reconstrução), cada inserção custa $O(1)$ amortizado.

\resp{hs-agrup} Na sondagem linear, uma chave cujo endereço cai em qualquer posição de um grupo contíguo de posições ocupadas vai parar na primeira posição livre depois do grupo. Assim, a posição logo após um grupo de $k$ ocupadas é escolhida com probabilidade $(k+1)/m$, em vez de $1/m$, e grupos grandes tendem a crescer ainda mais e a se fundir.

\resp{hs-duplo} A sequência $h_1(k)+i\,h_2(k)\pmod m$ só visita $m/\mathrm{mdc}(h_2(k),m)$ posições distintas. Com $m=10$ e $h_2(k)=5$, ela alterna entre apenas 2 posições, e a inserção pode falhar mesmo com a tabela quase vazia. Por isso usa-se $m$ primo, ou $m$ potência de 2 com $h_2(k)$ sempre ímpar.

\resp{hs-remencad} Use listas \emph{duplamente} encadeadas: com ponteiros para o anterior e o próximo, o nó é retirado da lista em $O(1)$, sem percorrê-la. (Sem o ponteiro para o nó, a remoção exige primeiro uma busca, de custo $\Theta(1+\alpha)$ esperado.)

\resp{hs-perfeito} No primeiro nível, uma função universal distribui as $n$ chaves em $m=n$ posições; a posição $j$ recebe $n_j$ chaves. No segundo nível, cada posição tem uma tabela própria de tamanho $n_j^2$ e uma função universal sorteada até não haver colisões (com tamanho $n_j^2$, isso acontece com probabilidade pelo menos $1/2$ a cada sorteio). A memória total esperada é $O(\sum_jn_j^2)=O(n)$, e cada busca consulta duas posições: $O(1)$ no pior caso.

\resp{hs-distintos} Insira cada valor numa tabela hash, contando apenas quando ele ainda não estava presente: $O(n)$ esperado. Em $O(n\log n)$ no pior caso: ordene (por exemplo, com mergesort) e conte as posições em que o valor muda.

\resp{hs-tamanho} São necessárias pelo menos $10^6/0{,}75\approx1\,333\,334$ posições. Para o método da divisão, escolha um primo pouco acima desse valor e longe de potências de 2, como $1\,333\,357$.

\resp{hs-busca} 31: $h(31)=5$, e a busca examina as posições $5,6,\dots,10$: 6 sondagens. 57: $h(57)=5$, e a busca atravessa o agrupamento inteiro, das posições 5 a 11, e só para na posição 12, vazia: 8 sondagens. 26: $h(26)=0$, posição vazia: 1 sondagem. Uma chave cujo endereço cai dentro do agrupamento paga pelo comprimento dele, mesmo quando não está na tabela.

\resp{hs-reinsere} $h(70)=70\bmod13=5$, ocupada pelo 18; a posição 6, marcada como removida, é reaproveitada, e o 70 fica nela, depois de 2 sondagens. Uma busca pelo 32 começa em $h(32)=6$, encontra o 70, passa pelo 59 na posição 7 e acha o 32 na posição 8: 3 sondagens. Entre a remoção e a inserção, era a marca de ``removida'' que fazia a busca pelo 32 continuar depois da posição 6.

\resp{hs-univsim} Com $(2,7)$: $13\bmod17=13$ e $13\bmod6=1$; $23\bmod17=6$ e $6\bmod6=0$; $35\bmod17=1$ e $1\bmod6=1$. Os endereços são 1, 0 e 1, e as chaves 3 e 14 colidem. Com $(11,3)$: $36\bmod17=2$, $91\bmod17=6$ e $157\bmod17=4$, que dão os endereços 2, 0 e 4, sem colisão. A colisão entre 3 e 14 depende da função sorteada.

\resp{hs-horner} ``casa'' (3, 1, 19, 1): $h=3$; $(3\cdot31+1)\bmod97=94$; $(94\cdot31+19)\bmod97=23$; $(23\cdot31+1)\bmod97=35$. ``saca'' (19, 1, 3, 1): $h=19$; $(19\cdot31+1)\bmod97=8$; $(8\cdot31+3)\bmod97=57$; $(57\cdot31+1)\bmod97=22$. Os endereços são 35 e 22: não há colisão, porque cada posição da palavra é multiplicada por uma potência diferente de $b$. Uma função que só somasse as letras daria $3+1+19+1=24$ para as duas, e todos os anagramas colidiriam.

\normalsize


\backmatter
% ================================================================
\chapter{Referências bibliográficas}
% ================================================================
\markboth{Referências bibliográficas}{}
\section*{Nota sobre as fontes}
Este livro foi elaborado a partir do material didático da disciplina de Algoritmos do IMPA Tech, ministrada pelo Prof.\ Uéverton Souza --- slides, notas de aula, monitorias, listas de exercícios, avaliações e propostas de trabalhos de pesquisa ---, e complementado com as obras listadas abaixo. Os exercícios marcados como \textsf{[Material de apoio]} vêm de listas de exercícios de outras disciplinas de análise de algoritmos, indicadas aos alunos como material complementar.

\section*{Livros}
\begin{enumerate}[label={[\arabic*]},leftmargin=2.2em]
  \item \textsc{Cormen}, T.\,H.; \textsc{Leiserson}, C.\,E.; \textsc{Rivest}, R.\,L.; \textsc{Stein}, C. \emph{Introduction to Algorithms}. 3.\ ed. Cambridge, MA: MIT Press, 2009.
  \item \textsc{Szwarcfiter}, J.\,L.; \textsc{Markenzon}, L. \emph{Estruturas de Dados e seus Algoritmos}. 3.\ ed. Rio de Janeiro: LTC, 2010.
  \item \textsc{Szwarcfiter}, J.\,L. \emph{Teoria Computacional de Grafos: os algoritmos}. Rio de Janeiro: Elsevier, 2018.
  \item \textsc{Kleinberg}, J.; \textsc{Tardos}, É. \emph{Algorithm Design}. Boston: Pearson/Addison-Wesley, 2006.
  \item \textsc{Oxley}, J.\,G. \emph{Matroid Theory}. 2.\ ed. Oxford: Oxford University Press, 2011.
  \item \textsc{Schrijver}, A. \emph{Combinatorial Optimization: Polyhedra and Efficiency}. Berlin: Springer, 2003.
\end{enumerate}

\section*{Material da disciplina e material de apoio}
\begin{enumerate}[label={[\arabic*]},leftmargin=2.2em,resume]
  \item \textsc{Souza}, U.\,S. \emph{Algoritmos}: slides, notas de aula, listas de exercícios, avaliações e propostas de pesquisa da disciplina. Rio de Janeiro: IMPA Tech.
  \item Material das monitorias da disciplina de Algoritmos do IMPA Tech (heaps binomiais e de Fibonacci, Union-Find, árvores AVL, programação dinâmica e revisão).
  \item Listas de exercícios de apoio: \emph{Análise e Projeto de Algoritmos} (Departamento de Computação, UFF/PURO); \emph{MAC5711 -- Análise de Algoritmos} (IME-USP); \emph{MC458 -- Projeto e Análise de Algoritmos I} (Unicamp).
  \item Problemas de juízes online citados nos exercícios: \emph{Beecrowd} (problemas 1191, 1195, 1200, 1201 e 2330), \emph{SPOJ} (problema SID) e \emph{Codeforces} (problema 1131F).
\end{enumerate}


\end{document}
